\documentclass[11pt,a4paper]{article}
\usepackage[margin=23mm,headheight=15pt]{geometry}
\usepackage{iftex}
\ifPDFTeX
  \usepackage[T2A]{fontenc}
  \usepackage[utf8]{inputenc}
  \usepackage[english,russian]{babel}
\else
  \usepackage{fontspec}
  \defaultfontfeatures{Ligatures=TeX}
  \usepackage{polyglossia}
  \setdefaultlanguage{russian}
  \setotherlanguage{english}
  \setmainfont{NotoSerif-Regular.ttf}[BoldFont=NotoSerif-Bold.ttf,ItalicFont=NotoSerif-Italic.ttf,BoldItalicFont=NotoSerif-BoldItalic.ttf]
  \setsansfont{NotoSans-Regular.ttf}[BoldFont=NotoSans-Bold.ttf,ItalicFont=NotoSans-Italic.ttf,BoldItalicFont=NotoSans-BoldItalic.ttf]
  \setmonofont{NotoSansMono-Regular.ttf}
  \newfontfamily\cyrillicfonttt{NotoSansMono-Regular.ttf}
  \newfontfamily\cyrillicfont{NotoSerif-Regular.ttf}[BoldFont=NotoSerif-Bold.ttf,ItalicFont=NotoSerif-Italic.ttf,BoldItalicFont=NotoSerif-BoldItalic.ttf]
  \newfontfamily\cyrillicfontsf{NotoSans-Regular.ttf}[BoldFont=NotoSans-Bold.ttf,ItalicFont=NotoSans-Italic.ttf,BoldItalicFont=NotoSans-BoldItalic.ttf]
\fi
\usepackage{amsmath,amssymb}
\usepackage{array,longtable,booktabs,tabularx}
\usepackage[table]{xcolor}
\usepackage{enumitem}
\usepackage{fancyhdr}
\usepackage{xurl}
\usepackage{hyperref}
\hypersetup{unicode=true,colorlinks=true,linkcolor=blue!45!black,urlcolor=blue!45!black,pdftitle={Рекомендации AHA/ACC/ADA/ASN 2026 по сердечно-сосудисто-почечно-метаболическому синдрому. Полный адаптированный перевод},pdfauthor={Неофициальный перевод Вербицкий С.В., orcid.org/0009-0008-5706-0174}}
\setlength{\parindent}{0pt}
\setlength{\parskip}{0.55\baselineskip}
\setlength{\emergencystretch}{3em}
\setlength{\LTpre}{0.5\baselineskip}
\setlength{\LTpost}{0.5\baselineskip}
\renewcommand{\arraystretch}{1.16}
\setlength{\tabcolsep}{4pt}
\setlist{nosep,topsep=5pt,leftmargin=*}
\newcolumntype{P}[1]{>{\raggedright\arraybackslash}p{#1}}
\newcommand{\src}[1]{\textsuperscript{[#1]}}
\newcommand{\editorial}[1]{\par\begingroup\small\color{black!70}\textbf{Примечание к переводу.} #1\par\endgroup}
\newcommand{\heading}[1]{\par\medskip\textbf{#1}\par}
\newcommand{\unitgfr}{мл/мин/1,73~м$^2$}
\pagestyle{fancy}
\fancyhf{}
\fancyhead[L]{\small Рекомендации AHA/ACC/ADA/ASN 2026}
\fancyhead[R]{\small}
\fancyfoot[L]{\scriptsize Неофициальный перевод}
\fancyfoot[R]{\thepage}
\setcounter{tocdepth}{3}
\makeatletter
\renewcommand{\@pnumwidth}{2em}
\makeatother

\begin{document}
\begin{titlepage}
\centering
{\large AHA / ACC / ADA / ASN\par}
\vspace{1cm}
{\LARGE\bfseries Рекомендации 2026 года по профилактике, выявлению, обследованию и ведению пациентов с сердечно-сосудисто-почечно-метаболическим синдромом\par}
\vspace{0.8cm}
{\large Доклад Объединённого комитета Американской коллегии кардиологов и Американской кардиологической ассоциации по клиническим практическим рекомендациям\par}
\vspace{1cm}
{\Large Адаптированный перевод на русский язык. Вербицкий С.В., orcid.org/0009-0008-5706-0174\par}
\vspace{0.4cm}
{\large\bfseries\par}
\vspace{0.4cm}
{\normalsize Включает вводные материалы, все основные разделы 1--9, таблицы 1--21, текстовую адаптацию рисунков 1--20, сведения о сотрудниках организаций и составах комитетов. Включены сведения о публикации, весь список литературы (1054 записи с переводами названий) и приложения 1--2.\par}
\vfill
\begin{flushleft}
\textbf{} \textit{Circulation}. 2026;154:e50--e158. DOI: \nolinkurl{10.1161/CIR.0000000000001453}. Дата: 28 июля 2026 года.

\textbf{Оригинальное название:}\par
{\small 2026 AHA/ACC/ADA/ASN Guideline for the Prevention, Detection, Evaluation, and Management of Cardiovascular-Kidney-Metabolic Syndrome: A Report of the American College of Cardiology/American Heart Association Joint Committee on Clinical Practice Guidelines.}

\textbf{Статус:} перевод не утверждён AHA, ACC, ADA или ASN и не проходил независимого медицинского рецензирования. Полнота охвата не означает независимого подтверждения точности всех медицинских формулировок; перед клиническим применением необходима сверка с оригиналом.
\end{flushleft}
\vspace{0.4cm}
{\small Дата подготовки текущей версии: 5 сентября 2026 года.}
\end{titlepage}

\section*{Редакционные принципы и границы версии}
\addcontentsline{toc}{section}{Редакционные принципы и границы версии}
Адаптация касается русского синтаксиса, последовательности размещения таблиц и удобства чтения. Она не означает адаптацию рекомендаций к российской системе здравоохранения, зарегистрированным в России показаниям или национальным клиническим рекомендациям.

Основной термин: \textbf{сердечно-сосудисто-почечно-метаболический синдром (ССПМ-синдром; в оригинале CKM syndrome)}. Английские сокращения сохранены там, где они необходимы для сопоставления с оригиналом, шкалами риска и исследованиями. Численные пороги, классы рекомендаций и уровни доказательности передаются без изменения. Значения в мг/дл и мг/г не заменяются пересчитанными значениями в единицах СИ. Десятичная точка в русском тексте заменена запятой.

Надстрочные номера в квадратных скобках воспроизводят ссылки оригинала; их нумерация относится к соответствующему разделу оригинальной статьи, а не к сквозному списку этого перевода. Список литературы воспроизведён полностью, от преамбулы до конца раздела 8, с русскими переводами названий и сохранением оригинальных библиографических записей. В исходном списке нет отдельного блока литературы к разделу 9. Повторяющиеся записи не объединены; выходные данные не дополнены сведениями из других источников. Таблица 1 представляет перечень публикаций, упомянутых в самом руководстве: их тексты сюда не добавляются.

Рисунки 1--20 переданы русскоязычными структурированными схемами и пояснениями, без воспроизведения декоративных иллюстраций. Таблицы перенесены ближе к первому содержательному обсуждению. Повторяющиеся журнальные колонтитулы и отметки о скачивании не воспроизводятся. Редакторские примечания явно отделены от текста авторов.

\textbf{Ограничение использования.} Перевод предназначен для работы с исходным руководством, а не для самостоятельного назначения лечения. Перед клиническим использованием формулировки, дозировки, противопоказания и показания следует сверять с оригиналом и инструкциями к лекарственным препаратам.

\tableofcontents
\clearpage

\section*{Сведения о разработке и авторах}
\addcontentsline{toc}{section}{Сведения о разработке и авторах}
Документ разработан совместно с Американской диабетической ассоциацией (ADA), Ассоциацией по вопросам ожирения --- подразделением ADA --- и Американским обществом нефрологии (ASN) и одобрен этими организациями.

\heading{Члены авторского комитета}
{\small
Chiadi E. Ndumele, MD, PhD, MHS, FAHA --- председатель; Fatima Rodriguez, MD, MPH, FAHA, FACC --- заместитель председателя; Dave L. Dixon, PharmD, FAHA, FACC, BCACP, CLS --- представитель Объединённого комитета; Sadiya S. Khan, MD, MSc, FAHA, FACC --- представитель Объединённого комитета; Debabrata Mukherjee, MD, MS, FAHA, FACC --- бывший член Объединённого комитета, входивший в его состав во время подготовки документа;

Mandeep Bajaj, MBBS, FAHA --- представитель ADA; Sripal Bangalore, MD, MHA, FAHA, FACC; Biykem Bozkurt, MD, PhD, FAHA, FACC; Khadijah Breathett, MD, MS, FAHA, FACC; Shoa L. Clarke, MD, PhD, FAHA; Ian H. de Boer, MD, MS, FAHA --- представитель ASN; David H. Ellison, MD, FASN, FAHA; Lorraine S. Evangelista, PhD, RN, FAHA; Sean P. Heffron, MD, MS, MSc, FAHA; Dhruv S. Kazi, MD, MSc, MS, FAHA, FACC; Ambar Kulshreshtha, MD, PhD, FAHA;

Ildiko Lingvay, MD, MPH, MSCS, FAHA --- представитель Ассоциации по вопросам ожирения; Cecilia C. Low Wang, MD, FAHA; Claudia A. Mercado --- представитель интересов пациентов; John Magaña Morton, MD, MPH, MHA; Ian J. Neeland, MD, FAHA; Neha Pagidipati, MD, MPH, FAHA, FACC; Tiffany M. Powell-Wiley, MD, MPH, FAHA; Janani Rangaswami, MD, FAHA; Goutham Rao, MD, FAHA; Nosheen Reza, MD, FAHA, FACC --- представитель JCPM; Anum Saeed, MD, FAHA; Wendy St. Peter, PharmD, FAHA;

J. Bradley Starks --- представитель интересов пациентов; Madeline Sterling, MD, MPH, FAHA; Amy W. Talbot, MPH --- бывший сотрудник совместной структуры AHA/ACC, работавший в ней в период подготовки документа; Andrew H. Tran, MD, MPH, MS, FAHA; Katherine R. Tuttle, MD, FASN, FAHA; Lisa B. VanWagner, MD, MSc, FAHA; Amanda R. Vest, MBBS, MPH, FAHA, FACC; Salim S. Virani, MD, PhD, FAHA, FACC.
}

Члены авторского комитета обязаны воздерживаться от голосования по разделам, к которым могут относиться их конкретные связи с индустрией; подробная информация приведена в приложении 1 оригинала. Состав комитета рецензентов и Объединённого комитета AHA/ACC по клиническим практическим рекомендациям приведён на странице e123 оригинала.

\heading{Библиографическая запись, рекомендуемая AHA}
{\small Ndumele CE, Rodriguez F, Dixon DL, Khan SS, Mukherjee D, Bajaj M, Bangalore S, Bozkurt B, Breathett K, Clarke SL, de Boer IH, Ellison DH, Evangelista LS, Heffron SP, Kazi DS, Kulshreshtha A, Lingvay I, Low Wang CC, Mercado CA, Morton JM, Neeland IJ, Pagidipati N, Powell-Wiley TM, Rangaswami J, Rao G, Reza N, Saeed A, St. Peter W, Starks JB, Sterling M, Talbot AW, Tran AH, Tuttle KR, VanWagner LB, Vest AR, Virani SS. 2026 AHA/ACC/ADA/ASN guideline for the prevention, detection, evaluation, and management of cardiovascular-kidney-metabolic syndrome: a report of the American College of Cardiology/American Heart Association Joint Committee on Clinical Practice Guidelines. \textit{Circulation}. 2026;154:e50--e158. DOI: \nolinkurl{10.1161/CIR.0000000000001453}.}

Подкаст и его расшифровка доступны в дополнительных материалах к оригинальной статье. Эти дополнительные материалы не входят в настоящий перевод.

© 2026 American College of Cardiology Foundation и American Heart Association, Inc.

\section*{Аннотация}
\addcontentsline{toc}{section}{Аннотация}
\textbf{Цель.} «Рекомендации AHA/ACC/ADA/ASN 2026 года по профилактике, выявлению, обследованию и ведению пациентов с сердечно-сосудисто-почечно-метаболическим синдромом» отменяют, заменяют и расширяют «Рекомендации AHA/ACC/TOS 2013 года по ведению взрослых с избыточной массой тела и ожирением». Документ предназначен прежде всего для клиницистов, оказывающих помощь пациентам на всём протяжении спектра ССПМ-синдрома --- состояния, характеризующегося взаимосвязями между метаболическими факторами риска, включая ожирение и сахарный диабет 2 типа, хронической болезнью почек и сердечно-сосудистыми заболеваниями.

\textbf{Методы.} С 29 октября 2024 года по 14 апреля 2025 года проводили всесторонний поиск клинических исследований, систематических обзоров и метаанализов, а также других доказательств, полученных в исследованиях с участием людей и опубликованных на английском языке начиная с 2015 года. Поиск выполняли в MEDLINE через PubMed, EMBASE, Кокрейновской библиотеке, ресурсах Агентства по исследованиям и качеству медицинской помощи США и других выбранных базах данных, относящихся к тематике руководства.

\textbf{Структура.} Задача этих клинических рекомендаций --- создать актуализируемый рабочий документ, отражающий современные знания о ССПМ-синдроме и адресованный всем практикующим кардиологам, эндокринологам, нефрологам, специалистам первичной помощи и другим врачам, ведущим таких пациентов.

\textbf{Ключевые слова:} научные заявления AHA; атеросклероз; фибрилляция предсердий; артериальное давление/артериальная гипертензия; сердечно-сосудистые заболевания; сердечно-сосудисто-почечно-метаболический синдром; коронарная болезнь; диабет; диагностика; сердечная недостаточность; хронические болезни почек; ведение пациентов с ССПМ-синдромом; MASLD; метаболический синдром; избыточная масса тела/ожирение; заболевание периферических артерий; почечная недостаточность; факторы риска; скрининг; вторичная профилактика; венозная тромбоэмболия; снижение массы тела.

\section*{Основные положения}
\addcontentsline{toc}{section}{Основные положения}
\begin{enumerate}[label=\arabic*.]
\item \textbf{Определение стадии ССПМ-синдрома.} Стадирование ССПМ-синдрома рекомендуется детям, подросткам и взрослым для предотвращения перехода к более поздним стадиям, выбора терапии в соответствии с абсолютным риском, уменьшения частоты сердечно-сосудистых событий и потери функции почек на протяжении всей жизни, а также для содействия обратному развитию синдрома за счёт изменения образа жизни и снижения массы тела.

\item \textbf{Количественная оценка риска с помощью уравнений PREVENT.} У лиц с риском сердечно-сосудистых заболеваний (ССЗ), соответствующих стадиям ССПМ 0--3, следует количественно оценивать риск с помощью уравнений PREVENT (Predicting Risk of Cardiovascular Disease EVENTS). Они позволяют определить 10- и 30-летний риск атеросклеротических ССЗ (АССЗ), сердечной недостаточности (СН) и всех ССЗ. Результаты PREVENT используются при определении стадии: прогнозируемый 10-летний риск всех ССЗ $\geq20\%$ является одним из критериев стадии 3. Прогнозируемый 10-летний риск ССЗ $\geq7{,}5\%$ дополнительно помогает определять приоритетность лекарственных вмешательств при ССПМ-синдроме.

\item \textbf{Регулярная оценка факторов ССПМ-риска.} Всем взрослым рекомендуется регулярно оценивать метаболические факторы риска и функцию почек. В отдельных подгруппах также рекомендуется целенаправленное обследование для выявления доклинической сердечной недостаточности (pre-HF), стеатотической болезни печени, ассоциированной с метаболической дисфункцией (MASLD), и обструктивного апноэ сна (ОАС).

\item \textbf{Оценка социальных детерминант здоровья и работа с ними.} Рекомендуется регулярно оценивать социальные детерминанты здоровья (SDOH), тесно связанные с развитием ССПМ-синдрома и его осложнений. Устранение или смягчение выявленных неблагоприятных социальных факторов --- важнейшая составляющая комплексной помощи при ССПМ-синдроме.

\item \textbf{Междисциплинарная помощь с выделенным координатором.} При сочетании сахарного диабета 2 типа (СД2), хронической болезни почек (ХБП) и ССЗ рекомендуется междисциплинарная модель помощи, ориентированная на пациента. Подчёркивается необходимость выделить специалиста, координирующего работу междисциплинарной команды, содействующего внедрению доказательных подходов и терапии в соответствии с клиническими рекомендациями (GDMT), а также поддерживающего пациентов и клиницистов.

\item \textbf{Выявление и лечение избыточной массы тела и ожирения.} Для характеристики риска, связанного с избытком жировой ткани, рекомендуется оценивать избыточную массу тела/ожирение и абдоминальное жироотложение с помощью как индекса массы тела (ИМТ), так и окружности талии. Коррекция избыточной массы тела и ожирения необходима для предотвращения прогрессирования и содействия обратному развитию ССПМ-синдрома. Основной акцент делается на поддержке изменений образа жизни; при необходимости дополнительно применяют лекарственную терапию ожирения и метаболическую/бариатрическую хирургию, чтобы облегчить снижение массы тела.

\item \textbf{Применение кардиопротективной сахароснижающей терапии при диабете.} Пациентам с СД2 и ССЗ либо повышенным риском ССЗ наряду с изменением образа жизни, контролем массы тела и факторов риска рекомендуется кардиопротективная сахароснижающая GDMT. Она включает ингибиторы натрий-глюкозного котранспортера 2 типа (SGLT2i), терапию на основе глюкагоноподобного пептида-1 (GLP-1) или их сочетание для улучшения сердечно-сосудистых и почечных исходов. Выбор препарата должен определяться сопутствующими состояниями: ХБП, АССЗ, СН, ожирением, выраженной гипергликемией и MASLD.

\item \textbf{Выявление ХБП и применение нефропротективных препаратов.} Для характеристики ХБП и выбора нефропротективной терапии, обеспечивающей одновременно сердечно-сосудистые и почечные преимущества, рекомендуется определять и расчётную скорость клубочковой фильтрации (рСКФ), и отношение альбумина к креатинину в моче (UACR). При сочетании ХБП с СД2 либо ХБП с альбуминурией препаратами первой линии должны быть ингибиторы ренин-ангиотензиновой системы (RASi) и SGLT2i. Если у пациентов с ХБП и СД2 альбуминурия сохраняется, для дополнительной защиты почек и сердечно-сосудистой системы следует добавить нестероидный антагонист минералокортикоидных рецепторов (MRA) или терапию на основе GLP-1.

\item \textbf{Коррекция ССПМ-факторов риска у пациентов с АССЗ.} При ведении пациентов с АССЗ следует уделять особое внимание сопутствующим компонентам ССПМ-синдрома: лечению ожирения посредством изменения образа жизни, лекарственной терапии и, когда это уместно, метаболической/бариатрической хирургии; применению кардиопротективных сахароснижающих препаратов при СД2; применению нефропротективных препаратов при ХБП. Цель --- снизить риск неблагоприятных сердечно-сосудистых событий и потери функции почек.

\item \textbf{Коррекция ССПМ-факторов риска у пациентов с СН.} Факторы ССПМ-риска должны учитываться при лечении СН. При СН со сниженной фракцией выброса (HFrEF) необходимо использовать сердечно-сосудистые и почечные преимущества RASi --- ингибиторов рецепторов ангиотензина и неприлизина (ARNI), ингибиторов ангиотензинпревращающего фермента (ACEi), блокаторов рецепторов ангиотензина II (ARB) --- и SGLT2i в составе четырёхкомпонентной терапии с бета-адреноблокаторами и стероидными MRA. При СН с умеренно сниженной (HFmrEF) или сохранённой (HFpEF) фракцией выброса SGLT2i применяют как GDMT первой линии; при ожирении или других факторах ССПМ-риска добавляют терапию на основе GLP-1. У пациентов с СД2 и ХБП следует рассмотреть нестероидные MRA для уменьшения неблагоприятных сердечно-сосудистых событий и потери функции почек.
\end{enumerate}

\section*{Преамбула}
\addcontentsline{toc}{section}{Преамбула}
Начиная с 1980 года Американская коллегия кардиологов (ACC) и Американская кардиологическая ассоциация (AHA) преобразуют научные данные в клинические практические рекомендации, направленные на улучшение сердечно-сосудистого здоровья. Эти рекомендации, основанные на систематических методах оценки и классификации доказательств, служат основой высококачественной сердечно-сосудистой помощи. Когда это применимо, в документ включают положения об экономической ценности вмешательств, опирающиеся на анализ соотношения затрат и эффективности. Методология таких положений изложена в заявлении AHA/ACC о методологии оценки стоимости и ценности.\src{1}

В публикации ACC/AHA об основных принципах и процессе разработки рекомендаций описаны лучшие подходы для клиницистов-кардиологов и приведены дополнительные сведения о методологии подготовки рекомендаций.\src{2} Согласование процессов одобрения лекарственных препаратов и медицинских устройств Управлением по контролю качества пищевых продуктов и лекарственных средств США (FDA) с методологией AHA/ACC подробнее обсуждается в документе о включении процессов FDA в разработку рекомендаций.\src{3}

Рекомендации описывают подходы, отвечающие потребностям пациентов в большинстве, но не во всех ситуациях, и не должны заменять клиническое суждение. ACC и AHA финансируют разработку и публикацию рекомендаций без коммерческой поддержки; члены авторских и рецензионных комитетов участвуют в этой работе на добровольной основе. Рекомендации отражают официальную позицию ACC и AHA. При подготовке некоторых руководств ACC и AHA сотрудничают с другими организациями.

\begin{flushright}
Catherine M. Otto, MD, FAHA, FACC\\
Председатель Объединённого комитета ACC/AHA\\
по клиническим практическим рекомендациям
\end{flushright}

\section{Введение}
\subsection{Методология и обзор доказательств}
Приведённые в руководстве рекомендации по возможности основаны на доказательствах. Первоначальный расширенный обзор данных проводился с октября 2024 года по апрель 2025 года. Он включал англоязычные публикации исследований с участием людей, индексированные в MEDLINE через PubMed, EMBASE, Кокрейновской библиотеке, ресурсах Агентства по исследованиям и качеству медицинской помощи США и других выбранных базах данных, соответствующих тематике руководства. По мере необходимости авторский комитет добавлял отдельные ключевые исследования, опубликованные по ноябрь 2025 года включительно. Итоговые таблицы доказательств обобщают данные, использованные комитетом для формулирования рекомендаций. Перечень отобранных и опубликованных в руководстве ссылок не является исчерпывающим.

\subsection{Организация авторского комитета}
В состав комитета вошли представители общей, профилактической и интервенционной кардиологии, эндокринологии, нефрологии, бариатрической хирургии, семейной медицины, внутренних болезней, педиатрии, гастроэнтерологии, сердечно-сосудистой эпидемиологии, экономики здравоохранения, сестринской практики расширенного уровня и клинической фармакологии, а также представители интересов пациентов. В комитете были представлены AHA, ACC, ADA, включая её направление по борьбе с ожирением, и Американское общество нефрологии.

\subsection{Связи с индустрией и другими организациями}
ACC и AHA придерживаются строгих правил и процедур, призванных исключить предвзятость и ненадлежащее влияние при разработке документов. Полный текст политики в отношении связей с индустрией и другими организациями (RWI) доступен в интернете. В приложении 1 оригинала приведены полные сведения о таких связях членов авторского комитета, в том числе относящихся к рассматриваемым вопросам.

\subsection{Комитет рецензентов}
Объединённый комитет назначил комитет для рецензирования руководства. В него вошли специалисты, выдвинутые ACC, AHA и организациями-партнёрами. Полные сведения о связях рецензентов с индустрией и другими организациями, включая связи, относящиеся к тематике документа, приведены в приложении 2 оригинала.

\subsection{Область применения руководства}
Руководство предназначено для специалистов клинической медицины и общественного здравоохранения. Клиницистам следует учитывать, что оно отменяет и заменяет опубликованные ранее «Рекомендации AHA/ACC/TOS 2013 года по ведению взрослых с избыточной массой тела и ожирением». В этом документе не даются общие рекомендации по следующим направлениям: контроль артериального давления (АД), хроническая коронарная болезнь, дислипидемия/холестерин крови, сердечная недостаточность и заболевание периферических артерий. Эти вопросы рассматриваются в других руководствах AHA/ACC, перечисленных в таблице 1.

Основная целевая аудитория --- клиницисты, ведущие пациентов на всём протяжении спектра ССПМ-синдрома, характеризующегося взаимосвязями метаболических факторов риска, включая ожирение и СД2, ХБП и ССЗ.\src{1} Цель документа --- дать краткие практические ориентиры по диагностике, стадированию, лечению и наблюдению таких пациентов. Представленные клинические рекомендации предназначены для взрослых ($\geq18$ лет); предполагается, что рекомендации для детей и подростков ($<18$ лет) станут предметом будущих документов.

Подчёркивается применение междисциплинарных моделей для согласованной, ориентированной на пациента помощи и важность устранения неблагоприятных социальных детерминант здоровья для комплексного ведения ССПМ-синдрома. Отдельное внимание уделяется изменению образа жизни и контролю массы тела, направленным на основную причину синдрома --- избыток и дисфункцию жировой ткани. Для более поздних стадий приводятся подходы к применению сердечно-сосудистой и нефропротективной терапии с доказанной пользой. Современные обновления в прогнозировании сердечно-сосудистого риска используются для более адресного назначения профилактических вмешательств.

Руководство содержит также положения об экономической ценности некоторых методов лечения ССПМ-синдрома, прежде всего оригинальных брендированных лекарственных препаратов. Хотя эти положения размещены рядом с соответствующими клиническими рекомендациями, они не предназначены для принятия решений о лечении конкретного пациента.

При подготовке документа авторский комитет рассмотрел ранее опубликованные руководства и другие относящиеся к теме документы о клинической практике. Таблица 1 содержит перечень публикаций AHA/ACC, признанных значимыми для этой работы. Она служит справочным ресурсом и позволяет не повторять уже имеющиеся рекомендации.

\heading{Таблица 1. Связанные публикации AHA/ACC}
{\small
\begin{longtable}{P{0.56\textwidth}P{0.27\textwidth}P{0.10\textwidth}}
\toprule
Название & Организации & Год / ссылка\\
\midrule
\endfirsthead
\toprule Название & Организации & Год / ссылка\\\midrule
\endhead
\multicolumn{3}{l}{\textbf{Клинические рекомендации}}\\
Первичная профилактика ССЗ & ACC/AHA & 2019\src{2}\\
Профилактика инсульта у пациентов, перенёсших инсульт или транзиторную ишемическую атаку & AHA/ASA & 2021\src{3}\\
Ведение пациентов с СН & AHA/ACC/HFSA & 2022\src{4}\\
Ведение пациентов с хронической коронарной болезнью & AHA/ACC/ACCP/ ASPC/NLA/PCNA & 2023\src{5}\\
Ведение пациентов с фибрилляцией предсердий & ACC/AHA/ACCP/HRS & 2023\src{6}\\
Ведение пациентов с заболеванием периферических артерий нижних конечностей & ACC/AHA/AACVPR/ APMA/ABC/SCAI/ SVM/SVN/SVS/ SIR/VESS & 2024\src{7}\\
Профилактика, выявление, обследование и ведение взрослых с повышенным АД & AHA/ACC/AANP/ AAPA/ABC/ACCP/ ACPM/AGS/AMA/ ASPC/NMA/PCNA/ SGIM & 2025\src{8}\\
Ведение пациентов с дислипидемией & ACC/AHA/AACVPR/ ABC/ACPM/ADA/ AGS/APhA/ASPC/ NLA/PCNA & 2026\src{9}\\
\midrule
\multicolumn{3}{l}{\textbf{Другие относящиеся к теме документы}}\\
Ценность примордиальной и первичной профилактики ССЗ & AHA & 2011\src{10}\\
Социальные детерминанты риска и исходов ССЗ & AHA & 2015\src{11}\\
Выявление и снижение риска ССЗ у женщин посредством сотрудничества с акушерами и гинекологами & AHA/ACOG & 2018\src{12}\\
Применение инструментов оценки риска при принятии решений о первичной профилактике АССЗ & AHA & 2019\src{13}\\
СД2 и СН & AHA/HFSA & 2019\src{14}\\
Сердечно-сосудистые аспекты ведения беременных & AHA & 2020\src{15}\\
Ожирение и ССЗ & AHA & 2021\src{16}\\
Обструктивное апноэ сна и ССЗ & AHA & 2021\src{17}\\
Life's Essential 8: восемь ключевых компонентов здоровья & AHA & 2022\src{18}\\
Неалкогольная жировая болезнь печени и сердечно-сосудистый риск & AHA & 2022\src{19}\\
Комплексная коррекция сердечно-сосудистых факторов риска у взрослых с СД2 & AHA & 2022\src{20}\\
Сердечно-сосудисто-почечно-метаболическое здоровье & AHA & 2023\src{1}\\
Краткий обзор научных данных и клинического ведения ССПМ-синдрома & AHA & 2023\src{21}\\
Новые уравнения прогнозирования абсолютного риска всех ССЗ с учётом ССПМ-здоровья & AHA & 2023\src{22}\\
Модели сердечно-сосудистой помощи, ориентированные на человека & AHA & 2023\src{23}\\
Роль первичной помощи в достижении целей Life's Essential 8 & AHA & 2024\src{24}\\
Возможности снижения риска ССЗ в послеродовом периоде после неблагоприятных исходов беременности & AHA & 2024\src{25}\\
Внедрение научных знаний об ожирении в клиническую практику & AHA & 2024\src{26}\\
Ведение взрослых с ожирением и СН & ACC & 2025\src{27}\\
Применение оценки риска для принятия решений о контроле АД при первичной профилактике ССЗ & AHA & 2025\src{28}\\
Первичная профилактика СН на основе оценки риска & AHA & 2025\src{29}\\
Медикаментозный контроль массы тела для оптимизации сердечно-сосудистого здоровья & ACC & 2026\src{30}\\
Роль физической активности в лечении ожирения и поддержании кардиометаболического здоровья & AHA & 2026\src{31}\\
\bottomrule
\end{longtable}
}

{\footnotesize
\textbf{Организации и обозначения в таблице 1.}
AACVPR --- Американская ассоциация сердечно-сосудистой и лёгочной реабилитации;
AANP --- Американская ассоциация практикующих медицинских сестёр;
AAPA --- Американская академия ассистентов врача;
ABC --- Ассоциация чернокожих кардиологов;
ACC --- Американская коллегия кардиологов;
ACCP --- Американская коллегия клинической фармации;
ACOG --- Американская коллегия акушеров и гинекологов;
ACPM --- Американская коллегия профилактической медицины;
ADA --- Американская диабетическая ассоциация;
AGS --- Американское гериатрическое общество;
AHA --- Американская кардиологическая ассоциация;
AMA --- Американская медицинская ассоциация;
APhA --- Американская ассоциация фармацевтов;
APMA --- Американская ассоциация подиатрической медицины;
ASA --- Американская ассоциация по изучению инсульта;
ASPC --- Американское общество профилактической кардиологии;
HFSA --- Американское общество по изучению сердечной недостаточности;
HRS --- Общество сердечного ритма;
NLA --- Национальная липидная ассоциация;
NMA --- Национальная медицинская ассоциация;
PCNA --- Ассоциация медицинских сестёр профилактической сердечно-сосудистой помощи;
SCAI --- Общество сердечно-сосудистой ангиографии и вмешательств;
SGIM --- Общество общей внутренней медицины;
SIR --- Общество интервенционной радиологии;
SVM --- Общество сосудистой медицины;
SVN --- Общество сосудистого сестринского дела;
SVS --- Общество сосудистой хирургии;
VESS --- Общество сосудистой и эндоваскулярной хирургии.

Клинические сокращения оригинала: AF --- фибрилляция предсердий; APO --- неблагоприятный исход беременности; BP --- АД; CCD --- хроническая коронарная болезнь; CKM --- сердечно-сосудисто-почечно-метаболический; CVD --- ССЗ; GYN --- гинекологи; HF --- СН; OB --- акушеры; OSA --- обструктивное апноэ сна; PAD --- заболевание периферических артерий; T2D --- СД2.
}

\subsection{Класс рекомендации и уровень доказательности}
Класс рекомендации (COR) отражает её силу и учитывает предполагаемую пользу в соотношении с риском. Уровень доказательности (LOE) характеризует качество научных данных в поддержку вмешательства с учётом типа, объёма и согласованности результатов клинических исследований и других источников (таблица 2).

\heading{Таблица 2. Применение класса рекомендации и уровня доказательности к клиническим стратегиям, вмешательствам, лечению и диагностическим исследованиям}
\textit{Редакция исходной таблицы: декабрь 2024 года.}

\heading{Класс (сила) рекомендации}
\heading{Класс 1 --- сильная рекомендация.} Польза $\ggg$ риск. Рекомендуемые формулировки: «рекомендуется»; «показано/полезно/эффективно/приносит пользу»; «следует выполнить/назначить». При сравнении эффективности: «лечение/стратегия A рекомендуется или показана предпочтительно перед лечением B»; «следует выбрать лечение A, а не B».
\heading{Класс 2a --- умеренная рекомендация.} Польза $\gg$ риск. Формулировки: «целесообразно»; «может быть полезно, эффективно или приносить пользу». При сравнении эффективности: «лечение/стратегия A, вероятно, рекомендуется или показана предпочтительно перед лечением B»; «целесообразно выбрать лечение A, а не B».
\heading{Класс 2b --- слабая рекомендация.} Польза $\geq$ риск. Формулировки: «может быть целесообразно»; «может быть рассмотрено»; «полезность/эффективность неизвестна, неясна, неопределённа или недостаточно установлена».
\heading{Класс 3: отсутствие пользы --- умеренная рекомендация.} Польза $=$ риск. Как правило, применим только при уровнях доказательности A или B. Формулировки: «не рекомендуется»; «не показано/не полезно/не эффективно/не приносит пользы»; «не следует выполнять/назначать».
\heading{Класс 3: вред --- сильная рекомендация.} Риск $>$ польза. Формулировки: «потенциально вредно»; «причиняет вред»; «связано с увеличением заболеваемости/смертности»; «не следует выполнять/назначать».

\heading{Уровень (качество) доказательности}
\heading{A.} Доказательства высокого качества, полученные более чем в одном рандомизированном контролируемом исследовании (РКИ); метаанализы РКИ высокого качества; одно или несколько РКИ, результаты которых подтверждаются высококачественными регистровыми исследованиями.
\heading{B-R --- рандомизированные данные.} Доказательства умеренного качества из одного или нескольких РКИ; метаанализы РКИ умеренного качества.
\heading{B-NR --- нерандомизированные данные.} Доказательства умеренного качества из одного или нескольких хорошо спланированных и выполненных нерандомизированных исследований, наблюдательных исследований или регистров; метаанализы таких исследований.
\heading{C-LD --- ограниченные данные.} Рандомизированные или нерандомизированные наблюдательные либо регистровые исследования с ограничениями дизайна или проведения; метаанализы таких исследований; физиологические исследования или исследования механизмов с участием людей.
\heading{C-EO --- мнение экспертов.} Согласованное мнение экспертов, основанное на клиническом опыте.

COR и LOE определяются независимо: любой класс рекомендации может сочетаться с любым уровнем доказательности. Уровень C не означает, что рекомендация слабая. Многие важные клинические вопросы, рассматриваемые в руководствах, не подходят для изучения в клинических испытаниях. Даже при отсутствии РКИ может существовать вполне определённый клинический консенсус о полезности или эффективности конкретного исследования либо лечения.

\textbf{Примечания исходной таблицы.} Необходимо указывать исход или результат вмешательства: улучшение клинического исхода, повышение диагностической точности либо дополнительную прогностическую информацию. Для рекомендаций сравнительной эффективности --- только COR 1 и 2a, LOE A и B --- исследования, обосновывающие сравнительные формулировки, должны содержать прямое сравнение оцениваемых методов лечения или стратегий. Методология оценки качества развивается; она включает стандартизированные, широко применяемые и предпочтительно валидированные инструменты градации доказательств, а для систематических обзоров --- участие комитета по оценке доказательств. EO --- мнение экспертов; LD --- ограниченные данные; NR --- нерандомизированные; R --- рандомизированные данные.


\section{Определения и классификация}
\subsection{Определения}
ССПМ-синдром представляет собой клинический синдром со стадиями от 1 до 4, которые описаны в разделе 2.3 «ССПМ-синдром и определения его стадий». Авторский комитет учитывает, что для обозначения компонентов синдрома, часто сопутствующих ему состояний и применяемых методов лечения используется множество распространённых терминов. Поэтому в этом разделе приводятся краткие определения.

\textbf{Ведение в соответствии с клиническими рекомендациями (GDMT).} Включает клиническую оценку, диагностические исследования, лекарственное лечение и лечебные процедуры. Для этих и всех других рекомендуемых схем лекарственной терапии читатель должен проверять дозы по инструкции к препарату, а также оценивать противопоказания и лекарственные взаимодействия. Рекомендации ограничены лекарствами, устройствами и методами лечения, разрешёнными к клиническому применению в США.

\heading{Дополнительные определения состояний, составляющих ССПМ-синдром}
\heading{Клинически выраженные АССЗ.} Включают острый коронарный синдром в анамнезе, инфаркт миокарда (ИМ), стабильную или нестабильную стенокардию, коронарную реваскуляризацию или реваскуляризацию других артерий, инсульт, транзиторную ишемическую атаку либо заболевание периферических артерий.
\heading{Клинически выраженная СН.} Сердечная недостаточность с клинической симптоматикой, включая СН со сниженной (HFrEF), умеренно сниженной (HFmrEF) и сохранённой (HFpEF) фракцией выброса.
\heading{Сахарный диабет 2 типа.} Сложное метаболическое расстройство, характеризующееся снижением чувствительности к действию инсулина (инсулинорезистентностью) и прогрессирующим нарушением его секреции, что приводит к гипергликемии.
\heading{Ожирение.} Избыток жировой ткани, представляющий риск для здоровья; диагностируется при ИМТ $\geq30$~кг/м$^2$. Абдоминальное ожирение диагностируют при окружности талии $\geq88$~см у женщин или $\geq102$~см у мужчин.
\heading{Хроническая болезнь почек.} Стойкое снижение клубочковой фильтрации, то есть рСКФ $<60$~\unitgfr, или альбуминурия, то есть UACR $\geq30$~мг/г, выявленные как минимум при двух измерениях с интервалом 3 месяца.
\heading{Артериальная гипертензия.} Повышение систолического АД $\geq130$~мм рт. ст. или диастолического АД $\geq80$~мм рт. ст. по среднему значению не менее двух тщательно выполненных измерений, полученных не менее чем при двух отдельных обследованиях.
\heading{Почечная недостаточность (kidney failure).} Хроническое снижение рСКФ $<15$~\unitgfr\ (ХБП стадии 5) или постоянная потребность в заместительной почечной терапии --- диализе. Ранее использовался термин «терминальная стадия болезни почек».
\heading{Предиабет.} Повышение глюкозы, сопряжённое с риском развития СД2. Определяется как глюкоза плазмы натощак 100--125~мг/дл, нарушение толерантности к глюкозе при пероральном глюкозотолерантном тесте с 75~г глюкозы --- значение через 2 часа 140--199~мг/дл --- или гликированный гемоглобин HbA1c 5,7--6,4\%.
\heading{Доклиническая сердечная недостаточность (pre-HF).} Наличие субклинического заболевания миокарда в виде структурных или функциональных нарушений, признаков повышения давления наполнения либо повышенных сердечных биомаркеров при отсутствии клинических симптомов СН.

\editorial{Термин pre-HF далее обозначает доклиническую стадию сердечной недостаточности; он не относится к заболеванию предсердий. Термин kidney failure используется в узком значении, заданном авторами, и не является синонимом любой стадии ХБП.}

\heading{Дополнительные определения состояний, часто встречающихся при ССПМ-синдроме}
\textbf{Стеатотическая болезнь печени, ассоциированная с метаболической дисфункцией (MASLD).} Характеризуется стеатозом печени на фоне метаболических факторов риска, например ожирения или СД2. Этот термин был недавно введён взамен термина «неалкогольная жировая болезнь печени».

\textbf{Факторы, усиливающие риск прогрессирования ССПМ-синдрома.} Клинические факторы --- например хронические воспалительные состояния или неблагоприятные исходы беременности в анамнезе, --- социальные или демографические факторы, а также семейный анамнез диабета или почечной недостаточности, связанные с более высоким риском прогрессирования ССПМ-синдрома и способные влиять на выбор профилактической терапии. Многие из них совпадают с усилителями риска, определёнными в междисциплинарном руководстве ACC/AHA 2026 года по дислипидемии,\src{1} однако существуют и различия, поскольку речь идёт о двух связанных, но не тождественных исходах: прогрессировании ССПМ-синдрома и АССЗ.

\heading{Образ жизни, лекарственная терапия и хирургические методы лечения ССПМ-синдрома}
\textbf{Терапия на основе GLP-1 с доказанной пользой.} Лекарственные средства, которые имитируют или усиливают действие инкретинового гормона GLP-1 либо действуют как двойные агонисты, одновременно имитирующие действие глюкозозависимого инсулинотропного полипептида. Их способность снижать риск подтверждена в исследованиях сердечно-сосудистых исходов. Также могут использоваться обозначения «инкретиновая терапия» или «терапия, связанная с инкретинами».

\textbf{Ингибиторы ренин-ангиотензиновой системы (RASi).} Класс препаратов, подавляющих ренин-ангиотензиновую систему; включает ACEi и ARB. Применяются для снижения АД, уменьшения протеинурии и предотвращения прогрессирования сердечно-сосудистых и почечных заболеваний.

\textbf{Ингибиторы ренин-ангиотензин-альдостероновой системы (RAASi).} Класс препаратов, включающий RASi, например ACEi и ARB, а также препараты, подавляющие альдостероновую систему, например MRA.

\textbf{Метаболическая и бариатрическая хирургия (MBS).} Группа операций --- например гастрошунтирование и рукавная резекция желудка, --- направленных на снижение массы тела и улучшение метаболического здоровья посредством изменения анатомии и физиологии желудочно-кишечного тракта, влияющего на энергетический баланс.

\textbf{Коррекция образа жизни.} Включает оценку исходных поведенческих привычек каждого человека, консультирование и поддержку здоровых изменений. Предполагает содействие питанию, полезному для сердца, регулярной физической активности, отказу от всех никотинсодержащих средств доставки, здоровым привычкам сна и поддержанию здоровой массы тела. Включает поведенческие компоненты программы AHA Life's Essential 8\src{2--4} и управление стрессом.

\subsection{Сокращения}
\editorial{Для удобства сопоставления с оригиналом слева сохранены его английские сокращения. В переводе также используются ССПМ, ССЗ, АССЗ, СН, ХБП, СД2, АД, ИМТ, рСКФ, РКИ, ИМ и ОАС.}
{\small
\begin{longtable}{P{0.18\textwidth}P{0.75\textwidth}}
\toprule Сокращение & Значение\\\midrule
\endfirsthead
\toprule Сокращение & Значение\\\midrule
\endhead
ACEi & Ингибиторы ангиотензинпревращающего фермента\\
APO & Неблагоприятный исход беременности\\
ARB & Блокаторы рецепторов ангиотензина II\\
ARNI & Ингибитор рецепторов ангиотензина и неприлизина\\
ASCVD & Атеросклеротическое сердечно-сосудистое заболевание\\
BNP & Натрийуретический пептид B-типа\\
BP & Артериальное давление\\
CAC & Кальций коронарных артерий; коронарный кальций\\
CCD & Хроническая коронарная болезнь\\
CI & Доверительный интервал\\
CKD & Хроническая болезнь почек\\
CKM & Сердечно-сосудисто-почечно-метаболический\\
CVD & Сердечно-сосудистое заболевание\\
DOAC & Прямые пероральные антикоагулянты\\
eGFR & Расчётная скорость клубочковой фильтрации\\
GDM & Гестационный сахарный диабет\\
GDMT & Ведение/терапия в соответствии с клиническими рекомендациями\\
GIP & Глюкозозависимый инсулинотропный полипептид\\
GLP-1 & Глюкагоноподобный пептид-1\\
GLP-RA & Агонист рецептора глюкагоноподобного пептида-1\\
HbA1c & Гликированный гемоглобин A1c\\
HDL-C & Холестерин липопротеинов высокой плотности (ХС ЛВП)\\
HF & Сердечная недостаточность\\
HFmrEF & СН с умеренно сниженной фракцией выброса\\
HFpEF & СН с сохранённой фракцией выброса\\
HFrEF & СН со сниженной фракцией выброса\\
HR & Отношение интенсивностей наступления события (hazard ratio)\\
Hs & Высокочувствительный\\
ICER & Инкрементальное соотношение затрат и эффективности\\
LDL-C & Холестерин липопротеинов низкой плотности (ХС ЛНП)\\
LSG & Лапароскопическая рукавная резекция желудка\\
LVEF & Фракция выброса левого желудочка\\
MACE & Большое неблагоприятное сердечно-сосудистое событие\\
MASLD & Стеатотическая болезнь печени, ассоциированная с метаболической дисфункцией\\
MBS & Метаболическая и бариатрическая хирургия\\
MI & Инфаркт миокарда\\
MRA & Антагонист минералокортикоидных рецепторов\\
nsMRA & Нестероидный антагонист минералокортикоидных рецепторов\\
NT-proBNP & N-концевой фрагмент прогормона натрийуретического пептида B-типа\\
OR & Отношение шансов\\
OSA & Обструктивное апноэ сна\\
PAD & Заболевание периферических артерий\\
PREVENT & Predicting Risk of Cardiovascular Disease Events --- прогнозирование риска сердечно-сосудистых событий\\
QALY & Год жизни с поправкой на её качество\\
RAASi & Ингибиторы ренин-ангиотензин-альдостероновой системы\\
RASi & Ингибиторы ренин-ангиотензиновой системы\\
RCT & Рандомизированное контролируемое исследование\\
RR & Относительный риск\\
RYGB & Гастрошунтирование по Ру\\
SDOH & Социальные детерминанты здоровья\\
SGLT2i & Ингибиторы натрий-глюкозного котранспортера 2 типа\\
T2D & Сахарный диабет 2 типа\\
UACR & Отношение альбумина к креатинину в моче\\
VKA & Антагонист витамина K\\
VTE & Венозная тромбоэмболия\\
\bottomrule
\end{longtable}
}

\subsection{ССПМ-синдром и определения его стадий}
\heading{Обзор}
Определение ССПМ-синдрома (таблица 3) отражает тесные взаимосвязи между метаболическими факторами риска, ХБП и ССЗ, а также их совместное, взаимно усиливающее влияние на заболеваемость и смертность, описанное в президентском консультативном документе AHA о ССПМ-здоровье.\src{1} Это определение служит отправной точкой для описания стадий синдрома, разработки инструментов скрининга и стратификации риска связанных с ним неблагоприятных исходов, а также для определения оптимальных клинических подходов, обеспечивающих своевременное начало доказательной профилактики и лечения.\src{2}

Межсистемное взаимодействие, подчёркнутое в определении, обосновывает отказ от изолированного ведения отдельных компонентов синдрома узкими специалистами. Оно подчёркивает важность скрининга ССПМ-синдрома как в первичной помощи, так и в соответствующих специализированных клиниках, а также ценность междисциплинарного сотрудничества. Определение признаёт неотъемлемую роль социальных детерминант здоровья и факторов образа жизни в укреплении и сохранении ССПМ-здоровья на протяжении всей жизни. Кроме того, единое определение облегчает взаимодействие научного сообщества и населения: оно помогает объяснить важность выявления ССПМ-синдрома и воздействия на его взаимосвязанные компоненты посредством изменения образа жизни, контроля массы тела и целевой лекарственной терапии для улучшения общественного здоровья.

Своевременное выявление компонентов ССПМ-синдрома, имеющих неблагоприятные клинические последствия, представляет важную возможность для профилактики, особенно с учётом наличия нескольких методов лечения с доказанной эффективностью. Стадии синдрома (рисунок 1) отражают его типичное патофизиологическое развитие: от избытка/дисфункции жировой ткани (стадия 1) к метаболическим факторам риска и/или ХБП (стадия 2), затем к субклиническому ССЗ либо его эквивалентам риска --- ХБП очень высокого риска или высокому расчётному риску ССЗ (стадия 3), --- и далее к клинически выраженному ССЗ при сопутствующих факторах ССПМ-риска (стадия 4).

Здоровый образ жизни и контроль массы тела --- основные способы сдерживания перехода к более поздним стадиям. Система стадирования направлена на предотвращение прогрессирования, включая такие неблагоприятные исходы, как ССЗ и почечная недостаточность. Она также подчёркивает более высокий абсолютный риск ССЗ на поздних стадиях, когда интенсификация лечения обеспечивает наибольшую чистую клиническую пользу.

Главное внимание уделяется ССЗ, поскольку именно они характеризуются наибольшей частотой и бременем смертности при ССПМ-синдроме. Однако ХБП также находится в центре внимания: последовательно возрастающим категориям риска ХБП по классификации KDIGO (Kidney Disease: Improving Global Outcomes; рисунок 2) соответствуют более поздние стадии ССПМ-синдрома, что отражает рост сердечно-сосудистого и почечного риска. Диагностические критерии стадирования у взрослых представлены в таблице 4.

\heading{Таблица 3. Определение ССПМ-синдрома}
\textbf{Полное определение.} ССПМ-синдром --- системное расстройство, характеризующееся патофизиологическими взаимодействиями метаболических факторов риска, ХБП и сердечно-сосудистой системы, которые приводят к полиорганной дисфункции и высокой частоте неблагоприятных сердечно-сосудистых исходов. Он охватывает как лиц с риском ССЗ из-за метаболических факторов риска, ХБП или их сочетания, так и пациентов с уже имеющейся ХБП или ССЗ, которые потенциально связаны с метаболическими факторами риска либо осложняют их течение. Вероятность ССПМ-синдрома и его неблагоприятных исходов дополнительно увеличивается при неблагоприятных условиях для здорового образа жизни и заботы о собственном здоровье, обусловленных общественной политикой, экономикой и окружающей средой.

\textbf{Упрощённое определение.} ССПМ-синдром --- нарушение здоровья, обусловленное взаимосвязями между болезнями сердца, болезнями почек, диабетом и ожирением и приводящее к неблагоприятным последствиям для здоровья.

\heading{Рисунок 1. Стадии ССПМ-синдрома: текстовая адаптация схемы}
\begin{enumerate}[label={},leftmargin=0pt,labelwidth=0pt,labelsep=0pt]
\item \textbf{Стадия 0: факторы риска отсутствуют.} Акцент на примордиальной профилактике --- предотвращении формирования факторов риска --- и сохранении сердечно-сосудистого здоровья.
\item \hfill $\Downarrow$ \hfill\mbox{}
\item \textbf{Стадия 1: избыток и/или дисфункция жировой ткани.} Избыточная масса тела/ожирение, абдоминальное ожирение, предиабет.
\item \hfill $\Downarrow$ \hfill\mbox{}
\item \textbf{Стадия 2: метаболические факторы риска, ХБП или их сочетание.} Артериальная гипертензия, гипертриглицеридемия, метаболический синдром, СД2, ХБП умеренного или высокого риска.
\item \hfill $\Downarrow$ \hfill\mbox{}
\item \textbf{Стадия 3: субклиническое ССЗ при ССПМ-синдроме.} Субклиническое АССЗ или pre-HF. Эквиваленты риска: ХБП очень высокого риска --- стадии G4--G5 либо соответствующая категория по карте KDIGO --- или высокий прогнозируемый 10-летний риск всех ССЗ $\geq20\%$ по PREVENT-CVD.
\item \hfill $\Downarrow$ \hfill\mbox{}
\item \textbf{Стадия 4: клинически выраженное ССЗ при ССПМ-синдроме.} Коронарная болезнь сердца, СН, фибрилляция предсердий, инсульт или заболевание периферических артерий.
\end{enumerate}

\textbf{Пояснение авторов к рисунку 1.} Адаптировано с разрешения по работе Ndumele и соавт.\src{1} © 2023 American Heart Association, Inc. Стадирование отражает прогрессирование патофизиологических нарушений и увеличение абсолютного риска ССЗ по всему спектру ССПМ-синдрома. К стадии 0 относятся лица без ССПМ-синдрома с нормальной массой тела, нормальными глюкозой, АД, липидами и функцией почек, без признаков субклинического или клинически выраженного ССЗ; приоритетами являются примордиальная профилактика и сохранение сердечно-сосудистого здоровья.

Стадия 1 включает избыток жировой ткани, её дисфункцию или их сочетание. Избыток жировой ткани устанавливают по массе тела либо абдоминальному ожирению; дисфункция проявляется предиабетом --- повышенной глюкозой или нарушением толерантности к ней. Стадия 2 включает метаболические факторы риска --- гипертриглицеридемию, гипертензию, метаболический синдром или СД2, --- ХБП умеренного или высокого риска либо их сочетание. Хотя гипертензия и ХБП обычно развиваются вследствие метаболических факторов риска, изогнутые стрелки исходной схемы показывают случаи их неметаболического происхождения; влияние на риск и лечебные подходы при этом сходны.

Стадия 3 включает субклиническое ССЗ при сопутствующих факторах ССПМ-риска --- избытке/дисфункции жировой ткани, метаболических факторах или ХБП, --- а также эквиваленты риска: ХБП очень высокого риска либо прогнозируемый 10-летний риск $\geq20\%$ по PREVENT. Стадия 4 включает клинически выраженные ССЗ --- коронарную болезнь сердца, СН, инсульт, заболевание периферических артерий или фибрилляцию предсердий --- в сочетании с факторами ССПМ-риска.

Сокращения исходной схемы: Afib --- фибрилляция предсердий; ASCVD --- АССЗ; CHD --- коронарная болезнь сердца; CKD --- ХБП; CKM --- ССПМ; CVD --- ССЗ; HF --- СН; PAD --- заболевание периферических артерий; PREVENT --- прогнозирование риска сердечно-сосудистых событий; KDIGO --- Kidney Disease: Improving Global Outcomes.

\heading{Рисунок 2. Карта KDIGO для классификации риска ХБП: табличная адаптация}
ХБП классифицируют по \textbf{причине, СКФ и альбуминурии}. Категории альбуминурии:
\begin{itemize}
\item \textbf{A1: нормальная или незначительно повышенная} --- UACR $<30$~мг/г, или $<3$~мг/ммоль.
\item \textbf{A2: умеренно повышенная} --- UACR 30--299~мг/г, или 3--29~мг/ммоль.
\item \textbf{A3: значительно повышенная} --- UACR $\geq300$~мг/г, или $\geq30$~мг/ммоль.
\end{itemize}

{\small
\begin{longtable}{P{0.09\textwidth}P{0.27\textwidth}P{0.10\textwidth}P{0.14\textwidth}P{0.14\textwidth}P{0.14\textwidth}}
\toprule Катег. & Функция почек & СКФ* & A1 & A2 & A3\\\midrule
\endfirsthead
\toprule Катег. & Функция почек & СКФ* & A1 & A2 & A3\\\midrule
\endhead
G1 & Нормальная или высокая & $\geq90$ & Низкий; С & Умеренный; Л & Высокий; ЛН\\
G2 & Незначительно сниженная & 60--89 & Низкий; С & Умеренный; Л & Высокий; ЛН\\
G3a & Незначительно или умеренно сниженная & 45--59 & Умеренный; Л & Высокий; Л & Очень высокий; ЛН\\
G3b & Умеренно или значительно сниженная & 30--44 & Высокий; Л & Очень высокий; ЛН & Очень высокий; ЛН\\
G4 & Значительно сниженная & 15--29 & Очень высокий; ЛН & Очень высокий; ЛН & Очень высокий; ЛН\\
G5 & Почечная недостаточность & $<15$ & Очень высокий; ЛН & Очень высокий; ЛН & Очень высокий; ЛН\\
\bottomrule
\end{longtable}
}
{\small *СКФ в \unitgfr. В каждой ячейке указаны категория риска и действие: С --- скрининг; Л --- лечение; ЛН --- лечение и направление к нефрологу. «Умеренный» соответствует обозначению оригинала «умеренно повышенный риск». Низкий риск означает отсутствие ХБП, если нет других маркеров заболевания почек.}

\textbf{Пояснение авторов к рисунку 2.} Модифицировано с разрешения KDIGO на условиях лицензии CC BY NC ND.\src{3} © 2022 KDIGO: Kidney Disease Improving Global Outcomes. Опубликовано Elsevier Inc. от имени Международного общества нефрологии. «Лечение» означает, что клиницисты должны начать соответствующее ведение пациентов с риском прогрессирования ХБП либо ССЗ как следствия ХБП. «Направление» означает консультацию нефрологической службы по мере необходимости с учётом местного порядка определения частоты наблюдения и сроков направления.

Под «причиной» понимается установленная клиницистом этиология ХБП. Большинство случаев ХБП при ССПМ-синдроме обусловлено диабетом, гипертензией и другими метаболическими факторами риска. Однако нефропротективные препараты, в частности ACEi/ARB и SGLT2i, продемонстрировали почечные и сердечно-сосудистые преимущества как при ХБП метаболического происхождения, так и при некоторых других её причинах. Поэтому система стадирования ССПМ применима почти ко всем пациентам с ХБП. Цвета исходного рисунка не соответствуют обозначениям класса рекомендации и уровня доказательности в таблице 2 AHA/ACC. CKD --- ХБП; CKM --- ССПМ; GFR --- расчётная скорость клубочковой фильтрации; KDIGO --- Kidney Disease: Improving Global Outcomes.

\heading{Таблица 4. Определения стадий ССПМ-синдрома и диагностические критерии у взрослых}
\heading{Стадия 0: факторы ССПМ-риска отсутствуют.} Нормальные ИМТ и окружность талии, нормогликемия, нормальное АД и липидный профиль, отсутствие ХБП, субклинического и клинически выраженного ССЗ.

\heading{Стадия 1: избыток или дисфункция жировой ткани.} Избыточная масса тела/ожирение, абдоминальное ожирение или дисфункция жировой ткани \textbf{при отсутствии других метаболических факторов риска и ХБП}:
\begin{itemize}
\item ИМТ $\geq25$~кг/м$^2$ либо $\geq23$~кг/м$^2$ у лиц азиатского происхождения; окружность талии $\geq88$~см у женщин и $\geq102$~см у мужчин либо, у лиц азиатского происхождения, соответственно $\geq80$ и $\geq90$~см; \textbf{или}
\item глюкоза крови натощак от 100 до 125~мг/дл либо HbA1c от 5,7 до 6,4\%\textsuperscript{*}.
\end{itemize}

\heading{Стадия 2: метаболические факторы риска, ХБП или их сочетание.} Метаболические факторы риска, ХБП умеренного или высокого риска либо их сочетание:
\begin{itemize}
\item артериальная гипертензия: систолическое АД $\geq130$~мм рт. ст., диастолическое АД $\geq80$~мм рт. ст. либо приём антигипертензивных препаратов;
\item гипертриглицеридемия $\geq150$~мг/дл;
\item метаболический синдром по критериям AHA/NHLBI\textsuperscript{\dag};
\item СД2: глюкоза крови натощак $\geq126$~мг/дл либо HbA1c $\geq6{,}5\%$;
\item и/или ХБП умеренного или высокого риска по KDIGO.
\end{itemize}

\heading{Стадия 3: субклиническое ССЗ при ССПМ-синдроме.} Субклиническое ССЗ --- субклинический коронарный атеросклероз или pre-HF --- при избытке/дисфункции жировой ткани, других метаболических факторах риска либо ХБП; \textbf{или эквивалент риска субклинического ССЗ}:
\begin{itemize}
\item субклинический коронарный атеросклероз преимущественно диагностируют при коронарном кальции с индексом Агатстона $\geq100$. Критериям также соответствуют случайно выявленный умеренный или выраженный коронарный кальциноз, субклинический коронарный атеросклероз по данным катетеризации коронарных артерий/КТ-ангиографии либо низкий лодыжечно-плечевой индекс без симптомов перемежающейся хромоты;
\item pre-HF диагностируют по повышенным сердечным биомаркерам: NT-proBNP $\geq125$~пг/мл; высокочувствительный тропонин T (hs-cTnT) $\geq14$~нг/л у женщин и $\geq22$~нг/л у мужчин; высокочувствительный тропонин I (hs-cTnI) $\geq10$~нг/л у женщин и $\geq12$~нг/л у мужчин; либо по эхокардиографическим параметрам\textsuperscript{\ddag}. Сочетание биомаркерных и эхокардиографических признаков указывает на наибольший риск СН.
\end{itemize}
Эквиваленты риска субклинического ССЗ:
\begin{itemize}
\item ХБП очень высокого риска: стадии G4--G5 либо категория очень высокого риска по классификации KDIGO\textsuperscript{§};
\item прогнозируемый 10-летний риск ССЗ $\geq20\%$ по PREVENT-CVD.
\end{itemize}

\heading{Стадия 4: клинически выраженное ССЗ при ССПМ-синдроме.} Клинически выраженное ССЗ --- коронарная болезнь сердца, СН, инсульт, заболевание периферических артерий, фибрилляция предсердий --- при избытке/дисфункции жировой ткани, других факторах ССПМ-риска или ХБП.
\begin{itemize}
\item \textbf{Стадия 4a:} почечная недостаточность отсутствует.
\item \textbf{Стадия 4b:} почечная недостаточность присутствует: рСКФ $<15$~\unitgfr\ либо потребность в постоянной заместительной почечной терапии.
\end{itemize}

\textbf{Примечания исходной таблицы 4.}

\textsuperscript{*} После беременности женщинам с гестационным диабетом следует проводить более интенсивный скрининг предиабета.

\textsuperscript{\dag} Метаболический синдром определяется при наличии трёх или более признаков: 1) окружность талии $\geq88$~см у женщин и $\geq102$~см у мужчин, а при азиатском происхождении --- соответственно $\geq80$ и $\geq90$~см; 2) ХС ЛВП $<40$~мг/дл у мужчин и $<50$~мг/дл у женщин; 3) триглицериды $\geq150$~мг/дл; 4) повышенное АД: систолическое $\geq130$~мм рт. ст. или диастолическое $\geq80$~мм рт. ст. и/или приём антигипертензивных препаратов; 5) глюкоза крови натощак $\geq100$~мг/дл.

\textsuperscript{\ddag} Порог коронарного кальция CAC $>100$ основан на высоком абсолютном риске ССЗ, влияющем на выбор лечения.

\textsuperscript{§} Умеренному или высокому риску ХБП по KDIGO соответствуют G1--G2 с A2--A3, G3a с A1--A2 либо G3b с A1; очень высокому риску --- G3a с A3, G3b с A2--A3 либо G4--G5. Альбуминурия A1: UACR $<30$~мг/г; A2: UACR от 30 до $<300$~мг/г; A3: UACR $\geq300$~мг/г.

\editorial{В оригинале знак \ddag{} стоит после упоминания эхокардиографических параметров, однако соответствующая сноска посвящена CAC. Кроме того, в основной строке указан порог CAC $\geq100$, а в сноске --- $>100$. Эти несогласованности переданы явно, без самостоятельного исправления клинического критерия.}

{\small Дополнительные обозначения таблицы 4: AHA/NHLBI --- Американская кардиологическая ассоциация/Национальный институт сердца, лёгких и крови США; CT --- компьютерная томография; DBP --- диастолическое АД; SBP --- систолическое АД; MetS --- метаболический синдром. Остальные сокращения раскрыты в разделе 2.2 и тексте таблицы.}

\subsubsection{Определения компонентов стадий ССПМ у детей и подростков}
\heading{Обзор}
Метаболические факторы риска и ХБП всё чаще появляются в возрасте до 18 лет, нередко сохраняются и прогрессируют во взрослом возрасте. Поэтому стадирование ССПМ показано на протяжении всей жизни. Пороговые значения факторов риска у детей и подростков имеют педиатрическую специфику. Из-за возрастных изменений роста для диагностики избыточной массы тела и ожирения у детей используют процентили ИМТ, а не фиксированные значения, применяемые у взрослых.\src{1} Конкретные критерии приведены в таблице 5, включая дополнительные процентильные пороги для ожирения 2-го и 3-го классов.

У детей младше 13 лет АД классифицируют по процентилям; в возрасте 13--17 лет применяют те же пороги АД, что и у взрослых. Поскольку данных для определения конкретного уровня АД в детстве, связанного с сердечно-сосудистыми исходами во взрослом возрасте, недостаточно, гипертензию определяют по нормативному распределению АД у здоровых детей.\src{2} У подростков соответствующий процентиль иногда превышает взрослые пороги, поэтому для упрощения практического применения у них приняты взрослые пороговые значения.\src{3}

Пороги патологических уровней липидов у детей приведены в таблице 5. Для диагностики дислипидемии предпочтительна липидограмма натощак, однако допустима и проба, взятая не натощак.\src{4} Критерии нарушений гликемии у детей совпадают с критериями для взрослых (таблица 5).\src{5}

\heading{Таблица 5. Диагностические критерии ожирения, артериальной гипертензии, дислипидемии и диабета у детей и подростков ($<18$ лет)}
\heading{Избыточная масса тела и ожирение\src{1}}
\begin{itemize}
\item \textbf{Избыточная масса тела:} ИМТ от 85-го до менее 95-го процентиля.
\item \textbf{Ожирение:} ИМТ $\geq95$-го процентиля.
\item \textbf{Тяжёлое ожирение:} ИМТ $\geq120\%$ от значения 95-го процентиля.
\item \textbf{Ожирение 2-го класса:} ИМТ от $120\%$ до менее $140\%$ от значения 95-го процентиля либо ИМТ от 35 до менее 40~кг/м$^2$ --- выбирают меньший порог с учётом возраста и пола.
\item \textbf{Ожирение 3-го класса:} ИМТ $\geq140\%$ от значения 95-го процентиля либо ИМТ $\geq40$~кг/м$^2$ --- выбирают меньший порог с учётом возраста и пола.
\end{itemize}

\heading{Артериальное давление\src{2}}
\textbf{Возраст младше 13 лет:}
\begin{itemize}
\item норма: менее 90-го процентиля;
\item повышенное АД: от 90-го до менее 95-го процентиля либо от 120/80~мм рт. ст. до менее 95-го процентиля --- выбирают меньший порог;
\item гипертензия 1-й стадии: от 95-го процентиля до менее чем «95-й процентиль + 12~мм рт. ст.» либо 130/80--139/89~мм рт. ст. --- выбирают меньший порог;
\item гипертензия 2-й стадии: $\geq$ «95-й процентиль + 12~мм рт. ст.» либо $\geq140/90$~мм рт. ст. --- выбирают меньший порог.
\end{itemize}

\textbf{Возраст 13 лет и старше:}
\begin{itemize}
\item норма: систолическое АД $<120$ и диастолическое $<80$~мм рт. ст.;
\item повышенное АД: систолическое 120--129 и диастолическое $<80$~мм рт. ст.;
\item гипертензия 1-й стадии: 130/80--139/89~мм рт. ст.;
\item гипертензия 2-й стадии: $\geq140/90$~мм рт. ст.
\end{itemize}

\heading{Дислипидемия\src{4}}
\textbf{Липидограмма натощак:}
\begin{itemize}
\item общий холестерин $\geq200$~мг/дл;
\item ХС ЛВП $<40$~мг/дл;
\item ХС ЛНП $\geq130$~мг/дл;
\item триглицериды: в возрасте 0--9 лет $\geq100$~мг/дл, в возрасте 10--19 лет $\geq130$~мг/дл;
\item холестерин не-ЛВП $\geq145$~мг/дл.
\end{itemize}
\textbf{Липидограмма не натощак:} холестерин не-ЛВП $\geq145$~мг/дл; ХС ЛВП $<40$~мг/дл.

\heading{Нарушения гликемии\src{5}}
\textbf{Предиабет --- любой из перечисленных критериев:}
\begin{itemize}
\item глюкоза плазмы натощак 100--125~мг/дл; \textbf{или}
\item HbA1c 5,7--6,4\%; \textbf{или}
\item глюкоза плазмы через 2 часа в пероральном глюкозотолерантном тесте с 75~г глюкозы 140--199~мг/дл.
\end{itemize}
\textbf{Диабет --- любой из перечисленных критериев:}
\begin{itemize}
\item глюкоза плазмы натощак $\geq126$~мг/дл; \textbf{или}
\item HbA1c $\geq6{,}5\%$; \textbf{или}
\item глюкоза плазмы через 2 часа в пероральном глюкозотолерантном тесте с 75~г глюкозы $\geq200$~мг/дл; \textbf{или}
\item глюкоза плазмы при случайном определении $\geq200$~мг/дл у пациента с классическими симптомами гипергликемии либо гипергликемическим кризом.
\end{itemize}
\editorial{Возрастной интервал 10--19 лет для триглицеридов сохранён из таблицы оригинала, хотя в её заголовке указан возраст младше 18 лет.}


\section{Обследование и диагностика}
\subsection{Диагностический подход к определению стадии ССПМ-синдрома}
\heading{Рекомендации по диагностическому подходу к стадированию}
Исследования, на которые опираются рекомендации, обобщены в таблице доказательств оригинального руководства.

{\small
\begin{longtable}{P{0.07\textwidth}P{0.11\textwidth}P{0.75\textwidth}}
\toprule COR & LOE & Рекомендация\\\midrule
\endfirsthead
\toprule COR & LOE & Рекомендация (продолжение)\\\midrule
\endhead
1 & B-NR & \textbf{1.} У детей и подростков ($<18$ лет), а также у взрослых ($\geq18$ лет) следует определять стадию ССПМ-синдрома путём оценки метаболических факторов риска, функции почек --- расчёт рСКФ с дополнительным определением UACR при стадии ССПМ $\geq2$ --- и наличия ССЗ. Это необходимо для предотвращения прогрессирования, содействия обратному развитию синдрома и персонализации лечения в соответствии с абсолютным риском ССЗ и ожидаемой чистой пользой терапии.\src{1--3}\\[5pt]
2a & B-NR & \textbf{2.} У взрослых с промежуточным 10-летним риском АССЗ --- PREVENT-ASCVD от 5 до $<10\%$ --- либо у отдельных взрослых с пограничным риском --- PREVENT-ASCVD от 3 до $<5\%$, --- когда решение о профилактической терапии остаётся неопределённым, выявление субклинического коронарного атеросклероза по коронарному кальцию (CAC) может быть полезным для выбора оптимальной профилактики ССЗ.\src{4--6}\\[5pt]
2a & B-NR & \textbf{3.} У взрослых с повышенным прогнозируемым 10-летним риском СН --- PREVENT-HF $\geq5\%$ --- обследование для выявления pre-HF с использованием сердечных биомаркеров: натрийуретических пептидов BNP и NT-proBNP, высокочувствительного сердечного тропонина (hs-cTn), --- может быть полезным для выбора дополнительных диагностических исследований, например визуализации сердца, и организации согласованной помощи для оптимальной профилактики СН.\src{7--10}\\
\bottomrule
\end{longtable}
}

\heading{Обзор}
Для определения стадии ССПМ-синдрома необходима оценка метаболического здоровья, функции почек и сердечно-сосудистых факторов риска. Это улучшает выявление часто нераспознанных либо бессимптомных компонентов синдрома. Определение стадии помогает направлять профилактику и лечение и способствует переходу к более ранней стадии благодаря существенным изменениям образа жизни или значительному снижению массы тела.

Важно выполнять стадирование, когда пациент клинически стабилен: это обеспечивает точность исходной оценки, воспроизводимость результатов и возможность интерпретировать изменения стадии в процессе наблюдения. В идеале антропометрическая оценка должна включать не только ИМТ, но и окружность талии, поскольку абдоминальное жироотложение играет центральную роль в ССПМ-синдроме. Другие прямые или антропометрические измерения жировой ткани, если они доступны, могут дополнительно уточнить ССПМ-риск.

На стадиях ССПМ 2--4 рекомендуется определять UACR в дополнение к рСКФ, чтобы всесторонне охарактеризовать ХБП и связанный с ней риск ССЗ. Это основано на наблюдаемой более высокой распространённости патологического UACR в популяциях, где уже присутствуют метаболические факторы риска или ХБП. Выявление субклинического ССЗ --- стадии 3, то есть субклинического коронарного атеросклероза либо pre-HF, --- может влиять на ведение посредством реклассификации риска ССЗ и уточнения прогноза. Это позволяет выбирать индивидуализированные, более интенсивные профилактические стратегии для лиц высокого риска.

\heading{Обоснование отдельных рекомендаций}
\textbf{1. Определение стадии ССПМ-синдрома.} Стадирование (таблица 4) позволяет выявлять и лечить пациентов на разных стадиях в соответствии с их абсолютным риском ССЗ и ожидаемой пользой вмешательств, снижающих риск. Интервенционные исследования, направленные на различные факторы ССПМ-риска, указывают, что более раннее выявление и начало вмешательства часто связаны с большей клинической пользой.

В 10-летнем продольном исследовании в Великобритании\src{11} оценивали влияние скрининговой программы, в рамках которой основные выявленные факторы риска ССЗ измеряли и регистрировали. Лицам с высоким риском предоставляли рекомендации по укреплению здоровья по стандартизированному сценарию. Скрининг был связан с улучшением уровней факторов риска --- АД и холестерина --- и поведения в отношении здоровья.

Проспективные когортные исследования также показали, что ожидаемая польза профилактической терапии зависит от совокупности факторов риска. В моделирующем исследовании с использованием нескольких когорт измерение клинических предикторов ССЗ, в том числе применяемых при стадировании ССПМ, позволяло оценивать индивидуальный прогноз и эффект лечения в отношении улучшения 10-летнего риска, риска на протяжении жизни и ожидаемой продолжительности жизни без ССЗ.\src{1} Оценка факторов риска, необходимая для стадирования ССПМ, позволяет количественно определить 10- и 30-летний риск АССЗ, СН и всех ССЗ с помощью уравнений PREVENT.\src{12}

\textbf{2. Оценка коронарного кальция.} Оценка CAC улучшает прогнозирование АССЗ\src{4,5} и отбор пациентов для интенсивной коррекции факторов риска АССЗ, в частности назначения статинов.\src{6} Данные MESA (Multi-Ethnic Study of Atherosclerosis --- многоэтническое исследование атеросклероза)\src{4} показали тесную связь CAC с 10-летним риском впервые возникшего АССЗ независимо от стандартных факторов риска, причём сходную в разных демографических группах.

По мере увеличения кальциевого индекса наблюдается, по-видимому, логарифмическое увеличение риска без убедительных признаков достижения плато. Например, у лиц с CAC $\geq1000$ риск ССЗ и смерти от всех причин выше, чем при CAC 400--999.\src{5}

В исследовании CorCal --- исследовании эффективности активной стратегии первичной сердечно-сосудистой профилактики со скринингом CAC или без него для предотвращения будущих больших неблагоприятных сердечно-сосудистых событий ---\src{6} 601 пациента первичной профилактики рандомизировали в группу выбора тактики на основе CAC ($n=302$) либо калькулятора риска ($n=299$). В группе CAC реклассификация риска была более выраженной. Через год в ней отмечались более высокая приверженность приёму статинов, более низкий уровень ХС ЛНП и сопоставимые либо меньшие расчётные затраты.

С учётом влияния CAC на реклассификацию риска его определение можно рассматривать в соответствии с междисциплинарным руководством ACC/AHA 2026 года по дислипидемии\src{13} у лиц с пограничным или промежуточным 10-летним риском --- PREVENT-ASCVD 3--10\%, --- когда решение о профилактической терапии остаётся неопределённым.

\editorial{В пояснительном абзаце оригинала указан диапазон «3--10\%». В формальной рекомендации 2 верхняя граница промежуточного риска обозначена как $<10\%$; оба фрагмента переданы с сохранением исходной записи.}

\textbf{3. Выявление доклинической СН с помощью биомаркеров и визуализации.} Повышение сердечных биомаркеров --- BNP, NT-proBNP\src{7,8} или hs-cTn\src{9} --- либо структурные или функциональные отклонения при визуализации сердца, например эхокардиографии,\src{10} позволяют выявлять pre-HF, связанную с повышенным риском СН. Одновременное наличие обоих типов признаков указывает на наибольший абсолютный риск.\src{14}

В исследовании STOP-HF (St. Vincents Screening to Prevent HF --- скрининг для профилактики СН)\src{7} у пациентов в возрасте $\geq40$ лет с хотя бы одним фактором риска ССЗ определение BNP, за которым при BNP $\geq50$~пг/мл следовали эхокардиография и направление для согласованного ведения, было связано с уменьшением шансов развития дисфункции левого желудочка или СН примерно на 45\%.

Крупные наборы данных из нескольких регионов показывают, что сердечные биомаркеры улучшают прогнозирование СН. Увеличение левого предсердия, гипертрофия левого желудочка, патологическая глобальная продольная деформация и повышенное отношение $E/e'$ по эхокардиографии являются независимыми предикторами впервые возникшей СН.\src{10}

Если pre-HF ранее не диагностирована, 10-летний риск PREVENT-HF $\geq5\%$ можно использовать как основание для избирательного исследования сердечных биомаркеров. Это связано с тем, что в популяционных когортах повышенные биомаркеры соответствовали сходному 10-летнему риску впервые возникшей СН --- 3--5\%.\src{15--18} Такой подход позволяет выявить лиц с повышенным риском и организовать согласованную профилактическую помощь; при необходимости можно рассмотреть визуализацию сердца для уточнения оценки риска.\src{15--18}

Доказательства также поддерживают определение сердечных биомаркеров на основании возраста и факторов ССПМ-риска --- например в возрасте $\geq40$ лет при наличии хотя бы одного фактора ССПМ-риска, как в STOP-HF.\src{7} Однако диагностическая результативность такого подхода ниже, чем стратегии, основанной на интегральной оценке риска СН по PREVENT-HF.\src{19,20} Дополнительная прогностическая ценность исследования биомаркеров у пациентов с далеко зашедшей ХБП или АССЗ неясна, поскольку в этих подгруппах повышенные биомаркеры встречаются очень часто.


\subsubsection{Продольная диагностическая оценка на всех стадиях ССПМ}
\heading{Рекомендации по диагностической оценке в динамике}
Исследования, обосновывающие рекомендации, обобщены в таблице доказательств оригинала.
{\small
\begin{longtable}{P{0.07\textwidth}P{0.11\textwidth}P{0.75\textwidth}}
\toprule COR & LOE & Рекомендация\\\midrule
\endfirsthead
\toprule COR & LOE & Рекомендация (продолжение)\\\midrule
\endhead
1 & B-NR & \textbf{1.} Всем взрослым с ССПМ-синдромом или риском его развития рекомендуется измерять ИМТ и окружность талии не реже одного раза в год для выявления риска перехода к более поздней стадии.\src{1,2}\\[5pt]
1 & B-NR & \textbf{2.} Всем взрослым без ССПМ-синдрома, то есть на стадии ССПМ 0, рекомендуется оценивать липиды, гликемию и функцию почек с расчётом рСКФ не реже одного раза в 5 лет, а АД --- не реже одного раза в год. Это необходимо для своевременного выявления факторов ССПМ-риска и оптимальной профилактики ССЗ.\src{3,4}\\[5pt]
1 & B-NR & \textbf{3.} Взрослым на стадии ССПМ 1 рекомендуется оценивать липиды, гликемию и функцию почек с расчётом рСКФ не реже одного раза в 2--3 года, а АД --- не реже одного раза в год, чтобы своевременно выявлять и корректировать факторы ССПМ-риска для оптимальной профилактики ССЗ.\src{5--9}\\[5pt]
1 & B-NR & \textbf{4.} Взрослым на стадии ССПМ 2 и выше рекомендуется оценивать липиды, гликемию и АД не реже одного раза в год для выбора оптимальной профилактики и ведения ССЗ.\src{10--13}\\[5pt]
1 & B-NR & \textbf{5.} Взрослым на стадии ССПМ 2 и выше рекомендуется определять и рСКФ, и UACR не реже одного раза в год для характеристики риска, связанного с ХБП, и выбора профилактики и ведения сердечно-сосудистых и почечных заболеваний.\src{14--16}\\
\bottomrule
\end{longtable}
}

\heading{Обзор}
Наблюдение за факторами ССПМ-риска в динамике важно для информирования клиницистов и пациентов о возникновении или прогрессировании синдрома. Более раннее выявление факторов риска и своевременное начало соответствующей терапии, уменьшающей риск, связанный с диабетом, гипертензией, дислипидемией и ХБП, ассоциируются с улучшением клинических исходов. В основе ССПМ-синдрома лежат избыток и дисфункция жировой ткани, особенно висцеральное жироотложение; избыток жировой ткани является модифицируемым фактором риска на всех стадиях синдрома.\src{1} Поэтому при наличии ССПМ-синдрома или риска его развития ИМТ и окружность талии следует оценивать ежегодно. В соответствии с действующими рекомендациями всем взрослым также рекомендуется ежегодно оценивать АД.\src{17}

Подход к наблюдению за другими факторами зависит от стадии ССПМ (рисунок 3). Рекомендуется регулярно оценивать липиды, гликемию и функцию почек, например рСКФ, причём частота обследований должна возрастать на более поздних стадиях, поскольку по мере развития синдрома увеличивается риск прогрессирования метаболических и почечных нарушений.\src{13,18--21} Начиная со стадии 2 ежегодная оценка должна включать оба показателя состояния почек --- рСКФ и UACR. Вместе они необходимы для надлежащей диагностики ХБП и дополняют друг друга в отношении прогноза и выбора лечения.\src{2,3}

\heading{Рисунок 3. Обследование взрослых в динамике в зависимости от исходной стадии ССПМ: табличная адаптация}
{\small
\begin{longtable}{P{0.16\textwidth}P{0.33\textwidth}P{0.44\textwidth}}
\toprule Исходная стадия & Ежегодно на всех стадиях & Дополнительные обследования\\\midrule
\endfirsthead
\toprule Исходная стадия & Ежегодно на всех стадиях & Дополнительные обследования\\\midrule
\endhead
0 & ИМТ, окружность талии и АД & Липиды, гликемия и рСКФ --- каждые 5 лет\\
1 & ИМТ, окружность талии и АД & Липиды, гликемия и рСКФ --- каждые 2--3 года\\
2--4 & ИМТ, окружность талии и АД & Липиды, гликемия, рСКФ и UACR --- ежегодно\\
\bottomrule
\end{longtable}
}
\textbf{Пояснение авторов к рисунку 3.} Систематическая оценка факторов ССПМ-риска позволяет своевременно выявлять их и количественно определять сердечно-сосудистый риск с помощью PREVENT. Частота и интенсивность обследований увеличиваются на более поздних стадиях. Всем пациентам следует ежегодно проводить антропометрические измерения и оценивать АД. Дополнительные исследования липидов, гликемии и рСКФ необходимо выполнять не реже одного раза в 5 лет на стадии 0 и чаще --- каждые 2--3 года --- на стадии 1. На стадии 2, когда присутствуют метаболические факторы риска, ХБП или их сочетание, эти обследования следует проводить ежегодно. Более высокая распространённость альбуминурии на стадии 2 также обосновывает ежегодное определение UACR для полноценной оценки ХБП и связанного с ней сердечно-сосудистого риска. BMI --- ИМТ; BP --- АД; CKM --- ССПМ; eGFR --- рСКФ; UACR --- отношение альбумина к креатинину мочи.

\heading{Обоснование отдельных рекомендаций}
\textbf{1. Антропометрическая оценка.} Коррекция избыточной массы тела и ожирения занимает центральное место в поддержании ССПМ-здоровья: избыток и дисфункция жировой ткани являются основными причинами формирования метаболических факторов риска и ХБП с последующим увеличением риска ССЗ.\src{2,22,23} Висцеральная жировая ткань способствует воспалению и инсулинорезистентности; возникающие межсистемные нарушения приводят к прогрессированию ССПМ-синдрома. Поэтому регулярная антропометрическая оценка на протяжении жизни важна для раннего выявления ССПМ-риска.\src{24,25}

Предотвращение увеличения массы тела с возрастом связано с меньшей вероятностью развития метаболических факторов риска.\src{2} Ежегодный скрининг избыточной массы тела --- ИМТ от 25 до $<30$~кг/м$^2$ --- и ожирения --- ИМТ $\geq30$~кг/м$^2$ --- поддерживается несколькими клиническими руководствами.\src{18,26} Для отдельных популяционных подгрупп, например американцев азиатского происхождения, предлагается порог избыточной массы тела ИМТ $\geq23$~кг/м$^2$, поскольку повышенный метаболический риск у них наблюдается уже при относительно небольшой массе тела.\src{27,28}

ИМТ имеет ограничения: он не отражает состав тела, в частности абдоминальное ожирение, определяющее ССПМ-риск. Совместное использование ИМТ и окружности талии даёт взаимодополняющую прогностическую информацию о будущих метаболических нарушениях, впервые возникающих ССЗ и смертности. Хотя отношение окружности талии к окружности бёдер характеризует метаболический и сердечно-сосудистый риск лучше, чем только окружность талии, сложность измерения и отсутствие стандартизации делают этот показатель менее удобным для клинической практики.\src{18,29--32}

\textbf{2. Обследование на стадии 0.} Возникновение гипертензии и повышение таких биомаркеров, как триглицериды и глюкоза, являются важнейшими индикаторами кардиометаболического и почечного здоровья.\src{4,33--35} Предотвращение развития гипертензии снижает риск в большей степени, чем лечение уже возникшей гипертензии.\src{36,37} В соответствии с действующими рекомендациями в идеале АД следует измерять при каждом визите; для большинства людей целевое значение составляет $<130/80$~мм рт. ст.\src{17}

Высокие триглицериды и низкий ХС ЛВП --- маркеры инсулинорезистентности; высокие триглицериды могут также причинно способствовать АССЗ.\src{38} Действующие руководства рекомендуют здоровым взрослым исследование липидов, включающее эти показатели, каждые 5 лет.\src{21} Обследование на предиабет и диабет с помощью глюкозы крови натощак, HbA1c или перорального глюкозотолерантного теста крайне важно ввиду тесной связи гипергликемии с сердечно-сосудистыми и почечными заболеваниями.\src{39,40}

Среди людей с нормальной массой тела гипергликемия чаще развивается у представителей небелых расовых групп, поэтому систематическое обследование всех взрослых независимо от ИМТ может способствовать более равноправному выявлению гипергликемии.\src{3} Раннее обнаружение снижения рСКФ также приносит клиническую пользу. На стадии ССПМ 0 оценка гликемии и функции почек не реже одного раза в 5 лет позволяет комплексно оценивать ССПМ-риск одновременно с исследованием липидов, сроки которого установлены рекомендациями для здоровых взрослых, и количественно определять риск ССЗ по PREVENT.

\textbf{3. Обследование на стадии 1.} У лиц со стадией ССПМ 1, обусловленной избыточной массой тела или ожирением, риск возникновения метаболических факторов и ХБП выше, чем у лиц с нормальной массой тела без факторов риска, соответствующих стадии 0.\src{5,6,41} Даже при отсутствии метаболических факторов избыток массы тела связан с повышенным сердечно-сосудистым риском, особенно риском СН.\src{6,42} Исследования «метаболически здорового ожирения» показали, что за 4--5 лет наблюдения метаболический синдром или диабет развивается приблизительно у 30--45\% таких лиц.\src{43,44}

Ввиду высокого краткосрочного риска диабета при избыточной массе тела и ожирении ведущие профессиональные общества и Рабочая группа США по профилактическим мероприятиям рекомендуют обследовать таких лиц на метаболический синдром и диабет каждые 3 года.\src{35,42,45} У взрослых с предиабетом риск впервые возникающего диабета очень высок,\src{7,8} что требует ещё более интенсивного наблюдения: ведущие общества рекомендуют ежегодное обследование на СД2.\src{45} Более интенсивное наблюдение необходимо и после гестационного диабета из-за повышенного риска СД2.\src{46}

Более раннее возникновение метаболических факторов связано с увеличением риска ССЗ и смертности, что подчёркивает необходимость их своевременного выявления и лечения.\src{47--49} Поэтому на стадии 1 оценивать метаболические факторы и рСКФ следует не реже одного раза в 2--3 года, чтобы своевременно выявлять факторы ССПМ-риска и при необходимости пересматривать стадию.

\textbf{4. Оценка метаболических факторов на стадии 2 и выше.} Избыточная масса тела, большая окружность талии, высокое АД, высокие триглицериды, низкий ХС ЛВП и гипергликемия связаны с увеличением риска ССЗ. Экспериментальные, генетические и интервенционные данные поддерживают причинную связь для всех перечисленных факторов, кроме низкого ХС ЛВП.\src{10--12,50} Выявление метаболического синдрома указывает на более высокий риск диабета и ССЗ, а также на большую вероятность дополнительных патофизиологических нарушений, способствующих ССЗ: системного воспаления, увеличения числа частиц ЛНП, протромботического состояния и эндотелиальной дисфункции.\src{50,51}

У пациентов с диабетом для оценки гликемии следует использовать HbA1c как показатель гликемического контроля за предшествующие 2--3 месяца, характеризующий риск микро- и макрососудистых осложнений.\src{52} Взаимосвязанность метаболических факторов при ССПМ-синдроме обусловливает их частое сочетание и необходимость комплексной оценки. Обследование как минимум ежегодно на стадии 2 позволяет рано интенсифицировать изменение образа жизни и при необходимости лекарственную терапию, чтобы корректировать патологические показатели и уменьшать риск сердечно-сосудистых событий и прогрессирования ХБП. Более своевременное начало лечения связано с улучшением сердечно-сосудистых клинических исходов при гипертензии, диабете и дислипидемии.\src{18,53}

\textbf{5. Совместная оценка рСКФ и UACR.} Диагностика ХБП включает определение обоих показателей; критерием служит хроническое снижение рСКФ $<60$~\unitgfr\ либо UACR $\geq30$~мг/г.\src{54} Большинство людей с ХБП не знают о своём диагнозе.\src{55} Большинство случаев ХБП обусловлено метаболическими факторами --- СД2 и гипертензией. Повышенная распространённость альбуминурии у лиц с этими факторами либо сниженной рСКФ обосновывает регулярное определение обоих показателей для выявления заболевания.\src{16,55}

У пациентов с уже установленной ХБП наличие альбуминурии влияет на выбор терапии для уменьшения связанного с ХБП риска (раздел 5.5.4 «Ведение пациентов с ХБП при стадиях ССПМ 2--3»). Как отражено в карте KDIGO (рисунок 2), рСКФ и UACR предоставляют взаимодополняющую и суммирующуюся прогностическую информацию о риске прогрессирования ХБП, ССЗ и смертности.\src{15} В клинической практике СКФ обычно рассчитывают по формулам на основе креатинина. Расчёт по креатинину и цистатину C вместе точнее, чем по любому из этих маркеров отдельно; однако доступность определения цистатина C ограничена.\src{14}

Для выявления ХБП на стадии ССПМ 2 рекомендуется определять рСКФ и UACR как минимум ежегодно. При далеко зашедшей ХБП можно рассмотреть более частую оценку --- каждые 3--6 месяцев --- для прогнозирования или реклассификации почечного и сердечно-сосудистого риска в соответствии с рекомендациями KDIGO.\src{54}

\subsection{Оценка социальных детерминант здоровья}
\heading{Рекомендации по оценке социальных детерминант здоровья}
Исследования, обосновывающие рекомендации, обобщены в таблице доказательств оригинала.
{\small
\begin{longtable}{P{0.07\textwidth}P{0.11\textwidth}P{0.75\textwidth}}
\toprule COR & LOE & Рекомендация\\\midrule
\endfirsthead
\toprule COR & LOE & Рекомендация (продолжение)\\\midrule
\endhead
1 & B-NR & \textbf{1.} Взрослым с ССПМ-синдромом или риском его развития ($\geq18$ лет) рекомендуется регулярный скрининг с помощью валидированного инструмента оценки SDOH (таблица 6) для реализации помощи, ориентированной на пациента.\src{1--7}\\[5pt]
1 & B-NR & \textbf{2.} При ССПМ-синдроме или риске его развития у детей и подростков ($<18$ лет) рекомендуется регулярно обследовать лиц, осуществляющих уход за ними, с помощью валидированного инструмента оценки SDOH таких лиц (таблица 7), чтобы обеспечить помощь, ориентированную на пациента.\src{8--10}\\
\bottomrule
\end{longtable}
}

\heading{Обзор}
Социальные детерминанты здоровья --- немедицинские факторы, влияющие на риск заболеваний и исходы, --- играют ключевую роль в ССПМ-здоровье на протяжении всей жизни.\src{11--14} Неблагоприятные SDOH на нескольких уровнях непосредственно влияют на поведение в отношении здоровья. У детей и подростков они способствуют неблагоприятным условиям ранней жизни в критические периоды развития и непосредственно воздействуют на сердечно-сосудистое здоровье во взрослом возрасте.\src{12} Многоуровневые социальные факторы также изменяют биологические пути, усиливающие кардиометаболический риск.\src{11,15} Посредством этих механизмов неблагоприятные SDOH приводят к формированию факторов ССПМ-риска.\src{16--20}

Неблагоприятные социальные условия ухудшают как способность заботиться о собственном здоровье, так и взаимодействие с системой здравоохранения. В конечном счёте это влияет на сердечно-сосудистую и общую смертность\src{11,14} и способствует неравенству в смертности от ССЗ.\src{21} Поэтому SDOH крайне важны для оценки и снижения риска ССПМ-синдрома и занимают центральное место в модели помощи.\src{5--7,20} Их значение также отражено включением индекса социальной депривации --- территориального показателя SDOH --- в уравнение PREVENT. Неблагоприятные социальные факторы, выявленные при клиническом скрининге, следует учитывать и корректировать в рамках междисциплинарной модели помощи (раздел 5.1).

С точки зрения пациентов скрининг социальных потребностей в целом приемлем, однако необходимо ясно объяснять связь SDOH с исходами ССПМ-синдрома, чтобы способствовать активному участию пациентов в скрининге и последующем направлении за помощью.\src{22,23} Medicare и Medicaid возмещают оценку SDOH каждые 6 месяцев в амбулаторных условиях и ежегодно при госпитализации; вместе с тем оценка социальных потребностей на каждом визите может лучше поддерживать профилактику и ведение ССПМ-синдрома.\src{24--26}

\heading{Обоснование отдельных рекомендаций}
\textbf{1. Скрининг у взрослых.} Регулярную оценку SDOH и связанных с ними социальных потребностей можно встроить в оказание помощи взрослым с ССПМ-синдромом, чтобы выявлять немедицинские факторы, на которые можно воздействовать в рамках помощи, ориентированной на пациента. Индивидуальные, территориальные и структурные SDOH и связанные с ними потребности можно оценивать с помощью валидированных инструментов.\src{1}

Инструменты для взрослых существенно различаются, но ключевые области оценки включают финансовые трудности, недостаточную обеспеченность продовольствием, нестабильность жилищных условий, безопасность и подверженность насилию, транспортные затруднения и потребности в коммунальных услугах (таблицы 6 и 7). Дополнительно можно оценивать поддержку семьи и социального окружения, поведение в отношении здоровья, медицинскую грамотность и психическое здоровье.\src{3} Немногочисленные интервенционные исследования показали улучшение АД, липидов, отказа от табака и потребления овощей и фруктов в результате выявления неблагоприятных SDOH и работы с ними.\src{27}

Инструменты скрининга можно интегрировать в электронную медицинскую документацию, чтобы уменьшить препятствия для их применения,\src{2,4} например трудности включения скрининга в привычный порядок приёма пациента. Другие сложности: недостаточная стандартизация методов скрининга и элементов данных, опасения по поводу конфиденциальности, особенно в небольших сельских сообществах, а также необходимость совершенствовать системы направления и вмешательства для удовлетворения социальных потребностей.\src{2,28}

\textbf{2. Скрининг лиц, осуществляющих уход за детьми и подростками.} Некоторые инструменты SDOH валидированы для всех возрастов, но разработаны и специальные инструменты для лиц, ухаживающих за детьми и подростками (таблица 7). Такой скрининг выявляет воздействие в детстве социальных факторов, способных влиять на ССПМ-здоровье на протяжении всей жизни.\src{9} Основные области оценки у лиц, ухаживающих за детьми с ССПМ-синдромом: экономическая стабильность, образование, семейная обстановка, здоровье и доступность медицинской помощи, а также условия района проживания.\src{9} Некоторые из этих инструментов прошли валидацию и проверку надёжности.\src{9}

Можно рассмотреть дополнительные инструменты индивидуального и территориального скрининга, специально оценивающие факторы, влияющие на развитие ССПМ-синдрома: сегрегацию по месту жительства, социальную уязвимость, индекс качества воздуха, условия для ходьбы и езды на велосипеде, медицинскую грамотность, дискриминацию, духовную сферу и другие факторы.\src{10} Для этой группы существуют особые препятствия: например клиницисты могут избегать скрининга из опасений, что работа с социальными потребностями нарушит конфиденциальность лиц, осуществляющих уход.\src{8}

\heading{Таблица 6. Инструменты скрининга SDOH у взрослых с ССПМ-синдромом или риском его развития}
\editorial{Каждый пункт соответствует строке исходной таблицы: название инструмента и все перечисленные в ней основные области оценки. Английские названия инструментов сохранены для их идентификации.}
\heading{Health Leads\src{29}}
Подверженность насилию; финансовые трудности; недостаточная обеспеченность продовольствием; нестабильность жилищных условий; социально-демо\-графические сведения; транспортные затруднения; потребности в коммунальных услугах.

\heading{Инструмент AHC HRSN Центра инноваций Medicare и Medicaid\src{30}}
Accountable Health Communities (AHC) Health-Related Social Needs Screening Tool: недостаточная обеспеченность продовольствием; нестабильность жилищных условий; безопасность в межличностных отношениях; транспортные проблемы; потребность в помощи с коммунальными услугами.

\heading{AAFP: The EveryONE Project\src{31}}
Уход за детьми; занятость; образование; финансы; питание; жильё; безопасность в межличностных отношениях; транспорт; коммунальные услуги.

\heading{PRAPARE Implementation and Action Toolkit\src{32}}
Семья и дом; деньги и ресурсы; социальное и эмоциональное здоровье; социально-демографические характеристики.

\heading{OCHIN: инструменты SDOH в электронных медицинских картах общественных центров здоровья\src{33}}
Образование; подверженность насилию; ограниченность финансовых ресурсов; недостаточная обеспеченность продовольствием; употребление алкоголя; расовая и этническая принадлежность; употребление табака и воздействие табачного дыма; недостаточная физическая активность; жилищная незащищённость; психическое здоровье --- депрессия, социальная изоляция, стресс.

\heading{HealthBegins Upstream Risks Screening Tool\src{34,35}\textsuperscript{*}}
Экономическая стабильность; образование; недостаточная обеспеченность продовольствием; жилищная незащищённость; социальное окружение и местное сообщество; район проживания и физическая среда; поведение в отношении здоровья --- физическая активность и характер питания.

\textbf{Примечания исходной таблицы 6.} Модифицировано с разрешения по работе Ndumele и соавт.\src{40} © 2023 American Heart Association, Inc. Z-коды Международной классификации болезней 10-го пересмотра (МКБ-10) предназначены для дополнительного кодирования и не являются кодами основного диагноза. Z55 --- проблемы, связанные с образованием и грамотностью; Z57 --- профессиональное воздействие факторов риска; Z58 --- проблемы, связанные с физической окружающей средой; Z59 --- проблемы, связанные с жилищными и экономическими обстоятельствами; Z60.9 --- неуточнённые проблемы, связанные с социальным окружением.

\textsuperscript{*} Инструмент ещё не валидирован либо сведения о валидности не указаны. AAFP --- Американская академия семейных врачей; HRSN --- социальные потребности, связанные со здоровьем; OCHIN --- Oregon Community Health Information Network, Орегонская информационная сеть общественного здоровья; PRAPARE --- Protocol for Responding to and Assessing Patients' Assets, Risks, and Experiences, протокол оценки ресурсов, рисков и опыта пациентов и реагирования на них; CKM --- ССПМ; SDOH --- социальные детерминанты здоровья.

\heading{Таблица 7. Инструменты скрининга SDOH у лиц, осуществляющих уход за детьми и подростками с ССПМ-синдромом или риском его развития}
\heading{HealthBegins Upstream Risks Screening Tool\src{34,35}\textsuperscript{*}}
Экономическая стабильность; образование; недостаточная обеспеченность продовольствием; жилищная незащищённость; социальное окружение и местное сообщество; район проживания и физическая среда; поведение в отношении здоровья --- физическая активность и характер питания.

\heading{Child ACE Tool\src{36}\textsuperscript{*}}
\textbf{Образование:} образование родителей. \textbf{Семейная обстановка:} подозрение на жестокое обращение с ребёнком; насилие со стороны интимного партнёра в семье; психическое заболевание у члена семьи; злоупотребление психоактивными веществами в семье; нахождение члена семьи в заключении; семейное положение родителей.

\heading{SEEK PSQ\src{37}}
\textbf{Экономическая стабильность:} недостаток продовольствия. \textbf{Семейная обстановка:} насилие между родителями со стороны интимного партнёра; депрессия или стресс у родителей; расстройство, связанное с употреблением психоактивных веществ или алкоголя, у родителей; употребление табака в доме; наличие огнестрельного оружия дома; помощь с уходом за ребёнком при необходимости. \textbf{Здоровье и медицинская помощь:} потребность в дымовом извещателе; потребность в контактных данных токсикологической службы.

\heading{IHELLP Screening Tool\src{38}}
\textbf{Экономическая стабильность:} недостаток продовольствия и нестабильность жилищных условий. \textbf{Образование:} опасения, связанные с образовательными потребностями ребёнка. \textbf{Семейная обстановка:} насилие в семье. \textbf{Здоровье и медицинская помощь:} опасения, связанные с медицинским страхованием ребёнка. \textbf{Условия проживания:} опасения относительно физического состояния жилья.

\heading{WE CARE Screening Tool\src{39}}
\textbf{Экономическая стабильность:} недостаток продовольствия; нестабильность жилищных условий; трудности с оплатой счетов; занятость родителей. \textbf{Образование:} образование родителей; отсутствие возможностей ухода за ребёнком. \textbf{Семейная обстановка:} насилие со стороны интимного партнёра в семье; симптомы депрессии у родителей; расстройство, связанное с употреблением алкоголя или психоактивных веществ, у члена семьи.

\textbf{Примечания исходной таблицы 7.} Модифицировано с разрешения по работе Ndumele и соавт.\src{40} © 2023 American Heart Association, Inc. Z-коды МКБ-10 используются для дополнительного кодирования, а не как коды основного диагноза: Z55 --- проблемы образования и грамотности; Z57 --- профессиональное воздействие факторов риска; Z58 --- проблемы физической окружающей среды; Z59 --- жилищные и экономические проблемы; Z60.9 --- неуточнённые проблемы социального окружения.

\textsuperscript{*} Инструмент ещё не валидирован либо сведения о валидности не указаны. ACE --- Adverse Childhood Experiences, неблагоприятный опыт детства; IHELLP --- Income, Housing, Education, Legal Status, Literacy, and Personal Safety, доход, жильё, образование, правовой статус, грамотность и личная безопасность; SEEK PSQ --- Safe Environment for Every Kid Parent Screening Questionnaire, скрининговый опросник родителей «Безопасная среда для каждого ребёнка»; WE CARE --- Well Child Care, Evaluation, Community Resources, Advocacy, Referral, Education, профилактическая помощь ребёнку, оценка, ресурсы сообщества, защита интересов, направление и обучение; CKM --- ССПМ; SDOH --- социальные детерминанты здоровья.


\section{Оценка риска}
\subsection{Количественная оценка риска ССЗ}
\editorial{В оглавлении оригинала для раздела 4.1 указана страница e64. Фактически раздел начинается на странице e67 (страница 18 исходного PDF); здесь использована фактическая пагинация.}
\heading{Рекомендации по количественной оценке риска ССЗ}
Исследования, обосновывающие рекомендации, обобщены в таблице доказательств оригинала.
{\small
\begin{longtable}{P{0.07\textwidth}P{0.11\textwidth}P{0.75\textwidth}}
\toprule COR & LOE & Рекомендация\\\midrule
\endfirsthead
\toprule COR & LOE & Рекомендация (продолжение)\\\midrule
\endhead
1 & B-NR & \textbf{1.} Взрослым в возрасте 30--79 лет без ССЗ --- коронарной болезни сердца, инсульта или СН --- рекомендуется рассчитывать 10-летний риск ССЗ и его компонентов, АССЗ и СН, с помощью PREVENT для количественной оценки риска, связанного с ССПМ-синдромом, и выбора профилактических стратегий.\src{1--4}\\[5pt]
2a & B-NR & \textbf{2.} У взрослых в возрасте 30--59 лет без ССЗ --- коронарной болезни сердца, инсульта или СН --- расчёт 30-летнего риска ССЗ и его компонентов, АССЗ и СН, с помощью PREVENT может быть полезен для количественной оценки риска, связанного с ССПМ-синдромом, и выбора профилактических стратегий.\src{1--4}\\
\bottomrule
\end{longtable}
}

\heading{Обзор}
Количественная оценка риска ССЗ занимает центральное место в обследовании и ведении ССПМ-синдрома: она используется для определения стадии (раздел 2.3), выбора обследований на субклиническое ССЗ (раздел 3.1) и лечебных рекомендаций, чтобы интенсивность профилактики соответствовала индивидуальному риску ССЗ (раздел 5.5.1).

Ожидаемое абсолютное снижение риска при применении конкретной терапии прямо пропорционально исходному прогнозируемому риску пациента. Поэтому стратегия, основанная на риске, эффективнее стратегии, учитывающей только уровни отдельных факторов, в отношении числа предотвращённых событий ССЗ, то есть числа пациентов, которых необходимо лечить для предотвращения одного события.\src{5} Этот вопрос обсуждается в дополняющем научном заявлении AHA/ACC.\src{6--8}

Уравнения PREVENT, доступные в онлайн-калькуляторе PREVENT, позволяют точно оценивать 10- и 30-летний риск ССЗ и его компонентов --- АССЗ и СН. Основным уравнением является PREVENT-CVD, поскольку при ССПМ-синдроме повышен риск и АССЗ, и СН. Количественные оценки PREVENT могут определять лечебные подходы к уменьшению сердечно-сосудистого риска, связанного с диабетом (раздел 5.5.1) и другими факторами, в частности гипертензией и дислипидемией.\src{9,10}

Количественная оценка должна направлять обсуждение между клиницистом и пациентом стратегий снижения риска, включая изменение образа жизни и лекарственную терапию. Она дополняет другие способы оценки риска: учёт индивидуальных сопутствующих заболеваний, усилителей ССПМ-риска (раздел 4.2) и выявление субклинического ССЗ (разделы 5.6.1 «Использование маркеров субклинического атеросклероза» и 5.6.2 «Доклиническая сердечная недостаточность»).

\heading{Обоснование отдельных рекомендаций}
\textbf{1. Оценка 10-летнего риска.} У взрослых 30--79 лет без ССЗ количественная оценка риска полезна для определения приоритетности начала терапии и повышения её результативности в отношении риска ССЗ, связанного с ССПМ-синдромом. После скрининга факторов ССЗ (раздел 3.1) бессимптомным взрослым этого возраста рекомендуется оценивать 10-летний риск с помощью PREVENT.\src{1,11} Эти уравнения используются для стадирования ССПМ, выбора обследований на субклиническое ССЗ и интенсификации профилактики при диабете, а также при дислипидемии и гипертензии, как описано в ранее выпущенных руководствах (таблица 8).

Базовые уравнения PREVENT включают традиционные факторы риска ССЗ и дополнительно введённые предикторы, значимые для ССПМ-синдрома, --- ИМТ и рСКФ. Расширенные модели включают UACR, HbA1c и индекс социальной депривации.\src{1,11} Внешняя валидация PREVENT проведена в разнообразной по составу группе из 3\,330\,085 взрослых жителей США. Для ССЗ продемонстрированы хорошая или отличная дискриминация: медиана C-статистики [межквартильный диапазон] у женщин 0,79 [0,76--0,81], у мужчин 0,76 [0,73--0,78]; а также точность: коэффициент наклона калибровки у женщин 1,03 [0,81--1,16], у мужчин 0,94 [0,81--1,13]. Сходные характеристики получены для отдельных компонентов --- АССЗ и СН. Тем не менее, как и у любой модели риска, возможны неточности, особенно в недостаточно представленных группах или при повышенном риске ССЗ вследствие усилителей ССПМ-риска.\src{1,12}

\textbf{2. Оценка 30-летнего риска.} У молодых взрослых 30-летний риск ССЗ можно рассчитывать по PREVENT одновременно с 10-летним риском и использовать при совместном принятии решений, особенно в консультировании по образу жизни. Польза оценки 10-летнего риска у молодых взрослых может быть ограничена, поскольку вероятность ССЗ в ближайшее десятилетие у них часто невелика.\src{13} Однако у многих из них 30-летний риск ССЗ и его компонентов повышен, несмотря на низкий 10-летний риск.\src{14--16} Поэтому расчёт 30-летнего риска по PREVENT может быть полезен.

Как абсолютный 30-летний риск, так и соответствующий возрастно-половой процентиль могут помочь реклассифицировать риск у молодых взрослых, особенно если он соответствует 75-му процентилю или выше, по аналогии с применением процентилей CAC. Референсные значения для расчёта процентилей 30-летнего риска доступны на странице калькулятора PREVENT профессионального сайта AHA.\src{17} Оценка 30-летнего риска позволяет, в частности, выявить молодых людей с несколькими факторами и повышенным долгосрочным риском, что может облегчить обсуждение риска между клиницистами и пациентами.\src{1,11}

Другие инструменты объяснения долгосрочного риска включают PREVENT-Risk Age: он определяет возраст, соответствующий прогнозируемому риску, для сравнения с хронологическим возрастом.\src{18} Оценка 30-летнего риска может повышать приверженность рекомендациям по образу жизни, а в некоторых случаях --- влиять на применение лекарственной терапии для предотвращения прогрессирования ССПМ-синдрома.

\heading{Таблица 8. Свод рекомендаций по применению уравнений PREVENT для конкретных исходов и порогов риска в первичной профилактике ССЗ}
\heading{PREVENT-CVD}
\textbf{Оцениваемый исход:} все ССЗ --- совокупность АССЗ и СН.
\begin{itemize}
\item \textbf{Стадирование ССПМ:} определить стадию 3 при 10-летнем риске PREVENT-CVD $\geq20\%$, что соответствует эквиваленту риска субклинического ССЗ. Источник: настоящее руководство AHA/ACC/ADA/ASN 2026 года.
\item \textbf{Лечение ССПМ:} рекомендовать начало терапии на основе GLP-1, SGLT2i или их сочетания при 10-летнем риске PREVENT-CVD $\geq7{,}5\%$, если имеются показания. Источник: настоящее руководство.
\item \textbf{Лечение гипертензии:} рекомендовать начало лекарственной терапии для интенсивного снижения АД при гипертензии 1-й стадии и 10-летнем риске PREVENT-CVD $\geq7{,}5\%$. Источник: руководство AHA/ACC 2025 года по повышенному АД.\src{9}
\end{itemize}

\heading{PREVENT-ASCVD}
\textbf{Оцениваемый исход:} АССЗ --- смертельная коронарная болезнь сердца, несмертельный ИМ, смертельный и несмертельный инсульт.
\begin{itemize}
\item \textbf{Стадирование ССПМ/выявление субклинического ССЗ:} рекомендовать обследование на субклинический атеросклероз при 10-летнем риске PREVENT-ASCVD от 3 до $<10\%$, если остаётся неопределённость относительно начала или интенсификации терапии.
\item \textbf{Лечение:} рекомендовать начало липидснижающей терапии при 10-летнем риске PREVENT-ASCVD $\geq5\%$.
\item Рассмотреть начало липидснижающей терапии при 10-летнем риске PREVENT-ASCVD от 3 до $<5\%$ после оценки усилителей риска, 30-летнего риска АССЗ или CAC, если сохраняется неопределённость.
\item Рассмотреть начало липидснижающей терапии при 30-летнем риске PREVENT-ASCVD $\geq10\%$.
\end{itemize}
Источник всех указанных применений PREVENT-ASCVD в этой таблице: руководство ACC/AHA 2026 года по дислипидемии.\src{10}

\heading{PREVENT-HF}
\textbf{Оцениваемый исход:} СН, включая СН со сниженной, умеренно сниженной и сохранённой фракцией выброса.

\textbf{Стадирование ССПМ/выявление субклинического ССЗ:} рекомендовать обследование на субклиническую СН и координацию помощи при 10-летнем риске PREVENT-HF $\geq5\%$. Источник: настоящее руководство AHA/ACC/ADA/ASN 2026 года.

\textbf{Примечание исходной таблицы 8.} См. дополняющее научное заявление AHA/ACC об оценке риска при ССПМ-синдроме.\src{8} ACC --- Американская коллегия кардиологов; ADA --- Американская диабетическая ассоциация; AHA --- Американская кардиологическая ассоциация; ASN --- Американское общество нефрологии; ASCVD --- АССЗ; CAC --- коронарный кальций; CHD --- коронарная болезнь сердца; CKM --- ССПМ; CVD --- ССЗ; GLP-1 --- глюкагоноподобный пептид-1; HF --- СН; MI --- ИМ; PREVENT --- прогнозирование риска сердечно-сосудистых событий; SGLT2i --- ингибиторы натрий-глюкозного котранспортера 2 типа.

\subsection{Усилители риска при ССПМ-синдроме}
\heading{Рекомендация по учёту усилителей риска}
Исследования, обосновывающие рекомендацию, обобщены в таблице доказательств оригинала.
{\small
\begin{longtable}{P{0.07\textwidth}P{0.11\textwidth}P{0.75\textwidth}}
\toprule COR & LOE & Рекомендация\\\midrule
2a & B-NR & \textbf{1.} У взрослых в возрасте 30--79 лет без ССЗ --- коронарной болезни сердца, инсульта или СН --- целесообразно учитывать усилители риска прогрессирования ССПМ-синдрома при принятии решений об интенсификации профилактических стратегий.\src{1--4}\\
\bottomrule
\end{longtable}
}

\heading{Обзор}
Усилители ССПМ-риска --- это факторы, позволяющие выявить человека с повышенной вероятностью перехода к более поздним стадиям и более высоким риском ССЗ или почечной недостаточности (рисунок 4).\src{1--5} После оценки риска ССЗ по PREVENT (раздел 4.1) и определения стадии ССПМ (раздел 2) усилители риска могут дополнительно уточнять вероятность прогрессирования, клинического ССЗ или ухудшения функции почек.\src{6,7}

Они охватывают широкий круг характеристик (таблица 9): хронические воспалительные заболевания в личном анамнезе; специфичные для пола факторы; нарушения сна; неблагополучие в сфере психического здоровья; семейный анамнез диабета или почечной недостаточности; более выраженное воздействие неблагоприятных SDOH; принадлежность к демографической подгруппе высокого риска.\src{5} Многие усилители ССПМ-риска совпадают с факторами, используемыми для профилактики АССЗ, но существуют и различия, обусловленные целью выявлять риск продвижения по стадиям ССПМ, а не только риск АССЗ.\src{8,9}

Кроме того, усилители ССПМ-риска вместе с 30-летним риском можно учитывать, если польза конкретной терапии для предотвращения прогрессирования остаётся неопределённой на основании стандартной оценки 10-летнего риска, уровней факторов риска или предпочтений пациента.\src{6}

\heading{Рисунок 4. Концептуальный подход к оценке и уточнению ССПМ-риска: текстовая адаптация}
\textbf{Шаг 1. Исходное обследование и стадирование.} Оценить факторы ССПМ-риска --- показатели жироотложения, АД, липиды, гликемию и функцию почек --- и неблагоприятные SDOH. Определить исходную категорию:
\begin{itemize}
\item стадия 1 --- избыток или дисфункция жировой ткани;
\item стадия 2 или 3 --- метаболические факторы риска и/или ХБП;
\item стадия 4 --- клинически выраженное ССЗ.
\end{itemize}

\textbf{Шаг 2. Уточнение стадии и оценка усилителей риска.} У лиц 30--79 лет со стадиями 1--3 рассчитать PREVENT-CVD, чтобы количественно определить 10-летний риск ССЗ и уточнить стадию. У лиц с метаболическими факторами или ХБП 10-летний риск $<20\%$ соответствует ветви стадии 2, а $\geq20\%$ --- стадии 3; субклиническое ССЗ или ХБП очень высокого риска также определяют стадию 3. Для стадий 1--3 оценить усилители ССПМ-риска (таблица 9), чтобы обсудить с пациентом риск и выбрать план лечения. При стадии 4 с имеющимся ССЗ PREVENT не показан.

\textbf{Шаг 3. Предотвращение прогрессирования.}
\begin{itemize}
\item Стадии 1--2: изменение образа жизни и/или лекарственное лечение по показаниям для первичной профилактики ССЗ.
\item Стадия 3: изменение образа жизни и лекарственное лечение по показаниям для первичной профилактики ССЗ.
\item Стадия 4: интенсивное изменение образа жизни и лекарственное лечение для предотвращения повторных событий ССЗ и почечной недостаточности.
\end{itemize}

\textbf{Пояснение авторов к рисунку 4.} Систематический подход начинается с оценки факторов риска и определения стадии. Последующий расчёт PREVENT количественно характеризует риск и помогает различать стадию 2 при 10-летнем риске $<20\%$ и стадию 3 при риске $\geq20\%$ у лиц с метаболическими факторами или ХБП. Наличие субклинического ССЗ либо ХБП очень высокого риска также указывает на стадию 3. PREVENT не показан при стадии 4 с существующим ССЗ. На стадиях 1--3 усилители риска прогрессирования дополняют обсуждение соотношения пользы и риска профилактической терапии между пациентом и клиницистом. Схема объединяет исходную оценку ССПМ, прогноз риска по PREVENT-CVD и индивидуализацию с помощью усилителей риска; знак сноски в оригинале отсылает к таблице 9. CKM --- ССПМ; CKD --- ХБП; CV --- сердечно-сосудистый; CVD --- ССЗ; pt --- пациент; SDOH --- социальные детерминанты здоровья.

\heading{Таблица 9. Усилители риска при ССПМ-синдроме}
\heading{Личный медицинский анамнез}
\begin{itemize}
\item Хронические воспалительные или аутоиммунные заболевания: например ревматоидный артрит, системная красная волчанка, ВИЧ/СПИД.\src{4,11}
\item Нарушения сна: например обструктивное апноэ сна.\src{16}
\item Неблагополучие в сфере психического здоровья: например депрессия или тревога.\src{17}
\item Специфичные для пола факторы: например преждевременная менопауза в возрасте $<40$ лет, неблагоприятные исходы беременности, синдром поликистозных яичников, эректильная дисфункция.\src{2}
\end{itemize}
\heading{Социальные детерминанты здоровья\src{1}}
Более низкий социально-экономический статус; большее бремя неблагоприятных индивидуальных и территориальных SDOH (раздел 3.2).
\heading{Воспаление\src{13}}
Повышенный высокочувствительный С-реактивный белок: $\geq2{,}0$~мг/л, если исследование выполнено.
\heading{Семейный анамнез\src{14}}
Диабет; почечная недостаточность.
\heading{Демографические характеристики\src{15}}
Группы повышенного риска: например лица, самостоятельно относящие себя по расовой или этнической принадлежности к выходцам из Южной Азии, коренным американцам или жителям островов Тихого океана\textsuperscript{*}.

\textbf{Примечания исходной таблицы 9.} Адаптировано с разрешения по работе Ndumele и соавт.\src{5} © 2023 American Heart Association, Inc. \textsuperscript{*} Механизмы увеличения риска могут включать социальные и/или генетические факторы, непропорционально воздействующие на людей определённого происхождения. CKM --- ССПМ; HIV/AIDS --- вирус иммунодефицита человека/синдром приобретённого иммунодефицита; SDOH --- социальные детерминанты здоровья; SLE --- системная красная волчанка.

\heading{Обоснование рекомендации}
\textbf{1. Индивидуализация профилактики.} Риск перехода к более поздним стадиям ССПМ значительно варьирует и связан с повышенным риском ССЗ, почечной недостаточности и смертности.\src{1--4} Поэтому учёт усилителей риска помогает пациенту и клиницисту совместно принимать решения о начале и интенсификации терапии для предотвращения прогрессирования синдрома и уменьшения риска ССЗ.\src{10}

Более интенсивная профилактика может быть полезна при конкретных заболеваниях, связанных с ухудшением ССПМ-патологии,\src{4,11} воздействии неблагоприятных SDOH,\src{1,12} повышенных биомаркерах воспаления, если они измерялись,\src{3,13} семейном анамнезе диабета или почечной недостаточности\src{14} либо принадлежности к демографической подгруппе с установленным повышенным риском прогрессирования (таблица 9).\src{15} Таким образом, усилители риска представляют собой доступные для оценки характеристики, позволяющие более полно оценить индивидуальный риск прогрессирования при их сочетании с исходным стадированием (раздел 3.1.1) и количественным прогнозом, например PREVENT (раздел 4.1).\src{6,7}


\section{Общие принципы помощи при ССПМ-синдроме}
\heading{Обзор}
Основные принципы ведения ССПМ-синдрома включают неосуждающий подход к лечению ожирения, комплексную коррекцию факторов риска, поддержку междисциплинарной организации помощи для уменьшения её фрагментации и включение мер по смягчению воздействия SDOH в обычное ведение.\src{1} Модель AHA Life's Essential 8 представляет набор показателей здорового образа жизни и оптимального контроля факторов риска, применимый на популяционном уровне и значимый на всех стадиях ССПМ-синдрома.\src{2--4}

Взаимосвязи метаболических факторов, ХБП и ССЗ важны для выбора лечения, в том числе методов с пользой для нескольких органов. Сочетание заболеваний неотъемлемо от ССПМ-синдрома и может служить ориентиром при принятии терапевтических решений для достижения максимальной клинической пользы (таблица 10). Междисциплинарные модели с выделенным координатором ССПМ-помощи могут содействовать практическому внедрению лечебных стратегий при разной клинической тяжести синдрома и в различных условиях оказания помощи.\src{5--9}

В регионах с ограниченными ресурсами модели помощи можно гибко адаптировать, учитывая различия в доступности узких специалистов и экономические условия.\src{5} Пациентам, у которых валидированные инструменты скрининга выявили больше неблагоприятных SDOH (таблицы 6 и 7), следует обеспечить доступ к ресурсам местного сообщества для социальной поддержки и оптимального ведения ССПМ-синдрома.\src{10}

\heading{Таблица 10. Лечебные подходы при сочетании компонентов ССПМ-синдрома}
\editorial{Исходная матрица преобразована в последовательные блоки. В каждом блоке сохранены три колонки оригинала --- ожирение, СД2 и ХБП, --- включая повторяющиеся сочетания. «Неприменимо» соответствует обозначению N/A; это не утверждение об отсутствии показаний к лечению заболевания.}

\heading{Базовое ведение основных факторов ССПМ-риска}
\begin{itemize}
\item \textbf{Ожирение:} изменение образа жизни; дополнительная лекарственная терапия ожирения либо метаболическая и бариатрическая хирургия по необходимости (раздел 5.4).
\item \textbf{СД2:} изменение образа жизни; контроль факторов риска; применение SGLT2i, терапии на основе GLP-1 или их сочетания у пациентов с повышенным риском событий ССЗ (раздел 5.5.1).
\item \textbf{ХБП:} RASi и SGLT2i как терапия первой линии при альбуминурической ХБП для снижения риска ССЗ и предотвращения потери функции почек (раздел 5.5.4).
\end{itemize}

\heading{Сопутствующее состояние: ожирение}
\begin{itemize}
\item \textbf{При основном факторе «ожирение»:} неприменимо.
\item \textbf{При СД2:} при ожирении II класса и выше отдавать приоритет терапии на основе GLP-1 для лечения СД2 и уменьшения риска ССЗ.
\item \textbf{При ХБП:} целенаправленное снижение массы тела уменьшает альбуминурию.
\end{itemize}

\heading{Сопутствующее состояние: СД2}
\begin{itemize}
\item \textbf{При ожирении:} при ожирении II класса и выше отдавать приоритет терапии на основе GLP-1 для лечения СД2 и уменьшения риска ССЗ.
\item \textbf{При основном факторе «СД2»:} неприменимо.
\item \textbf{При ХБП:} RASi и SGLT2i --- терапия первой линии для уменьшения сердечно-сосудистого и почечного риска; при диабетической болезни почек с альбуминурией, сохраняющейся на терапии первой линии, добавление агониста рецептора GLP-1 и/или финеренона.
\end{itemize}

\heading{Сопутствующее состояние: альбуминурическая ХБП}
\begin{itemize}
\item \textbf{При ожирении:} целенаправленное снижение массы тела уменьшает альбуминурию.
\item \textbf{При СД2:} RASi и SGLT2i --- терапия первой линии для уменьшения сердечно-сосудистого и почечного риска; при диабетической болезни почек с альбуминурией, сохраняющейся на терапии первой линии, добавление агониста рецептора GLP-1 и/или финеренона.
\item \textbf{При основном факторе «ХБП»:} неприменимо.
\end{itemize}

\heading{Сопутствующее состояние: MASLD (раздел 7.2)}
\begin{itemize}
\item \textbf{При ожирении:} рекомендуется снижение массы тела; при MASLD с фиброзом рекомендуется агонист рецептора GLP-1.
\item \textbf{При СД2:} снижение массы тела, если это уместно; при MASLD рекомендуется агонист рецептора GLP-1 ввиду высокого риска фиброза.
\item \textbf{При ХБП:} неприменимо.
\end{itemize}

\heading{Сопутствующее состояние: субклинический коронарный атеросклероз (раздел 6.2)}
\begin{itemize}
\item \textbf{При ожирении:} возможное более выраженное абсолютное снижение риска при применении показанной профилактической терапии.
\item \textbf{При СД2:} возможное более выраженное абсолютное снижение риска при применении показанной профилактической терапии.
\item \textbf{При ХБП:} возможное более выраженное абсолютное снижение риска при применении показанной профилактической терапии.
\end{itemize}

\heading{Сопутствующее состояние: доклиническая СН (раздел 5.6.2)}
\begin{itemize}
\item \textbf{При ожирении:} значительное снижение массы тела может уменьшать риск развития СН.
\item \textbf{При СД2:} отдавать приоритет SGLT2i как основному кардиопротективному сахароснижающему средству для профилактики СН.
\item \textbf{При ХБП:} применять SGLT2i для профилактики СН; при диабетической болезни почек наличие pre-HF служит доводом в пользу рассмотрения добавления агониста рецептора GLP-1 или финеренона.
\end{itemize}

\heading{Сопутствующее состояние: АССЗ (раздел 6.2)}
\begin{itemize}
\item \textbf{При ожирении:} комплексная помощь команды по контролю массы тела для её снижения и уменьшения ССПМ-риска; терапия на основе GLP-1 для уменьшения MACE и сердечно-сосудистой смерти.
\item \textbf{При СД2:} изменение образа жизни, контроль факторов риска и применение SGLT2i, терапии на основе GLP-1 либо их сочетания при АССЗ для уменьшения MACE и сердечно-сосудистой смерти.
\item \textbf{При ХБП:} применение SGLT2i для уменьшения MACE, сердечно-сосудистой смерти и потери функции почек.
\end{itemize}

\heading{Сопутствующее состояние: HFpEF (раздел 6.3)}
\begin{itemize}
\item \textbf{При ожирении:} терапия на основе GLP-1 рекомендуется для уменьшения событий СН; рассмотреть физические тренировки и питание с дефицитом калорий.
\item \textbf{При СД2:} SGLT2i должны быть кардиопротективной сахароснижающей терапией первой линии для уменьшения событий СН; дополнительно рассмотреть терапию на основе GLP-1 при наличии диабета и ожирения.
\item \textbf{При ХБП:} SGLT2i в составе стандартной GDMT СН уменьшают события СН и потерю функции почек; дополнительно рассмотреть финеренон при диабетической болезни почек.
\end{itemize}

\heading{Сопутствующее состояние: HFrEF (раздел 6.3)}
\begin{itemize}
\item \textbf{При ожирении:} соотношение риска и пользы терапии на основе GLP-1 при HFrEF неясно, особенно при отсутствии АССЗ.
\item \textbf{При СД2:} SGLT2i должны быть кардиопротективной сахароснижающей терапией первой линии для уменьшения событий СН и сердечно-сосудистой смерти.
\item \textbf{При ХБП:} SGLT2i и ARNI/RASi в составе стандартной GDMT СН уменьшают события СН, сердечно-сосудистую смерть и потерю функции почек.
\end{itemize}

\textbf{Пояснение авторов к таблице 10.} Таблица описывает особенности ведения при сочетании основных факторов ССПМ-риска --- ожирения, диабета и ХБП --- с другими сопутствующими состояниями. Верхняя строка исходной таблицы отражает базовое ведение основных факторов. Последующие строки показывают, как сочетание соответствующего фактора с другим заболеванием должно влиять на клиническую тактику. Подробные рекомендации по ведению ожирения, СД2 и ХБП при наличии и отсутствии ССЗ приведены в последующих разделах.

ARNI --- ингибитор рецепторов ангиотензина и неприлизина; ASCVD --- АССЗ; CKD --- ХБП; CKM --- ССПМ; CV --- сердечно-сосудистый; CVD --- ССЗ; DKD --- диабетическая болезнь почек; GDMT --- терапия в соответствии с рекомендациями; GLP-1 --- глюкагоноподобный пептид-1; GLP-1 RA --- агонист рецептора GLP-1; HF --- СН; HFrEF --- СН со сниженной фракцией выброса; HFpEF --- СН с сохранённой фракцией выброса; MACE --- большие неблагоприятные сердечно-сосудистые события; MASLD --- стеатотическая болезнь печени, ассоциированная с метаболической дисфункцией; N/A --- неприменимо; RASi --- ингибиторы ренин-ангиотензиновой системы; SGLT2i --- ингибиторы натрий-глюкозного котранспортера 2 типа; T2D --- СД2.

\subsection{Междисциплинарная помощь}
\heading{Рекомендации по междисциплинарной помощи}
Исследования, обосновывающие рекомендации, обобщены в таблице доказательств оригинала.
{\small
\begin{longtable}{P{0.07\textwidth}P{0.11\textwidth}P{0.75\textwidth}}
\toprule COR & LOE & Рекомендация\\\midrule
\endfirsthead
\toprule COR & LOE & Рекомендация (продолжение)\\\midrule
\endhead
1 & B-R & \textbf{1.} Взрослым со стадиями ССПМ 2--4 и сочетанием как минимум двух из следующих состояний: диабет, ХБП и ССЗ, --- рекомендуется помощь междисциплинарной команды с выделенным координатором ССПМ-помощи. Это способствует ведению межсистемных нарушений, включая вмешательства по изменению образа жизни и оптимизацию GDMT.\src{1--5}\\[5pt]
1 & C-EO & \textbf{2.} При ССПМ-синдроме у взрослых следует уделять приоритетное внимание уменьшению клинического воздействия неблагоприятных SDOH для оптимизации комплексной помощи.\\
\bottomrule
\end{longtable}
}

\heading{Обзор}
У пациентов с ССПМ-синдромом часто имеются несколько заболеваний, требующих участия разных клиницистов, что может приводить к фрагментации помощи.\src{6} Включение клинических подходов нескольких дисциплин в единую ССПМ-модель и координация действий членов команды --- обязательные компоненты комплексной помощи. В президентском консультативном документе AHA о ССПМ-здоровье описаны два междисциплинарных подхода: 1) модель помощи, ориентированная на ценность для пациента (value-based care), и 2) модель, ориентированная на объём оказываемой помощи (volume-based care; рисунок 5).\src{7}

В модели value-based care междисциплинарная команда, включающая координатора ССПМ-помощи и представителей нескольких специальностей, разрабатывает протоколы ведения пациентов с сочетанием как минимум двух компонентов --- диабета, ХБП и ССЗ. Координатор обеспечивает междисциплинарную коммуникацию (таблицы 11 и 12). В модели volume-based care пациентов высокого риска направляют к соответствующим узким специалистам, а координатор помогает им ориентироваться в помощи нескольких клиницистов.

Координатор также поддерживает другие важные составляющие оптимального ведения: комплексное управление лекарственной терапией, содействие изменению образа жизни и координацию помощи. В обеих моделях помощь может оказываться дистанционно или очно. Воздействие неблагоприятных SDOH, выявленных при систематической оценке (раздел 3.2), следует смягчать путём включения в команду работников общественного здравоохранения, социальных работников, специалистов по сопровождению случая или навигаторов пациента, а также привлечения ресурсов местного сообщества.\src{6,7}

\heading{Рисунок 5. Применение междисциплинарных моделей при ССПМ-синдроме: текстовая адаптация}
По горизонтальной оси исходной схемы показан спектр мультиморбидности при ССПМ-синдроме; по вертикальной --- абсолютный риск неблагоприятных событий ССЗ. По мере увеличения числа и тяжести сочетанных заболеваний повышаются риск и потребность в специализированной координации.
\begin{enumerate}[label=\arabic*.]
\item \textbf{Меньшее бремя факторов ССПМ-риска.} Поддержка изменения образа жизни и GDMT, направленной на факторы риска.
\item \textbf{Сочетание как минимум двух состояний: СД2, ХБП, ССЗ.} Модель value-based care с поддержкой координатора ССПМ-помощи и междисциплинарной команды.
\item \textbf{Сочетание как минимум двух состояний: СД2 высокого риска, ХБП очень высокого риска, ССЗ.} Модель volume-based care с направлением к узким специалистам и навигационной поддержкой координатора.
\end{enumerate}
\textbf{Определения в сноске рисунка.} СД2 высокого риска: HbA1c $\geq9\%$ либо микрососудистые осложнения СД2. ХБП очень высокого риска определяется по карте KDIGO (рисунок 2).

\textbf{Пояснение авторов к рисунку 5.} Концептуальная модель адаптирует междисциплинарную помощь к профилю мультиморбидности и соответствующему абсолютному риску ССЗ. У большинства пациентов с меньшим бременем факторов риска клиницисты должны поддерживать и поощрять изменение образа жизни и начало рекомендованной терапии с учётом профиля ССПМ-риска и расчётного риска ССЗ. Меньшая группа с сочетанием как минимум двух из состояний СД2, ХБП и ССЗ имеет более высокий абсолютный риск; для неё рекомендуется value-based care с координатором и виртуальной междисциплинарной командой. Ещё меньшая группа с сочетанием как минимум двух из состояний СД2 высокого риска, ХБП очень высокого риска или клинически выраженного ССЗ имеет наибольший абсолютный риск. Для неё рекомендуется междисциплинарная стратегия volume-based care с направлением к узким специалистам и навигационной поддержкой координатора. CKD --- ХБП; CKM --- ССПМ; CVD --- ССЗ; GDMT --- терапия в соответствии с рекомендациями; T2D --- СД2.

\heading{Обоснование отдельных рекомендаций}
\textbf{1. Междисциплинарная команда и координатор.} Такие команды крайне важны для воздействия на взаимозависимые патологические процессы и многоуровневые неблагоприятные SDOH, участвующие в развитии синдрома.\src{6} Несколько РКИ показали, что междисциплинарная работа при сочетании компонентов ССПМ улучшает внедрение терапии, изменение образа жизни или оба направления.\src{1--4,8--10}

Длительное наблюдение в исследовании STENO-2 (Intensified Multifactorial Intervention in Patients With T2D and Microalbuminuria --- интенсивное многофакторное вмешательство при СД2 и микроальбуминурии) показало, что многофакторное вмешательство команды из врача, медицинской сестры и диетолога, направленное на контроль факторов риска посредством изменения образа жизни и целевой лекарственной терапии, улучшало исход по общей смертности у пациентов с СД2, избыточной массой тела/ожирением и альбуминурией.\src{4}

В исследовании COORDINATE-Diabetes (Coordinating Cardiology Clinics Randomized Trial of Interventions to Improve Outcomes-Diabetes --- рандомизированное исследование координации кардиологических клиник для улучшения исходов при диабете) многокомпонентное вмешательство с акцентом на сотрудничество кардиологов и специалистов по диабету улучшало назначение доказательной терапии пациентам с СД2 и АССЗ.\src{1}

Метаанализ 19 РКИ показал, что команды с представителями как минимум трёх медицинских дисциплин существенно уменьшали сердечно-сосудистые факторы риска у пациентов с диабетом в первичной помощи.\src{2} Несколько наблюдательных клинических исследований также поддерживают междисциплинарный подход при ССПМ-синдроме.\src{11--13} Рандомизированные данные обосновывают участие фармацевтов,\src{3,8,14--16} медицинских сестёр,\src{15} навигаторов без профессиональной лицензии\src{3} и работников общественного здравоохранения\src{5,15} в командной помощи, начале доказательной терапии и последующем наблюдении. Ряд исследований указывает на экономическую эффективность междисциплинарной помощи в долгосрочной перспективе.\src{17--19}

\textbf{2. Смягчение воздействия неблагоприятных SDOH.} Эти факторы повышают риск ССЗ и смерти от всех причин при ССПМ-синдроме.\src{20} Работа с ними, вероятно, улучшает качество комплексной помощи ввиду их многоуровневого влияния на осложнения. Она может улучшать клинические исходы, однако убедительных данных, подтверждающих эту гипотезу, пока недостаточно.

В небольшом рандомизированном перекрёстном исследовании 44 взрослых с диабетом и недостаточной обеспеченностью продовольствием доставка питания, адаптированного к медицинским потребностям, улучшала качество рациона и обеспеченность продовольствием и уменьшала гипогликемию.\src{21} Другое, нерандомизированное проспективное исследование показало связь скрининга неблагоприятных SDOH и попыток их коррекции с небольшим улучшением АД и липидных показателей, но не HbA1c.\src{22}

Вмешательство работника общественного здравоохранения у госпитализированных пациентов с низким социально-экономическим статусом увеличивало доступность первичной помощи после выписки и уменьшало повторные госпитализации в течение 30 дней.\src{23} Ещё одно небольшое РКИ показало, что координация помощи под руководством социального работника уменьшала частоту 30-дневной повторной госпитализации на 22\% у пожилых стационарных пациентов с высоким риском повторного поступления.\src{24}

Поддержка, направленная на неблагоприятные SDOH, вероятно, полезна и сопряжена с минимальным риском, но настоятельно необходимы дальнейшие исследования для подтверждения клинической пользы и экономической эффективности и определения оптимальных вмешательств.

\heading{Таблица 11. Междисциплинарная команда ССПМ-помощи}
{\small
\begin{longtable}{P{0.28\textwidth}P{0.65\textwidth}}
\toprule Категория помощи & Участники команды\\\midrule
\endfirsthead
\toprule Категория помощи & Участники команды\\\midrule
\endhead
Общемедицинская помощь & Специалист первичной помощи\textsuperscript{*}; координатор ССПМ-помощи\textsuperscript{\dag}; клинический фармацевт; медицинская сестра; диетолог\\
Специализированная помощь & Специалисты по кардиологии\textsuperscript{*}, эндокринологии\textsuperscript{*}, нефрологии\textsuperscript{*} и гепатологии\textsuperscript{*}; специалист по обучению пациентов с диабетом\\
Социальная поддержка & Социальный работник; работник общественного здравоохранения\\
\bottomrule
\end{longtable}
}
\textsuperscript{*} Это могут быть врачи со степенью MD или DO либо медицинские специалисты расширенной практики. \textsuperscript{\dag} Роль координатора ССПМ-помощи могут также выполнять специалисты определённого профиля, например по обучению пациентов с диабетом или координации помощи при СН. CKM --- ССПМ.

\heading{Таблица 12. Координированные услуги для оптимального ведения ССПМ-синдрома}
\heading{Комплексное управление лекарственной терапией\textsuperscript{*}}
\begin{itemize}
\item Обучение пациента компонентам ССПМ-синдрома и соответствующим препаратам для расширения его возможностей участвовать в лечении.
\item Сверка списка принимаемых лекарственных средств.
\item Начало, титрация и мониторинг доказательной ССПМ-терапии.
\item Работа с побочными эффектами и необоснованным прекращением приёма препаратов.
\item Совместное принятие решений при отмене лекарств, к которым нет показаний.
\item Оценка приверженности и консультирование для её повышения.
\item Содействие доступу к терапии, включая программы помощи пациентам.
\end{itemize}
\heading{Содействие изменению образа жизни}
Обучение и сопровождение изменений питания и физической активности, а также отказа от табака.
\heading{Координация помощи}
\begin{itemize}
\item Обеспечение коммуникации между клиницистами пациента.
\item Информирование всей команды о трудностях начала доказательной терапии или соблюдения назначений.
\item Обеспечение понимания пациентом сложных планов лечения.
\end{itemize}
\heading{Навигационная поддержка пациента}
\begin{itemize}
\item Обеспечение понимания роли каждого участника команды и полного плана помощи.
\item Содействие коммуникации пациента с клиницистами.
\item Информирование пациента о контрольных исследованиях и приёмах.
\item Содействие доступу к программам социальной помощи.
\end{itemize}
\textsuperscript{*} Комплексное управление лекарственной терапией --- стандарт практики, предусматривающий индивидуальную оценку всех принимаемых средств: рецептурных и безрецептурных лекарств, альтернативных лекарственных средств, витаминов и пищевых добавок. Для каждого средства проверяют наличие показаний и обоснованность применения, эффективность при заболеваниях пациента, безопасность с учётом сопутствующей патологии и других препаратов, а также возможность приёма пациентом в соответствии с индивидуальным планом. CKM --- ССПМ.

\subsection{Стадия ССПМ 0}
\heading{Обзор}
Стадия 0 соответствует отсутствию ССПМ-синдрома вследствие отсутствия факторов риска: нормальные ИМТ и окружность талии, нормогликемия, нормальное АД, нормальный липидный профиль и отсутствие ХБП, субклинического либо клинически выраженного ССЗ.\src{1} Поскольку факторы риска чаще возникают с возрастом, стадия 0 наиболее распространена, хотя и не исключительно, среди детей, подростков и молодых взрослых. Мужчины реже, чем женщины, не имеют факторов ССПМ-риска; аналогичное различие наблюдается у чернокожих взрослых по сравнению с белыми.\src{2}

К сожалению, 90--95\% взрослых в США находятся на стадиях ССПМ 1--4. Это подчёркивает необходимость эффективнее поддерживать идеальное сердечно-сосудистое здоровье с раннего возраста.\src{2--4} Продольные данные показывают, что большинство людей среднего возраста со стадией 0 за 10 лет переходят к более поздней стадии, что связано с более высоким риском ССЗ и смертности.\src{5}

Поэтому основное направление на стадии 0 --- примордиальная профилактика, то есть предотвращение появления факторов ССПМ-риска. Особое внимание уделяется предотвращению избытка и дисфункции жировой ткани: возникновение метаболических факторов от детского и молодого взрослого возраста до среднего возраста преимущественно связано с возрастной прибавкой массы тела.\src{6} Рекомендации по примордиальной профилактике обсуждаются в руководстве ACC/AHA 2019 года по первичной профилактике.\src{7}

\subsection{Профилактика ССЗ при ССПМ-синдроме}
\heading{Обзор}
На стадиях ССПМ 1--3 одна из главных задач --- предотвращение перехода к более поздним стадиям и клинически выраженному ССЗ. Это также благоприятно влияет на ХБП, MASLD и другие полиорганные проявления синдрома.\src{1} Необходимо применять общие принципы ведения (раздел 5) и корректировать ключевые факторы --- ожирение, СД2 и ХБП, --- чтобы уменьшить вероятность событий ССЗ и почечной недостаточности (рисунок 6). Кроме того, АД и липиды следует контролировать и корректировать согласно действующим клиническим рекомендациям.\src{2--6} Последующие разделы 5.4 «Стадия ССПМ 1», 5.5 «Стадия ССПМ 2» и 5.6 «Стадия ССПМ 3» подробно описывают лечение пациентов на соответствующих стадиях.

\heading{Рисунок 6. Принципы ведения стадий ССПМ 1--3 при ожирении, СД2 или ХБП: текстовая адаптация}
\heading{Применять общие принципы ССПМ-помощи}
\begin{itemize}
\item Укреплять сердечно-сосудистое здоровье с помощью модели Life's Essential 8.
\item При сочетании компонентов --- СД2, ХБП, выраженного субклинического ССЗ --- организовать междисциплинарную помощь, включая выделенного координатора.
\item Работать с неблагоприятными SDOH с участием работников общественного здравоохранения и социальных работников, привлекая ресурсы местного сообщества.
\item Контролировать АД и липиды согласно действующим рекомендациям.
\end{itemize}

\heading{Выбирать ведение с учётом профиля риска пациента: ожирение}
\begin{itemize}
\item Консультировать по снижению массы тела без осуждения.
\item Использовать изменения образа жизни и поведенческие вмешательства для снижения массы тела не менее чем на 5--10\% и уменьшения риска прогрессирования ССПМ-синдрома.
\item Если целевое снижение массы тела не достигнуто, рассмотреть терапию на основе GLP-1 либо метаболическую и бариатрическую хирургию для уменьшения массы тела и риска прогрессирования синдрома.
\end{itemize}

\heading{Выбирать ведение с учётом профиля риска пациента: СД2}
\begin{itemize}
\item Изменение образа жизни и контроль массы тела для улучшения гликемического контроля.
\item Комплексная коррекция факторов риска посредством изменения образа жизни и целевой лекарственной терапии.
\item При повышенном риске ССЗ --- 10-летний PREVENT-CVD $\geq7{,}5\%$ либо возраст $>50$ лет с другими факторами риска --- начать терапию на основе GLP-1 или SGLT2i для уменьшения событий ССЗ и сердечно-сосудистой смертности.
\item Если HbA1c на 0,5--1 процентного пункта выше целевого, применять метформин совместно с SGLT2i или терапией на основе GLP-1.
\end{itemize}

\heading{Выбирать ведение с учётом профиля риска пациента: ХБП}
\begin{itemize}
\item \textbf{При СД2 или альбуминурии:} использовать RASi и SGLT2i как терапию первой линии для улучшения сердечно-сосудистых и почечных исходов.
\item \textbf{При СД2 и альбуминурии, сохраняющейся несмотря на терапию первой линии:} рассмотреть добавление терапии на основе GLP-1 или финеренона для улучшения сердечно-сосудистых и почечных исходов.
\end{itemize}

\textbf{Пояснение авторов к рисунку 6.} Ведение пациентов с ожирением, диабетом и/или ХБП на стадиях ССПМ 1--3 должно сочетать общие принципы помощи с дополнительными лечебными подходами к отдельным факторам риска, показанными на схеме. Если HbA1c существенно превышает целевое значение --- на $\geq0{,}5$--1 процентного пункта, --- метформин можно назначать одновременно с SGLT2i, что особенно важно с учётом ограниченного сахароснижающего действия этих препаратов, либо с препаратом на основе GLP-1. Альтернативно метформин можно добавить позднее в зависимости от предпочтений пациента.

CHW --- работник общественного здравоохранения; CKD --- ХБП; CKM --- ССПМ; CV --- сердечно-сосудистый; CVD --- ССЗ; DM --- диабет; GLP-1 --- глюкагоноподобный пептид-1; RASi --- ингибиторы ренин-ангиотензиновой системы; SDOH --- социальные детерминанты здоровья; SGLT2i --- ингибиторы натрий-глюкозного котранспортера 2 типа; T2D --- СД2.


\subsection{Стадия ССПМ 1}
\heading{Обзор}
Стадия ССПМ 1 определяется как избыточная масса тела/ожирение, абдоминальное ожирение или дисфункция жировой ткани при отсутствии других метаболических факторов риска, ХБП и субклинического либо клинически выраженного ССЗ.\src{1} Совместное измерение ИМТ и окружности талии улучшает характеристику состава тела и метаболического риска. Повышение обоих показателей (таблица 4) выделяет группу, которой следует уделять приоритетное внимание при снижении массы тела. Другие методы --- отношение окружности талии к окружности бёдер, биоимпедансный анализ и двухэнергетическая рентгеновская абсорбциометрия (DEXA) --- также лучше характеризуют состав тела, чем один ИМТ; однако акцент сделан на окружности талии ввиду доступности и простоты измерения.

Определение дисфункции жировой ткани включает предиабет, отражающий метаболические нарушения, способствующие инсулинорезистентности даже при отсутствии избытка жировой ткани. Цель ведения стадии 1 --- предотвратить появление дополнительных метаболических факторов посредством изменения образа жизни и снижения массы тела. Контроль массы тела должен быть многокомпонентным: консультирование по образу жизни с добавлением лекарственной терапии и/или метаболической и бариатрической хирургии по необходимости. Целевое снижение массы тела не менее чем на 5--10\% направлено на уменьшение ССПМ-риска; более выраженное снижение связано с большей пользой.

Лицам с предиабетом независимо от ИМТ следует уделять приоритетное внимание при изменении образа жизни, дополняя его снижением массы тела при наличии показаний. Если гипергликемия сохраняется и прогрессирует несмотря на эти меры, для предотвращения диабета можно рассмотреть лекарственную терапию, например метформин.\src{2}

\subsubsection{Общий подход к ведению пациентов с ожирением}
\heading{Рекомендации по общему подходу к ведению ожирения}
Исследования, обосновывающие рекомендации, обобщены в таблице доказательств оригинала.
{\small
\begin{longtable}{P{0.07\textwidth}P{0.11\textwidth}P{0.75\textwidth}}
\toprule COR & LOE & Рекомендация\\\midrule
\endfirsthead
\toprule COR & LOE & Рекомендация (продолжение)\\\midrule
\endhead
1 & A & \textbf{1.} Взрослым с избыточной массой тела или ожирением рекомендуется изменение образа жизни как стратегия первой линии для снижения массы тела не менее чем на 5--10\% от исходной.\src{1--3}\\[5pt]
1 & B-NR & \textbf{2.} При консультировании взрослых с избыточной массой тела или ожирением рекомендуется начинать обсуждение снижения массы тела без осуждения, чтобы повысить результативность снижения веса.\src{4}\\[5pt]
2a & B-NR & \textbf{3.} У взрослых с ожирением интегрированные междисциплинарные команды по контролю массы тела могут быть эффективны для усиления ориентированного на пациента подхода к снижению веса.\src{2}\\
\bottomrule
\end{longtable}
}

\heading{Обзор}
Контроль массы тела --- основа помощи при ССПМ-синдроме. Поддержание здоровой массы связано с меньшим риском прогрессирования, а выраженное снижение веса способствует уменьшению риска ССЗ и обратному развитию синдрома. Ожирение --- сложное хроническое состояние с многофакторной этиологией, возникающее при взаимодействии социальных, поведенческих и биологических факторов.\src{5} После формирования ожирения намеренное снижение массы тела запускает множество компенсаторных механизмов, влияющих на обмен веществ, голод и насыщение; они часто приводят к повторному набору веса.\src{6} Поэтому для оптимального контроля массы тела необходимо длительное взаимодействие с пациентом.

Для стадий ССПМ 1--3 появляется всё больше эффективных возможностей снижения массы тела. Первая линия --- изменение образа жизни через изменение поведения; дополнительно рассматривают всё более эффективную лекарственную терапию ожирения и метаболическую/бариатрическую хирургию (рисунок 7). Единого оптимального подхода для всех пациентов на любой стадии не существует. Выбор следует индивидуализировать с учётом выраженности ожирения, стадии ССПМ, доступных ресурсов и предпочтений пациента. Стратегии, помогающие ориентироваться среди лечебных возможностей, например интегрированные многопрофильные команды, важны для помощи, ориентированной на пациента.

\heading{Рисунок 7. Ведение ожирения при стадиях ССПМ 1--3: текстовая адаптация алгоритма}
\textbf{Исходная группа:} пациенты с ССПМ-синдромом и ожирением без ССЗ.
\begin{enumerate}[label=\arabic*.]
\item Консультирование по снижению массы тела без осуждения --- \textbf{COR 1}.
\item Поддержка изменения поведения и образа жизни --- \textbf{COR 1}.
\item Оценка достаточности снижения массы тела: достигнуто ли снижение не менее чем на 5--10\%?
\begin{itemize}
\item \textbf{Да:} постоянная поддержка изменения образа жизни для длительного удержания достигнутого результата --- \textbf{COR 1}.
\item \textbf{Нет:} терапия на основе GLP-1 --- \textbf{COR 2a} --- либо метаболическая и бариатрическая хирургия --- \textbf{COR 2a}.
\end{itemize}
\item После хирургического лечения --- постоянная поддержка изменения образа жизни для длительного удержания снижения массы тела, \textbf{COR 1}.
\item При терапии на основе GLP-1 повторно оценить достаточность снижения массы тела (не менее 5--10\%). Если результат достаточен --- продолжать поддержку образа жизни для удержания результата, \textbf{COR 1}. Если недостаточен --- обеспечить повышение дозы, сменить препарат на основе GLP-1 и/или направить к специалисту по контролю массы тела, \textbf{COR 1}.
\end{enumerate}
\textbf{Условные обозначения оригинала:} COR 1, 2a, 2b, 3 «отсутствие пользы» и 3 «вред» соответствуют классам рекомендации из таблицы 2; в данном алгоритме задействованы классы 1 и 2a. CKM --- ССПМ; GLP-1 --- глюкагоноподобный пептид-1.

\heading{Обоснование отдельных рекомендаций}
\textbf{1. Изменение образа жизни как первая линия.} Избыточная масса тела и ожирение --- ведущие проблемы общественного здравоохранения, затрагивающие более 70\% взрослых в США.\src{7} Изменение образа жизни с поведенческой поддержкой для длительного изменения питания и физической активности остаётся терапией первой линии. Оно обеспечивает не только умеренное снижение массы тела на 5--10\%, но и благоприятно влияет на несколько метаболических факторов независимо от снижения веса.\src{8--11}

Кроме того, изменение образа жизни воздействует на патофизиологические процессы, связанные с избытком жировой ткани и инсулинорезистентностью и увеличивающие риск ССЗ: системное воспаление, эндотелиальную дисфункцию, атерогенную дислипидемию и гиперкоагуляцию. Одной лекарственной коррекции сопутствующих ожирению заболеваний может быть недостаточно для воздействия на эти процессы. Однако хронический характер ожирения и компенсаторные изменения обмена веществ, голода и насыщения способствуют возврату веса.\src{6} Поэтому для оптимального результата необходимы частое и длительное взаимодействие с пациентом и поведенческая поддержка.

\textbf{2. Обсуждение без осуждения.} Консультирование по снижению массы тела --- важный первый шаг в поддержке пациента. Но из-за широко распространённых стигмы и предубеждений в отношении веса клинические обсуждения часто имеют осуждающий характер. Они чрезмерно акцентируют личную ответственность по сравнению со сложными социальными и биологическими факторами, участвующими в развитии ожирения и возврате веса после его снижения.\src{12} Столкновение со стигматизацией по поводу веса может независимо ассоциироваться с худшими клиническими исходами.\src{13}

Исследования показывают, что после неосуждающего обсуждения лица с ожирением чаще успешно достигают снижения массы тела более чем на 10\%, чем после разговора, воспринимаемого как осуждающий.\src{14} Кроме того, столкновение с предубеждениями снижает вовлечённость в помощь по поводу ожирения. STOP Obesity Alliance предложил инструментарий для эффективного и неосуждающего обсуждения, исходящий из представления об ожирении как хроническом клиническом состоянии с многофакторной этиологией, существенно влияющем на здоровье и качество жизни.\src{15} Инструментарий \textit{Weight Can't Wait --- Guide for the Management of Obesity in the Primary Care Setting} также включает стратегии перехода от обсуждения к действиям, поддерживающим усилия пациента.\src{15}

\textbf{3. Междисциплинарная помощь.} За последние годы возможности лечения ожирения существенно расширились.\src{16} Помимо специалистов по лечению ожирения, команда может включать специалистов по питанию, кинезиологов, фармацевтов и специалистов по психическому здоровью; такая поддержка особенно эффективна для изменения образа жизни и снижения массы тела.\src{17,18} Лекарственная терапия, особенно на основе GLP-1, эффективна как дополнение, а междисциплинарная работа помогает обеспечить доступность и надлежащее применение препаратов.\src{19}

Метаболическая и бариатрическая хирургия --- наиболее мощный метод снижения массы тела --- стала безопаснее благодаря менее инвазивным вариантам операций.\src{20,21} Включение соответствующих специалистов в команды по контролю массы тела может улучшить доступ к хирургическому лечению. Участие многопрофильной команды также рекомендовано хирургическими руководствами для оптимизации до- и послеоперационной помощи.\src{22}

Ответ на разные методы снижения веса варьирует; выбор дополнительно зависит от массы тела, предыдущего опыта снижения веса, профиля ССПМ-риска, доступности терапии и предпочтений пациента. Имеются данные, что интегрированные команды, объединяющие специалистов по лечению ожирения и бариатрической хирургии с другими участниками, поддерживающими снижение веса и оказывающими социальную помощь, улучшают ориентированное на пациента ведение, помогая ориентироваться в расширяющемся выборе методов.\src{23,24}

\subsubsection{Интенсивное изменение образа жизни для снижения массы тела}
\heading{Рекомендации по интенсивному изменению образа жизни}
Исследования, обосновывающие рекомендации, обобщены в таблице доказательств оригинала.
{\small
\begin{longtable}{P{0.07\textwidth}P{0.11\textwidth}P{0.75\textwidth}}
\toprule COR & LOE & Рекомендация\\\midrule
\endfirsthead
\toprule COR & LOE & Рекомендация (продолжение)\\\midrule
\endhead
1 & A & \textbf{1.} Взрослых с избыточной массой тела или ожирением клиницисты должны консультировать не реже одного раза в год о пользе снижения массы тела не менее чем на 5--10\% для уменьшения ССПМ-риска.\src{1,2}\\[5pt]
1 & A & \textbf{2.} Направление взрослых с избыточной массой тела или ожирением на очные либо дистанционные многокомпонентные поведенческие вмешательства полезно для снижения массы тела и ССПМ-риска.\src{1,3,4}\\[5pt]
2a & B-R & \textbf{3.} При существенных препятствиях для участия в многокомпонентных поведенческих программах самостоятельное использование цифровых решений для здоровья может быть полезно взрослым с избыточной массой тела или ожирением для снижения веса и ССПМ-риска.\src{5--7}\\[5pt]
1 & A & \textbf{4.} Взрослым с избыточной массой тела или ожирением следует рекомендовать изменение образа жизни с умеренным ограничением калорийности и регулярной физической активностью для снижения веса и сердечно-сосудистого риска.\src{8}\\[5pt]
1 & B-NR & \textbf{5.} Взрослым, успешно снизившим массу тела, следует рекомендовать долгосрочные программы её контроля, включающие акцент на повышении физической активности, чтобы удерживать достигнутое снижение веса.\src{4}\\
\bottomrule
\end{longtable}
}

\heading{Обзор}
AHA/ACC и несколько профессиональных обществ рекомендуют интенсивные многокомпонентные поведенческие вмешательства под руководством клиницистов, умеющих эффективно начинать консультирование по снижению веса.\src{9--11} Их компоненты должны включать изменение питания для дефицита 500--750~ккал/сут, не менее 150 минут аэробной активности умеренной или высокой интенсивности в неделю и не менее 14 месяцев поведенческой терапии. Последняя включает самоконтроль питания, физической активности и массы тела, а также решение возникающих проблем под руководством специалиста.\src{9,10} Такие вмешательства обеспечивают снижение массы тела в среднем на 8\%\src{10} --- величину, для которой показано улучшение метаболического риска.\src{12}

\editorial{В оригинале на странице e77, включая изображение страницы, указано именно «не менее 14 месяцев» поведенческой терапии. Эта длительность сохранена без замены числом занятий или другим сроком.}

Однако по ряду причин эти результаты могут завышать ожидаемое снижение веса в обычной практике. В исследования обычно включают людей с большей мотивацией, чем у большинства пациентов. Удержание сниженного веса часто затруднено.\src{13} Доступность интенсивных программ ограничена, особенно в сельских и неблагополучных городских районах; стоимость и страховое покрытие нередко препятствуют участию.\src{14} Эффективные программы часто реализуются многопрофильными командами, а не только врачами.\src{15}

Успех программ следует оценивать не только по изменению веса, но и по трём критериям: доступность, достаточно низкая стоимость для эффективного предоставления и возможность длительного участия, поскольку ожирение является хроническим состоянием. Технологические решения потенциально могут соответствовать этим критериям.\src{16}

\heading{Обоснование отдельных рекомендаций}
\textbf{1. Консультирование и значение снижения веса.} Консультирование врачом рекомендуется для вовлечения пациента в снижение массы тела. Структурированное консультирование может повышать мотивацию и активное участие.\src{17} Систематический обзор мотивационного интервьюирования в первичной помощи показал, что некоторые вмешательства приводили к умеренному снижению веса на 5\%.\src{2} При обсуждении следует подчёркивать пользу для ССПМ-риска. Приоритетная группа --- лица с одновременным повышением ИМТ и окружности талии.

Умеренное снижение массы тела на 5--10\% у пациентов с избыточной массой тела/ожирением и диабетом связано, среди прочего, с улучшением гликемии, АД, ХС ЛВП и триглицеридов. Это может уменьшать риск впервые возникающих факторов ССПМ-риска, таких как диабет и гипертензия.\src{18} Намеренное снижение веса улучшает контроль уже существующих факторов: АД, гликемии, дислипидемии, альбуминурии и системного воспаления.\src{19} Большее снижение веса связано с более выраженным улучшением профиля ССПМ-риска,\src{20} а при снижении на $\geq10\%$ возможно и некоторое уменьшение риска ССЗ.\src{21}

\textbf{2. Многокомпонентные вмешательства.} Такие программы сочетают групповые занятия, руководство клинициста, конкретные изменения питания и поддержку физической активности.\src{1} Наиболее эффективны они при поведенческой поддержке, обеспечивающей длительную приверженность изменениям образа жизни.\src{1} Рандомизированные исследования показали значимое снижение массы тела с улучшением метаболических факторов и функции почек. Важно, что очный и дистанционный форматы продемонстрировали сходную эффективность.\src{22}

Систематический обзор 2024 года показал: поведенческие вмешательства разной длительности, независимо от конкретного сочетания поведенческих, пищевых и двигательных компонентов, связаны со снижением массы тела приблизительно на 5\% за примерно 12 месяцев у пациентов с избыточной массой тела/ожирением и СД2.\src{23} Успешные программы предусматривали частую обратную связь и поддержку, благоприятно влияющие на приверженность.\src{24} В нескольких исследованиях сопутствующие ССЗ и ОАС не препятствовали успешному снижению веса.\src{9} Несмотря на убедительные доказательства снижения веса и улучшения метаболического и почечного риска, расширение доступности программ в регионах с ограниченными ресурсами и для людей с низкой медицинской грамотностью либо другими препятствиями остаётся сложной задачей.

\textbf{3. Самостоятельное использование цифровых решений.} Цифровые вмешательства включают разнообразные приложения для ограничения калорийности, физической активности, самоконтроля, обучения и мотивации. Преимущественно это приложения для смартфонов; некоторые программы используют носимые устройства. Большинство разработано частными компаниями, но многие недороги либо имеют бесплатные ограниченные версии. Лишь немногие тщательно изучены или сопоставлены с альтернативами. Тем не менее по этой теме опубликовано много систематических обзоров с качественным анализом и метаанализов.\src{25--29}

Кокрейновский обзор 2024 года показал, что отдельные приложения связаны с умеренным краткосрочным снижением массы тела: $-2{,}6$~кг через 6--8 месяцев, но через 12 месяцев относительная польза невелика. Значительная неоднородность содержания и подходов затрудняет выводы. Высокая частота прекращения использования также остаётся серьёзной проблемой.\src{30} И всё же для многих пациентов с избыточной массой тела/ожирением и ССПМ-синдромом недорогие цифровые решения могут быть единственной альтернативой.

Эта область развивается: недавно появились цифровые помощники по здоровью на основе искусственного интеллекта и аналогичные приложения. Хотя имеются обнадёживающие данные об исходах,\src{31} эффективность таких подходов ещё предстоит систематически изучить.

\textbf{4. Питание и физическая активность.} Умеренное ограничение калорийности создаёт устойчивый энергетический дефицит; для снижения массы тела доступны разные схемы питания. Систематический обзор и сетевой метаанализ 121 клинического исследования показали, что несколько вариантов, включая низкожировое и низкоуглеводное питание, умеренно влияют на вес --- приблизительно на 4~кг --- и АД.\src{8} К сожалению, в большинстве исследований через 12 месяцев преимущества для веса и факторов риска исчезают, вероятно, из-за снижения приверженности.

В научном заявлении описано соответствие популярных диет рекомендациям AHA по питанию. Особенно заметные кардиометаболические преимущества имеют DASH, средиземноморский, пескетарианский и вегетарианский типы питания, снижающие риск гипертензии, СД2 и ССЗ.\src{32} Напротив, кетогенные диеты с очень низким содержанием углеводов и высоким содержанием насыщенных жиров вызывают снижение веса, но приводят к выраженной гиперхолестеринемии с потенциальным повышением сердечно-сосудистого риска. Осведомлённость клинициста об оптимальных для кардиометаболического здоровья схемах питания помогает давать точные и эффективные рекомендации.

Питание с ограничением времени приёма пищи, наиболее популярной формой которого является интервальное голодание, продемонстрировало определённую эффективность для снижения веса.\src{33,34} Однако исследования показывают, что оно не эффективнее ограничения калорийности в отношении снижения массы тела или улучшения сердечно-сосудистого риска.\src{35,36} Одна физическая активность вряд ли приведёт к клинически значимому снижению веса, но её сочетание с изменением питания обеспечивает большее долгосрочное снижение, чем одно изменение питания.\src{37,38}

\textbf{5. Удержание достигнутого результата.} Повторный набор веса --- сложное явление, включающее поведенческие, физиологические и средовые факторы.\src{39} Регистры снижения веса дают важные сведения об успешном удержании результата. Систематический обзор регистровых данных выявил несколько полезных стратегий, включая регулярный завтрак и наличие здоровой пищи дома.\src{40} Наиболее последовательно с удержанием результата связано повышение физической активности; несколько исследований показали, что её более высокий уровень ассоциируется с большей вероятностью сохранения сниженного веса.\src{38,41}

В РКИ после 8-недельного низкокалорийного вмешательства, приведшего к снижению веса у пациентов с ожирением, физические упражнения через 12 месяцев были эффективнее отсутствия вмешательства по поддержанию результата: $-4{,}1$~кг (95\% ДИ от $-7{,}8$ до $-0{,}4$). Однако они уступали лираглутиду: $-7{,}8$~кг (95\% ДИ от $-10{,}1$ до $-3{,}1$), и сочетанию упражнений с лираглутидом: $-9{,}5$~кг (95\% ДИ от $-13{,}1$ до $-5{,}9$).\src{42}


\subsubsection{Лекарственная терапия ожирения для контроля массы тела при стадии ССПМ 1}
\heading{Рекомендации по лекарственной терапии ожирения}
Исследования, обосновывающие рекомендации, обобщены в таблице доказательств оригинала.
{\small
\begin{longtable}{P{0.07\textwidth}P{0.11\textwidth}P{0.75\textwidth}}
\toprule COR & LOE & Рекомендация\\\midrule
\endfirsthead
\toprule COR & LOE & Рекомендация (продолжение)\\\midrule
\endhead
2a & A & \textbf{1.} У взрослых со стадией ССПМ 1, то есть избыточной массой тела/ожирением при ИМТ $\geq27$~кг/м$^2$ без других метаболических факторов, ХБП и субклинического либо клинически выраженного ССЗ, добавление терапии на основе GLP-1 с доказанной пользой к структурированному вмешательству по изменению образа жизни может быть полезным для снижения веса, улучшения гликемии и профиля ССПМ-риска.\src{1--3}\\[5pt]
\multicolumn{2}{P{0.18\textwidth}}{Экономическая оценка} & \textbf{2. Экономически неэффективно; высокая степень уверенности.} Для взрослых со стадией ССПМ 1 добавление терапии на основе GLP-1 с доказанной пользой к структурированному изменению образа жизни, согласно прогнозу, не является экономически эффективным при ценах в США 2025 года.\\[5pt]
2b & A & \textbf{3.} У взрослых со стадией ССПМ 1 для снижения массы тела может быть целесообразна лекарственная терапия ожирения, не основанная на GLP-1.\src{1}\\
\bottomrule
\end{longtable}
}
\textbf{Примечание оригинала об экономической оценке.} Эти положения предназначены для решений на уровне населения и системы здравоохранения, а не для непосредственного определения лечения конкретного пациента.

\heading{Обзор}
Раздел посвящён лекарственной терапии при избыточной массе тела/ожирении без других компонентов ССПМ. Рекомендации по препаратам на основе GLP-1 при дополнительных состояниях приведены далее: при СД2 --- в разделе 5.5.1, ХБП --- 5.5.4, АССЗ --- 6.2, СН --- 6.3, MASLD --- 7.2. Под терапией на основе GLP-1 понимается более широкий класс средств, включающих агонизм рецептора GLP-1, в том числе тирзепатид --- агонист рецепторов GLP-1 и глюкозозависимого инсулинотропного полипептида (GIP).

Появление терапии на основе GLP-1 существенно расширило возможности медикаментозного контроля массы тела. Доступны три препарата --- лираглутид, семаглутид и тирзепатид; многие другие находятся в разработке. Семаглутид и тирзепатид обеспечивают снижение массы тела вплоть до 15--21\% и улучшают профиль ССПМ-риска. Несмотря на пользу, эти препараты часто вызывают желудочно-кишечные нежелательные явления, а после отмены достигнутые преимущества утрачиваются.\src{1,4} Стоимость также остаётся препятствием для многих пациентов.

FDA одобрило в США несколько препаратов для лечения ожирения, не основанных на GLP-1: орлистат, фентермин, фентермин/топирамат и налтрексон/бупропион. Они снижают вес и несколько улучшают факторы ССПМ-риска, но эффект умерен по сравнению с GLP-1-препаратами, а побочные эффекты могут ограничивать применение. Клинические эффекты и противопоказания одобренных препаратов обобщены в таблице 13.

\heading{Обоснование отдельных рекомендаций}
\textbf{1. Терапия на основе GLP-1.} Многочисленные рандомизированные исследования лираглутида, семаглутида и тирзепатида с контролем плацебо на фоне изменения образа жизни показали эффективность снижения веса у лиц без диабета.\src{5} Снижение массы тела более чем на 5\% и значительное уменьшение окружности талии встречались очень часто. Снижение АД происходило параллельно уменьшению веса; улучшались также липиды и высокочувствительный С-реактивный белок. Гликемия последовательно улучшалась; имеются данные о предотвращении прогрессирования в диабет при инсулинорезистентности.\src{6}

Более новые препараты --- семаглутид и тирзепатид --- эффективнее снижают вес и, возможно, лучше переносятся, чем лираглутид.\src{1--3,7} Исследование SURMOUNT-5 (исследование тирзепатида у участников с ожирением либо избыточной массой тела и связанными с весом сопутствующими заболеваниями) показало превосходство тирзепатида над семаглутидом по общему снижению массы тела и удержанию результата.\src{8} В настоящее время доступен один пероральный GLP-1-препарат --- семаглутид. Его пероральная форма, по-видимому, несколько менее эффективна для снижения веса, чем подкожная.\src{9} Как обсуждается в разделе 8, при недостаточном снижении или повторном наборе веса на GLP-1-терапии показано направление к специалистам по контролю массы тела для рассмотрения альтернативного препарата либо метаболической и бариатрической хирургии.

\textbf{2. Экономическая оценка.} Моделирующее исследование сравнивало семаглутид и тирзепатид с одним изменением образа жизни с позиции здравоохранения США и на горизонте всей жизни.\src{10} В подгруппе с ожирением без сопутствующих диабета, гипертензии или установленного ССЗ инкрементальное соотношение затрат и эффективности (ICER) составляло для семаглутида 1\,012\,000 долл. США на дополнительный QALY при чистой цене 8412 долл. США, а для тирзепатида --- 449\,000 долл. США на QALY при чистой цене 6236 долл. США.

Другая модель для когорты более высокого риска без диабета, но с такими состояниями, как гипертензия, показала, что семаглутид при годовой стоимости 13\,618 долл. США по прейскурантной цене не будет экономически эффективным по сравнению с одним стандартным изменением образа жизни: ICER составлял 237\,000 долл. США на дополнительный QALY.\src{11} ICER был ещё выше при сравнении семаглутида с другими препаратами против ожирения --- например 469\,000 долл. США на QALY по сравнению с фентермином/топираматом --- или лираглутида с одним изменением образа жизни: 483\,000 долл. США на QALY.

За предшествующий год стоимость семаглутида снизилась вследствие переговоров, скидок и прямых продаж потребителям. Предварительное обновление прежнего анализа указывает на улучшение экономической эффективности в этой подгруппе более высокого риска, использующей препарат для контроля веса.\src{12}

\textbf{3. Препараты, не основанные на GLP-1.} В США доступны несколько таких средств; обычно они меньше влияют на вес и профиль ССПМ-риска, чем GLP-1-терапия. Наибольшее снижение веса связано с фентермином/топираматом,\src{13} затем следуют бупропион/налтрексон\src{14--16} и орлистат.\src{17,18} Все три улучшают гликемию; фентермин/топирамат и орлистат также снижают АД.

Несмотря на эффективность фентермина/топирамата, присутствие стимулятора вызывает опасения при использовании у особенно восприимчивых к осложнениям или побочным эффектам пациентов. Фентермин отдельно также применяют при ожирении, но часто краткосрочно, в течение 12 недель, из-за опасений сердечно-сосудистых эффектов и возможности злоупотребления. Бупропион/налтрексон связан с небольшим повышением АД.\src{1,4} Поэтому, как обсуждается далее в разделе 6.3.1, этих средств в целом следует избегать при ССЗ ввиду возможного увеличения сердечно-сосудистого риска. Как и при GLP-1-терапии, желудочно-кишечные нежелательные явления часты при бупропионе/налтрексоне и особенно при орлистате.

\heading{Таблица 13. Средние изменения массы тела и метаболического риска, частые побочные эффекты и предостережения для одобренных препаратов против ожирения}
\editorial{Широкая таблица представлена отдельными блоками препаратов. В каждом блоке последовательно сохранены доза из оригинала, среднее максимальное снижение веса с поправкой на плацебо, максимальное снижение веса, длительность исследования, метаболические эффекты, частые нежелательные явления и предостережения. Это не схема назначения: частота приёма в исходной таблице не указана.}

\heading{Тирзепатид, 15 мг\src{1,3,19--21}\textsuperscript{\dag}}
\textbf{Снижение веса:} с поправкой на плацебо 15\%; максимальное 21\%. \textbf{Длительность:} 72 недели. \textbf{Уменьшение метаболических показателей с поправкой на плацебо к завершению РКИ:} систолическое АД на 6,5~мм рт. ст.; диастолическое АД на 3~мм рт. ст.; HbA1c на 0,4\%; триглицериды на 19\%. \textbf{Частые побочные эффекты:} тошнота, диарея, запор. \textbf{Предостережения:} избегать при множественной эндокринной неоплазии II типа и панкреатите в анамнезе.

\heading{Семаглутид, 2,4 мг\src{2,5,22--24}}
\textbf{Снижение веса:} с поправкой на плацебо 12\%; максимальное 15\%. \textbf{Длительность:} 68 недель. \textbf{Уменьшение показателей с поправкой на плацебо:} систолическое АД на 4,5~мм рт. ст.; диастолическое АД на 2~мм рт. ст.; HbA1c на 0,3\%; триглицериды на 16\%. \textbf{Частые побочные эффекты:} тошнота, диарея, рвота, запор. \textbf{Предостережения:} избегать при множественной эндокринной неоплазии II типа и панкреатите в анамнезе.

\heading{Лираглутид, 3 мг\src{1,19,25}}
\textbf{Снижение веса:} с поправкой на плацебо 5\%; максимальное 8\%. \textbf{Длительность:} 56 недель. \textbf{Уменьшение показателей с поправкой на плацебо:} систолическое АД на 3~мм рт. ст.; HbA1c на 0,3\%; триглицериды на 8~мг/дл. \textbf{Частые побочные эффекты:} учащение пульса, тошнота, диарея, запор. \textbf{Предостережения:} избегать при множественной эндокринной неоплазии II типа и панкреатите в анамнезе.

\heading{Фентермин/топирамат, 15 мг/92 мг\src{1,13,17,19}}
\textbf{Снижение веса:} с поправкой на плацебо 8\%; максимальное 13\%. \textbf{Длительность:} 56 недель. \textbf{Уменьшение показателей с поправкой на плацебо:} систолическое АД на 3~мм рт. ст.; диастолическое АД на 1~мм рт. ст.; триглицериды на 20~мг/дл. \textbf{Частые побочные эффекты:} учащение пульса, головокружение, сухость во рту, запор. \textbf{Предостережения:} ограничивать дозу при СКФ $<50$; избегать при коронарной болезни; резкая отмена связана с риском судорог.

\heading{Налтрексон/бупропион, 32 мг/360 мг\src{1,14,19}}
\textbf{Снижение веса:} с поправкой на плацебо 4\%; максимальное 8\%. \textbf{Длительность:} 56 недель. \textbf{Уменьшение показателей с поправкой на плацебо:} триглицериды на 9~мг/дл; HbA1c на 1\% при СД2. \textbf{Частые побочные эффекты:} повышение систолического АД на 2~мм рт. ст., головная боль, тошнота, запор. \textbf{Предостережения:} ограничивать дозу при СКФ $<50$; контролировать АД; избегать при гипертензии и коронарной болезни; противопоказано лицам с риском судорог.

\heading{Орлистат, 120 мг\src{1,18,19}}
\textbf{Снижение веса:} с поправкой на плацебо 3\%; максимальное 10\%. \textbf{Длительность:} 52 недели\textsuperscript{\ddag}. \textbf{Эффекты с поправкой на плацебо:} снижение систолического АД на 2,5~мм рт. ст.; возможная помощь в предотвращении прогрессирования в диабет. \textbf{Частые побочные эффекты:} подтекание кала, стеаторея, неотложные позывы к дефекации. \textbf{Предостережения:} рекомендуется дополнительный приём жирорастворимых витаминов.

\heading{Фентермин, 15--37,5 мг\src{26--28}}
\textbf{Снижение веса:} с поправкой на плацебо 3--5\%; максимальное 10\%. \textbf{Длительность:} 28 недель. \textbf{Показатели метаболического эффекта с поправкой на плацебо:} систолическое АД 1,3~мм рт. ст.; частота сердечных сокращений 5 уд/мин; триглицериды 10\%. \textbf{Частые побочные эффекты:} сухость во рту, бессонница, головная боль, тревога, возбуждение, гипертензия (редко). \textbf{Предостережения:} избегать при неконтролируемой гипертензии и ССЗ, расстройстве вследствие злоупотребления психоактивными веществами в анамнезе, глаукоме.

\textbf{Примечания исходной таблицы 13.} Показатели снижения веса предполагают максимальную одобренную или переносимую дозу во время РКИ. \textsuperscript{\dag} Включены исходы у пациентов без диабета. \textsuperscript{\ddag} Максимальное снижение веса отмечено к концу первого года при соблюдении «диеты для снижения веса»; исследование продолжалось ещё год с «диетой для поддержания веса», в течение которого наблюдался повторный набор.

A1c --- HbA1c; BP --- АД; CAD --- коронарная болезнь; CVD --- ССЗ; DBP --- диастолическое АД; DM --- сахарный диабет; GFR --- СКФ; HTN --- гипертензия; RCT --- РКИ; SBP --- систолическое АД; T2D --- СД2; TG --- триглицериды.
\editorial{В строке фентермина исходная колонка озаглавлена как снижение метаболических факторов, но отдельные значения, включая ЧСС 5 уд/мин, приведены без знака направления; эта запись сохранена без самостоятельного уточнения. В предостережениях при СКФ $<50$ единицы в исходной таблице не указаны. Изменения HbA1c сохранены в процентах, как в оригинале; речь идёт об абсолютной разнице процентных значений.}

\subsubsection{Хирургические вмешательства для снижения массы тела при стадиях ССПМ 1--3}
\heading{Рекомендации по хирургическому снижению массы тела}
Исследования, обосновывающие рекомендации, обобщены в таблице доказательств оригинала.
{\small
\begin{longtable}{P{0.07\textwidth}P{0.11\textwidth}P{0.75\textwidth}}
\toprule COR & LOE & Рекомендация\\\midrule
\endfirsthead
\toprule COR & LOE & Рекомендация (продолжение)\\\midrule
\endhead
2a & A & \textbf{1.} При стадиях ССПМ 1--3 и ожирении с недостаточным снижением веса на фоне изменения образа жизни, с дополнительной лекарственной терапией ожирения или без неё, метаболическая и бариатрическая хирургия может быть полезна для снижения массы тела и сдерживания прогрессирования синдрома.\src{1--8}\\[5pt]
\multicolumn{2}{P{0.18\textwidth}}{Экономическая оценка} & \textbf{2. Экономически эффективно; высокая степень уверенности.} При тяжёлом ожирении бариатрическая хирургия экономически эффективна по сравнению с медикаментозным лечением независимо от исходного наличия диабета.\\[5pt]
2a & B-NR & \textbf{3.} После метаболической и бариатрической операции при значительном возврате веса --- $\geq25\%$ от общей потерянной массы --- терапия на основе GLP-1 может быть полезна для повторного снижения веса и ведения сопутствующих заболеваний.\src{9--13}\\
\bottomrule
\end{longtable}
}
\textbf{Примечание оригинала.} Экономические положения предназначены для решений на популяционном и системном уровнях, а не для непосредственного определения лечения отдельных пациентов.

\heading{Обзор}
Термин «метаболическая и бариатрическая хирургия» (MBS) объединяет операции по снижению веса, клиническая ценность которых выходит за пределы изменения массы тела: они уменьшают сопутствующую патологию и долгосрочную смертность. Это давно эффективная возможность для тех, кто не может снизить вес другими способами. Гастрошунтирование по Ру (RYGB) и лапароскопическая рукавная резекция желудка (LSG) составляют 90\% вмешательств.\src{9}

Хотя GLP-1-терапия обещает значимое снижение веса и ССПМ-риска, долгосрочных данных всё ещё недостаточно, а прекращение лечения с возвратом веса встречается часто. До настоящего времени MBS в целом демонстрировала превосходство по длительному снижению веса и уменьшению сердечно-сосудистого риска, заболеваемости и смертности; это может измениться с накоплением данных о GLP-1.\src{14--16} MBS в целом безопасна,\src{17} при этом хирургические и лекарственные подходы не исключают друг друга. GLP-1-терапию можно применять до операции для снижения веса либо после неё при повторном наборе; при недостаточном эффекте лекарств MBS можно предложить как альтернативу. Для выбора хирургических и медикаментозных вариантов уместно совместное принятие решений, в идеале при поддержке междисциплинарной команды.

\heading{Обоснование отдельных рекомендаций}
\textbf{1. Эффективность хирургии.} MBS эффективна при ИМТ $\geq35$~кг/м$^2$ независимо от сопутствующих заболеваний и может быть полезна при СД2 с пороговым ИМТ 30~кг/м$^2$.\src{6} Имеются данные о пользе у пациентов азиатского происхождения и при ИМТ $<30$~кг/м$^2$.\src{3} Средняя потеря избыточной массы больше после RYGB (56\%), чем после LSG (46\%),\src{2,18} и различия сохраняются через 10 лет.\src{4} MBS благоприятно влияет на все факторы ССПМ-риска, а также на ОАС и MASLD.

В РКИ GATEWAY (Gastric Bypass to Treat Obese Patients With Steady Hypertension --- гастрошунтирование при ожирении и стойкой гипертензии) MBS обеспечивала более частую ремиссию гипертензии и уменьшение потребности в антигипертензивных средствах по сравнению с медикаментозным лечением; в оригинале приведено сопоставление 2,4\% против 46,9\%.\src{7,19} Исследования STAMPEDE (Surgical Treatment and Medications Potentially Eradicate Diabetes Efficiently) и ARMMS-T2D (Alliance of Randomized Trials of Medicine versus Metabolic Surgery in T2D) показали более частую ремиссию диабета после MBS, чем при медикаментозной терапии/изменении образа жизни; в ARMMS-T2D GLP-1-терапию получали 40\% участников.\src{6,8} Ремиссия факторов ССПМ-риска выраженнее после RYGB, чем после LSG.\src{20,21}

Наблюдательные исследования с подбором сопоставимых групп также выявили снижение смертельных и несмертельных событий ССЗ на 42\% и общей смертности на 49\% при сравнении MBS с обычной помощью.\src{1,5} Крупные регистры показывают частоту больших хирургических осложнений $<5\%$ и смертность $<0{,}1\%$. Американская коллегия хирургов предлагает пациентам онлайн-калькулятор бариатрического риска для поддержки совместного выбора операции.

\textbf{2. Экономическая оценка.} Модели сравнивали LSG и RYGB с одной медикаментозной терапией при тяжёлом ожирении, с исходным диабетом или без него.\src{22,23} Обобщение данных разных исследований показало: ICER зависит от популяции и вмешательства, но последовательно остаётся ниже порога ACC/AHA 120\,000 долл. США на дополнительный QALY.

Например, в одном исследовании за 5 лет RYGB давало дополнительно 0,44 QALY при дополнительных затратах 20\,633 долл. США, то есть ICER 46\,877 долл. США на QALY в долларах 2020 года.\src{23} ICER обратно зависел от исходного ИМТ и тяжести диабета: при большем ИМТ и более тяжёлом диабете хирургия была экономически привлекательнее. Он также снижался с длительностью наблюдения, если преимущества сохранялись долго.

Важная оговорка: эти исследования сравнивали операции с прежними препаратами против ожирения. Недавно опубликованная модель предположила, что GLP-1-терапия при годовой стоимости 19\,881 долл. США не будет экономически эффективна по сравнению с RYGB.\src{24} Однако исследование не содержит достаточных подробностей для проверки калибровки и предполагает исходную цену препарата значительно выше стоимости 2023 года. Поэтому сравнительная эффективность бариатрической хирургии и высокоинтенсивной GLP-1-терапии остаётся неопределённой.

\textbf{3. Возврат веса после операции.} Обычно MBS обеспечивает отличное снижение массы, однако у 10\% пациентов после RYGB и 30\% после LSG результат недостаточен --- менее 20\% общей массы тела.\src{25} Частичное восстановление веса после минимума, достигаемого через 12--18 месяцев, распространено.\src{26} В проспективной когорте 1406 взрослых после MBS при наблюдении 6,6 года у пациентов с увеличением веса после операции более чем на 20\% отмечалось клинически значимое ухудшение качества жизни (42\%), а у 35,3\% --- ухудшение диабета.\src{27} В PCORnet Bariatric Study средний возврат веса составлял 7,2~кг после RYGB и 8,2~кг после LSG.\src{25}

До недавнего времени возможности коррекции недостаточного снижения и повторного набора были ограничены; одним из вариантов являлась повторная операция.\src{28} Несколько наблюдательных исследований показали, что GLP-1-терапия уменьшает послеоперационный возврат веса.\src{10--12,29} Эффективные препараты включают лираглутид, семаглутид и тирзепатид --- агонист GIP/GLP-1. Такое применение соответствует пониманию ожирения как хронического состояния, нередко требующего возобновления лечения на протяжении жизни.\src{30}

\subsubsection{Ведение после гестационного диабета}
\heading{Рекомендации по ведению после гестационного диабета}
Исследования, обосновывающие рекомендации, обобщены в таблице доказательств оригинала.
{\small
\begin{longtable}{P{0.07\textwidth}P{0.11\textwidth}P{0.75\textwidth}}
\toprule COR & LOE & Рекомендация\\\midrule
\endfirsthead
\toprule COR & LOE & Рекомендация (продолжение)\\\midrule
\endhead
1 & B-R & \textbf{1.} Женщинам с диагностированным гестационным сахарным диабетом (ГСД) рекомендуется пероральный глюкозотолерантный тест через 4--12 недель после родов для скрининга предиабета и диабета и выявления повышенного риска СД2 и прогрессирования ССПМ.\src{1--4}\\[5pt]
1 & C-LD & \textbf{2.} Женщинам с ГСД рекомендуется консультирование медицинским специалистом для информирования о повышенном риске СД2 и перехода к более поздним стадиям ССПМ.\src{5}\\[5pt]
1 & B-R & \textbf{3.} Женщинам после родов с ГСД в анамнезе и предиабетом рекомендуется оптимизировать гликемию и метаболические факторы посредством изменения образа жизни или метформина для снижения риска впервые возникающего СД2 и прогрессирования ССПМ.\src{6--9}\\
\bottomrule
\end{longtable}
}

\heading{Обзор}
ГСД означает нарушение толерантности к глюкозе или углеводам, выявленное либо проявившееся во время беременности. Ввиду повышенного риска СД2 после ГСД подход к последующей оценке гликемии, липидов и функции почек --- не реже одного раза в 2--3 года --- должен соответствовать подходу при стадии ССПМ 1 (раздел 3.1.1).

Нелеченый ГСД затрагивает 2--10\% родов в США и связан с вероятностью перехода в СД2 35--60\% в течение 10--20 лет.\src{10} Ожирение, чрезмерная прибавка веса при беременности и сохранение прибавленного веса после родов являются факторами риска ГСД, с последующим десятикратным риском СД2 и двукратным риском ССЗ. Это подчёркивает значение ранней послеродовой диагностики и ведения: скрининга СД2, своевременного начала помощи для повышения осведомлённости о риске и профилактике, консультирования по лактации и здоровому образу жизни, лекарственных и немедикаментозных стратегий снижения дальнейшего риска.

Лицам группы риска следует изменять образ жизни и снижать вес для уменьшения вероятности СД2 и прогрессирования ССПМ.\src{11} Необходимы усиленное наблюдение, диагностика и раннее вмешательство. Однако существующие прогностические модели СД2 после ГСД имеют методологические недостатки, а лекарственные подходы, способствующие изменению образа жизни у женщин с ГСД, требуют дальнейшего изучения.

\heading{Обоснование отдельных рекомендаций}
\textbf{1. Послеродовой скрининг.} Данные поддерживают пероральный глюкозотолерантный тест с 75~г глюкозы через 4--12 недель после родов для выявления диабета и предиабета.\src{1} Альтернативы включают раннее тестирование до выписки, чтобы уменьшить вероятность пропуска скрининга через 4--12 недель,\src{2,3} глюкозу плазмы натощак через 6--12 недель либо HbA1c после 13 недель.\src{12} При использовании глюкозы натощак и HbA1c ($\geq7{,}5\%$) в раннем послеродовом периоде сообщается о более высокой выявляемости нарушенной толерантности к глюкозе по сравнению с критериями глюкозотолерантного теста.\src{1}

\editorial{Порог HbA1c $\geq7{,}5\%$ в предыдущем предложении воспроизведён из данного абзаца оригинала. Он не заменяет диагностические пороги диабета и предиабета, приведённые авторами в таблицах 4 и 5.}

Лишь немногие исследования сравнивали глюкозу натощак и/или HbA1c с глюкозотолерантным тестом у женщин после ГСД. Один HbA1c, измеренный рано или поздно после родов, показал недостаточную чувствительность для выявления нарушенной толерантности.\src{4} Непрерывный мониторинг глюкозы не является рутинным компонентом помощи всем женщинам с ГСД,\src{13} но интерес к его применению во время и после беременности растёт.\src{14} После нормального раннего послеродового глюкозотолерантного теста можно рассматривать длительное тестирование с интервалами 1--3 года для минимизации риска диабета у матери.\src{15} Поддерживается также ежегодное определение HbA1c.\src{12}

\textbf{2. Консультирование.} Традиционно оно ограничивалось одним визитом через 4--6 недель после родов. Однако новые данные указывают, что обсуждение снижения ССПМ-риска у женщин высокого риска должно начинаться до родов и продолжаться до завершения послеродовых визитов матери и ребёнка.\src{16--18} Включение консультирования до, во время и после беременности существенно улучшало обследование гликемии и уменьшало риски и осложнения ССПМ в этой популяции.\src{5} Консультирование должно отвечать потребностям и опасениям матерей, повышать осведомлённость о прогрессировании ССПМ и помогать активно управлять рисками, например сохранением веса после родов.

Консультирование по питанию либо его сочетанию с упражнениями эффективно способствует снижению веса после ГСД.\src{6,19--21} Во время лактации уменьшаются инсулинорезистентность, глюкоза крови, серотонин и пролиферация бета-клеток. Лактация также уменьшает окислительный стресс поджелудочной железы и риск СД2 после ГСД; её следует обсуждать до, во время и после беременности.\src{22--24} Наконец, женщинам с ГСД необходима биопсихосоциальная оценка послеродовой депрессии, связанной с ССПМ-риском, для раннего выявления, направления и управления риском.\src{25}

\textbf{3. Профилактика СД2 после ГСД.} При ГСД в анамнезе и предиабете после родов необходимо оптимизировать гликемию и метаболические факторы посредством изменения образа жизни либо метформина.\src{6--9} Апостериорный анализ DPP (Diabetes Prevention Program) показал, что 12-месячное вмешательство по образу жизни или метформин улучшали гликемический контроль и снижали риск СД2 на 50--60\% по сравнению с контролем у женщин с ГСД.\src{19} Интенсивное изменение образа жизни сохраняло преимущества через 10 лет. Планирование целей, самоконтроль, уменьшение жиров и калорийности и 150 минут умеренных упражнений в неделю снижали ССПМ-риск.\src{6}

Последующие исследования расширили применение послеродовых вмешательств и получили сопоставимые результаты.\src{20,26,27} Программа для женщин с избыточной массой тела после ГСД, включавшая обучение, контроль веса, шагомер, телефонную и текстовую поддержку, обеспечила большее снижение веса через 6 месяцев: $-3{,}9$~кг (стандартное отклонение 7,0) против $-0{,}7$~кг (стандартное отклонение 3,8) в контроле.\src{8} Наблюдательные исследования связывают послеродовое питание DASH или средиземноморского типа с меньшей частотой гипертензии и СД2.\src{9,28,29} Пиоглитазон уменьшает риск СД2 у латиноамериканских женщин после ГСД, но повышает риск прибавки веса, отёков и переломов.\src{30,31}


\subsection{Стадия ССПМ 2}
\heading{Обзор}
Стадия ССПМ 2 определяется наличием метаболических факторов риска, ХБП умеренного или высокого риска либо их сочетанием при отсутствии субклинического и клинически выраженного ССЗ. Профилактика на этой стадии основывается на принципах стадий 0 и 1, но смещает акцент на своевременную коррекцию метаболических факторов и ХБП для предотвращения клинического ССЗ и почечной недостаточности.

Изменение образа жизни остаётся основой профилактики, однако лекарственное лечение метаболических факторов и/или ХБП также занимает центральное место в уменьшении риска прогрессирования. Как описано далее, помимо терапии гипертензии\src{1} и гипертриглицеридемии\src{2} по действующим руководствам, подчёркивается применение кардиопротективных сахароснижающих средств при СД2 и нефропротективных препаратов при ХБП, уменьшающих неблагоприятные сердечно-сосудистые и почечные события.

10-летний риск ССЗ $\geq7{,}5\%$ по PREVENT-CVD дополнительно помогает выбирать пациентов для лекарственной терапии и должен обсуждаться с пациентом. Снижение веса посредством изменения образа жизни, препаратов против ожирения или бариатрической хирургии важно для предотвращения прогрессирования и содействия переходу к более ранней стадии ССПМ.

\subsubsection{Ведение пациентов с СД2 при стадиях ССПМ 2--3}
\heading{Рекомендации по ведению СД2 при стадиях ССПМ 2--3}
Исследования, обосновывающие рекомендации, обобщены в таблице доказательств оригинала.
{\small
\begin{longtable}{P{0.07\textwidth}P{0.11\textwidth}P{0.75\textwidth}}
\toprule COR & LOE & Рекомендация\\\midrule
\endfirsthead
\toprule COR & LOE & Рекомендация (продолжение)\\\midrule
\endhead
1 & A & \textbf{1.} Взрослым со стадиями ССПМ 2--3 и СД2 рекомендуется изменение образа жизни, включающее изменение поведения, здоровое питание и рекомендованный уровень физической активности, при необходимости с дополнительными лекарственными и/или хирургическими подходами, для достижения и поддержания здоровой массы тела и улучшения гликемического контроля.\src{1--5}\\[5pt]
1 & B-R & \textbf{2.} Взрослым со стадиями ССПМ 2--3 и СД2 рекомендуется интенсивное многофакторное вмешательство: изменение образа жизни и целевая лекарственная коррекция гипергликемии, гипертензии, дислипидемии и альбуминурии для уменьшения риска ССЗ, смертности и почечных осложнений.\src{6--8}\\[5pt]
1 & A & \textbf{3.} У взрослых со стадиями ССПМ 2--3, СД2 и повышенным риском ССЗ --- 10-летний PREVENT-CVD $\geq7{,}5\%$ --- план лечения должен включать SGLT2i или терапию на основе GLP-1 с доказанной пользой для уменьшения сердечно-сосудистых событий и смертности.\src{9--12}\\[5pt]
\multicolumn{2}{P{0.18\textwidth}}{Экономическая оценка} & \textbf{4. Экономически эффективно; умеренная степень уверенности.} У взрослых со стадиями ССПМ 2--3, СД2 и повышенным риском ССЗ лечение SGLT2i при ценах США 2025 года, согласно прогнозу, экономически эффективно по сравнению с обычной помощью.\\[5pt]
\multicolumn{2}{P{0.18\textwidth}}{Экономическая оценка} & \textbf{5. Экономически неэффективно; умеренная степень уверенности.} У взрослых со стадиями ССПМ 2--3, СД2 и повышенным риском ССЗ GLP-1-терапия при ценах США 2025 года, согласно прогнозу, не является экономически эффективной по сравнению с обычной помощью.\\[5pt]
2a & A & \textbf{6.} У взрослых со стадиями ССПМ 2--3, СД2 и повышенным риском ССЗ при HbA1c на 0,5--1 процентного пункта выше индивидуальной цели метформин в сочетании с кардиопротективной сахароснижающей терапией --- GLP-1-препаратом или SGLT2i --- может быть эффективен для достижения целевой гликемии.\src{13--15}\\[5pt]
2b & B-NR & \textbf{7.} У взрослых со стадиями ССПМ 2--3, СД2 и повышенным риском ССЗ и/или несколькими факторами ССПМ-риска можно рассмотреть сочетание GLP-1-терапии с доказанной пользой и SGLT2i для дополнительного уменьшения сердечно-сосудистых событий по сравнению с одним кардиопротективным сахароснижающим препаратом.\src{16--21}\\
\bottomrule
\end{longtable}
}
\textbf{Примечание оригинала.} Экономические положения предназначены для решений на уровне населения и системы здравоохранения и не должны непосредственно определять лечение конкретного пациента.

\heading{Обзор}
СД2 --- частое осложнение избыточной массы тела и ожирения и один из основных факторов риска ССЗ.\src{22} Большинство пациентов с СД2 имеют дополнительные метаболические факторы, включая избыток жировой ткани, гипертензию и дислипидемию. Изменение образа жизни крайне важно для гликемического контроля, снижения веса и улучшения метаболических факторов и связанных патофизиологических нарушений.

Оптимальное ведение гликемии и сопутствующих состояний уменьшает микро- и макрососудистые осложнения за счёт достижения комплексных целей контроля факторов риска и применения антиагрегантов при наличии показаний.\src{6,7,23,24} Цели включают HbA1c $<7{,}0\%$ для небеременных взрослых без значимой гипогликемии и АД $<130/80$~мм рт. ст.\src{23,25,26} Менее строгая цель HbA1c, например $<8{,}0\%$, может быть уместна при ограниченной ожидаемой продолжительности жизни либо когда вред более интенсивного контроля может превышать пользу.\src{26}

Поскольку большинство лиц с диабетом в возрасте $\geq40$ лет имеют как минимум промежуточный риск АССЗ, рекомендуется статинотерапия умеренной или высокой интенсивности с рассмотрением добавления эзетимиба при более высоком риске для снижения ХС ЛНП не менее чем на 50\%.\src{23,27,28} При повышенном прогнозируемом риске ССЗ рекомендуется SGLT2i или GLP-1-терапия; выбор определяется такими сопутствующими состояниями, как ХБП, pre-HF, ожирение, выраженная гипергликемия и MASLD (рисунок 8).\src{9--12,23,26}

\heading{Рисунок 8. Ведение СД2 при стадиях ССПМ 2--3: текстовая адаптация алгоритма}
\textbf{Исходная группа:} ССПМ-синдром с СД2 без ССЗ.
\begin{enumerate}[label=\arabic*.]
\item Рекомендовать и поддерживать изменение образа жизни --- \textbf{COR 1}.
\item Обеспечить многокомпонентный контроль факторов риска --- \textbf{COR 1}.
\item Оценить, повышен ли риск. В сноске рисунка повышенный риск определён как 10-летний PREVENT-CVD $\geq7{,}5\%$ либо возраст $\geq50$ лет при СД2 и дополнительных факторах ССПМ-риска.
\begin{itemize}
\item \textbf{Нет:} стандартная сахароснижающая лекарственная терапия.
\item \textbf{Да:} добавить кардиопротективное сахароснижающее средство --- \textbf{COR 1}.
\end{itemize}
\item Выбор кардиопротективного препарата:
\begin{itemize}
\item \textbf{GLP-1-терапия, COR 1:} приоритет при тяжёлом ожирении, MASLD или выраженной гипергликемии. Совместно применять метформин, если он ещё не назначен и HbA1c превышает цель более чем на 0,5--1 процентного пункта, --- \textbf{COR 2a}.
\item \textbf{SGLT2i, COR 1:} приоритет при ХБП или pre-HF. Совместно применять метформин при тех же условиях --- \textbf{COR 2a}.
\end{itemize}
\item При нескольких факторах ССПМ-риска или высоком прогнозируемом риске --- комбинированная GLP-1-терапия и SGLT2i, \textbf{COR 2b}.
\end{enumerate}
Легенда исходной схемы включает COR 1, 2a, 2b, 3 «отсутствие пользы» и 3 «вред»; в алгоритме используются 1, 2a и 2b. CKM --- ССПМ; CVD --- ССЗ; GLP-1 --- глюкагоноподобный пептид-1; HbA1c --- гликированный гемоглобин; HF --- СН; MASLD --- стеатотическая болезнь печени, ассоциированная с метаболической дисфункцией; SGLT2i --- ингибиторы натрий-глюкозного котранспортера 2 типа; T2D --- СД2.
\editorial{В рисунке 8 возрастной критерий записан как $\geq50$ лет, тогда как в рисунке 6 --- $>50$ лет. Оба порога сохранены по соответствующим фрагментам оригинала.}

\heading{Обоснование отдельных рекомендаций}
\textbf{1. Контроль массы тела.} Он имеет многочисленные преимущества при СД2, включая снижение HbA1c, уменьшение стеатоза печени и улучшение сердечно-сосудистых факторов.\src{26,29} Когортные и интервенционные исследования подтверждают пользу изменения образа жизни, регулярной физической активности и сбалансированного питания для веса, гликемии и метаболических показателей.\src{30--32}

Препараты против ожирения, особенно на основе GLP-1, полезны как дополнение у отдельных пациентов с СД2 и ИМТ $\geq27$~кг/м$^2$. Среди доступных препаратов наибольшую эффективность в снижении веса и метаболических преимуществах имеют тирзепатид --- двойной агонист GIP и GLP-1 --- и семаглутид --- агонист GLP-1.\src{4,5} Данные также поддерживают MBS при сопутствующем ожирении.\src{3} Помимо улучшения веса, гликемии и ССПМ-факторов в клинических исследованиях, наблюдательные и сопоставленные когортные исследования выявили снижение ССЗ, сердечно-сосудистой и общей смертности при СД2.\src{33--35} Поэтому MBS следует рассматривать при СД2 и ИМТ $\geq30$~кг/м$^2$, а для американцев азиатского происхождения --- $\geq27{,}5$~кг/м$^2$. В этой популяции её рекомендуют Американское общество метаболической и бариатрической хирургии и ADA.\src{29,36}

\textbf{2. Интенсивное многофакторное вмешательство.} Большинство пациентов с СД2 имеют несколько дополнительных факторов ССЗ --- избыток абдоминальной жировой ткани, гипертензию и дислипидемию --- и соответствуют критериям метаболического синдрома.\src{37} Помимо его диагностических компонентов, риск увеличивают воспаление, эндотелиальная дисфункция, повышенная концентрация атерогенных липопротеинов и протромботическое состояние; на них необходимо воздействовать изменением образа жизни и снижением веса.

В STENO-2 у пациентов с СД2, избыточной массой тела/ожирением и альбуминурией интенсивное многофакторное вмешательство включало изменение поведения --- питание, упражнения, отказ от курения --- и многокомпонентную лекарственную терапию: сахароснижающие препараты, ACEi или ARB, липидснижающую терапию и аспирин. Оптимизация факторов риска была связана с уменьшением событий ССЗ, общей и сердечно-сосудистой смертности по сравнению с обычной терапией.\src{6,7} Преимущества нарастали на протяжении 8 лет исследования, указывая на эффективность постоянного поддержания снижения риска. Через 21 год наблюдения пациенты, получавшие интенсивную терапию первоначально в течение 7,8 года, жили дольше; увеличение продолжительности жизни было связано с временем без впервые возникшего ССЗ.\src{38}

\textbf{3. Кардиопротективная сахароснижающая терапия.} При СД2 и высоком риске ССЗ SGLT2i и/или GLP-1-терапия с доказанной сердечно-сосудистой пользой рекомендуются независимо от HbA1c и применения метформина.\src{9,12} Оба класса уменьшают MACE и сердечно-сосудистую смертность в многочисленных исследованиях.\src{39,40} SGLT2i особенно влияют на госпитализации по поводу СН, а GLP-1-терапия сильнее влияет на впервые возникающие события АССЗ.

Исследования при СД2 включали пациентов с ССЗ либо дополнительными факторами помимо диабета, то есть популяцию повышенного риска. С учётом распределения риска в исследованиях, абсолютного снижения риска и числа пациентов, которых необходимо лечить, соответствующего порогам других руководств, для приоритетного назначения SGLT2i или GLP-1-терапии рекомендуется порог 10-летнего PREVENT-CVD $\geq7{,}5\%$.\src{41} Дополняющий подход --- приоритет лечения лиц $\geq50$ лет с СД2 и дополнительными факторами. Однако наблюдательные данные указывают на возможную сердечно-сосудистую пользу предпочтительного применения этих классов и при умеренном риске.\src{11}

У молодых и средневозрастных взрослых без повышенного 10-летнего риска дополнительным ориентиром при обсуждении лекарственной терапии может быть высокий для возраста и пола 30-летний риск. В оригинале он обозначен «$\geq75\%$ для возраста и пола». С учётом различий механизмов и клинических эффектов SGLT2i можно отдавать приоритет при ХБП и/или pre-HF, а GLP-1-терапии --- при ожирении II класса и выше, выраженной гипергликемии (HbA1c $\geq9\%$) и/или MASLD. Нежелательные эффекты и способы их смягчения приведены в таблицах 14 и 15.

\editorial{Запись «$\geq75\%$ для возраста и пола» воспроизведена из данного абзаца. В разделе 4.1 авторы обсуждают соответствующий возрастно-половой процентиль 30-летнего риска и порог 75-го процентиля; эту запись не следует смешивать с абсолютной вероятностью ССЗ 75\%.}

\textbf{4. Экономическая оценка SGLT2i.} Предыдущие исследования с позиции здравоохранения США показывали, что добавление SGLT2i при диабете экономически эффективно по сравнению с прежней стандартной терапией и экономит затраты либо экономически эффективно по сравнению с ингибиторами дипептидилпептидазы-4. Цены в этих моделях были несколько ниже оптовой закупочной стоимости SGLT2i: 3500--5700 долл. США в год против примерно 6500 долл. США в год в 2025 году. С этой оговоркой исследования указывают на вероятную экономическую эффективность при пороге 120\,000 долл. США на дополнительный QALY.\src{42,43} Она может дополнительно улучшиться при снижении цен в результате переговоров Medicare или появления генерических SGLT2i.

\textbf{5. Экономическая оценка GLP-1-терапии.} Несколько исследований с позиции здравоохранения США оценивали добавление этого класса, включая его более старые препараты, к стандартной помощи при диабете.\src{44--48} В целом они показывали отсутствие экономической эффективности по сравнению с обычной помощью, особенно при использовании в первой линии. Однако ICER широко варьировал в зависимости от стоимости и препарата сравнения.\src{49} Необходима комплексная экономическая оценка новых GLP-1-препаратов, включая их сочетание с SGLT2i, для уточнения ценности этих дорогостоящих методов.

\textbf{6. Дополнительный гликемический контроль.} Снижение сердечно-сосудистого риска приоритетно, но контроль гликемии дополнительно необходим для уменьшения микрососудистых осложнений. SGLT2i умеренно снижают HbA1c --- на 0,5--1 процентного пункта. GLP-1-терапия влияет сильнее, однако недостаточная титрация или отмена из-за побочных эффектов и стоимости часто ограничивают результат. Поэтому при HbA1c более чем на 0,5--1 процентного пункта выше цели можно рассмотреть дополнительное сахароснижающее средство вместе с SGLT2i или GLP-1-препаратом.

Особое внимание уделяется метформину: он благоприятно влияет на HbA1c и вес, имеет очень низкий риск гипогликемии при отсутствии инсулинотерапии, широко доступен и безопасно сочетается с кардиопротективными сахароснижающими средствами.\src{13--15,26} Сердечно-сосудистая польза SGLT2i и GLP-1-терапии не зависит от метформина, и их следует рассматривать в первой линии.\src{14,15} Пиоглитазон можно сочетать с кардиопротективной терапией для достижения гликемических целей; он может уменьшать риск АССЗ, но связан с увеличением риска СН.\src{50} Тактика контроля гликемии должна соответствовать рекомендациям ADA.\src{26}

\textbf{7. Сочетание SGLT2i и GLP-1-терапии.} Недостаточно интервенционных данных о сравнительном влиянии комбинации на сердечно-сосудистые и почечные исходы по отношению к каждому классу отдельно. Хотя оба класса уменьшают MACE и сердечно-сосудистую смертность,\src{39,40} их механизмы и преимущества несколько различаются.\src{16} Клинические исследования не выявили неоднородности эффекта каждого класса в зависимости от приёма другого.\src{17--19,21} Наблюдательные данные и модели предполагают дополнительную сердечно-сосудистую пользу сочетания.\src{19,20}

Такой подход можно рассмотреть у отдельных пациентов с СД2 высокого риска. Дополнительными основаниями могут быть влияние GLP-1-терапии на ожирение, выраженную гипергликемию или MASLD и влияние SGLT2i на ХБП или высокий риск СН. Однако величина клинической пользы комбинации не определена в клинических испытаниях. Возможные преимущества необходимо сопоставлять с полипрагмазией, стоимостью и предпочтениями пациента.

\heading{Таблица 14. Нежелательные эффекты SGLT2i при СД2: профилактические и лечебные меры}
\heading{Риск диабетического кетоацидоза при дефиците инсулина}
При СД2 встречается значительно реже. При подозрении прекратить препарат, незамедлительно обследовать и лечить. Учитывать предрасполагающие факторы, например кетогенную диету, и клинические проявления, включая эугликемический кетоацидоз. Уменьшать риск с помощью заранее составленного плана действий на период заболевания. Отменять перед плановой операцией, например за 3--4 дня, при критическом состоянии или длительном голодании.

\heading{Грибковые инфекции половых органов}
Уменьшать риск соблюдением гигиены; избегать применения у лиц высокого риска. При выраженной гипергликемии применять с осторожностью либо избегать для уменьшения инфекционного риска.

\heading{Некротизирующий фасциит промежности --- гангрена Фурнье}
Редкое осложнение; при подозрении требуется незамедлительное лечение.

\heading{Уменьшение внутрисосудистого объёма жидкости}
Оценивать объёмный статус и АД, особенно при заболевании или голодании. При необходимости рассматривать уменьшение дозы других препаратов, снижающих объём жидкости; контролировать функцию почек после начала лечения.

\textbf{Примечание исходной таблицы 14.} Адаптировано и воспроизведено с разрешения Комитета ADA по профессиональной практике в диабетологии.\src{26} © 2026 American Diabetes Association. DKA --- диабетический кетоацидоз; SGLT2i --- ингибиторы натрий-глюкозного котранспортера 2 типа; T2D --- СД2.

\heading{Таблица 15. Нежелательные эффекты GLP-1-терапии при СД2: профилактические и лечебные меры}
\heading{Желудочно-кишечные нежелательные явления}
Предупредить об их возможности; объяснить изменения питания для смягчения симптомов; при желудочно-кишечных затруднениях рассмотреть более медленную титрацию. Не рекомендуется при гастропарезе.

\heading{Опухоли C-клеток щитовидной железы}
Выявлены у грызунов; значение для человека не установлено. Противопоказано при личном или семейном анамнезе медуллярного рака щитовидной железы и при синдроме множественной эндокринной неоплазии 2 типа.

\heading{Панкреатит}
Сообщалось об остром панкреатите, но причинная связь не установлена. Не начинать лечение при высоком риске панкреатита и прекращать при подозрении на него.

\heading{Острое заболевание желчного пузыря}
При подозрении на желчнокаменную болезнь или холецистит обследовать желчный пузырь; избегать применения у лиц группы риска.

\heading{Илеус}
Сообщалось о случаях кишечной непроходимости; уровень риска надёжно не установлен.

\heading{Влияние на одновременно принимаемые пероральные препараты вследствие замедленного опорожнения желудка}
Во время титрации может нарушаться всасывание пероральных препаратов, включая контрацептивы. Пациенткам, принимающим пероральные контрацептивы, следует рекомендовать переход на непероральный метод либо добавление барьерного метода на 4 недели после начала лечения и на 4 недели после каждого повышения дозы.

\heading{Осложнения диабетической ретинопатии при её наличии в анамнезе}
Тщательно наблюдать пациентов высокого риска: пожилых и с большей длительностью СД2 ($\geq10$ лет).

\heading{Неартериитная передняя ишемическая нейропатия зрительного нерва (NAION)}
Редкое осложнение; отслеживать признаки NAION при офтальмологических обследованиях.

\heading{Острое повреждение почек}
Контролировать функцию почек у пациентов с ХБП, сообщающих о тяжёлых желудочно-кишечных нежелательных реакциях.

\heading{Риск лёгочной аспирации}
Объяснять отмену перед хирургическими вмешательствами для уменьшения возможной аспирации при общей анестезии или глубокой седации.

\textbf{Примечание исходной таблицы 15.} Адаптировано и воспроизведено с разрешения Комитета ADA по профессиональной практике в диабетологии.\src{26} © 2026 American Diabetes Association. GI --- желудочно-кишечный; GLP-1 --- глюкагоноподобный пептид-1; T2D --- СД2.
\editorial{Таблицы 14 и 15 передают обобщённые формулировки исходника. Они не заменяют индивидуальную оценку и инструкцию конкретного препарата; указания о пероральной контрацепции и периоперационной отмене не дополнены в переводе самостоятельными уточнениями.}

\subsubsection{Гипертриглицеридемия}
\heading{Обзор}
Гипертриглицеридемия определяется как триглицериды натощак $\geq150$~мг/дл либо не натощак $\geq175$~мг/дл.\src{1,2} Некоторые исследования менделевской рандомизации указывают на причинную связь повышенных триглицеридов с АССЗ.\src{3} Действующие руководства также рассматривают стойкое повышение триглицеридов как усилитель риска АССЗ.\src{2,4--6} Оно связано с большей частотой ХБП и более быстрым снижением рСКФ, возможно, через альбуминурию.\src{7,8}

При выявлении гипертриглицеридемии необходимо установить и устранить обратимые вторичные причины: нездоровый образ жизни, гипотиреоз, неконтролируемую гипергликемию, выраженную альбуминурию и влияние лекарств.\src{1,6--8} Первоначальные меры по руководствам --- изменение образа жизни и статинотерапия при пограничном, промежуточном или более высоком прогнозируемом риске. При сохранении гипертриглицеридемии на максимально переносимых статинах либо дополнительных сердечно-сосудистых факторах для уменьшения риска АССЗ можно рассмотреть икозапент этил.\src{9}

Тяжёлая гипертриглицеридемия ($\geq500$~мг/дл) связана с ещё большим риском АССЗ и повышенным риском панкреатита. Помимо снижения ХС ЛНП, меры могут включать диету с очень низким содержанием жиров, отказ от рафинированных углеводов и алкоголя, направление к диетологу и добавление фибрата либо рецептурных омега-3 жирных кислот для снижения триглицеридов и предотвращения острого панкреатита.\src{1}

\subsubsection{Артериальная гипертензия}
\heading{Обзор}
Гипертензия --- важный компонент ССПМ-синдрома, способствующий прогрессирующему поражению органов-мишеней, включая ухудшение функции почек и ССЗ.\src{1,2} На стадии ССПМ 2 она способствует развитию ХБП, а ХБП, в свою очередь, увеличивает риск и тяжесть гипертензии. Поэтому эффективный контроль АД необходим для уменьшения заболеваемости и смертности.

Польза для снижения АД показана для ограничения натрия, снижения веса, увеличения физической активности и полезных для сердца схем питания, таких как DASH или средиземноморская.\src{3,4} Действующие рекомендации поддерживают сочетание изменения образа жизни с антигипертензивными препаратами при среднем АД $\geq140/90$~мм рт. ст. либо при среднем АД $\geq130/80$~мм рт. ст. в сочетании с клиническим ССЗ, перенесённым инсультом, диабетом, ХБП или 10-летним PREVENT-CVD $\geq7{,}5\%$.\src{5} Общая цель лечения --- АД $<130/80$~мм рт. ст.

К препаратам первой линии относятся RASi, диуретики тиазидного типа и длительно действующие дигидропиридиновые блокаторы кальциевых каналов. Для более быстрого контроля и лучшей приверженности предпочтительно начинать с комбинации препаратов первой линии в одной таблетке. RASi обеспечивают контроль АД, нефро- и кардиопротекцию\src{6} и поэтому предпочтительны при ХБП с альбуминурией. Подробные рекомендации приведены в руководстве AHA/ACC 2025 года по артериальному давлению.\src{5}

\subsubsection{Ведение пациентов с ХБП при стадиях ССПМ 2--3}
\heading{Рекомендации}
Исследования, на которых основаны рекомендации, обобщены в таблице доказательств оригинала.

{\small
\begin{longtable}{P{0.07\textwidth}P{0.11\textwidth}P{0.75\textwidth}}
\toprule COR & LOE & Рекомендация\\\midrule
\endfirsthead
\toprule COR & LOE & Рекомендация (продолжение)\\\midrule
\endhead
1 & B-R & \textbf{1.} У взрослых со стадиями ССПМ 2--3, ХБП и СД2 либо ХБП без СД2, но с UACR $\geq30$~мг/г, при рСКФ $\geq30$~\unitgfr\ рекомендуется RASi, например ACEi или ARB, в максимально переносимой дозе для замедления утраты функции почек и уменьшения риска ССЗ.\src{1--4}\\[5pt]
1 & A & \textbf{2.} У взрослых со стадиями ССПМ 2--3, ХБП и СД2 либо ХБП без СД2, но с UACR $\geq200$~мг/г, при рСКФ $\geq20$~\unitgfr\ рекомендуется SGLT2i для замедления утраты функции почек и уменьшения риска госпитализации по поводу СН и сердечно-сосудистой смертности.\src{5--10}\\[5pt]
2a & B-R & \textbf{3.} У взрослых со стадиями ССПМ 2--3 и ХБП без СД2, но с UACR от $\geq30$ до 199~мг/г, можно рассмотреть SGLT2i для замедления утраты функции почек и уменьшения риска госпитализации по поводу СН и сердечно-сосудистой смертности.\src{6,10}\\[5pt]
\multicolumn{2}{P{0.18\textwidth}}{Экономическая оценка} & \textbf{4. Экономически эффективно; высокая степень уверенности.} У взрослых со стадиями ССПМ 2--3 и ХБП с СД2 или без него лечение SGLT2i при ценах США 2025 года, согласно прогнозу, экономически эффективно по сравнению с отсутствием SGLT2i.\\[5pt]
1 & A & \textbf{5.} У взрослых со стадиями ССПМ 2--3, ХБП, СД2 и UACR $\geq30$~мг/г, сохраняющимся несмотря на применение ACEi/ARB и SGLT2i с учётом переносимости, при рСКФ $\geq25$~\unitgfr\ рекомендуется добавить нестероидный антагонист минералокортикоидных рецепторов (nsMRA) с доказанной нефро- и кардиопротективной эффективностью для уменьшения риска утраты функции почек, почечной недостаточности и ССЗ.\src{11--15}\\[5pt]
\multicolumn{2}{P{0.18\textwidth}}{Экономическая оценка} & \textbf{6. Экономически неэффективно; умеренная степень уверенности.} У взрослых со стадиями ССПМ 2--3, ХБП, СД2 и UACR $\geq30$~мг/г лечение финереноном при ценах США 2025 года, согласно прогнозу, не является экономически эффективным по сравнению с обычной помощью.\\[5pt]
1 & B-R & \textbf{7.} У взрослых со стадиями ССПМ 2--3, ХБП, СД2 и UACR $\geq100$~мг/г, сохраняющимся несмотря на применение ACEi/ARB и SGLT2i с учётом переносимости, рекомендуется терапия на основе GLP-1 с доказанной нефро- и кардиопротективной эффективностью для уменьшения риска утраты функции почек, почечной недостаточности и ССЗ.\src{16--18}\\[5pt]
\multicolumn{2}{P{0.18\textwidth}}{Экономическая оценка} & \textbf{8. Не определена; недостаточно доказательств.} У взрослых со стадиями ССПМ 2--3, ХБП, СД2 и UACR $\geq100$~мг/г, сохраняющимся несмотря на применение ACEi/ARB и SGLT2i с учётом переносимости, экономическая эффективность добавления терапии на основе GLP-1 не определена.\\
\bottomrule
\end{longtable}
}
\textbf{Примечание оригинала.} Экономические положения предназначены для решений на уровне населения и системы здравоохранения, а не для непосредственного определения лечения отдельных пациентов.

\heading{Обзор}
Главная цель ведения пациентов с ХБП на стадиях ССПМ 2--3 --- уменьшить риск утраты функции почек и больших неблагоприятных сердечно-сосудистых событий (MACE). Интенсивная коррекция сопутствующих факторов риска, включая гипертензию и дислипидемию, посредством изменения образа жизни и лекарственной терапии уменьшает почечный и сердечно-сосудистый риск. Рандомизированные контролируемые исследования (РКИ) указывают, что аторвастатин может замедлять прогрессирование ХБП.\src{19,20} Центральный компонент ведения ХБП --- препараты, которые предотвращают утрату функции почек и одновременно обеспечивают сердечно-сосудистую пользу (рисунок 9).

При ХБП и СД2 либо ХБП и альбуминурии (UACR $\geq30$~мг/г) RASi (ACEi или ARB) и SGLT2i рекомендуются как терапия первой линии благодаря длительному опыту подтверждённой эффективности и безопасности и высокой доступности; для SGLT2i наиболее убедительные данные при отсутствии СД2 получены при UACR $\geq200$~мг/г. Финеренон и/или препарат на основе GLP-1 с доказанной пользой следует рассматривать как дополнительную терапию у пациентов высокого риска с СД2, ХБП и сохраняющейся альбуминурией на фоне первой линии.

Сердечно-сосудистая польза этих препаратов в исследованиях ХБП главным образом связана с уменьшением госпитализаций по поводу СН, MACE и сердечно-сосудистой смерти. Этот подход соответствует общему принципу: наиболее интенсивное лечение предназначено для людей с наибольшим абсолютным риском. Поэтому сочетание ХБП с СД2 либо альбуминурией определяет ключевые группы с высоким исходным абсолютным риском, получающие пользу от терапии. Популяционные исследования показывают, что у подавляющего большинства таких лиц 10-летний PREVENT-CVD составляет $\geq7{,}5\%$.\src{21}

Альбуминурия --- ключевой ориентир для ведения ХБП (рисунки 10 и 11): порог UACR $\geq30$~мг/г используется для RASi, SGLT2i и финеренона, последнего --- при СД2; порог UACR $\geq100$~мг/г --- для терапии на основе GLP-1 при СД2. Данные исследований при ХБП и СД2 показывают дополнительное снижение альбуминурии при сочетании SGLT2i с финереноном по сравнению с каждым препаратом отдельно; почти все участники получали фоновую терапию RASi.\src{22} Однако исследований сердечно-сосудистой пользы сочетания более двух из перечисленных четырёх классов нефропротективных препаратов мало. При этом значимой неоднородности снижения риска в зависимости от применения других нефропротективных средств не наблюдалось. Моделирование указывает на существенное увеличение общей выживаемости и выживаемости без больших сердечно-сосудистых и почечных событий при сочетании нескольких нефропротективных препаратов по сравнению с одними RASi.\src{23}

\heading{Рисунок 9. Ведение ХБП при стадиях ССПМ 2--3: текстовая адаптация}
\textbf{Исходная ситуация:} ССПМ-синдром с ХБП --- рСКФ $<60$ и/или UACR $\geq30$~мг/г --- без ССЗ. В исходном блоке рисунка единицы рСКФ не напечатаны; в тексте руководства используется \unitgfr.

\textbf{Ветвь без СД2.}
\begin{itemize}
\item Определить наличие альбуминурии: UACR $\geq30$~мг/г.
\item Если альбуминурии нет: «нефропротективная терапия не требуется» --- формулировка соответствующего блока оригинальной схемы.
\item Если UACR $\geq200$~мг/г: RASi + SGLT2i (COR 1).
\item Если UACR от $\geq30$ до 199~мг/г: RASi (COR 1) и SGLT2i (COR 2a).
\end{itemize}

\textbf{Ветвь с СД2.}
\begin{enumerate}
\item RASi + SGLT2i (COR 1).
\item Проверить, сохраняется ли альбуминурия. Если нет --- продолжить текущую терапию и повторить UACR через 1 год.
\item Если альбуминурия сохраняется --- добавить нефропротективную терапию (COR 1): агонист рецептора GLP-1 (GLP-1 RA; COR 1) либо nsMRA (COR 1). GLP-1 RA отдаётся приоритет при неконтролируемой гипергликемии, ожирении или MASLD.
\item Проверять UACR каждые 3--6 месяцев (COR 2a).
\item Если альбуминурия больше не сохраняется --- продолжить текущую терапию; если сохраняется --- рассмотреть сочетание GLP-1 RA + nsMRA.
\end{enumerate}

\textbf{Примечание авторов к рисунку 9.} На момент подготовки руководства семаглутид --- единственный GLP-1 RA, для которого показано улучшение почечных и сердечно-сосудистых исходов при СД2 и ХБП с альбуминурией.

CKD --- ХБП; CVD --- ССЗ; eGFR --- рСКФ; GLP-1 --- глюкагоноподобный пептид-1; GLP-1 RA --- агонист его рецептора; MASLD --- стеатотическая болезнь печени, ассоциированная с метаболической дисфункцией; nsMRA --- нестероидный антагонист минералокортикоидных рецепторов; RASi --- ингибитор ренин-ангиотензиновой системы; SGLT2i --- ингибитор натрий-глюкозного котранспортера 2 типа; UACR --- отношение альбумина к креатинину мочи. Легенда схемы обозначает классы 1, 2a, 2b, 3 «нет пользы» и 3 «вред»; в действующих блоках этого рисунка используются классы 1 и 2a.

\editorial{Рисунок 9 является укрупнённым алгоритмом. Его блок об отсутствии потребности в нефропротективной терапии относится к показанной ветви без СД2 и без альбуминурии; он не заменяет общие меры ведения ХБП. Пороговые значения рСКФ для начала лечения и UACR для добавления отдельных препаратов приведены в формальных рекомендациях 1--8 и не повторены во всех блоках схемы.}

\heading{Рисунок 10. Кардио- и нефропротективная терапия при ХБП без СД2 в зависимости от UACR: текстовая адаптация}
На горизонтальной оси оригинала отмечены UACR 0, 30, 100 и 200~мг/г.
\begin{itemize}
\item \textbf{RASi:} при ХБП без СД2 UACR $\geq30$~мг/г служит основанием для назначения с целью кардио- и нефропротекции.
\item \textbf{SGLT2i:} при ХБП без СД2 UACR $\geq30$~мг/г обеспечивает умеренную доказательную поддержку применения, а UACR $\geq200$~мг/г --- убедительную поддержку кардио- и нефропротективной пользы.
\end{itemize}
\textbf{Пояснение авторов.} При отсутствии СД2 UACR $\geq30$~мг/г является показанием к нефропротективной терапии. RASi рекомендуются всем пациентам с UACR $\geq30$~мг/г. Убедительные доказательства пользы SGLT2i для сердца и почек имеются при UACR $\geq200$~мг/г; при UACR 30--199~мг/г доказательства более умеренные. CKD --- ХБП; CVD --- ССЗ; RASi --- ингибиторы ренин-ангиотензиновой системы; SGLT2i --- ингибиторы натрий-глюкозного котранспортера 2 типа; T2D --- СД2; UACR --- отношение альбумина к креатинину мочи.

\heading{Рисунок 11. Кардио- и нефропротективная терапия при ХБП с СД2 в зависимости от UACR: текстовая адаптация}
На горизонтальной оси оригинала отмечены UACR 0, 30, 100 и 200~мг/г.
\begin{itemize}
\item \textbf{RASi:} показаны всем пациентам с ХБП и СД2 независимо от UACR.
\item \textbf{SGLT2i:} показаны всем пациентам с ХБП и СД2 независимо от UACR.
\item \textbf{nsMRA:} при ХБП и СД2 на первой линии (RASi и SGLT2i) UACR $\geq30$~мг/г служит основанием для добавления nsMRA с целью кардио- и нефропротекции.
\item \textbf{GLP-1 RA:} при ХБП и СД2 на первой линии (RASi и SGLT2i) UACR $\geq100$~мг/г служит основанием для добавления GLP-1 RA с целью кардио- и нефропротекции.
\end{itemize}
\textbf{Пояснение авторов.} При СД2 RASi и SGLT2i рекомендуются как первая линия всем пациентам с ХБП. Если на этой терапии сохраняется альбуминурия, UACR $\geq30$~мг/г является показанием к добавлению nsMRA, а UACR $\geq100$~мг/г --- к добавлению GLP-1 RA. CKD --- ХБП; GLP-1 --- глюкагоноподобный пептид-1; nsMRA --- нестероидный антагонист минералокортикоидных рецепторов; RASi --- ингибиторы ренин-ангиотензиновой системы; SGLT2i --- ингибиторы натрий-глюкозного котранспортера 2 типа; T2D --- СД2; UACR --- отношение альбумина к креатинину мочи.

\heading{Обоснование отдельных рекомендаций}
\textbf{1. RASi.} Многочисленные РКИ и их метаанализы показали, что RASi уменьшают риск утраты функции почек и почечной недостаточности при ХБП с СД2 либо альбуминурией.\src{1--4} В исследовании RENAAL (снижение частоты конечных точек при инсулиннезависимом диабете с помощью лозартана) лозартан значимо уменьшал первичную комбинированную конечную точку --- удвоение исходного сывороточного креатинина, терминальная почечная недостаточность или смерть. Он также уменьшал частоту удвоения креатинина, терминальной почечной недостаточности и первой госпитализации по поводу СН, но не влиял на смертность у пациентов с СД2 и альбуминурией.\src{24} Аналогично метаанализ РКИ при диабете выявил снижение сердечно-сосудистых событий на RASi по сравнению с плацебо.\src{25} Важно, что в этих метаанализах польза RASi была менее выраженной при сравнении с активным контролем --- другими антигипертензивными препаратами --- как для замедления ХБП, так и для профилактики ССЗ.\src{4,25}

\textbf{2. SGLT2i при СД2 либо выраженной альбуминурии.} Исследования и метаанализы при ХБП с СД2 либо ХБП без СД2, но с UACR $\geq200$~мг/г, показали значимое снижение риска утраты функции почек, почечной недостаточности, сердечно-сосудистой смерти и госпитализации по поводу СН по сравнению с плацебо.\src{5--7} Польза SGLT2i для почечных и сердечно-сосудистых событий наблюдалась независимо от фонового применения RASi\src{26} или GLP-1 RA,\src{10} что указывает на самостоятельный и дополнительный эффект SGLT2i. В исследования обычно включали лиц с рСКФ $\geq20$~\unitgfr, что поддерживает начало терапии в этой популяции; лечение продолжали вплоть до почечной недостаточности, при этом сохранялись признаки почечной и сердечной пользы. Сердечно-сосудистый эффект в целом определялся главным образом уменьшением сердечно-сосудистой смерти, особенно смерти от СН и внезапной сердечной смерти, без значимого влияния на инфаркт миокарда или инсульт.\src{27}

\textbf{3. SGLT2i при UACR 30--199 мг/г без СД2.} В совместном метаанализе SMART-C при альбуминурии, определённой как UACR $\geq30$~мг/г, SGLT2i уменьшали MACE: отношение рисков (HR) 0,90; 95\% ДИ 0,84--0,96; и сердечно-сосудистую смерть: HR 0,80; 95\% ДИ 0,72--0,88.\src{10} Однако доля участников без диабета была небольшой, и все они происходили из EMPA-Kidney --- исследования защиты сердца и почек эмпаглифлозином. В нём для пациентов с UACR 30--199 дополнительно требовалась рСКФ 20--45~\unitgfr.

В EMPA-Kidney снижение риска первичной конечной точки --- прогрессирование болезни почек или сердечно-сосудистая смерть --- на эмпаглифлозине по сравнению с плацебо было более выраженным при более высоком UACR: при UACR $\geq30$--300 HR 0,91 (95\% ДИ 0,65--1,26), а при UACR $\geq300$ HR 0,67 (95\% ДИ 0,58--0,78).\src{6} Вместе с тем после первоначального снижения рСКФ её дальнейшее падение было медленнее на эмпаглифлозине во всех ключевых подгруппах, включая UACR $<30$~мг/г.\src{6} Поэтому SGLT2i можно рассматривать при ХБП без СД2 с UACR 30--199~мг/г, но доказательства совокупной почечной и сердечно-сосудистой пользы здесь менее убедительны, чем при UACR $\geq200$~мг/г.

\textbf{4. Экономическая оценка SGLT2i.} Исследования добавления SGLT2i к прежней стандартной помощи при диабете прогнозировали экономическую эффективность такого подхода с позиции системы здравоохранения США. В модели диабетической болезни почек добавление эмпаглифлозина к стандартной терапии было экономически эффективным по сравнению с одной стандартной терапией.\src{28} В модели использовалась стоимость 3200 долл. США в год --- приблизительно 50\% скидки от оптовой закупочной цены. Она существенно ниже цен в других моделях, но выше согласованной Medicare цены, вводимой в 2026 году.

В высококачественной модели дапаглифлозина при недиабетической ХБП препарат увеличивал ожидаемую продолжительность жизни на 2 года, дисконтированные QALY --- с 6,75 до 8,06 --- и уменьшал число новых случаев почечной недостаточности, требующей заместительной почечной терапии, на протяжении жизни когорты. При цене 4416 долл. США в год добавление дапаглифлозина к стандартной помощи прогнозировалось как экономически эффективное на пожизненном горизонте: ICER 60\,000 долл. США на дополнительный QALY.\src{29} В совокупности эти данные дают высокую степень уверенности в экономической эффективности SGLT2i при ХБП.

\textbf{5. Финеренон.} Исследования nsMRA финеренона у пациентов с ХБП, СД2 и альбуминурией, сохранявшейся несмотря на RASi, при рСКФ $\geq25$~\unitgfr\ показали значимую почечную и сердечно-сосудистую пользу.\src{11,12,30--32} В FIDELIO-DKD (эффективность и безопасность финеренона при СД2 и диабетической болезни почек) и FIGARO-DKD (снижение сердечно-сосудистой смертности и заболеваемости финереноном при диабетической болезни почек) финеренон по сравнению с плацебо значимо уменьшал альбуминурию, утрату функции почек, почечную недостаточность, смерть от болезни почек, СН и сердечно-сосудистую смертность.\src{11--13,30--32}

Неоднородности эффекта в зависимости от применения SGLT2i или GLP-1 RA не выявлено, хотя статистическая мощность для этих подгрупп была недостаточной. В исследования включали только пациентов с исходным калием $\leq4{,}8$~ммоль/л; прекращение лечения из-за гиперкалиемии чаще происходило на финереноне, чем на плацебо.\src{13} Предшествующее применение одновременно RASi и SGLT2i не было обязательным условием исследований финеренона при ХБП. Однако, как отмечено выше, длительный опыт эффективности и безопасности RASi и SGLT2i определяет их место в первой линии, а добавление финеренона рекомендуется для подгруппы высокого риска с СД2 и сохраняющейся альбуминурией на первой линии.

\textbf{6. Экономическая оценка финеренона.} В анализе добавления финеренона к стандартной помощи при диабете и ХБП использовалась годовая стоимость 7300 долл. США. ICER относительно стандартной терапии составил 135\,000 долл. США на дополнительный QALY.\src{33} Следовательно, при пороге 120\,000 долл. США на QALY применение по этому показанию прогнозируется как экономически неэффективное. В вероятностном анализе чувствительности примерно 35\% моделирований дали ICER $<120\,000$ долл. США на QALY. Стоимость годового курса финеренона по Federal Supply Schedule в августе 2025 года составляла 7600 долл. США.\src{34}

\textbf{7. Терапия на основе GLP-1.} Метаанализы исследований GLP-1 RA в подгруппе СД2 с ХБП показали уменьшение риска снижения функции почек, почечной недостаточности, MACE, госпитализации по поводу СН, сердечно-сосудистой и общей смерти по сравнению с плацебо независимо от применения SGLT2i.\src{16} В FLOW --- исследовании семаглутида в сравнении с плацебо при СД2 и ХБП --- семаглутид подкожно по 1,0~мг еженедельно уменьшал риск больших неблагоприятных почечных исходов: диализ, трансплантация, впервые возникшее снижение рСКФ $<15$~\unitgfr, снижение рСКФ на $\geq50\%$ или смерть от почечных причин. Он также уменьшал MACE, госпитализации по поводу СН, сердечно-сосудистую и общую смерть при СД2 и ХБП с альбуминурией.\src{17}

Как и финеренон, семаглутид рекомендуется в качестве дополнительной терапии пациентам высокого риска с СД2, ХБП и сохраняющейся альбуминурией --- $\geq100$~мг/г согласно критериям исследования --- несмотря на первую линию RASi и SGLT2i. С учётом дополнительных физиологических и клинических преимуществ GLP-1 RA, описанных в других разделах руководства, семаглутиду можно отдавать приоритет перед финереноном при неконтролируемой гипергликемии, ожирении или MASLD.

\textbf{8. Экономическая оценка семаглутида.} Авторы не выявили экономических оценок на основе FLOW при диабете и ХБП с сопутствующим АССЗ или без него. Выраженная польза в FLOW предполагает возможность существенного улучшения здоровья населения при широком применении по этому показанию. Однако высокая стоимость семаглутида\src{34} и большое число подходящих для терапии пациентов подчёркивают срочную необходимость строгой экономической оценки.

\subsection{Стадия ССПМ 3}
Стадия ССПМ 3 включает популяцию первичной профилактики высокого риска с субклиническим ССЗ --- АССЗ или СН, --- либо эквивалентами риска: ХБП очень высокого риска по классификации KDIGO или высоким прогнозируемым риском ССЗ (10-летний PREVENT-CVD $>20\%$).\src{1--3} Порог 20\% для высокого 10-летнего риска PREVENT-CVD соответствует частоте событий ССЗ 2\% в год, сопоставимой с обычно наблюдаемой у пациентов с уже имеющимся ССЗ.

Принципы ведения стадии 2 применяются и на стадии 3.\src{1,2,4} Однако существенно больший абсолютный риск обосновывает более интенсивное изменение образа жизни и лекарственное лечение для предотвращения ближайших клинических событий; ожидаемое абсолютное снижение риска больше, чем на стадии 2. Хотя исследования специально не набирали пациентов по признаку стадии ССПМ 3, испытания SGLT2i,\src{5,6} GLP-1 RA\src{7--9} и финеренона,\src{10--13} включавшие сопоставимые по абсолютному риску популяции, показывают малое число пациентов, которых необходимо лечить для предотвращения одного неблагоприятного сердечно-сосудистого или почечного события.

На стадии 3 целесообразность комбинированной лекарственной терапии можно рассматривать более настойчиво: польза может превышать вред, связанный с полипрагмазией, стоимостью и нежелательными эффектами. Моделирование указывает, что сочетания ССПМ-терапии могут давать дополнительную защиту сердца и почек.\src{6}
\editorial{В данном абзаце оригинала напечатано PREVENT-CVD $>20\%$, а затем обсуждается порог 20\%. Знак не заменён редактором на $\geq$; критерии стадирования необходимо сопоставлять с соответствующей таблицей руководства.}

\subsubsection{Использование маркеров субклинического атеросклероза}
\heading{Рекомендация}
Исследования, на которых основана рекомендация, обобщены в таблице доказательств оригинала.

{\small
\begin{longtable}{P{0.07\textwidth}P{0.11\textwidth}P{0.75\textwidth}}
\toprule COR & LOE & Рекомендация\\\midrule
2a & B-NR & \textbf{1.} У взрослых со стадией ССПМ 3 вследствие клинически значимого коронарного кальция (индекс CAC $>100$ либо умеренный или выраженный CAC при качественной оценке) начало или усиление профилактической терапии, показанной при ССПМ-синдроме, может быть полезным для уменьшения сердечно-сосудистых событий.\src{1--4}\\
\bottomrule
\end{longtable}
}

\heading{Обзор}
Субклинический атеросклероз, например повышенный CAC, связан с увеличенным сердечно-сосудистым риском при наличии факторов ССПМ-риска, таких как диабет и ХБП.\src{5--7} CAC отражает суммарную нагрузку атеросклеротическими бляшками и позволяет реклассифицировать риск при добавлении к многофакторным уравнениям.\src{8} Наличие и выраженность CAC могут определять начало или усиление гиполипидемической терапии (раздел 7.1 «Контроль ХС ЛНП при ССПМ-синдроме»), а также других профилактических методов лечения ССПМ-синдрома.

Другие показатели атеросклероза --- лодыжечно-плечевой индекс, толщина комплекса интима--медиа сонных артерий --- показали полезность в оценке риска,\src{9--11} но доказательства их применения для выбора клинической тактики ограничены. Нулевой CAC в общей популяции связан с низким риском событий в кратко- и среднесрочной перспективе, однако при ССПМ-синдроме его информативность может быть ограничена. Особенно это относится к молодым взрослым с ускоренным накоплением CAC вследствие нескольких взаимосвязанных факторов риска и высоким долгосрочным риском ССЗ. Это указывает на необходимость раннего начала коррекции факторов риска.\src{12--16} При ХБП CAC также связан с повышенным сердечно-сосудистым риском и требует интенсивной его коррекции;\src{17} интерпретацию следует индивидуализировать у пациентов на диализе или после трансплантации.

\heading{Обоснование рекомендации}
\textbf{1.} Исследования, использующие оценки относительного снижения риска из РКИ и абсолютную частоту событий из эпидемиологических исследований CAC, показывают, что умеренный или выраженный CAC выявляет пациентов с высоким абсолютным риском ССЗ, наиболее вероятно получающих пользу от нескольких профилактических методов лечения ССПМ-синдрома: статинов, терапии на основе GLP-1, интенсивного снижения АД и икозапента этила при гипертриглицеридемии.\src{1--4} Хотя применительно к CAC это ещё не изучено, те же принципы, вероятно, относятся к SGLT2i, RASi и nsMRA.

CAC $>100$ выявляет пациентов, наиболее вероятно получающих пользу от профилактики благодаря высокому абсолютному риску и, следовательно, малому числу нуждающихся в лечении для предотвращения одного события. При более высоком CAC это число ещё меньше. Альтернативно CAC $\geq75$-го перцентиля для возраста и пола может указывать на повышенный риск у женщин и в более молодых популяциях даже при меньшем абсолютном CAC.\src{18,19}

Важно, что эти работы являются моделированием результатов РКИ применительно к эпидемиологическим данным, где начало профилактической терапии в ходе наблюдения влияет на частоту событий.\src{20} Они не определяют последовательность добавления препаратов. Случайно обнаруженный умеренный или выраженный коронарный кальциноз на КТ грудной клетки без ЭКГ-синхронизации обычно соответствует CAC $\geq100$, имеет сходную прогностическую ценность и также выявляет пациентов с ССПМ-синдромом, наиболее вероятно получающих пользу от профилактической терапии.\src{21--24}

\subsubsection{Доклиническая сердечная недостаточность}
\heading{Рекомендации}
Исследования, на которых основаны рекомендации, обобщены в таблице доказательств оригинала.

{\small
\begin{longtable}{P{0.07\textwidth}P{0.11\textwidth}P{0.75\textwidth}}
\toprule COR & LOE & Рекомендация\\\midrule
\endfirsthead
\toprule COR & LOE & Рекомендация (продолжение)\\\midrule
\endhead
1 & B-R & \textbf{1.} У взрослых со стадией ССПМ 3 вследствие доклинической СН рекомендуются интенсивное изменение образа жизни и интенсивный контроль факторов риска для улучшения ССПМ-здоровья и предотвращения перехода к клинически выраженной СН.\src{1,2}\\[5pt]
1 & B-NR & \textbf{2.} У взрослых со стадией ССПМ 3 вследствие доклинической СН в сочетании с СД2 или ХБП рекомендуется SGLT2i как терапия первой линии для профилактики СН.\src{3--6}\\[5pt]
2a & B-NR & \textbf{3.} У взрослых со стадией ССПМ 3 вследствие доклинической СН с диабетом, ХБП и UACR $\geq100$~мг/г целесообразно добавлять к SGLT2i терапию на основе GLP-1 с доказанной пользой для профилактики СН.\src{7--10}\\[5pt]
2a & B-NR & \textbf{4.} У взрослых со стадией ССПМ 3 вследствие доклинической СН с диабетом, ХБП и UACR $\geq30$~мг/г целесообразно добавлять nsMRA к SGLT2i для профилактики СН.\src{7--10}\\
\bottomrule
\end{longtable}
}

\heading{Обзор}
Доклиническая, или субклиническая, СН определяется повышением сердечных биомаркеров либо признаками сердечной дисфункции при визуализации в отсутствие симптомов СН (таблица 16) и входит в стадию ССПМ 3. Она весьма распространена уже к среднему возрасту, особенно при факторах ССПМ-риска. Наибольший риск впервые возникшей СН отмечается при сочетании отклонений биомаркеров и эхокардиографии.\src{11,12} Однако исследований, специально включавших пациентов с диагностированной доклинической СН и нацеленных на её профилактику, мало; основная доказательная база --- анализы post hoc.\src{13}

Главная цель ведения --- уменьшить прогрессирование до клинической СН (рисунок 12). Как отмечено в разделе 3.1 «Диагностический подход к стадированию ССПМ», исследование сердечных биомаркеров можно выполнять при риске PREVENT-HF $\geq5\%$. При доклинической СН визуализация сердца уточняет абсолютный риск и может дополнительно служить основанием для направления на согласованное ведение с целью предотвращения СН. Профилактика должна включать изменение образа жизни, усиление контроля факторов риска и применение доказательной ССПМ-терапии при наличии показаний. Помимо лечения факторов риска, при систолической дисфункции левого желудочка рекомендуются бета-блокаторы и ACEi/ARB.\src{14,15}

Рекомендации по отдельным состояниям, способным приводить к доклинической СН, обсуждаются в других документах: наследственный риск кардиомиопатии,\src{16} кардиотоксические воздействия,\src{17} врождённые пороки сердца,\src{18} а также коронарная болезнь,\src{19} клапанные пороки\src{20} и фибрилляция предсердий (ФП).\src{21}

\heading{Таблица 16. Определение доклинической СН}
\textbf{Эхокардиографические и биомаркерные пороги.}
{\small
\begin{longtable}{P{0.35\textwidth}P{0.59\textwidth}}
\toprule Категория & Порог\\\midrule
\endfirsthead
\toprule Категория & Порог (продолжение)\\\midrule
\endhead
Структурное заболевание сердца & Индекс объёма левого предсердия (LAVI) $\geq29$~мл/м$^2$\\
 & Индекс массы левого желудочка (LVMI) $>116$~г/м$^2$ у мужчин; $>95$~г/м$^2$ у женщин\\
 & Относительная толщина стенки (RWT) $>0{,}42$\\
 & Толщина стенки левого желудочка $\geq12$~мм\\
 & Фракция выброса левого желудочка (LVEF) $<50\%$\\
 & Глобальная продольная деформация (GLS) $<16\%$\\
 & Септальная скорость $e'$ $<7$~см/с\\
 & Скорость трикуспидальной регургитации (TR) $>2{,}8$~м/с\\
 & Скорость трикуспидальной регургитации (TR) $>2{,}8$~м/с --- повторная строка оригинала\\[4pt]
Неинвазивная визуализация для оценки давления наполнения & Расчётное систолическое давление в лёгочной артерии $>35$~мм рт. ст.\\
 & Среднее $E/e'$ $\geq15$\\[4pt]
Сердечные биомаркеры\textsuperscript{†} & BNP $\geq35$~пг/мл\textsuperscript{‡}\\
 & NT-proBNP $\geq125$~пг/мл\textsuperscript{†}\\
\bottomrule
\end{longtable}
}
\textbf{Примечания исходной таблицы 16.} Адаптировано с разрешения Heidenreich и соавторов.\src{17} © 2022 American Heart Association, Inc. и American College of Cardiology Foundation.

*Доклиническая СН, также называемая стадией B СН в руководстве 2022 года, определяется отсутствием симптомов и признаков СН при наличии одного из следующего: (1) структурное заболевание сердца; (2) неинвазивные признаки повышенного давления наполнения; (3) факторы риска с повышенными сердечными биомаркерами при отсутствии другого диагноза, объясняющего повышение, например острого коронарного синдрома, ХБП\textsuperscript{†}, тромбоэмболии лёгочной артерии или миоперикардита.

\textsuperscript{†}В руководстве ACC/AHA 2022 года по СН указаны только BNP/NT-proBNP; универсальное определение СН, напротив, включает стойкое повышение тропонина выше 99-го перцентиля в нормальной референсной популяции.

\textsuperscript{‡}Указанные пороги натрийуретических пептидов могут иметь меньшую специфичность, особенно у пожилых пациентов, при ФП или ХБП.

BNP --- натрийуретический пептид B-типа; GLS --- глобальная продольная деформация; LAVI --- индекс объёма левого предсердия; LV --- левый желудочек; LVEF --- его фракция выброса; LVMI --- индекс его массы; RWT --- относительная толщина стенки.
\editorial{Повтор строк скорости TR, отсутствие отдельного названия категории у части эхокардиографических показателей и размещение знаков сносок †/‡ воспроизводят исходную таблицу. Повтор не означает дополнительного диагностического критерия.}

\heading{Рисунок 12. Обследование и ведение при доклинической СН и ССПМ-синдроме: текстовая адаптация}
\begin{enumerate}
\item Выявить ССПМ-синдром без симптомной клинической СН, но с повышенным ближайшим риском её развития: PREVENT-HF $\geq5\%$.
\item Определить сердечные биомаркеры --- натрийуретические пептиды и/или высокочувствительный тропонин --- для выявления доклинической СН.
\item Если сердечные биомаркеры повышены: оптимизировать контроль факторов риска; рассмотреть визуализацию сердца для уточнения риска; начать профилактическое лечение для уменьшения риска СН.
\end{enumerate}
\textbf{Примечание авторов.} Высокочувствительный тропонин можно рассматривать при ожирении, когда натрийуретические пептиды могут быть низкими даже при сердечной дисфункции и повышенном риске СН. Обследование на доклиническую СН при ССПМ-синдроме показано при повышенном прогнозируемом риске: PREVENT-HF $\geq5\%$. CKM --- ССПМ; HF --- СН; PREVENT --- Predicting Risk of Cardiovascular Disease Events, прогнозирование риска сердечно-сосудистых событий.

\heading{Обоснование отдельных рекомендаций}
\textbf{1. Образ жизни и факторы риска.} На всех стадиях ССПМ имеются убедительные доказательства улучшения ССПМ-здоровья при интенсивном изменении образа жизни и контроле факторов риска в рамках согласованной помощи с лекарственной терапией.\src{22--24} Крупных рандомизированных исследований поведенческих вмешательств при доклинической СН мало. В небольшом одноцентровом исследовании 46 пациентов с доклинической СН --- толщина межжелудочковой перегородки $>11$~мм и либо NT-proBNP $>40$~пг/мл, либо hs-cTnT $>0{,}6$~пг/мл --- высокоинтенсивные тренировки в течение года улучшали механику миокарда по сравнению с контролем.\src{1}

Анализ post hoc Look AHEAD (действия для здоровья при диабете), включавшего взрослых 45--76 лет с избыточной массой тела или ожирением и диабетом, выявил значимое снижение hs-cTn через 1 и 4 года, связанное с существенным снижением массы тела ($\geq5\%$), по сравнению с минимальным снижением ($<5\%$).\src{2} В той же популяции интенсивное изменение образа жизни было связано с меньшим риском СН с сохранённой фракцией выброса (HFpEF) у лиц с высоким исходным NT-proBNP либо со стабильным или снижающимся NT-proBNP в течение года.\src{2,25} Наблюдательные когорты также показывают связь большей физической активности с меньшим риском перехода к СН при повышенных сердечных биомаркерах.\src{26}

\textbf{2. SGLT2i.} При диабете получены убедительные доказательства первичной профилактики СН с помощью SGLT2i. Метаанализ трёх исследований --- EMPA-REG, CANVAS (оценка сердечно-сосудистых эффектов канаглифлозина) и DECLARE-TIMI 58 (влияние дапаглифлозина на сердечно-сосудистые события) --- включал 30\,431 пациента с СД2 без СН и показал уменьшение впервые возникшей госпитализации по поводу СН на 21\%: HR 0,79; 95\% ДИ 0,71--0,88.\src{6}

Эти исследования специально не включали пациентов с доклинической СН и не анализировали такую подгруппу. Однако последующие биомаркерные анализы выявили высокие медианные NT-proBNP и hs-cTn; более чем у 50\% участников уровни соответствовали доклинической СН.\src{27--30} В небольшом исследовании 44 пациентов с диабетом и нарушенной глобальной продольной деформацией рандомизация на эмпаглифлозин сопровождалась значимым улучшением деформации: увеличение на 2\% через 6 месяцев.\src{31} Сходные результаты получены в других небольших механистических исследованиях SGLT2i.\src{32--34}

\textbf{3. GLP-1 RA.} Метаанализ пяти исследований при диабете --- четырёх более старых, менее мощных GLP-1 RA и одного семаглутида --- не выявил значимого уменьшения госпитализаций по поводу СН.\src{35} Более поздний метаанализ десяти исследований GLP-1 RA оценил снижение на 13\%: HR 0,87; 95\% ДИ 0,80--0,94.\src{36}

В анализах post hoc SURPASS 4 --- сравнения тирзепатида раз в неделю с инсулином гларгином раз в сутки у пациентов с СД2 и повышенным сердечно-сосудистым риском --- у участников с ССЗ или высоким его риском, доклинической СН по повышенному NT-proBNP и без госпитализаций по поводу СН в анамнезе NT-proBNP значимо снижался от исходного уровня к 52-й неделе.\src{37}

В единственном специализированном исследовании GLP-1 RA при диабете и альбуминурической ХБП (FLOW) семаглутид значимо уменьшал события СН во всей выборке: HR 0,74; 95\% ДИ 0,58--0,94. Сходное снижение впервые возникшей СН отмечено у пациентов без исходной СН: HR 0,68; 95\% ДИ 0,50--0,91.\src{8,9} Для предотвращения одного события СН за 3 года требовалось лечить 49 пациентов.\src{9} При диабете и ХБП моделирование предполагает дополнительную пользу добавления GLP-1 RA к SGLT2i: снижение госпитализаций по поводу СН на 43\% (HR 0,57; 95\% ДИ 0,47--0,70), если предполагать пожизненное применение терапии.\src{38}

\textbf{4. nsMRA.} В объединённых анализах FIDELITY при диабете и альбуминурической ХБП финеренон уменьшал СН на 21\%: HR 0,79; 95\% ДИ 0,66--0,92.\src{7} Сходная польза отмечалась без СН в анамнезе\src{10} и при исходной гипертрофии левого желудочка.\src{39} Эффект сохранялся с фоновыми SGLT2i\src{40} или GLP-1 RA\src{41} и без них.

Эти данные указывают на целесообразность второй линии --- GLP-1 RA с доказанной сердечно-сосудистой пользой либо nsMRA --- для профилактики СН при диабете и альбуминурической ХБП после начала SGLT2i либо при противопоказаниях к SGLT2i.\src{42} Однако в исследованиях GLP-1 RA и nsMRA при диабете и ХБП большая доля участников имела исходное АССЗ: 46\% в FIDELITY и 23\% в FLOW. Модель добавления nsMRA к SGLT2i предполагала последовательную дополнительную пользу со снижением госпитализаций по поводу СН на 50\% (HR 0,50; 95\% ДИ 0,39--0,64), но основывалась на пожизненном применении и, вероятно, в значительной мере отражала нефропротекцию.\src{38}

\section{Ведение пациентов с ССЗ при ССПМ-синдроме}
\subsection{Стадия ССПМ 4}
\heading{Обзор}
Стадия ССПМ 4 определяется наличием клинического ССЗ --- АССЗ, СН или ФП --- у людей с избытком или дисфункцией жировой ткани, другими метаболическими факторами риска либо ХБП. Дополнительно выделяют отсутствие почечной недостаточности (стадия 4a) и её наличие (стадия 4b), что может влиять на тактику ведения.

Вследствие сложных разнонаправленных взаимосвязей факторы ССПМ-риска усугубляют неблагоприятные исходы при установленном ССЗ. На стадии 4 ведение направлено на оптимизацию вторичной профилактики ССЗ с учётом этих факторов. Для всех пациентов подчёркивается важность интенсивной комплексной коррекции образа жизни и факторов риска, достижения и поддержания здоровой массы тела и оптимального контроля клинических факторов риска.

Следующие разделы содержат рекомендации для стадии 4 с акцентом на сочетания ССЗ с ожирением, СД2 и ХБП. Повышенный риск сердечно-сосудистых и почечных событий и смерти подчёркивает необходимость согласованной междисциплинарной помощи. Важно учитывать дополнительные руководства: по контролю гипертензии\src{1} и холестерина\src{2} для вторичной профилактики АССЗ;\src{3} по терапии, соответствующей клиническим рекомендациям (GDMT), при хронической коронарной болезни,\src{3} заболевании периферических артерий,\src{4} инсульте\src{5} и СН;\src{6} а также по стратегиям контроля частоты сокращений и ритма при ФП.\src{7}

\subsection{Ведение пациентов с АССЗ}
\heading{Обзор}
Цель ведения АССЗ --- одного из фенотипов стадии ССПМ 4 --- улучшить функциональное состояние и качество жизни, оптимизировать вторичную профилактику и предотвратить повторные события АССЗ (рисунок 13).\src{1} Оптимальное ведение хронической коронарной болезни,\src{2} заболевания периферических артерий\src{3} и цереброваскулярной болезни\src{4} изложено в соответствующих руководствах AHA/ACC.

В них подчёркивается применение антиагрегантов и высокоинтенсивной либо максимально переносимой терапии статинами в первой линии. При наличии показаний рассматривают дополнительные гиполипидемические препараты для более интенсивного снижения ХС ЛНП либо коррекции триглицеридов.\src{5} При заболевании периферических артерий дополнительно рассматривают ривароксабан с низкой дозой аспирина для предотвращения MACE и событий со стороны конечностей у лиц без повышенного риска кровотечений.\src{3}

При сопутствующей гипертензии показан интенсивный контроль АД с целью $<130/80$~мм рт. ст. Выбор антигипертензивного препарата определяется сопутствующими состояниями: например, ACEi/ARB отдают приоритет при диабете или ХБП.\src{6} Кроме того, при АССЗ часто присутствуют ожирение (раздел 6.2.1), диабет (раздел 6.2.2) и/или ХБП (раздел 6.2.3), требующие особых подходов, подробно описанных далее.

\heading{Рисунок 13. Общая схема ведения АССЗ на стадии ССПМ 4: текстовая адаптация}
\heading{Применять первую линию GDMT при АССЗ}
\begin{itemize}
\item Аспирин; при наличии показаний --- дополнительный антиагрегант, антикоагулянт либо оба.
\item Максимально переносимая терапия статинами; при наличии показаний --- дополнительные или альтернативные препараты: эзетимиб, бемпедоевая кислота, ингибиторы PCSK9.
\end{itemize}

\heading{Применять общие принципы ведения ССПМ-синдрома}
\begin{itemize}
\item Укреплять сердечно-сосудистое здоровье на основе Life's Essential 8.
\item Организовать междисциплинарную помощь, включая выделенного координатора ССПМ-помощи.
\item Работать с неблагоприятными социальными детерминантами здоровья (SDOH), привлекая работников общественного здравоохранения (CHW), социальных работников и ресурсы сообщества.
\item Контролировать АД согласно действующим рекомендациям.
\end{itemize}

\heading{Выбирать ведение с учётом профиля риска: ожирение}
\begin{itemize}
\item Интенсивно изменять образ жизни для снижения массы тела и улучшения профиля ССПМ-риска.
\item Начать терапию на основе GLP-1 совместно с консультированием по образу жизни для уменьшения событий ССЗ.
\item Если цели снижения массы тела не достигнуты, рассмотреть метаболическую и бариатрическую хирургию для улучшения профиля ССПМ-риска и уменьшения событий ССЗ.
\end{itemize}

\heading{Выбирать ведение с учётом профиля риска: СД2}
\begin{itemize}
\item Поддерживать изменение образа жизни для улучшения профиля ССПМ-риска.
\item Начать терапию на основе GLP-1 или SGLT2i для уменьшения событий ССЗ и сердечно-сосудистой смертности.
\item Можно рассматривать совместное применение терапии на основе GLP-1 и SGLT2i.
\end{itemize}

\heading{Выбирать ведение с учётом профиля риска: ХБП}
\begin{itemize}
\item При АССЗ с СД2 либо альбуминурией использовать RASi и SGLT2i как первую линию для улучшения сердечно-сосудистых и почечных исходов.
\item При АССЗ с СД2 и сохраняющейся альбуминурией, несмотря на первую линию, рассмотреть добавление терапии на основе GLP-1 или финеренона для улучшения сердечно-сосудистых и почечных исходов.
\end{itemize}

\textbf{Пояснение авторов к рисунку 13.} Ведение пациентов с АССЗ на стадии ССПМ 4 должно сочетать первую линию GDMT при АССЗ, общие принципы ведения ССПМ-синдрома и дополнительные лечебные подходы к отдельным факторам ССПМ-риска, показанные на схеме. CHW --- работник общественного здравоохранения; CKM --- ССПМ; CVD --- ССЗ; GDMT --- лечение согласно клиническим рекомендациям; GLP-1 --- глюкагоноподобный пептид-1; RASi --- ингибиторы ренин-ангиотензиновой системы; SDOH --- социальные детерминанты здоровья; SGLT2i --- ингибиторы натрий-глюкозного котранспортера 2 типа; T2D --- СД2.

\subsubsection{Стадия ССПМ 4 с ожирением и АССЗ}
\heading{Рекомендации}
Исследования, на которых основаны рекомендации, обобщены в таблице доказательств оригинала.

{\small
\begin{longtable}{P{0.07\textwidth}P{0.11\textwidth}P{0.75\textwidth}}
\toprule COR & LOE & Рекомендация\\\midrule
\endfirsthead
\toprule COR & LOE & Рекомендация (продолжение)\\\midrule
\endhead
1 & A & \textbf{1.} У взрослых со стадией ССПМ 4, избыточной массой тела или ожирением и АССЗ рекомендуется интенсивное многокомпонентное поведенческое вмешательство по изменению образа жизни для снижения массы тела и улучшения факторов ССПМ-риска.\src{1--5}\\[5pt]
1 & B-R & \textbf{2.} У взрослых со стадией ССПМ 4, избыточной массой тела или ожирением (ИМТ $\geq27$~кг/м$^2$) и АССЗ рекомендуется терапия на основе GLP-1 с доказанной сердечно-сосудистой пользой в дополнение к консультированию по здоровому питанию и регулярной физической активности для уменьшения риска сердечно-сосудистых событий.\src{6,7}\\[5pt]
\multicolumn{2}{P{0.18\textwidth}}{Экономическая оценка} & \textbf{3. Экономически эффективно; умеренная степень уверенности.} У взрослых со стадией ССПМ 4, избыточной массой тела или ожирением (ИМТ $\geq27$~кг/м$^2$) и АССЗ терапия на основе GLP-1 с доказанной сердечно-сосудистой пользой при ценах США 2025 года, согласно прогнозу, экономически эффективна по сравнению с обычной помощью.\\[5pt]
2a & B-NR & \textbf{4.} У взрослых со стадией ССПМ 4, ожирением (ИМТ $\geq30$~кг/м$^2$) и АССЗ, не достигших целей снижения массы тела изменением образа жизни с максимально переносимой лекарственной терапией или без неё, метаболическая и бариатрическая хирургия (MBS) может быть полезной для снижения массы тела как минимум на 5--10\%, улучшения ССПМ-здоровья и уменьшения риска сердечно-сосудистых событий.\src{8,9}\\[5pt]
3: вред & B-R & \textbf{5.} У взрослых со стадией ССПМ 4, избыточной массой тела или ожирением (ИМТ $\geq27$~кг/м$^2$) и АССЗ лечение налтрексоном/бупропионом либо препаратами, содержащими фентермин, потенциально вредно, поскольку они могут повышать АД и частоту сердечных сокращений.\src{10,11}\\
\bottomrule
\end{longtable}
}
\textbf{Примечание оригинала.} Экономические положения предназначены для решений на уровне населения и системы здравоохранения, а не для непосредственного определения лечения отдельных пациентов.

\heading{Обзор}
При избыточной массе тела или ожирении и установленном АССЗ (стадия ССПМ 4) лечение избытка массы тела может включать интенсивное изменение образа жизни, лекарственную терапию и бариатрическую хирургию (рисунок 14). Подход определяется профилем ССПМ-риска, личными предпочтениями и доступностью лечения. Снижение массы тела примерно на 5\% при интенсивном изменении образа жизни существенно улучшает несколько метаболических факторов риска у пациентов с ожирением.\src{12,13}

Вмешательства для снижения массы тела можно проводить во время или сразу после кардиореабилитации при АССЗ с избыточной массой тела или ожирением, хотя эффективность этой стратегии требует дополнительных доказательств.\src{14} Возможности снижения массы тела одним изменением образа жизни ограничены; поэтому для уменьшения массы тела и риска ССЗ при АССЗ может потребоваться добавление лекарственной терапии ожирения.

В SELECT (влияние семаглутида на сердечно-сосудистые исходы при избыточной массе тела или ожирении) подкожный семаглутид у пациентов с ИМТ $\geq27$~кг/м$^2$ и АССЗ без СД2 значимо уменьшал массу тела и MACE;\src{6} часть сердечно-сосудистой пользы не зависела от снижения массы тела.\src{8,9,15} При ИМТ $\geq30$~кг/м$^2$ и АССЗ направление на бариатрическую операцию может быть полезным, если интенсивное изменение образа жизни и лекарственная терапия ожирения не дают клинически значимого снижения массы тела: наблюдательные исследования с подобранными группами сравнения показывают существенное улучшение профиля ССПМ-риска и исходов ССЗ.\src{8,9,15}

\heading{Рисунок 14. Ведение сочетания ожирения и АССЗ при ССПМ-синдроме: текстовая адаптация}
\begin{enumerate}
\item Исходная ситуация --- ожирение и АССЗ.
\item Консультировать по снижению массы тела без осуждения, обсуждая изменение образа жизни.
\item Применять многокомпонентное интенсивное вмешательство по изменению образа жизни (COR 1) \textbf{и} терапию на основе GLP-1 (COR 1).
\item Оценить достаточность снижения массы тела: в блоке оригинала указано $>5$--10\%.
\item Если снижение достаточное --- постоянно поддерживать изменение образа жизни и приверженность лекарствам для сохранения результата (COR 1).
\item Если снижение недостаточное --- рассмотреть метаболическую и бариатрическую хирургию (COR 2a), затем постоянно поддерживать изменение образа жизни для сохранения сниженной массы тела (COR 1).
\end{enumerate}

\textbf{Примечание авторов к рисунку 14.} Сердечно-сосудистые преимущества семаглутида при ожирении и АССЗ, по-видимому, в значительной мере не зависят от достигнутого снижения массы тела. Это может обосновывать продолжение терапии на основе GLP-1 даже при недостаточном снижении массы тела. ASCVD --- АССЗ; CKM --- ССПМ; COR --- класс рекомендации; GLP-1 --- глюкагоноподобный пептид-1. Легенда оригинальной схемы включает классы 1, 2a, 2b, 3 «нет пользы» и 3 «вред»; действующие блоки отмечены классами 1 и 2a.

\heading{Обоснование отдельных рекомендаций}
\textbf{1. Изменение образа жизни.} Систематический обзор вмешательств на основе образа жизни, большинство из которых оценивалось в РКИ, показал более высокую вероятность клинически значимого снижения массы тела на 5\% при поведенческих вмешательствах по сравнению с контролем.\src{2} Look AHEAD при избыточной массе тела/ожирении и СД2, с ССЗ и без него, не выявило влияния интенсивного изменения образа жизни на сердечно-сосудистые события. Возможные объяснения --- умеренное снижение массы тела и широкое применение статинов в контрольной группе.\src{16}

Несмотря на отсутствие снижения событий ССЗ, изменение поведения, более здоровое питание и физическая активность, приводящие к снижению массы тела на 5--10\%, улучшают факторы ССПМ-риска, включая ремиссию СД2.\src{17,18} Польза наблюдалась в разных популяциях и географических условиях.\src{19} Метаанализ исследований первичной и вторичной профилактики ССЗ показал связь средиземноморского питания с уменьшением риска событий ССЗ на 38\%.\src{4}

К сожалению, изменение образа жизни не устраняет компенсаторные нейрометаболические механизмы, активируемые при снижении массы тела: например, уменьшение лептина жировой ткани и увеличение желудочного грелина. Они часто способствуют возврату массы тела из-за усиления голода и изменения метаболизма. Поэтому нередко нужны дополнительные методы лечения для более значительного и устойчивого снижения массы тела.\src{3}

\textbf{2. Семаглутид.} В SELECT\src{6} включили 17\,604 участника с ИМТ $\geq27$~кг/м$^2$ и установленным АССЗ --- перенесённым инфарктом миокарда, инсультом или симптомным заболеванием периферических артерий, --- без СД2, терминальной болезни почек или диализа. Семаглутид 2,4~мг уменьшал риск первичной комбинированной конечной точки --- сердечно-сосудистая смерть, нефатальный инфаркт или нефатальный инсульт --- на 20\% по сравнению с плацебо.

За 4 года наблюдения семаглутид также сопровождался средним снижением массы тела на 9\%, значимым улучшением кардиометаболических показателей, системного воспаления и уменьшением сердечно-сосудистых событий независимо от исходного HbA1c.\src{3,20} Важно, что снижение риска ССЗ возникало до клинически значимого снижения массы тела, а уменьшение событий было сходным независимо от величины достигнутого снижения массы тела.\src{21} Это указывает на наличие части сердечно-сосудистых преимуществ, не связанных с массой тела. Хотя семаглутид отменяли чаще, главным образом из-за желудочно-кишечных эффектов, общая частота серьёзных нежелательных явлений была ниже, чем на плацебо.

\textbf{3. Экономическая оценка семаглутида.} Модель оценивала семаглутид для вторичной профилактики АССЗ в национально репрезентативной когорте взрослых США, подходящих для лечения.\src{22} По сравнению с одной обычной помощью его добавление примерно 4 миллионам подходящих взрослых без диабета могло предотвратить около 360\,000 MACE. При чистой цене США 2023 года 8604 долл. США ICER составил 148\,100 долл. США на дополнительный QALY; 95\% интервал неопределённости 127\,100--173\,400.

В анализе чувствительности при дополнительном снижении годовой стоимости на 18\%, до 7055 долл. США, семаглутид становился экономически эффективным при пороге 120\,000 долл. США на QALY (рисунок 15). Средняя годовая стоимость за предшествующий год снизилась, а цена прямой покупки потребителем за собственные средства --- 5988 долл. США --- уже ниже пороговой цены экономической эффективности.

\heading{Рисунок 15. ICER семаглутида для вторичной профилактики ССЗ у взрослых США с избыточной массой тела/ожирением без диабета: текстовая адаптация}
\textbf{Оси графика:} по горизонтали --- годовая стоимость семаглутида в долларах США; по вертикали --- инкрементальное соотношение «затраты--эффективность» (ICER), долл. США на дополнительный год жизни с поправкой на качество (QALY). ICER увеличивается с ростом цены препарата.

На линии обозначены четыре варианта цены США: чистая цена для лечения диабета; цена прямой покупки за собственные средства для лечения ожирения; чистая цена для лечения ожирения; прейскурантная цена для лечения ожирения. Пунктиром показан порог экономической эффективности 120\,000 долл. США на QALY. Точные численные значения, приведённые в подписи оригинала, переданы далее; приблизительные координаты остальных точек не превращены в новые точные оценки.

{\small
\begin{longtable}{P{0.49\textwidth}P{0.44\textwidth}}
\toprule Порог экономической эффективности, долл. США на QALY & Годовая цена семаглутида, при которой достигается порог, долл. США\\\midrule
50\,000 & 3295\\
100\,000 & 5940\\
120\,000 & 7055\\
150\,000 & 8710\\
\bottomrule
\end{longtable}
}

\textbf{Пояснение авторов.} При чистой цене США 2023 года, после скидок и возвратов, 8604 долл. США ICER инъекционного семаглутида относительно обычной помощи составляет 148\,100 долл. США на QALY. Вертикальные линии указывают цены достижения порогов из таблицы. Если бы текущая цена самостоятельной оплаты --- 5988 долл. США в год, то есть цена для пациентов, оплачивающих лечение сами, а не через страхование, --- была доступна всем, терапия была бы экономически эффективной для вторичной профилактики ССЗ: ICER 99\,610 долл. США на QALY.

Адаптировано с разрешения Hennessy и соавторов.\src{22} © 2026 American Medical Association. Все права защищены, включая права на анализ текста и данных, обучение ИИ и аналогичные технологии. CVD --- ССЗ; ICER --- инкрементальное соотношение «затраты--эффективность»; QALY --- год жизни с поправкой на качество; USD --- доллары США.
\editorial{Цены и слова «текущая», «за предшествующий год» относятся к исходному руководству и его экономическим данным, а не к независимо проверенным ценам на дату перевода. График не используется для пересчёта стоимости лечения конкретного пациента.}

\textbf{4. Метаболическая и бариатрическая хирургия.} При АССЗ и ожирении (ИМТ $\geq30$~кг/м$^2$), если цели снижения массы тела не достигнуты изменением образа жизни и лекарствами, MBS, включая желудочное шунтирование по Ру (RYGB) или лапароскопическую рукавную гастрэктомию (LSG), может приносить пользу. Она обеспечивает значительное снижение массы тела --- до 30--35\% --- и улучшает несколько факторов ССПМ-риска.\src{23--25}

В анализе пациентов с перенесённым инфарктом из регистра SWEDEHEART (в исходнике название раскрыто как «Интенсивное раннее и устойчивое снижение холестерина не-ЛВП после инфаркта миокарда и прогноз») после MBS сердечно-сосудистая смерть и повторный инфаркт наблюдались реже, чем у подобранных контролей, но различия по инсульту отсутствовали: HR смерти 0,45 (95\% ДИ 0,29--0,70); инфаркта 0,24 (0,14--0,41); инсульта 0,91 (0,38--2,20).\src{8}

В общей популяции частота осложнений MBS низкая, а результаты сопоставимы с лапароскопической холецистэктомией.\src{26} Однако абсолютный риск сердечно-сосудистых и внесердечных периоперационных осложнений существенно выше при АССЗ, чем без него.\src{27,28} В отсутствие рандомизированных данных соотношение абсолютной пользы и риска определено не полностью, хотя у пациентов со стабильным АССЗ и низким хирургическим риском MBS, вероятно, безопасна и эффективна. Нужны долгосрочные данные о мальабсорбции, камнях в почках и других осложнениях при АССЗ.\src{29}

\textbf{5. Потенциально вредные препараты.} Систематический обзор и метаанализ 154 плацебо-контролируемых РКИ лекарственной терапии ожирения при избыточной массе тела или ожирении показал, что налтрексон связан со значимым повышением систолического и диастолического АД как отдельно, так и вместе с бупропионом. При монотерапии взвешенная средняя разность систолического АД составила 2,70~мм рт. ст. (95\% ДИ 2,00--3,70), диастолического --- 4,00~мм рт. ст. (3,49--4,51). При сочетании с бупропионом соответствующие разности составляли 2,01~мм рт. ст. (1,26--2,76) и 1,25~мм рт. ст. (0,52--1,97).

Фентермин при отдельном применении сопровождался значимо более высокой частотой сердечных сокращений: взвешенная средняя разность 4,20; 95\% ДИ 0,46--7,94.\src{11} С учётом этих неблагоприятных сердечно-сосудистых эффектов налтрексон/бупропион и фентерминсодержащие средства не рекомендуются пациентам с ССЗ.\src{10}
\editorial{В оригинале для разности частоты сердечных сокращений при фентермине ошибочно напечатана единица «mm Hg» (мм рт. ст.). Числа сохранены; единица не заменена без оговорки на уд/мин. Эту строку следует сверять с цитируемым исследованием перед клиническим использованием.}

\subsubsection{Стадия ССПМ 4 с СД2 и АССЗ}
\heading{Рекомендации}
Исследования, на которых основаны рекомендации, обобщены в таблице доказательств оригинала.

{\small
\begin{longtable}{P{0.07\textwidth}P{0.11\textwidth}P{0.75\textwidth}}
\toprule COR & LOE & Рекомендация\\\midrule
\endfirsthead
\toprule COR & LOE & Рекомендация (продолжение)\\\midrule
\endhead
1 & A & \textbf{1.} Взрослым со стадией ССПМ 4, СД2 и АССЗ рекомендуется индивидуальный план изменения образа жизни с акцентом на полезное для сердца питание и физическую активность для улучшения гликемического контроля, достижения и поддержания здоровой массы тела и улучшения других факторов риска АССЗ.\src{1,2}\\[5pt]
1 & A & \textbf{2.} Взрослым со стадией ССПМ 4, СД2 и АССЗ рекомендуется SGLT2i либо терапия на основе GLP-1 с доказанной сердечно-сосудистой пользой для уменьшения риска сердечно-сосудистых событий и сердечно-сосудистой смертности.\src{3}\\[5pt]
\multicolumn{2}{P{0.18\textwidth}}{Экономическая оценка} & \textbf{3. Экономически эффективно; умеренная степень уверенности.} У взрослых со стадией ССПМ 4, СД2 и АССЗ лечение SGLT2i при ценах США 2025 года, согласно прогнозу, экономически эффективно по сравнению с обычной помощью.\\[5pt]
\multicolumn{2}{P{0.18\textwidth}}{Экономическая оценка} & \textbf{4. Не определена; недостаточно доказательств.} У взрослых со стадией ССПМ 4, СД2 и АССЗ экономическая эффективность терапии на основе GLP-1 при ценах США 2025 года не определена.\\[5pt]
2a & C-LD & \textbf{5.} Взрослым со стадией ССПМ 4, СД2 и АССЗ сочетание SGLT2i с терапией на основе GLP-1 может быть полезно для улучшения сердечно-сосудистых исходов.\src{4}\\
\bottomrule
\end{longtable}
}
\textbf{Примечание оригинала.} Экономические положения предназначены для решений на уровне населения и системы здравоохранения, а не для непосредственного определения лечения отдельных пациентов.

\heading{Обзор}
Полезное для сердца питание и регулярная физическая активность остаются основой лечения при СД2 и АССЗ; их польза для профиля ССПМ-риска доказана.\src{5--12} Важен также контроль курения, АД и холестерина.\src{13,14} Изменение образа жизни и целевая лекарственная терапия для комплексного контроля факторов риска уменьшали сердечно-сосудистые события и смертность при СД2, избыточной массе тела/ожирении и микроальбуминурии.\src{15}

Кардиопротективная сахароснижающая терапия --- SGLT2i или препараты на основе GLP-1 --- занимает центральное место при СД2 и АССЗ (рисунок 16). Оба класса уменьшают MACE и смертность от ССЗ; основное влияние SGLT2i связано с событиями СН, а терапии на основе GLP-1 --- с атеросклеротическими событиями.\src{3,16,17} При симптомном заболевании периферических артерий семаглутид увеличивает дистанцию ходьбы.\src{18}

Терапия на основе GLP-1 часто вызывает значительное снижение массы тела и улучшает профиль ССПМ-риска и MASLD, тогда как SGLT2i заметно влияет на ХБП.\src{19--30} Оба класса следует рассматривать как первую линию при сочетании СД2 и АССЗ.\src{31} Выбор первого препарата индивидуализируют с учётом профиля ССПМ-риска, наличия MASLD (раздел 7.2), приоритетного типа снижения сердечно-сосудистого риска, предпочтений пациента и стоимости.

Наконец, при выраженной гипергликемии SGLT2i или препаратов на основе GLP-1 может быть недостаточно для контроля гликемии. Тогда возможны другие сахароснижающие средства; метформин обладает превосходным профилем эффективности и безопасности при низком риске гипогликемии.

\heading{Рисунок 16. Ведение сочетания СД2 и АССЗ при ССПМ-синдроме: текстовая адаптация}
\begin{enumerate}
\item Исходная ситуация --- СД2 и АССЗ.
\item Рекомендовать и поддерживать изменение образа жизни (COR 1).
\item Добиваться комплексного контроля факторов риска с использованием разных методов (COR 1).
\item Добавить к схеме контроля гликемии кардиопротективный сахароснижающий препарат (COR 1):
\begin{itemize}
\item терапию на основе GLP-1 (COR 1), отдавая ей приоритет при тяжёлом ожирении, неконтролируемой гипергликемии или MASLD;
\item либо SGLT2i (COR 1), отдавая ему приоритет при ХБП либо доклинической/клинической СН.
\end{itemize}
\item Начать комбинированную терапию на основе GLP-1 и SGLT2i (COR 2a).
\item Совместно использовать другие сахароснижающие средства, например метформин, если HbA1c превышает индивидуальную цель более чем на 0,5--1 процентного пункта.
\end{enumerate}
ASCVD --- АССЗ; CKM --- ССПМ; COR --- класс рекомендации; GLP-1 --- глюкагоноподобный пептид-1; HbA1c --- гликированный гемоглобин; HF --- СН; MASLD --- стеатотическая болезнь печени, ассоциированная с метаболической дисфункцией; SGLT2i --- ингибиторы натрий-глюкозного котранспортера 2 типа; T2D --- СД2. Легенда оригинала включает классы 1, 2a, 2b, 3 «нет пользы» и 3 «вред»; в блоках используются 1 и 2a.
\editorial{Пункт о комбинированной терапии передаёт повелительную формулировку блока схемы, но отмечен COR 2a; в формальной рекомендации 5 сказано «может быть полезно», а не «рекомендуется всем». Порог HbA1c в схеме напечатан как «$>0{,}5$--1\% above individual goal» и передан как разница в процентных пунктах.}

\heading{Обоснование отдельных рекомендаций}
\textbf{1. Питание и физическая активность.} Полезная для сердца модель питания --- ключевое вмешательство при СД2. Средиземноморское питание, DASH, вегетарианские и веганские модели помогают снижать массу тела и улучшать гликемический контроль.\src{1,2} Проспективные когорты показали существенно меньшую вероятность событий ССЗ и сердечно-сосудистой смерти у взрослых с СД2, придерживающихся здорового питания.\src{32} Однако РКИ интенсивного изменения образа жизни при СД2, где 14\% участников имели ССЗ, не показало уменьшения событий АССЗ, несмотря на первоначально успешное снижение массы тела.\src{33}

Снижение массы тела --- важный компонент лечения СД2; при необходимости питание следует адаптировать для достижения клинически значимого результата. Разработка подходящего плана питания требует времени и усилий; оптимальна помощь зарегистрированного диетолога-нутрициолога или программы обучения диабету. Регулярная физическая активность, включая сокращение времени сидения, аэробные нагрузки и силовые тренировки, связана с улучшением нескольких факторов ССПМ-риска --- HbA1c, окружности талии, доли жировой массы, АД --- и показателей смертности.\src{5--12} Поэтому необходимо прикладывать все усилия для внедрения индивидуально подобранного, полезного для сердца образа жизни у взрослых со стадией ССПМ 4, СД2 и АССЗ.

\textbf{2. Кардиопротективные сахароснижающие средства.} При СД2 и АССЗ SGLT2i и терапия на основе GLP-1 значимо уменьшают MACE и сердечно-сосудистую смертность, способствуют снижению массы тела\src{34} и замедляют болезнь почек (раздел 5.5.4).\src{17,35,36} SGLT2i, по-видимому, преимущественно уменьшают впервые возникшую и ухудшающуюся СН.\src{17,37--44} GLP-1 RA преимущественно уменьшают атеросклеротические события --- инфаркт и инсульт.\src{45--48}

Оба класса улучшают почечные исходы.\src{49--51} Терапия на основе GLP-1 дополнительно обеспечивает значительное снижение массы тела и улучшение гликемии, особенно семаглутид и тирзепатид,\src{19--25} а также уменьшает фиброз печени при MASH (раздел 7.2).\src{52--58} Приоритет SGLT2i либо GLP-1-терапии определяется индивидуально: необходимой выраженностью контроля гликемии и снижения массы тела, ХБП и текущей рСКФ, другими факторами ССПМ-риска, MASLD или MASH, стоимостью и предпочтениями пациента. SGLT2i не следует начинать при рСКФ $<20$~\unitgfr. Благодаря различным механизмам совместное применение обоих классов может давать большее снижение сердечно-сосудистого риска.

\textbf{3. Экономическая оценка SGLT2i.} Несколько исследований с позиции здравоохранения США оценивали добавление SGLT2i к прежней стандартной терапии диабета.\src{59,60} Различия моделей и ключевых входных параметров не позволяют напрямую сравнивать результаты, однако исследования прогнозировали экономическую эффективность SGLT2i относительно одной стандартной помощи. Моделируемая популяция не ограничивалась установленным АССЗ, но результаты были устойчивыми к разным допущениям об исходном риске MACE, поэтому, вероятно, применимы и при АССЗ. Экономическая ценность может дополнительно улучшиться при снижении цен вследствие переговоров Medicare\src{61} либо появления генерических форм SGLT2i.

\textbf{4. Экономическая оценка GLP-1-терапии.} Недавняя модель показала, что при текущих ценах США препараты на основе GLP-1 по сравнению с обычной помощью не являются экономически эффективными для взрослых с диабетом и избыточной массой тела/ожирением.\src{62} Авторы не выявили высококачественных исследований подгруппы диабета с АССЗ. Её риск событий выше, чем в общей популяции диабета, и дополнительная польза может быть больше. Однако из-за отсутствия доказательств экономическая эффективность в этой подгруппе признана неопределённой.

\textbf{5. Сочетание SGLT2i и GLP-1-терапии.} Исследования совместного применения показывают улучшение HbA1c, АД и массы тела. Заранее запланированные и последующие подгрупповые анализы предполагают дополнительное снижение сердечно-сосудистого риска.\src{63--69} Метаанализ РКИ SGLT2i показал сердечно-сосудистую и почечную пользу независимо от фоновой GLP-1-терапии.\src{66} Аналогично метаанализ GLP-1-терапии показал уменьшение MACE независимо от фонового SGLT2i.\src{68}

Метаанализ рандомизированных и нерандомизированных данных показал улучшение сердечно-сосудистых исходов на комбинации, но также большую частоту гипогликемии;\src{64} сходные результаты получены в популяционном когортном исследовании.\src{70} Другой метаанализ РКИ SGLT2i или GLP-1-терапии, несмотря на улучшение профиля ССПМ-риска, не смог оценить влияние комбинации на сердечно-сосудистые исходы и смертность из-за недостатка данных.\src{65}

Наконец, заранее запланированные анализы исследования перорального сема\-глутида\src{4} показали независимость наблюдаемой пользы от сопутствующего SGLT2i, который получали 26,9\% исходно и 48,9\% в какой-либо момент исследования. Поэтому сочетание SGLT2i и терапии на основе GLP-1 можно рассматривать для улучшения гликемического контроля, особенно при добавлении GLP-1-терапии к SGLT2i, контроля факторов ССПМ-риска, сердечно-сосудистых и почечных исходов.

\subsubsection{Стадия ССПМ 4 с ХБП и АССЗ}
\heading{Обзор}
Руководства по АССЗ относят пациентов с ХБП к группе высокого риска повторных событий АССЗ.\src{1--3} Взаимосвязь двунаправленная: события АССЗ также увеличивают риск последующей почечной недостаточности.\src{4} Как и при ХБП на стадиях ССПМ 2--3, при АССЗ и ХБП рекомендуется терапия статинами для уменьшения сердечно-сосудистого риска.\src{5--8}

ХБП повышает риск острого повреждения почек и кровотечения при коронарографии и вмешательствах, поэтому необходимость этих процедур следует тщательно взвешивать. При стабильной коронарной болезни и ХБП, когда может потребоваться реваскуляризация, следует обратиться к разделу 9.3 руководства ACC/AHA/SCAI 2021 года по коронарной реваскуляризации.\src{9}

Нефропротективная терапия ХБП улучшает и почечные, и сердечно-сосудистые исходы, причём польза сходна при наличии и отсутствии АССЗ. Поэтому ведение ХБП при АССЗ аналогично её ведению на стадиях ССПМ 2--3 (раздел 5.5.4). RASi и SGLT2i --- первая линия благодаря широкому кругу подходящих пациентов и обширным доказательствам эффективности и безопасности. Добавление терапии на основе GLP-1 либо финеренона следует рассматривать при ХБП, АССЗ и СД2 с остаточной альбуминурией, сохраняющейся несмотря на первую линию. Выбор препарата зависит от профиля ССПМ-риска и других сопутствующих состояний.

\subsection{Ведение пациентов с СН}
\heading{Обзор}
Пациентам со стадией ССПМ 4 и СН показана стандартная терапия согласно клиническим рекомендациям (GDMT), дополненная коррекцией факторов ССПМ-риска, включая ожирение, диабет и ХБП. При СН со сниженной фракцией выброса (HFrEF) первая линия GDMT включает четыре компонента: RASi --- ингибитор рецепторов ангиотензина и неприлизина (ARNI), ACEi либо ARB, --- бета-блокатор, антагонист минералокортикоидных рецепторов (MRA) и SGLT2i.\src{1,2} Исследования показывают снижение сердечно-сосудистой смертности и госпитализаций по поводу СН уже в течение 30 дней после начала лечения, что подчёркивает важность раннего вмешательства.\src{3--5}

После назначения четырёхкомпонентной терапии пациентам, идентифицирующим себя как чернокожие, рекомендуются гидралазин/изосорбида динитрат для уменьшения сердечно-сосудистой смертности и госпитализаций по поводу СН.\src{1,6} При задержке жидкости диуретики облегчают признаки и симптомы застоя.\src{1} Имплантируемые кардиовертеры-дефибрилляторы и сердечная ресинхронизирующая терапия при наличии показаний играют важную роль в ведении HFrEF.\src{1}

При СН с умеренно сниженной (HFmrEF) или сохранённой фракцией выброса (HFpEF) первая линия включает SGLT2i\src{2} и диуретики при застое.\src{1} Доказательства пользы ARNI, ARB и стероидных MRA при HFpEF ограничены анализами post hoc и вторичными анализами, поэтому в руководствах США им присвоен класс 2b.\src{1} Вместе с тем появляются данные об эффективности nsMRA финеренона при HFmrEF/HFpEF.\src{7} Бета-блокаторы не рекомендуются при HFpEF из-за отсутствия пользы, если нет иных показаний, например ФП или симптомной коронарной болезни.\src{1}

Несколько ССПМ-состояний влияют на заболеваемость и смертность при СН.\src{8} Подходы к ожирению, диабету и ХБП представлены в следующих подразделах и на рисунке 17. За конкретными рекомендациями по лечению СН следует обращаться к руководству AHA/ACC по СН.\src{1}

\heading{Рисунок 17. Общая схема ведения СН на стадии ССПМ 4: текстовая адаптация}
\heading{Применять первую линию GDMT при СН}
\begin{itemize}
\item HFrEF: бета-блокатор, стероидный MRA, ARNI/ACEi/ARB, SGLT2i и диуретики по необходимости.
\item HFpEF: SGLT2i и диуретики по необходимости.
\end{itemize}
\heading{Применять общие принципы ведения ССПМ-синдрома}
\begin{itemize}
\item Укреплять сердечно-сосудистое здоровье на основе Life's Essential 8.
\item Организовать междисциплинарную помощь, включая выделенного координатора.
\item Работать с неблагоприятными SDOH с участием работников общественного здравоохранения и социальных работников, используя ресурсы сообщества.
\item Контролировать АД и липиды согласно действующим рекомендациям.
\end{itemize}
\heading{Выбирать ведение с учётом профиля риска: ожирение}
\begin{itemize}
\item При HFpEF начать терапию на основе GLP-1 для улучшения функциональных возможностей, симптомов и исходов СН.
\item При HFpEF рассмотреть физические тренировки и питание с дефицитом калорий для улучшения функциональных возможностей.
\item При HFrEF польза снижения массы тела и начала GLP-1-терапии остаётся неопределённой.
\end{itemize}
\heading{Выбирать ведение с учётом профиля риска: СД2}
\begin{itemize}
\item При СН отдавать приоритет SGLT2i как первой линии кардиопротективной сахароснижающей терапии.
\item При HFpEF и других факторах ССПМ-риска рассмотреть добавление GLP-1-терапии для улучшения симптомов и уменьшения неблагоприятных ССПМ-исходов.
\end{itemize}
\heading{Выбирать ведение с учётом профиля риска: ХБП}
\begin{itemize}
\item При СН и рСКФ $\geq30$ использовать ARNI/ACEi/ARB для улучшения исходов СН и почечных исходов.
\item При СН и рСКФ $\geq20$ использовать SGLT2i для улучшения исходов СН и почечных исходов.
\item При HFmrEF/HFpEF, СД2, рСКФ $\geq25$ и UACR $\geq30$~мг/г рассмотреть финеренон для улучшения исходов СН и почечных исходов.
\end{itemize}
\textbf{Пояснение авторов.} Ведение ССПМ-синдрома при СН на стадии 4 должно сочетать первую линию GDMT для СН, общие принципы ССПМ-помощи и дополнительные подходы к отдельным факторам риска, как показано на схеме.

ACEi --- ингибиторы ангиотензинпревращающего фермента; ARNI --- ингибиторы рецепторов ангиотензина и неприлизина; CHW --- работник общественного здравоохранения; CKM --- ССПМ; CVD --- ССЗ; eGFR --- рСКФ; GDMT --- терапия согласно клиническим рекомендациям; GLP-1 --- глюкагоноподобный пептид-1; HF --- СН; HFpEF --- СН с сохранённой фракцией выброса; HFrEF --- СН со сниженной фракцией выброса; MRA --- антагонист минералокортикоидных рецепторов; SDOH --- социальные детерминанты здоровья; SGLT2i --- ингибиторы натрий-глюкозного котранспортера 2 типа; T2D --- СД2.
\editorial{В блоках рисунка 17 единицы рСКФ не напечатаны; в руководстве применяется \unitgfr. Общая формулировка блока ARNI/ACEi/ARB не заменяет уточнение фенотипа СН и показаний в формальных рекомендациях раздела 6.3.3.}

\subsubsection{Стадия ССПМ 4 с ожирением и СН}
\heading{Рекомендации}
Исследования, на которых основаны рекомендации, обобщены в таблице доказательств оригинала.
{\small
\begin{longtable}{P{0.07\textwidth}P{0.11\textwidth}P{0.75\textwidth}}
\toprule COR & LOE & Рекомендация\\\midrule
\endfirsthead
\toprule COR & LOE & Рекомендация (продолжение)\\\midrule
\endhead
1 & A & \textbf{1.} У взрослых со стадией ССПМ 4, ожирением и симптомной HFpEF показана терапия на основе GLP-1 с доказанной сердечно-сосудистой пользой для улучшения профиля ССПМ-риска, функциональных возможностей и симптомов СН, а также предотвращения событий ухудшения СН.\src{1--3}\\[5pt]
\multicolumn{2}{P{0.18\textwidth}}{Экономическая оценка} & \textbf{2. Не определена; недостаточно доказательств.} У взрослых со стадией ССПМ 4, ожирением и симптомной HFpEF экономическая эффективность добавления к GDMT терапии на основе GLP-1 с доказанной сердечно-сосудистой пользой не определена.\\[5pt]
2a & B-R & \textbf{3.} У взрослых со стадией ССПМ 4, ожирением и HFpEF сочетание физических тренировок с питанием с дефицитом калорий может быть полезным для улучшения функциональных возможностей.\src{4--6}\\[5pt]
2b & B-NR & \textbf{4.} У отдельных взрослых со стадией ССПМ 4, ожирением и симптомной HFrEF можно рассмотреть лечение ожирения для улучшения функциональных возможностей и обеспечения возможности трансплантации сердца, если по остальным критериям пациент подходит для неё.\src{7--11}\\
\bottomrule
\end{longtable}
}
\textbf{Примечание оригинала.} Экономические положения предназначены для решений на уровне населения и системы здравоохранения, а не для непосредственного определения лечения отдельных пациентов.

\heading{Обзор}
Для улучшения клинических исходов СН приоритет следует отдавать доказательной GDMT у всех пациентов без противопоказаний независимо от ожирения.\src{12} При HFrEF основа лечения --- RASi (ARNI, ACEi либо ARB), бета-блокатор, MRA и SGLT2i. При HFmrEF или HFpEF рекомендуется SGLT2i (рисунок 18).\src{12}

Помимо стандартной GDMT, лечение ожирения становится потенциальной возможностью улучшить функциональные возможности и симптомы СН. Особенно это относится к HFpEF, которая эпидемиологически сильнее связана с ожирением, чем HFrEF, и при которой уже завершены исследования лечения ожирения.\src{13,14} Питание с дефицитом калорий и аэробные тренировки у пожилых с HFpEF дают умеренные, но клинически важные улучшения массы тела и функциональных возможностей.\src{4}

Немногочисленные ретроспективные исследования бариатрической хирургии при исходной СН предполагают среднесрочную пользу: снижение госпитализаций по поводу СН и общей смертности за 4--5 лет наблюдения.\src{15,16} Лечение связанной с ожирением HFpEF семаглутидом или тирзепатидом улучшает оцениваемое пациентом состояние здоровья и качество жизни, функциональные возможности по дистанции 6-минутной ходьбы (6MWD) и снижает частоту событий ухудшения СН, включая госпитализации.\src{1,3,17,18}

\heading{Рисунок 18. Ведение ожирения и СН в зависимости от фенотипа СН: текстовая адаптация}
\textbf{Исходная ситуация:} ожирение и СН.
\begin{itemize}
\item \textbf{HFrEF:} бета-блокаторы (COR 1), стероидные MRA (COR 1), ARNI/ACEi/ARB (COR 1) и SGLT2i (COR 1). Эти четыре компонента отмечены в оригинале серым фоном как GDMT для СН.
\item \textbf{HFpEF:} SGLT2i (COR 1; в серой зоне GDMT), терапия на основе GLP-1 (COR 1) и физические тренировки с питанием с дефицитом калорий (COR 2a).
\end{itemize}
\textbf{Примечание авторов.} Адаптировано с разрешения Heidenreich и соавторов.\src{12} © 2022 American Heart Association, Inc. и American College of Cardiology Foundation. За исключением класса рекомендации SGLT2i при HFpEF, обновление которого ожидается в следующей редакции руководства по СН 2022 года, GDMT соответствует этому руководству для соответствующего фенотипа СН и выделена серым.

COR --- класс рекомендации; GDMT --- терапия согласно клиническим рекомендациям; GLP-1 --- глюкагоноподобный пептид-1; HF --- СН; HFpEF --- СН с сохранённой фракцией выброса; HFrEF --- СН со сниженной фракцией выброса; SGLT2i --- ингибиторы натрий-глюкозного котранспортера 2 типа. Легенда оригинала включает классы 1, 2a, 2b, 3 «нет пользы» и 3 «вред»; в блоках использованы 1 и 2a.

\heading{Обоснование отдельных рекомендаций}
\textbf{1. Терапия на основе GLP-1.} В STEP-HFpEF --- исследовании эффекта семаглутида при ожирении и СН с сохранённой фракцией выброса --- у взрослых с HFpEF (ФВ ЛЖ $\geq45\%$), ожирением без СД2 и клиническим суммарным баллом Канзасского опросника кардиомиопатии (KCCQ-CSS) $<90$ семаглутид значимо уменьшал массу тела и улучшал оценку состояния при СН по KCCQ-CSS: изменение от исходного уровня 16,6 против 8,7 балла; $P<0{,}001$.\src{1,19} Также улучшались 6MWD, NT-proBNP и С-реактивный белок. STEP-HFpEF-DM включало пациентов со связанной с ожирением HFpEF и СД2 и дало сходные результаты.\src{2}

Участники STEP-HFpEF и STEP-HFpEF-DM, уже получавшие петлевые диуретики, имели наибольшее улучшение KCCQ-CSS на семаглутиде; семаглутид также сопровождался меньшей потребностью в дозе диуретиков.\src{20} Эти наблюдения были расширены в SUMMIT --- исследовании тирзепатида при HFpEF и ожирении. Тирзепатид обеспечивал значительное снижение массы тела, улучшение KCCQ-CSS (19,5 против 12,7 балла; $P<0{,}001$) и снижение первичной комбинированной конечной точки на 38\%: сердечно-сосудистая смерть, ухудшение СН с госпитализацией либо усилением лекарственной терапии СН.\src{3}

\textbf{2. Экономическая оценка.} Хотя экономическая эффективность GLP-1 RA и двойных агонистов GLP-1/GIP оценивалась для снижения массы тела, исследований именно при HFpEF на основе STEP-HFpEF и SUMMIT не было. Поэтому экономическая эффективность по этому показанию остаётся неопределённой из-за недостатка доказательств.

\textbf{3. Физические тренировки и питание.} SECRET --- исследование ограничения калорий и тренировок при СН с нормальной фракцией выброса --- оценивало 20 недель гипокалорийного питания, контролируемых тренировок либо их сочетания у 100 пожилых пациентов с HFpEF и ожирением; средний возраст 67 лет.\src{4} Пиковое потребление кислорода ($VO_2$, мл/кг/мин) увеличивалось, а масса тела уменьшалась наиболее выраженно при сочетании вмешательств; в группе диеты отмечалось небольшое снижение безжировой массы.

В SECRET-II --- исследовании непереносимости нагрузки у пожилых с HFpEF --- 88 пожилых пациентов со связанной с ожирением HFpEF рандомизировали на 20 недель ограничения калорий и аэробных тренировок с силовыми тренировками либо без них.\src{5} В обеих группах масса тела значимо снижалась, а пиковый $VO_2$ улучшался. Силовые тренировки не уменьшали потерю безжировой массы, но увеличивали силу мышц ног и качество скелетных мышц.\src{5}

Одногрупповая 15-недельная прагматическая программа контроля массы тела с еженедельными консультациями специалистов по физической нагрузке, диетологов и специалистов по поведению обеспечила среднее снижение массы тела на 8,1~кг ($-6{,}7\%$) при HFpEF, связанной с ожирением.\src{6} Несколько других исследований подтверждают безопасность и эффективность тренировок при HFpEF, но не ограничивались пациентами с ожирением.\src{21--23}

\textbf{4. Лечение ожирения при HFrEF.} Соотношение пользы и риска и подходящие стратегии остаются неопределёнными.\src{24} Одни физические упражнения мало влияют на массу тела при HFrEF.\src{25} Небольшие исследования изменения образа жизни предполагают эффективность и безопасность снижения массы тела,\src{8--10} а пилотное исследование орлистата ($n=21$) показало улучшение функциональных возможностей.\src{7}

Однако LIVE --- исследование лираглутида при стабильной хронической СН --- и FIGHT --- исследование функционального эффекта GLP-1 при лечении СН --- показали численно более высокую частоту событий ССЗ на лираглутиде, с наибольшим риском при функциональном классе NYHA III--IV.\src{26} В SELECT у взрослых с ИМТ $\geq27$~кг/м$^2$ и АССЗ без СД2 семаглутид уменьшал риск ССЗ в подгруппе HFrEF, но 90\% этих пациентов имели NYHA I--II. Поэтому при далеко зашедшей HFrEF соотношение пользы и риска GLP-1 RA неясно; при повышении дозы рекомендуется тщательное наблюдение.\src{27}

Безопасность метаболической и бариатрической хирургии (MBS) при HFrEF также не установлена.\src{28} Эти пробелы особенно важны, когда тяжёлое ожирение препятствует трансплантации сердца при далеко зашедшей HFrEF. Появляющиеся наблюдательные данные показывают, что хирургическое снижение массы тела может помочь подготовить к трансплантации пациентов с ожирением, получающих длительную поддержку устройством вспомогательного кровообращения левого желудочка.\src{11,29}

\subsubsection{Стадия ССПМ 4 с СД2 и СН}
\heading{Рекомендации}
Исследования, на которых основаны рекомендации, обобщены в таблице доказательств оригинала.
{\small
\begin{longtable}{P{0.07\textwidth}P{0.11\textwidth}P{0.75\textwidth}}
\toprule COR & LOE & Рекомендация\\\midrule
\endfirsthead
\toprule COR & LOE & Рекомендация (продолжение)\\\midrule
\endhead
1 & A & \textbf{1.} У пациентов со стадией ССПМ 4, СД2 и СН следует отдавать приоритет SGLT2i как первой линии кардиопротективной сахароснижающей терапии для уменьшения сердечно-сосудистой смерти и госпитализаций по поводу СН.\src{1--3}\\[5pt]
\multicolumn{2}{P{0.18\textwidth}}{Экономическая оценка} & \textbf{2. Экономически эффективно; высокая степень уверенности.} У пациентов со стадией ССПМ 4, СД2 и симптомной СН с ФВ ЛЖ $\leq40\%$ добавление SGLT2i к прежней GDMT при ценах США 2025 года, согласно прогнозу, экономически эффективно по сравнению с прежней GDMT: ICER $<120\,000$ долл. США на дополнительный QALY.\\[5pt]
\multicolumn{2}{P{0.18\textwidth}}{Экономическая оценка} & \textbf{3. Не определена; противоречивые доказательства.} У пациентов со стадией ССПМ 4, СД2 и симптомной HFpEF экономическая эффективность добавления SGLT2i при ценах США 2025 года не определена.\\[5pt]
2a & B-R & \textbf{4.} У пациентов со стадией ССПМ 4, СД2, HFpEF и другими факторами ССПМ-риска добавление терапии на основе GLP-1 с доказанной сердечно-сосудистой пользой к базовой терапии SGLT2i может быть полезным для улучшения симптомов СН и уменьшения неблагоприятных ССПМ-исходов.\src{4--6}\\[5pt]
2a & B-NR & \textbf{5.} У пациентов со стадией ССПМ 4, СД2, стабильной СН, рСКФ $\geq30$~\unitgfr\ и HbA1c выше индивидуальной цели добавление метформина к SGLT2i может быть полезным для достижения целевой гликемии.\src{7--10}\\
\bottomrule
\end{longtable}
}
\textbf{Примечание оригинала.} Экономические положения предназначены для решений на уровне населения и системы здравоохранения, а не для непосредственного определения лечения отдельных пациентов.

\heading{Обзор}
СД2 существенно увеличивает риск СН. Если нынешние эпидемиологические тенденции сохранятся, СН станет ведущим сердечно-сосудистым осложнением СД2.\src{11} Поэтому СД2 и СН часто сосуществуют при ССПМ-синдроме; особенно выражено сочетание СД2 с HFpEF. Рекомендации на рисунке 19 акцентируют терапию, улучшающую ССПМ-исходы при стадии 4 с СД2 и СН. Стандартная GDMT для СН является основой у всех пациентов.

Подгрупповые анализы крупных исследований и несколько РКИ при СД2 и HFpEF показывают эффективность SGLT2i для уменьшения событий СН и сердечно-сосудистой смертности; польза имеется и при HFrEF, и при HFpEF.\src{12,13} Следовательно, SGLT2i должны быть приоритетной первой линией кардиопротективных сахароснижающих средств у всех пациентов с СД2 и СН.

Хотя RASi отдельно не обсуждаются в этом подразделе, их основополагающую роль нельзя упускать: помимо пользы при СН, они дают сосудистую пользу при СД2. Добавление GLP-1-терапии при СД2, ожирении и HFpEF улучшает симптомы и функциональные возможности.\src{4,5} Добавление её к SGLT2i также позволяет корректировать дополнительные факторы ССПМ-риска.

\heading{Рисунок 19. Ведение СД2 и СН в зависимости от фенотипа СН: текстовая адаптация}
\textbf{Исходная ситуация:} СД2 и СН.
\begin{itemize}
\item \textbf{HFrEF:} бета-блокаторы, стероидные MRA, ARNI/ACEi/ARB и SGLT2i; каждый компонент --- COR 1.
\item \textbf{HFpEF:} SGLT2i (COR 1). Затем оценить наличие ожирения, неконтролируемой гипергликемии, MASLD, АССЗ или ХБП с альбуминурией. Если их нет --- продолжить текущую терапию. Если присутствуют --- рассмотреть GLP-1-терапию (COR 2a).
\end{itemize}
\textbf{Примечание авторов.} Адаптировано с разрешения Heidenreich и соавторов.\src{34} © 2022 American Heart Association, Inc. и American College of Cardiology Foundation. Серым выделена GDMT для соответствующего фенотипа СН согласно руководству 2022 года: четыре класса при HFrEF и SGLT2i при HFpEF.

ACEi --- ингибиторы АПФ; ARB --- блокаторы рецепторов ангиотензина II; ARNI --- ингибиторы рецепторов ангиотензина и неприлизина; ASCVD --- АССЗ; CKD --- ХБП; COR --- класс рекомендации; GDMT --- терапия согласно клиническим рекомендациям; GLP-1 --- глюкагоноподобный пептид-1; HF --- СН; HFpEF --- СН с сохранённой фракцией выброса; HFrEF --- СН со сниженной фракцией выброса; MASLD --- стеатотическая болезнь печени, ассоциированная с метаболической дисфункцией; MRA --- антагонист минералокортикоидных рецепторов; SGLT2i --- ингибиторы натрий-глюкозного котранспортера 2 типа; T2D --- СД2. Легенда содержит классы 1, 2a, 2b, 3 «нет пользы» и 3 «вред»; использованы 1 и 2a.

\heading{Обоснование отдельных рекомендаций}
\textbf{1. SGLT2i.} Показана широкая эффективность для уменьшения госпитализаций по поводу СН и сердечно-сосудистой смерти на стадии ССПМ 4 при СД2 и СН с любой фракцией выброса. Метаанализ EMPEROR-Reduced (50\% с СД2) и DAPA-HF (45\% с СД2) дал HR 0,74 (95\% ДИ 0,65--0,84) для сердечно-сосудистой смерти либо первой госпитализации по поводу СН при СД2 и HFrEF.\src{1} В анализе post hoc DECLARE-TIMI 58 при СД2 и HFrEF ($n=671$, 3,9\% выборки) дапаглифлозин также уменьшал сердечно-сосудистую смерть и госпитализации по поводу СН: HR 0,62; 95\% ДИ 0,45--0,86.\src{2}

Более поздний метаанализ DELIVER и EMPEROR-Preserved показал, что дапаглифлозин и эмпаглифлозин уменьшают комбинированную частоту сердечно-сосудистой смерти и госпитализаций по поводу СН более чем на 20\% при СД2 --- примерно 45\% и 49\% участников соответственно --- и ФВ ЛЖ $>40\%$.\src{3} Благодаря пользе при HFrEF, HFmrEF и HFpEF SGLT2i должны быть первой линией кардиопротективных сахароснижающих средств при диабете и СН, если нет противопоказаний.

\textbf{2. Экономическая оценка при HFrEF.} Две высококачественные имитационные модели оценивали добавление SGLT2i к прежней GDMT у взрослых с симптомной HFrEF с позиции здравоохранения США на пожизненном горизонте.\src{14,15} Обе отдельно анализировали исходный диабет. Годовая стоимость препаратов 4192--5684 долл. США соответствовала Federal Supply Schedule в апреле 2025 года: 5200--5600 долл. США.\src{16}

В одной модели добавление SGLT2i давало 0,70 дополнительного QALY (95\% интервал неопределённости 0,23--1,20) при дополнительных затратах 46\,500 долл. США (32\,700--49\,700). ICER составлял 66\,800 долл. США на QALY (53\,800--116\,600), и стратегия была экономически эффективной в 95\% из 10\,000 вероятностных моделирований. Дополнительно сообщались ICER 57\,300 долл. США на дополнительный год жизни (44\,800--123\,800) и 53\,400 долл. США на дополнительный год жизни при равной оценке его ценности (42\,800--97\,200; показатель equal value of life-years gained). Во второй модели ICER при диабете составлял 79\,700 долл. США на QALY; оценка без диабета была сходной.

\textbf{3. Экономическая оценка при HFpEF.} Две высококачественные модели добавления SGLT2i к прежней GDMT при симптомной HFpEF не представили отдельных анализов по диабету.\src{17,18} Поскольку в РКИ при HFpEF эффект SGLT2i не различался в зависимости от диабета, возможно, и экономическая эффективность от него не зависит, как при HFrEF. Прежние исследования с позиции здравоохранения США и пожизненным горизонтом показывали малую вероятность экономической эффективности при HFpEF, но ICER зависел от влияния на смертность.

Одна модель использовала EMPEROR-Preserved, где влияние на сердечно-сосу\-дистую смертность не достигало статистической значимости: HR 0,91; 95\% ДИ 0,76--1,09. В основном анализе без предполагаемого снижения сердечно-сосудистой смертности при годовой цене 3920 долл. США ICER составлял 437\,400 долл. США на QALY; пациенты с диабетом и без него были объединены.\src{18} Если предполагалось уменьшение сердечно-сосудистой смертности на 9\%, ICER снижался до 174\,053 долл. США на QALY.

Вторая модель объединяла EMPEROR-Preserved и DELIVER и вероятностно учитывала снижение сердечно-сосудистой смертности из опубликованного метаанализа: HR 0,88; 95\% ДИ 0,77--1,00. При годовой цене 4506 долл. США\src{17} ICER относительно прежней GDMT составлял 141\,200 долл. США на QALY. Эти анализы предполагают отсутствие экономической эффективности при текущих ценах: ICER $\geq120\,000$ долл. США на QALY. Однако они объединяли лиц с диабетом и без него, а при диабете возможна дополнительная сердечно-сосудистая и почечная польза. В свете противоречивых данных экономическая эффективность при диабете и HFpEF признана неопределённой.

\textbf{4. Дополнительная GLP-1-терапия.} Как обсуждалось в разделе 6.2.1 «Стадия ССПМ 4 с ожирением и АССЗ», STEP-HFpEF-DM оценивало GLP-1-терапию при HFpEF, ожирении и СД2. Подкожный семаглутид значимо снижал массу тела, улучшал KCCQ-CSS и 6MWD относительно плацебо.\src{4} В SUMMIT половина участников имела СД2; тирзепатид уменьшал массу тела, улучшал KCCQ-CSS и 6MWD и снижал события ухудшения СН на 38\%, без различий эффекта по наличию СД2.\src{5}

Доказательств пользы GLP-1 при HFpEF становится больше; эффективность и безопасность при HFrEF менее ясны из-за некоторого увеличения риска неблагоприятных исходов в небольших исследованиях.\src{19--21} Интервенционных данных по сочетанию SGLT2i и GLP-1-терапии мало, но в целом эффект GLP-1 не различался при наличии и отсутствии SGLT2i.\src{22} В испытаниях также последовательно показано улучшение факторов ССПМ-риска на GLP-1-терапии. Поэтому её целесообразно добавлять к SGLT2i при СД2, HFpEF и дополнительных факторах риска или связанных состояниях: ожирении (раздел 6.3.1), ХБП с альбуминурией (раздел 6.3.3), MASLD (раздел 7.2), АССЗ (раздел 6.2.2).
\editorial{Ссылка на раздел 6.2.1 в первом предложении этого обоснования воспроизводит оригинал; развёрнутое обсуждение STEP-HFpEF-DM в данном переводе находится в разделе 6.3.1.}

\textbf{5. Метформин и гликемические цели.} Данные о метформине при СН предшествуют современной GDMT и полностью наблюдательные. Многочисленные наблюдательные исследования поддерживают прямую линейную связь гликемии с развитием и событиями СН.\src{23--25} Однако при СД2 и уже имеющейся СН смертность имеет U-образную связь с гликемией,\src{26--28} что подчёркивает потенциальный вред чрезмерно агрессивного снижения глюкозы. Поддержание HbA1c в пределах 7--8\% считается подходящим при СД2 и СН.\src{29}

Нерандомизированные сообщения в разных возрастах и по всему спектру СН --- амбулаторная хроническая СН, впервые диагностированная СН, недавняя госпитализация --- не подтверждают вреда метформина, хотя при декомпенсированной СН он противопоказан.\src{9,30} Поэтому при необходимости дополнительного контроля гликемии целесообразно добавлять метформин к SGLT2i, особенно с учётом умеренного сахароснижающего эффекта SGLT2i, при ССПМ-синдроме с СД2 и стабильной СН. Дозу следует уменьшать в зависимости от функции почек:\src{31}
\begin{itemize}
\item рСКФ $\geq60$~\unitgfr: обычная доза;
\item рСКФ 45--59~\unitgfr: 1000--1500~мг/сут;
\item рСКФ 30--44~\unitgfr: $\leq1000$~мг/сут;
\item рСКФ $<30$~\unitgfr: противопоказан.
\end{itemize}
За дополнительными указаниями по гликемии авторы рекомендуют обращаться к руководствам ADA.\src{32,33}

\subsubsection{Стадия ССПМ 4 с ХБП и СН}
\heading{Рекомендации}
Исследования, на которых основаны рекомендации, обобщены в таблице доказательств оригинала.
{\small
\begin{longtable}{P{0.07\textwidth}P{0.11\textwidth}P{0.75\textwidth}}
\toprule COR & LOE & Рекомендация\\\midrule
\endfirsthead
\toprule COR & LOE & Рекомендация (продолжение)\\\midrule
\endhead
1 & A & \textbf{1.} У взрослых со стадией ССПМ 4, ХБП, рСКФ $\geq30$~\unitgfr\ и HFrEF рекомендуется начать ARNI либо другой RASi, если начать ARNI невозможно, для уменьшения риска сердечно-сосудистой смерти или госпитализации по поводу СН и утраты функции почек.\src{1--8}\\[5pt]
\multicolumn{2}{P{0.18\textwidth}}{Экономическая оценка} & \textbf{2. Экономически эффективно; умеренная степень уверенности.} У взрослых со стадией ССПМ 4, ХБП, рСКФ $\geq30$~\unitgfr\ и HFrEF лечение ARNI, согласно прогнозу, экономически эффективно по сравнению с терапией RASi.\\[5pt]
1 & A & \textbf{3.} У взрослых со стадией ССПМ 4, ХБП и СН с любой фракцией выброса при рСКФ $\geq20$~\unitgfr\ рекомендуется начать SGLT2i для уменьшения сердечно-сосудистой смертности, госпитализаций по поводу СН и, возможно, утраты функции почек.\src{9--12}\\[5pt]
\multicolumn{2}{P{0.18\textwidth}}{Экономическая оценка} & \textbf{4. Экономически эффективно; умеренная степень уверенности.} У взрослых со стадией ССПМ 4, ХБП и симптомной СН с ФВ ЛЖ $\leq40\%$ добавление SGLT2i к прежней GDMT при ценах США 2025 года, согласно прогнозу, экономически эффективно по сравнению с прежней GDMT.\\[5pt]
\multicolumn{2}{P{0.18\textwidth}}{Экономическая оценка} & \textbf{5. Не определена; противоречивые доказательства.} У взрослых со стадией ССПМ 4, ХБП и симптомной HFpEF экономическая эффективность добавления SGLT2i к прежней GDMT при ценах США 2025 года не определена.\\[5pt]
2a & B-R & \textbf{6.} У взрослых со стадией ССПМ 4, ХБП, СД2, UACR $\geq30$~мг/г и СН с ФВ ЛЖ $>40\%$ (HFmrEF и HFpEF) при рСКФ $\geq25$~\unitgfr\ целесообразно начать нестероидный MRA для уменьшения риска госпитализации по поводу СН и утраты функции почек.\src{13}\\[5pt]
\multicolumn{2}{P{0.18\textwidth}}{Экономическая оценка} & \textbf{7. Не определена; недостаточно доказательств.} У взрослых со стадией ССПМ 4, ХБП с рСКФ $\geq25$~\unitgfr, СД2 и СН с ФВ ЛЖ $>40\%$ (HFmrEF и HFpEF) экономическая эффективность добавления nsMRA к остальной GDMT не определена.\\[5pt]
2b & B-R & \textbf{8.} У взрослых со стадией ССПМ 4, ХБП и HFrEF при рСКФ $>30$~\unitgfr\ может быть целесообразным применение новых пероральных калийсвязывающих препаратов для уменьшения риска гиперкалиемии и обеспечения возможности ингибирования РААС.\src{14--17}\\[5pt]
\multicolumn{2}{P{0.18\textwidth}}{Экономическая оценка} & \textbf{9. Не определена; недостаточно доказательств.} У взрослых со стадией ССПМ 4, ХБП и HFrEF при рСКФ $>30$~\unitgfr\ экономическая эффективность добавления новых пероральных калийсвязывающих препаратов к GDMT неясна.\\
\bottomrule
\end{longtable}
}
\textbf{Примечание оригинала.} Экономические положения предназначены для решений на уровне населения и системы здравоохранения, а не для непосредственного определения лечения отдельных пациентов.
\editorial{В экономических пунктах 4 и 5 исходной таблицы напечатано «adding an SGLT2i prior to GDMT». По содержанию обоснований 4 и 5 речь идёт о добавлении к прежней GDMT, а не о назначении раньше остальной базовой терапии; это контекстное прочтение явно принято в переводе.}

\heading{Обзор}
СН часто выявляется или развивается вместе с ХБП. При HFrEF установлены четыре основных компонента лекарственной терапии, улучшающих выживаемость и уменьшающих госпитализации по поводу СН (рисунок 20).\src{18--20} Из этих четырёх классов --- шести, если пациент идентифицирует себя как чернокожий, --- два обладают нефропротективными свойствами: (1) RASi, среди которых ARNI превосходит ACEi/ARB по пользе при СН, и (2) SGLT2i.\src{13,21--25} Следует избегать совместного применения более одного RASi из ACEi/ARB/ARNI: это может ухудшать функцию почек и вызывать ангионевротический отёк, гиперкалиемию и гипотензию.\src{26}

При ХБП с HFpEF либо HFmrEF наиболее убедительные данные о нефропротекции и уменьшении комбинированных исходов смертности и госпитализации по поводу СН имеются для SGLT2i.\src{27} Финеренон также можно рассматривать в подгруппе с СД2, ХБП и HFmrEF/HFpEF без гиперкалиемии на основании клинической пользы в разных исследованиях.

Эти рекомендации дополняют существующие рекомендации классов 1 и 2 руководства ACC/AHA по СН\src{18} и предлагают дополнительные меры против риска, связанного с сочетанием ХБП и СН. Пороги рСКФ для начала лекарств отражают критерии включения в клинические исследования. Продолжение терапии после снижения рСКФ ниже порога начала лечения рассматривается отдельно в разделе 6.5 «Лечение ССПМ-синдрома при далеко зашедшей ХБП».

\heading{Рисунок 20. Ведение ХБП и СН в зависимости от фенотипа СН: текстовая адаптация}
\textbf{Исходная ситуация:} ХБП и СН.
\heading{Ветвь HFrEF}
\begin{itemize}
\item Для улучшения исходов СН: бета-блокаторы (COR 1) и стероидные MRA (COR 1).
\item Для улучшения исходов СН и почечных исходов: ARNI/ACEi/ARB при рСКФ $>30$ на момент начала (COR 1) и SGLT2i при рСКФ $>20$ на момент начала (COR 1).
\end{itemize}
\heading{Ветвь HFmrEF/HFpEF}
\begin{itemize}
\item Для улучшения исходов СН и почечных исходов: SGLT2i при рСКФ $>20$ на момент начала (COR 1).
\item Проверить наличие СД2 и ХБП с альбуминурией. Если нет --- продолжить текущую терапию; если есть --- добавить nsMRA (COR 2a).
\end{itemize}
\textbf{Примечание авторов.} Адаптировано с разрешения Heidenreich и соавторов.\src{18} © 2022 American Heart Association, Inc. и American College of Cardiology Foundation. Серым выделена GDMT для соответствующего фенотипа СН согласно руководству 2022 года: бета-блокаторы, стероидные MRA, ARNI/ACEi/ARB и SGLT2i при HFrEF; SGLT2i при HFmrEF/HFpEF.

CKD --- ХБП; eGFR --- рСКФ; GDMT --- терапия согласно клиническим рекомендациям; HF --- СН; HFmrEF --- СН с умеренно сниженной фракцией выброса; HFpEF --- СН с сохранённой фракцией выброса; HFrEF --- СН со сниженной фракцией выброса; nsMRA --- нестероидный антагонист минералокортикоидных рецепторов; SGLT2i --- ингибиторы натрий-глюкозного котранспортера 2 типа; T2D --- СД2. Легенда содержит классы 1, 2a, 2b, 3 «нет пользы» и 3 «вред»; использованы 1 и 2a.
\editorial{В рисунке 20 напечатаны строгие пороги рСКФ $>30$ и $>20$, тогда как формальные рекомендации 1 и 3 используют $\geq30$ и $\geq20$. Оба варианта переданы без неоговорённого выравнивания. Единицы рСКФ в блоках не указаны; в тексте используются \unitgfr. Порог рСКФ и UACR для nsMRA следует сверять с рекомендацией 6.}

\heading{Обоснование отдельных рекомендаций}
\textbf{1. RASi и ARNI.} При HFrEF RASi улучшали исходы СН и уменьшали риск ухудшения болезни почек в нескольких исследованиях.\src{1--8,28} ARNI даёт большее снижение комбинированной частоты сердечно-сосудистой смерти и госпитализаций по поводу СН, чем другие RASi, поэтому при доступности предпочтителен при HFrEF.\src{2} Польза показана при рСКФ $>30$~\unitgfr. Большинство исследований исключало ХБП G4--G5; данные для этих стадий ограничены небольшими анализами post hoc.\src{1,5,22}

В UK HARP 3 (британское исследование защиты сердца и почек III) при ХБП значимого различия почечных эффектов ARNI и ARB не выявлено, но СН имела лишь небольшая часть участников.\src{29} Ограниченные данные допускают сходную пользу ARNI при HFpEF и ХБП G2--G3b.\src{30} В ранних исследованиях ACEi и ARB при СН функцию почек систематически не оценивали, а пациентов с креатинином $>2{,}5$~мг/дл обычно исключали.\src{28}

RASi могут сопровождаться умеренным острым снижением функции почек, обычно гемодинамического происхождения. Оно не свидетельствует о вреде и может сочетаться с лучшими долгосрочными исходами.\src{23,31,32} Перед прекращением терапии следует рассмотреть другие причины ухудшения функции почек.

\textbf{2. Экономическая оценка ARNI.} По меньшей мере два высококачественных анализа сравнивали ARNI с RASi при HFrEF NYHA II--IV с позиции здравоохранения США на пожизненном горизонте.\src{33,34} ICER составлял приблизительно 50\,000 долл. США на QALY в долларах 2015 года. После поправки на инфляцию около 35\% между 2015 и 2025 годами ICER возрастает примерно до 68\,000 долл. США на QALY, оставаясь заметно ниже порога 120\,000. Это предполагает экономическую эффективность ARNI при симптомной HFrEF.

Однако сохраняются два источника неопределённости. Во-первых, ни один анализ не изучал подгруппу ХБП. При ХБП риск событий СН во время наблюдения выше, а в PARADIGM-HF относительный эффект не зависел от рСКФ. Поэтому абсолютная польза может быть больше, если её не нивелирует увеличение нежелательных явлений. Это справедливо для ранней и умеренной ХБП, но может не распространяться на далеко зашедшую ХБП, когда госпитализации из-за перегрузки объёмом могут не предотвращаться ARNI.

Во-вторых, экономическая эффективность чувствительна к цене. Стоимость ARNI значительно выросла со времени анализов, но в 2026 году ожидается её снижение для получателей Medicare Part D применительно к годовому курсу. В совокупности результаты предполагают вероятную экономическую эффективность ARNI при стадии ССПМ 4 с ХБП и СН, особенно при ценах США 2026 года.

\textbf{3. SGLT2i.} Начало SGLT2i эффективно уменьшает сердечно-сосудистую смертность, госпитализации по поводу СН и снижение рСКФ при стадии ССПМ 4, ХБП (СКФ $>20$~\unitgfr) и HFrEF, особенно при альбуминурии.\src{9--12} При рСКФ $>20$~\unitgfr\ и HFmrEF/HFpEF также показано уменьшение комбинированной сердечно-сосудистой смертности и госпитализаций, но при тяжёлом снижении функции почек --- ХБП G4 --- данных меньше; наличие или отсутствие альбуминурии недостаточно изучено в исследованиях СН.\src{9--12}

В анализе EMPEROR-Preserved (HFpEF/HFmrEF) и EMPEROR-Reduced (HFrEF) эмпаглифлозин сходным образом улучшал исходы СН и замедлял снижение рСКФ во всех категориях риска KDIGO. Для начала SGLT2i при рСКФ $<20$~\unitgfr\ нужны дополнительные доказательства.\src{12}

\textbf{4. Экономическая оценка SGLT2i при HFrEF.} Две высококачественные модели прогнозировали экономическую эффективность добавления SGLT2i к прежней GDMT при симптомной HFrEF с позиции здравоохранения США на пожизненном горизонте.\src{35,36} ICER составлял 68\,000--84\,000 долл. США на QALY при годовой цене 4192--5684 долл. США, сопоставимой с Federal Supply Schedule в апреле 2025 года --- 5200--5600 долл. США.\src{37} В вероятностных анализах ICER был $<120\,000$ долл. США на QALY более чем в 95\% моделирований в одном исследовании\src{35} и примерно в 70\% в другом.\src{36}

Ни одно исследование отдельно не анализировало ХБП. При ХБП риск АССЗ и событий СН выше, а SGLT2i дополнительно замедляют ухудшение функции почек. В ключевых РКИ относительное снижение больших неблагоприятных сердечно-сосудистых событий не зависело от ХБП; поэтому абсолютная польза при HFrEF с ХБП, вероятно, больше, чем без ХБП. Оговорка: это вероятнее для ранней и умеренной ХБП и менее определённо для далеко зашедшей. Поэтому допустимо экстраполировать экономическую эффективность при HFrEF в целом на эту более рискованную подгруппу с потенциально большей абсолютной пользой. Ожидается дальнейшее улучшение экономической эффективности при снижении цен из-за появления генериков и переговоров Medicare.\src{38}

\textbf{5. Экономическая оценка SGLT2i при HFpEF.} Две высококачественные модели с позиции здравоохранения США и пожизненным горизонтом прогнозировали, что добавление SGLT2i к прежней GDMT при симптомной HFpEF не является экономически эффективным; ICER чувствителен к влиянию на смертность.\src{39,40} В модели EMPEROR-Preserved эмпаглифлозин при годовой цене 3920 долл. США давал ICER 437\,400 долл. США на QALY.\src{40} Вторая модель вероятностно учитывала снижение сердечно-сосудистой смертности из метаанализа EMPEROR-Preserved и DELIVER: при годовой цене 4506 долл. США ICER составлял 141\,200 долл. США на QALY.\src{39}

Однако эти анализы объединяли участников с ХБП и без неё и непосредственно не учитывали почечную пользу SGLT2i. В целом они предполагают отсутствие экономической эффективности при текущих ценах: ICER $\geq120\,000$ долл. США на QALY.

\textbf{6. Финеренон.} Исследования при диабетической болезни почек показали значимое уменьшение утраты функции почек и впервые возникшей СН; большая часть эффекта в отношении СН, вероятно, реализуется через нефропротекцию.\src{41,42} В описании FINEARTS-HF --- исследования эффективности и безопасности финеренона по сравнению с плацебо при СН --- оригинал сообщает об уменьшении комбинированной конечной точки «смерть от всех причин и госпитализация по поводу СН» у пациентов с «HFmrEF/HFrEF».\src{43}

Объединённый метаанализ исследований финеренона показал меньшее\src{28} снижение рСКФ,\src{13} однако в одном FINEARTS-HF уменьшения первичной почечной конечной точки не было.\src{44} Вероятное объяснение --- отсутствие более выраженной ХБП: средняя исходная рСКФ 62~\unitgfr\ и средний UACR $<20$~мг/г. Совокупность исследований предполагает сердечно-сосудистую и почечную пользу финеренона при альбуминурической диабетической болезни почек и HFmrEF/HFpEF, хотя нужна более специализированная работа в этой популяции. Рекомендации по началу финеренона пока не распространяются на рСКФ $<25$~\unitgfr\ или калий $>5$~ммоль/л: таких пациентов исключали из исследований.\src{13}
\editorial{В приведённом предложении о FINEARTS-HF оригинал использует «HFmrEF/HFrEF» и смерть от всех причин; далее в том же абзаце и в формальной рекомендации 6 обсуждаются HFmrEF/HFpEF и ФВ ЛЖ $>40\%$. Эти несогласованности отмечены, а не исправлены скрыто. Надстрочная ссылка 28 после слова «меньшее» также присутствует в исходном PDF. Для клинического применения описание конечной точки и популяции требуется сверить с первичной публикацией исследования.}

\textbf{7. Экономическая оценка nsMRA.} Авторы не выявили строгих анализов добавления nsMRA к GDMT с позиции здравоохранения США. Учитывая стоимость финеренона по Federal Supply Schedule в апреле 2025 года --- 7612 долл. США за годовой курс\src{37} --- и потенциально большое число подходящих пациентов, изучение экономической эффективности и бюджетного влияния должно стать исследовательским приоритетом.

\textbf{8. Калийсвязывающие препараты.} Новые пероральные средства, такие как патиромер, могут обеспечивать более полное применение ингибиторов РААС при разных рСКФ за счёт уменьшения риска гиперкалиемии. Однако в исследования с сопутствующей СН включали мало пациентов с рСКФ $<30$~\unitgfr.\src{14,15,45} Натрия циркония циклосиликат может повышать риск событий СН.\src{16} Текущие данные не определяют, выше ли этот риск при нём, чем при патиромере.\src{46}

\textbf{9. Экономическая оценка калийсвязывающих средств.} Авторы не выявили строгих анализов патиромера или натрия циркония циклосиликата при стадии ССПМ 4 с ХБП и СН с позиции здравоохранения США. Предотвращение гиперкалиемии и переносимость более высоких доз ингибиторов РААС могут приносить значительную пользу, но препараты в США дороги. Патиромер стоит на 24--28\% больше, чем препарат натрия циркония, названный в этой строке оригинала «sodium zirconium silicate». Экономическая эффективность и бюджетное влияние этих средств --- срочный исследовательский приоритет.

\subsection{Фибрилляция предсердий}
\heading{Обзор}
Прогрессирующие последствия избытка и дисфункции жировой ткани увеличивают риск впервые возникшей ФП. ФП, в свою очередь, связана с повышенным риском нескольких компонентов стадии ССПМ 4: инсульта, инфаркта миокарда, СН, ХБП и заболевания периферических артерий. Риск ФП повышают компоненты стадий 1--3: избыточная масса тела/ожирение, гипертензия, СД2 и ХБП.

На стадиях 1--3 профилактика ФП включает здоровые изменения образа жизни, контроль массы тела, гликемии и АД, прекращение курения, ограничение алкоголя или отказ от него, лечение нарушений дыхания во сне.\src{1} Данные о предотвращении ФП с помощью ACEi, ARB, статинов, SGLT2i и GLP-1-терапии ограничены.\src{1} MRA связаны с уменьшением впервые возникшей и рецидивирующей ФП.\src{2}

При имеющейся ФП --- одном из компонентов стадии ССПМ 4 --- для уменьшения её бремени, рецидивов и прогрессирования рекомендуются изменение образа жизни и лечение избыточной массы тела/ожирения (разделы 5.4.1--5.4.4), гипертензии (5.5.3), СД2 (5.5.1) и ХБП (5.5.4). При ФП и ИМТ $\geq40$~кг/м$^2$ без бариатрической операции в анамнезе прямые пероральные антикоагулянты (DOAC), а именно апиксабан и ривароксабан, могут быть вариантами тромбопрофилактики. Профилактика и ведение ФП всесторонне обсуждаются в руководстве ACC/AHA/ACCP/HRS 2023 года по диагностике и лечению ФП.\src{1}

\subsection{Лечение ССПМ-синдрома при далеко зашедшей ХБП}
\heading{Рекомендация}
Исследования, на которых основана рекомендация, обобщены в таблице доказательств оригинала.
{\small
\begin{longtable}{P{0.07\textwidth}P{0.11\textwidth}P{0.75\textwidth}}
\toprule COR & LOE & Рекомендация\\\midrule
2a & B-R & \textbf{1.} У взрослых с ХБП, получающих нефропротективную терапию ради почечной и сердечно-сосудистой пользы, при снижении рСКФ ниже порога начала конкретного препарата целесообразно продолжать лечение, если оно безопасно переносится.\src{1--8}\\
\bottomrule
\end{longtable}
}
\heading{Обзор}
Как подчёркивает карта риска KDIGO (рисунок 2), ХБП G4 и G5 обусловливают очень высокий риск неблагоприятных почечных и сердечно-сосудистых исходов. Доказательства ССПМ-терапии для этой группы ограничены: большинство данных получено в подгрупповых анализах РКИ с недостаточной статистической мощностью. В целом они указывают на безопасность препаратов при G4--G5; некоторые предполагают сохранение благоприятных эффектов. При этом несколько РКИ нефропротективных средств показывают, что продолжение лечения после снижения рСКФ ниже порога начала может сопровождаться сохраняющейся сердечно-сосудистой пользой.

При почечной недостаточности --- терминальной болезни почек --- на диализе ССЗ часты и являются ведущей причиной смерти; их фенотипы имеют особенности.\src{9} При хроническом диализе особенно выражены СН и перегрузка внеклеточной жидкостью; более частый диализ потенциально может их уменьшать.\src{10} Кальцификация средней оболочки артерий способствует сосудистой болезни, что отчасти может объяснять отрицательные результаты испытаний статинов у диализных пациентов.\src{11}

Профилактика и лечение ССЗ на диализе представляют особые трудности из-за специфической патофизиологии, рисков и вопросов безопасности, а также малого числа клинических исследований. Применение ССПМ-терапии после трансплантации почки изучается и здесь не рассматривается.

\heading{Обоснование рекомендации}
\textbf{1.} Как описано в разделах 5.5.4, 6.2.3 и 6.3.3, RASi, SGLT2i, финеренон и GLP-1-терапия сохраняют функцию почек и связаны с улучшением почечных и сердечно-сосудистых исходов. Пороги рСКФ для начала: RASi $\geq30$~\unitgfr, SGLT2i $\geq20$~\unitgfr, финеренон $\geq25$~\unitgfr.

В нескольких крупных исследованиях разрешалось продолжать исследуемые препараты после снижения рСКФ ниже порога начала и даже после начала диализа; польза сохранялась без явного увеличения нежелательных эффектов.\src{3,8} В подгруппе ХБП G4 метаанализа FIDELITY при диабете направление снижения риска сердечно-сосудистой конечной точки --- сердечно-сосудистая смерть, нефатальный инфаркт, нефатальный инсульт или госпитализация по поводу СН --- на финереноне было сходным, но статистически незначимым и требует дальнейшего изучения.\src{6} Гиперкалиемия при ХБП G4 на финереноне встречалась чаще.\src{6}

Недавний метаанализ предполагает, что ингибиторы РААС при далеко зашедшей ХБП уменьшают частоту почечной недостаточности, требующей заместительной терапии, но не смерти.\src{7} Хотя контроль АД, липидов и гликемии является центральным компонентом ССПМ-помощи, при рСКФ $<15$~\unitgfr\ (ХБП G5) целевые значения не определены и зависят от клинического суждения и совместного принятия решений с пациентом.

\section{Важные дополнительные вопросы}
\subsection{Контроль ХС ЛНП при ССПМ-синдроме}
\heading{Обзор}
ССПМ-синдром способствует ускоренному атерогенезу и повышенному риску событий АССЗ.\src{1,2} Снижение ХС ЛНП изменением образа жизни и лекарствами занимает центральное место в первичной и вторичной профилактике. При гипертриглицеридемии также следует рассматривать ХС не-ЛВП или аполипопротеин B как более надёжные показатели количества атерогенных частиц, чем ХС ЛНП в этой подгруппе.\src{3,4}

На стадиях ССПМ 1--2 наличие СД2 или ХБП у взрослых $\geq40$ лет либо промежуточного прогнозируемого риска АССЗ (PREVENT-ASCVD 5--9,9\%) обосновывает применение статинов и других средств снижения ХС ЛНП для уменьшения событий АССЗ. При пограничном риске (PREVENT-ASCVD 3--4,9\%) гиполипидемическую терапию можно рассматривать в рамках обсуждения риска с учётом его усилителей.

На стадии 3 наличие и выраженность субклинического коронарного атеросклероза помогают определять цель ХС ЛНП. Согласно руководству ACC/AHA 2026 года по дислипидемии, при CAC 100--999 единиц Агатстона (AU) либо случайно выявленном умеренном или выраженном CAC рекомендуется ХС ЛНП $<70$~мг/дл. При очень выраженном CAC $\geq1000$~AU рекомендуется более интенсивная цель $<55$~мг/дл. VESALIUS-CV --- исследование эволокумаба при высоком сердечно-сосудистом риске без перенесённых инфаркта или инсульта --- показало сердечно-сосудистую пользу сходного интенсивного снижения ХС ЛНП при диабете высокого риска, коронарном атеросклерозе или их сочетании без инфаркта или инсульта в анамнезе.\src{5}

При установленном АССЗ (стадия ССПМ 4) такие компоненты, как диабет и ХБП, связаны с повышенным риском повторных событий; пациенты часто относятся к группе АССЗ очень высокого риска. Им показано интенсивное снижение ХС ЛНП до $<55$~мг/дл с помощью высокоинтенсивной терапии статинами, эзетимиба и при необходимости дополнительных препаратов: моноклональных антител к PCSK9, инклисирана либо бемпедоевой кислоты. Конкретные рекомендации приведены в руководстве ACC/AHA 2026 года по дислипидемии.\src{4}

\subsection{Стеатотическая болезнь печени, ассоциированная с метаболической дисфункцией}
\heading{Рекомендации}
Исследования, на которых основаны рекомендации, обобщены в таблице доказательств оригинала.
{\small
\begin{longtable}{P{0.07\textwidth}P{0.11\textwidth}P{0.75\textwidth}}
\toprule COR & LOE & Рекомендация\\\midrule
\endfirsthead
\toprule COR & LOE & Рекомендация (продолжение)\\\midrule
\endhead
1 & B-NR & \textbf{1.} У взрослых с ССПМ-синдромом и диабетом либо $\geq2$ кардиометаболическими факторами риска рекомендуется рассчитывать индекс Fibrosis-4 (FIB-4) каждые 1--2 года для оценки риска фиброза печени, связанного с MASLD.\src{1--3}\\[5pt]
2a & C-LD & \textbf{2.} У взрослых со стадией ССПМ 1 вследствие предиабета целесообразно рассчитывать FIB-4 каждые 2--3 года для оценки риска фиброза печени, связанного с MASLD.\src{4}\\[5pt]
1 & B-R & \textbf{3.} У взрослых с ССПМ-синдромом и MASLD рекомендуется изменение образа жизни --- физическая активность и питание для снижения массы тела, при необходимости дополненные лекарственной терапией ожирения и/или MBS, --- для уменьшения тяжести MASLD, риска её прогрессирования и улучшения профиля ССПМ-риска.\src{5--7}\\[5pt]
1 & B-R & \textbf{4.} У взрослых с ССПМ-синдромом и MASLD в сочетании с СД2 либо MASH с признаками клинически значимого фиброза печени показана терапия на основе GLP-1 с доказанной пользой для улучшения гистологической картины MASLD.\src{8--11}\\
\bottomrule
\end{longtable}
}

\heading{Обзор}
MASLD, ранее называвшаяся неалкогольной жировой болезнью печени, --- весьма распространённое хроническое прогрессирующее и угрожающее жизни заболевание.\src{12--14} Стеатогепатит, ассоциированный с метаболической дисфункцией (MASH, ранее NASH), является тяжёлой формой MASLD, ведущей причиной цирроза и рака печени и одним из основных показаний к её трансплантации.\src{14,15} MASLD --- патофизиологическое следствие эктопического накопления жира, воспаления и инсулинорезистентности, лежащих в основе ССПМ-синдрома.\src{16,17} При ней также высока частота ССЗ, поэтому нужны скрининг факторов ССПМ-риска и профилактика.

Стадия фиброза лучше всего прогнозирует неблагоприятные исходы MASLD.\src{1,18} Оценку риска начинают с FIB-4; при необходимости последовательно выполняют дальнейшие исследования для исключения либо подтверждения клинически значимого фиброза --- стадии $\geq2$.\src{19--21} Изменение образа жизни со снижением массы тела имеет решающее значение: снижение на $\geq10\%$ приводит к улучшению и ремиссии MASLD.\src{22,23} Метаболическая и бариатрическая хирургия вызывает ремиссию даже на поздних стадиях.\src{6,24,25}

При MASH и фиброзе стадии 2 или 3 GLP-1 RA семаглутид связан с уменьшением стеатогепатита и фиброза.\src{11,26} Ресметиром, агонист бета-рецептора тиреоидных гормонов, также одобрен для улучшения гистологической картины при MASH и фиброзе стадии 2--3. Статины при MASLD безопасны и эффективны и должны применяться согласно рекомендациям.\src{12,19,20,27} При клинически значимом фиброзе ($\geq2$) целесообразно направление к гастроэнтерологу или гепатологу для обсуждения медикаментозного и хирургического снижения массы тела, направленной на печень лекарственной терапии и последующего наблюдения.\src{12,19,20}

\heading{Обоснование отдельных рекомендаций}
\textbf{1. Скрининг при диабете и множественных факторах риска.} Метаанализы и объединённые исследования показывают наибольший риск MASLD и её прогрессирования при СД2 и/или более чем двух кардиометаболических факторах, например ожирении, гиперлипидемии, гипертензии.\src{14,28--30} По NHANES III распространённость и тяжесть MASLD увеличивались с числом метаболических сопутствующих состояний, особенно при СД2. Сочетание MASLD и СД2 имело наибольший риск смерти.\src{29,30}

В общей популяции FIB-4 --- показатель из доступных клинических данных (возраст) и лабораторных показателей (АСТ, АЛТ, тромбоциты) --- экономически эффективен, предсказывает клинические исходы и позволяет эффективно исключать клинически значимый фиброз (таблица 17).\src{31,32} Переход FIB-4 из низкого риска в промежуточный или высокий может использоваться для оценки прогрессирования MASLD.\src{33} При кардиометаболических факторах скрининг MASLD высокого риска рекомендуется каждые 1--2 года при СД2 либо $\geq2$ факторах и каждые 2--3 года при $<2$ факторах без СД2.\src{12,19,20,34}
\editorial{Начало обоснования 1 использует «более чем два» фактора, тогда как формальная рекомендация и заключительное предложение --- «два или более». Различие исходных формулировок сохранено.}

\heading{Таблица 17. Повышенный риск MASLD}
{\small
\begin{longtable}{P{0.17\textwidth}P{0.21\textwidth}P{0.54\textwidth}}
\toprule Возраст, годы & FIB-4\src{20} & Последующее ведение\src{55}\\\midrule
\endfirsthead
\toprule Возраст, годы & FIB-4 & Последующее ведение (продолжение)\\\midrule
\endhead
$<35$ & \multicolumn{2}{P{0.75\textwidth}}{Не применяется}\\[4pt]
35--64 & $<1{,}3$ & Плановое наблюдение\\
 & 1,3--2,67 & Направить на VCTE или ELF\\
 & $>2{,}67$ & Направить к гепатологу\\[4pt]
$\geq65$ & $<2{,}0$ & Плановое наблюдение\\
 & 2,0--2,67 & Направить на VCTE или ELF\\
 & $>2{,}67$ & Направить к гепатологу\\
\bottomrule
\end{longtable}
}
\textbf{Примечание исходной таблицы 17.} Расчёт индекса фиброза печени:
\[
\mathrm{FIB}\text{-}4 =
\frac{\text{возраст}\times\mathrm{AST}}
{\text{число тромбоцитов}\times\sqrt{\mathrm{ALT}}}.
\]
FIB-4 не следует применять у пациентов с острым заболеванием; см. таблицу 18. ALT --- АЛТ, аланинаминотрансфераза; AST --- АСТ, аспартатаминотрансфераза; ELF --- расширенный тест фиброза печени (Enhanced Liver Fibrosis); VCTE --- вибрационно-контролируемая транзиентная эластография.
\editorial{Формула воспроизводит оригинал, в котором единицы входных лабораторных показателей не указаны. Для практического расчёта необходимо использовать валидированный калькулятор с указанными единицами; произвольная подстановка результатов в иных единицах недопустима.}

\textbf{2. Скрининг при предиабете.} MASLD --- ведущая причина печёночной недостаточности в США и частое сопутствующее состояние на всём спектре ССПМ.\src{15,35} Она также является независимым фактором сердечно-сосудистой заболеваемости и смертности.\src{36} Наибольший риск при MASLD имеют пациенты с выраженным фиброзом печени.\src{37,38} MASLD часто изучалась в контексте диабета,\src{39,40} но данные указывают на риск и при предиабете;\src{4} эндокринологические руководства признают предиабет состоянием высокого риска MASLD.\src{34}

Популяция предиабета велика и неоднородна; специфических доказательств для предиабета без других факторов риска мало. При имеющихся данных целесообразен скрининг выраженного фиброза у взрослых со стадией ССПМ 1 вследствие предиабета. Американская ассоциация по изучению болезней печени рекомендует FIB-4 как предпочтительный неинвазивный скрининговый тест благодаря высокой отрицательной прогностической ценности.\src{20} Интерпретация и дальнейшие действия зависят от возраста и результата (таблица 17). Вторичная оценка может включать VCTE или ELF; при отклонениях следует направлять к гастроэнтерологу/гепатологу (таблица 18).\src{20}

\heading{Таблица 18. Средства вторичной оценки риска клинически значимого фиброза при риске MASLD}
{\small
\begin{longtable}{P{0.25\textwidth}P{0.11\textwidth}P{0.18\textwidth}P{0.36\textwidth}}
\toprule Метод & Результат & Риск & Действие\\\midrule
\endfirsthead
\toprule Метод & Результат & Риск & Действие (продолжение)\\\midrule
\endhead
VCTE* & $<8$ & Низкий & Стандартное наблюдение за риском MASLD каждые 1--3 года\textsuperscript{†} с расчётом FIB-4\\[4pt]
 & 8--12 & Промежуточный & Направление к гастроэнтерологу/гепатологу\\[4pt]
 & $>12$ & Высокий & Направление к гастроэнтерологу/гепатологу\\[4pt]
ELF\textsuperscript{‡}\src{55} & $<7{,}7$ & Низкий & Стандартное наблюдение за риском MASLD каждые 1--3 года\textsuperscript{†} с расчётом FIB-4\\[4pt]
 & 7,7--9,8 & Промежуточный & Направление к гастроэнтерологу/гепатологу\\[4pt]
 & $>9{,}8$ & Высокий & Направление к гастроэнтерологу/гепатологу\\
\bottomrule
\end{longtable}
}
\textbf{Примечания исходной таблицы 18.}

*VCTE --- ультразвуковой метод, выполняемый непосредственно в клинике. Требуется голодание не менее 3 часов. Телосложение влияет на точность; метод не рекомендуется при ИМТ $\geq40$~кг/м$^2$.

\textsuperscript{†}FIB-4 рассчитывается как $(\text{возраст}\times\mathrm{AST})/(\text{число тромбоцитов}\times\sqrt{\mathrm{ALT}})$. FIB-4 рекомендуется каждые 1--2 года при СД2 либо $\geq2$ факторах ССПМ-риска; при стадии ССПМ 1 вследствие предиабета целесообразен интервал 2--3 года.

\textsuperscript{‡}ELF --- показатель на основе биомаркеров крови; образец направляется в референсную лабораторию.

ALT --- АЛТ; AST --- АСТ; CKM --- ССПМ; GI --- гастроэнтерология; MASLD --- стеатотическая болезнь печени, ассоциированная с метаболической дисфункцией.
\editorial{В исходной таблице 18 единицы результатов VCTE не напечатаны; числа перенесены без самостоятельного добавления единиц. Общие интервалы 1--3 года конкретизируются сноской †.}

\textbf{3. Снижение массы тела и физическая активность.} Метаанализ показал связь различных методов снижения массы тела с клинически значимым улучшением биомаркеров болезни печени при MASLD, хотя долгосрочных исследований клинических исходов не хватает.\src{5} При избыточной массе тела/ожирении снижение общей массы на $\geq10\%$ может обеспечить ремиссию MASLD и улучшение воспаления, фиброза и кардиометаболических факторов.\src{5,22} Даже снижение на 3--5\% уменьшает жир в печени, в том числе при исходно нормальной массе тела.\src{5,41}

GLP-1-терапия может эффективно дополнять лечение для достижения целей массы тела, с дополнительными гистологическими преимуществами, описанными далее.\src{8,9,42} MBS в крупных когортах связана с меньшим риском впервые возникших больших неблагоприятных печёночных исходов и MACE.\src{6,24,25}

В метаанализе физические упражнения, даже без диетического вмешательства или снижения массы тела, уменьшали стеатоз и показатели биомаркеров печени; в той же фразе оригинала сообщается об уменьшении кардиореспираторной тренированности.\src{43} Тренировки улучшают качество жизни и ряд кардиометаболических показателей даже при далеко зашедшей болезни печени.\src{44,45} Наблюдательные исследования связывают регулярную физическую активность с меньшим риском фиброза, цирроза, рака печени и общей смерти.\src{46--48} Малоподвижное поведение --- независимый предиктор MASLD и связано с высоким риском её прогрессирования.\src{49,50}
\editorial{В предложении о метаанализе физических упражнений оригинал грамматически относит слово «reduced» также к «cardiorespiratory fitness». Это не согласуется с обычным смыслом улучшения тренированности и отмечено как неоднозначность исходного текста; направление эффекта не исправлено без оговорки.}

\textbf{4. Терапия на основе GLP-1 и другие препараты.} GLP-1-терапия улучшает массу тела, инсулинорезистентность и неинвазивные показатели печени, а также благоприятно влияет на сердечно-сосудистые и почечные исходы.\src{51} В исследовании фазы IIb по подбору дозы семаглутида разрешение MASH при фиброзе стадий 1--3 было значительно чаще на дозе 0,4~мг, чем на плацебо.\src{8} ESSENCE --- исследование семаглутида при неалкогольном стеатогепатите --- при MASH и фиброзе стадии 2 или 3 показало улучшение стеатогепатита и фиброза относительно плацебо.\src{11,52}

В подисследовании рандомизированного открытого исследования фазы III SURPASS-3 с параллельными группами тирзепатид, агонист рецепторов GLP-1/GIP, одобренный для СД2, значимо уменьшал содержание жира в печени по сравнению с инсулином деглудек.\src{9} Исследования «случай--контроль» при циррозе и СД2 связывают применение GLP-1 RA с улучшением печёночных исходов.\src{53,54}

Среди других препаратов ресметиром стал первым одобренным FDA средством для MASLD, у которого улучшение гистологической картины не зависит от снижения массы тела.\src{26} При клинически значимом фиброзе рекомендуется участие гепатолога для выбора направленной на печень терапии и последующего мониторинга.

\subsection{Венозная тромбоэмболия}
\heading{Обзор}
ССПМ-синдром связан с протромботическим состоянием, способным привести к венозной тромбоэмболии (ВТЭ).\src{1} При острой ВТЭ риск рецидива ступенчато увеличивается с числом компонентов ССПМ-синдрома.\src{2--4} Абдоминальная жировая ткань --- один из наиболее сильных факторов риска ВТЭ.\src{3,5--7} Доказательства связи диабета с ВТЭ неоднозначны,\src{8--15} но более низкая рСКФ, высокий UACR и более поздние стадии ХБП связаны с развитием ВТЭ и худшими исходами.\src{16--20}

Лечение и профилактика ВТЭ включают парентеральные антикоагулянты --- например нефракционированный или низкомолекулярный гепарин --- либо пероральные: DOAC или антагонист витамина K (VKA). Преимущества DOAC перед VKA --- быстрое начало действия, предсказуемый эффект, более широкое терапевтическое окно, отсутствие необходимости рутинного мониторинга и меньше взаимодействий с пищей и лекарствами.\src{21} Появляющиеся данные предполагают, что DOAC также могут предотвращать сердечно-сосудистые события при ВТЭ.\src{22--24}

При ХБП, ожирении и после MBS антикоагуляция сложнее из-за изменённой фармакокинетики.\src{21} Здесь рассматривается острая ВТЭ при ожирении или ХБП без метаболической и бариатрической операции.\src{25} Специфическое ведение тромбоэмболии лёгочной артерии (ТЭЛА), включая пациентов с ожирением или ХБП, изложено в руководстве AHA/ACC 2026 года по острой ТЭЛА.\src{26}

\textbf{Ожирение.} Проспективных РКИ, сравнивающих отдельные DOAC с низкомолекулярным гепарином/VKA при ВТЭ, нет. Обычно DOAC и VKA сравнивают внутри весовых категорий\src{24,27} либо между ними.\src{28,29} Большинство исследований объединяет несколько DOAC, а не оценивает препараты отдельно. При ИМТ $\geq30$~кг/м$^2$ объединённое применение DOAC было связано со снижением комбинированных событий инсульта/системной эмболии/инфаркта/общей смерти на 20\% и больших кровотечений на 44\% по сравнению с VKA при сходной профилактике ВТЭ.\src{24}

При ожирении III степени метаанализы показали большую эффективность и безопасность DOAC, чем VKA, преимущественно в исследованиях ривароксабана или апиксабана.\src{30--33} При более тяжёлом ожирении --- ИМТ $\geq50$~кг/м$^2$ или масса $>150$~кг --- данных мало.\src{27,34}

\textbf{ХБП.} Высоки риски как образования ВТЭ, так и кровотечения на антикоагулянтах. Почечный клиренс DOAC различается: дабигатран 80\%, эдоксабан 50\%, ривароксабан 35\%, апиксабан 27\%.\src{21} В исследования ВТЭ фазы III не включали пациентов с рСКФ $<30$~\unitgfr\ или на заместительной почечной терапии. При ХБП G3 (рСКФ 30--59~\unitgfr) подгрупповые анализы предполагают, что DOAC не уступают VKA по рецидивам ВТЭ или большим кровотечениям.\src{35--38} Метаанализы ХБП G2--G3 (рСКФ 30--89~\unitgfr) показывают уменьшение больших кровотечений на большинстве DOAC по сравнению с VKA при сходной профилактике ВТЭ.\src{39--42}

При терминальной болезни почек на диализе два высококачественных наблюдательных исследования базы US Renal Data System связали апиксабан с меньшей частотой больших кровотечений и сходной смертностью относительно варфарина.\src{43,44} Метаанализ наблюдательных исследований ХБП G4--G5 показал меньшие риски рецидива ВТЭ и кровотечения на апиксабане при сходной общей смертности.\src{45} Данных о ривароксабане, эдоксабане и дабигатране при далеко зашедшей ХБП крайне мало.\src{35,46}

\subsection{Обструктивное апноэ сна}
\heading{Рекомендации}
Исследования, на которых основаны рекомендации, обобщены в таблице доказательств оригинала.
{\small
\begin{longtable}{P{0.07\textwidth}P{0.11\textwidth}P{0.75\textwidth}}
\toprule COR & LOE & Рекомендация\\\midrule
\endfirsthead
\toprule COR & LOE & Рекомендация (продолжение)\\\midrule
\endhead
2a & C-LD & \textbf{1.} У взрослых с ССПМ-синдромом и ожирением целесообразна ежегодная оценка симптомов апноэ сна, включая валидированные скрининговые инструменты, для облегчения выявления обструктивного апноэ сна (OSA).\src{1,2}\\[5pt]
1 & B-R & \textbf{2.} У взрослых с ССПМ-синдромом и OSA стратегия лечения должна включать вмешательства по контролю массы тела при наличии показаний в дополнение к постоянному положительному давлению в дыхательных путях (CPAP) и другим методам, чтобы улучшить симптомы OSA и уменьшить риск прогрессирования стадии ССПМ.\src{3--6}\\
\bottomrule
\end{longtable}
}

\heading{Обзор}
По когортному исследованию 2007--2010 годов распространённость OSA, определяемого индексом апноэ--гипопноэ (AHI) $\geq5$ событий/ч, среди взрослых США 30--70 лет оценивается в 33,9\% у мужчин и 17,4\% у женщин.\src{7} Консервативное моделирование прогнозирует значительный рост к 2050 году: до 55\% мужчин и 35\% женщин в возрасте 30--69 лет.\src{8}

ССЗ и OSA имеют много общих факторов ССПМ-риска. Поэтому OSA чаще встречается при ССПМ-синдроме и ожирении, чем в общей популяции, а ССПМ-синдром --- чаще среди людей с OSA. OSA вызывает фрагментацию и недостаток сна. Последствия связаны как с прямым ухудшением функционирования и качества жизни, так и с активацией патофизиологических механизмов, способствующих прогрессированию ССПМ-синдрома: ССЗ, цереброваскулярных событий, СД2, гипертензии.\src{9,10}

Лечение OSA улучшает симптомы, качество жизни и АД.\src{11} Однако РКИ CPAP в сравнении с обычной помощью не показали уменьшения смертности, сердечно-сосудистых или почечных событий.\src{11--14} Тяжёлое OSA связано с повышенной общей смертностью;\src{9,10} сердечно-сосудистая смертность среди лиц с OSA в США увеличивалась на протяжении последних двух десятилетий.\src{8}

\heading{Обоснование отдельных рекомендаций}
\textbf{1. Скрининг.} Американская академия медицины сна рекомендует ежегодный скрининг взрослых из отдельных групп высокого риска: ожирение, СН, ФП, резистентная гипертензия, нарушение толерантности к глюкозе или СД2, ночные нарушения ритма, инсульт, лёгочная гипертензия, подготовка к бариатрической операции и коронарная болезнь.\src{1} Многие из этих факторов характерны для ССПМ-синдрома, при котором распространённость OSA высока. Поэтому, несмотря на ограниченные доказательства, скрининг при ССПМ-синдроме и ожирении целесообразен.

Возможны целенаправленный опрос о симптомах, валидированные анкеты либо скрининговые устройства для апноэ. Часто применяются Берлинский опросник, STOP-BANG и STOP.\src{15,16} Необходимо учитывать их ограничения и сохранять клиническую настороженность: в отдельных популяциях, особенно у женщин и пациентов с ССЗ, они могут работать хуже.\src{2,15}

При клиническом подозрении для подтверждения диагноза следует выполнить полисомнографию либо домашнее исследование апноэ технически подходящим устройством.\src{15,17} Полисомнография предпочтительна при значимых кардиореспираторных заболеваниях, возможной слабости дыхательных мышц вследствие нервно-мышечных болезней, гиповентиляции в бодрствовании или подозрении на гиповентиляцию во сне, хроническом приёме опиоидов, инсульте в анамнезе либо тяжёлой бессоннице.\src{15,17}

\textbf{2. Снижение массы тела и лечение OSA.} При избытке массы тела её снижение более чем на 15\% существенно улучшает OSA и уменьшает AHI. Такая степень снижения массы тела также улучшает факторы ССПМ-риска и потенциально влияет на события ССЗ.

В SURMOUNT-OSA --- исследовании тирзепатида при OSA ---\src{5} тирзепатид с титрацией подкожной дозы до 15~мг по сравнению с плацебо улучшал средний AHI на $-20{,}0$ событий/ч (95\% ДИ от $-25{,}8$ до $-14{,}2$) у лиц без CPAP и на $-23{,}8$ событий/ч (от $-29{,}6$ до $-17{,}9$) у получавших CPAP. На тирзепатиде снижались высокочувствительный С-реактивный белок, АД и гипоксическая нагрузка, обусловленная апноэ сна.

MBS, вызывающая снижение массы тела более чем на 20\%, также улучшает OSA и связана с уменьшением сердечно-сосудистых и почечных событий, сердечно-сосудистой и общей смертности у людей с OSA.\src{3,18} В Sleep AHEAD --- исследовании апноэ сна в рамках Look AHEAD ---\src{4} у пациентов с СД2, ожирением и OSA, снизивших массу тела на $\geq10$~кг посредством поведенческого вмешательства, AHI существенно улучшался.

Учитывая доказанную пользу значимого снижения массы тела для OSA и исходов ССПМ-синдрома, контроль массы тела должен быть основой лечения OSA при ССПМ-синдроме.\src{19} Рекомендуется направление к специалисту по медицине сна для обсуждения CPAP и других методов уменьшения симптомов. CPAP улучшает контроль АД, но уменьшение событий ССЗ не доказано.\src{15}

\subsection{Беременность и ССПМ-здоровье}
\heading{Рекомендации}
Исследования, на которых основаны рекомендации, обобщены в таблице доказательств оригинала.
{\small
\begin{longtable}{P{0.07\textwidth}P{0.11\textwidth}P{0.75\textwidth}}
\toprule COR & LOE & Рекомендация\\\midrule
\endfirsthead
\toprule COR & LOE & Рекомендация (продолжение)\\\midrule
\endhead
\multicolumn{3}{P{0.93\textwidth}}{\textbf{Ведение ССПМ до родов}}\\[4pt]
1 & B-NR & \textbf{1.} Лицам с ССПМ-синдромом, планирующим беременность, рекомендуется оптимизировать массу тела, функцию почек, гликемию и контроль АД посредством изменения образа жизни и подходящей лекарственной терапии для уменьшения риска неблагоприятных исходов беременности (APO) и оптимизации ССПМ-здоровья после родов.\src{1--4}\\[5pt]
1 & B-R & \textbf{2.} Лицам с диабетом, планирующим беременность, рекомендуется оптимизировать гликемический контроль с целью HbA1c $<6{,}5\%$ посредством изменения образа жизни и при необходимости лекарств для уменьшения риска APO и оптимизации ССПМ-здоровья после родов.\src{5--7}\\[5pt]
1 & C-LD & \textbf{3.} Лица, планирующие беременность, со стадиями ССПМ 2--4 должны получать помощь междисциплинарной команды для уменьшения риска APO и оптимизации ССПМ-здоровья после родов.\src{8--11}\\[5pt]
\multicolumn{3}{P{0.93\textwidth}}{\textbf{Ведение ССПМ после родов}}\\[4pt]
1 & B-NR & \textbf{4.} Лицам, перенёсшим APO, рекомендуется скрининг факторов ССПМ-риска --- АД, липидов, гликемии, ХБП, ИМТ и окружности талии --- как минимум один раз в течение первого года после родов, а также консультирование по образу жизни для выбора мер оптимизации ССПМ-здоровья.\src{12--15}\\[5pt]
1 & B-NR & \textbf{5.} Лицам, перенёсшим APO, рекомендуется переход от послеродового наблюдения к длительному ведению в первичном звене для мониторинга факторов ССПМ-риска и оптимизации ССПМ-здоровья.\src{16,17}\\
\bottomrule
\end{longtable}
}

\heading{Обзор}
Беременность сопровождается многочисленными физиологическими сосудистыми и метаболическими изменениями, поддерживающими рост плода. Однако нарушение их регуляции с избыточной инсулинорезистентностью, накоплением жира, гиперкоагуляцией, ремоделированием сердца и снижением сосудистого сопротивления\src{12,18} может приводить к APO. К APO относятся гестационный диабет, гипертензивные расстройства беременности, малые размеры плода или новорождённого для гестационного возраста, низкая масса тела при рождении, отслойка плаценты и преждевременные роды. Их факторы риска пересекаются с факторами ССПМ-синдрома.\src{19}

Для минимизации риска APO женщинам с ССПМ-синдромом, планирующим беременность, следует пройти консультирование до зачатия\src{20--22} и оптимизировать факторы ССПМ-риска изменением образа жизни и при необходимости лекарствами. При стадиях ССПМ 2--4 помощь должна оказываться многопрофильной межпрофессиональной командой.

APO также связаны с повышенным последующим риском факторов ССПМ --- гипертензии, диабета, дислипидемии --- и ССЗ, включая СН, коронарную болезнь, заболевание периферических артерий и сосудистую деменцию.\src{12,23--27} Для уменьшения последующих неблагоприятных сердечно-сосудистых событий после APO целесообразны скрининг и коррекция факторов риска ССЗ в течение года после родов. Рекомендуется междисциплинарный подход с участием первичного звена. Конкретные рекомендации по АД и дислипидемии приведены в руководствах AHA/ACC 2025 года по высокому АД и ACC/AHA 2026 года по дислипидемии.\src{28}

\heading{Обоснование отдельных рекомендаций}
\textbf{1. Оптимизация факторов риска до беременности.} Неблагоприятные факторы ССПМ-риска до беременности --- избыточная масса тела/ожирение, гипертензия, гипертриглицеридемия, диабет, ХБП --- связаны с повышенным риском APO.\src{18,29} Наблюдательные когорты поддерживают их оптимизацию изменением образа жизни, лекарствами и/или хирургическим лечением до беременности.\src{29--31}

В когорте 226\,958 беременностей снижение массы тела до беременности на 10\% было связано с уменьшением риска APO примерно на 10\%.\src{1} Систематические обзоры и метаанализы выявили меньший риск APO после бариатрической операции до беременности по сравнению с отсутствием операции.\src{2,32} Оптимальное питание и физическая активность до беременности также связаны с меньшим риском APO.\src{3}

В наблюдательных исследованиях более здоровое питание, включая средиземноморскую модель, было связано с меньшими шансами гипертензивных расстройств беременности: отношение шансов (OR) 0,6; 95\% ДИ 0,4--0,9.\src{33,34} Метаанализ показал связь большей физической активности до беременности с меньшим риском гестационного диабета: OR 0,45; 95\% ДИ 0,28--0,75.\src{35} В CHAP --- исследовании хронической гипертензии и беременности --- оптимизация АД во время беременности до $<140/90$~мм рт. ст. сопровождалась меньшим риском неблагоприятных материнских и плодовых исходов.\src{36} Хотя лекарства могут улучшать факторы ССПМ-риска до беременности, препаратов, противопоказанных во время беременности, следует избегать, когда это безопасно возможно.

\textbf{2. Диабет при планировании беременности.} Женщины с диабетом имеют больший риск APO, чем женщины без диабета или с диабетом, возникающим во время беременности.\src{18} Риск напрямую связан с гипергликемией до и во время беременности.\src{5} Наблюдательные исследования и ADA поддерживают оптимизацию HbA1c до беременности до $<6{,}5\%$.\src{5,12} Инсулин является стандартом контроля глюкозы при беременности и при диабете 1 типа. При СД2 на неинсулиновых сахароснижающих средствах инсулин можно начать до беременности.\src{37}

РКИ изменения образа жизни до беременности с оценкой APO при диабете немного.\src{6,7} В РКИ 199 взрослых с избыточной массой тела/ожирением и ГСД в анамнезе 16-недельное вмешательство по изменению образа жизни до беременности, продолжавшееся до зачатия, давало большее снижение массы тела: 4,8~кг (95\% ДИ 3,4--6,0) против 0,7~кг (от $-0{,}9$ до 2,3) в контроле. Доля снизивших массу тела на $\geq5\%$ составила 50,0\% (17/34) против 13,6\% (3/22).\src{7} Однако значимого различия рецидивов ГСД не было: 57,9\% во вмешательстве против 44,0\% в контроле.

До беременности следует определить HbA1c, креатинин и UACR. При планировании беременности целевое АД должно соответствовать междисциплинарному руководству 2025 года. При гипертензии, ХБП и диабете необходимо пересмотреть и заменить тератогенные препараты, например ACEi и ARB.\src{5,28}

\textbf{3. Межпрофессиональная помощь.} Планирующих беременность женщин со стадией ССПМ $\geq2$ должна обследовать и наблюдать межпрофессиональная команда, которая может включать акушера, специалиста по медицине матери и плода, эндокринолога, нефролога, кардиолога и зарегистрированного диетолога. Состав определяется индивидуальным профилем риска.\src{38} РКИ мало, но наблюдательные данные, особенно при диабете и ХБП, поддерживают этот подход.

В проспективной когорте помощь до беременности при диабете была связана с лучшей гликемией и меньшим риском APO, чем у не посещавших такую программу: OR 0,2; 95\% ДИ 0,05--0,89.\src{8} Метаанализы и когорты показывают повышенный риск APO при ХБП.\src{9--11,39} Поэтому женщинам с ХБП требуется тщательное дородовое наблюдение межпрофессиональной командой, включая нефрологов и специалистов по медицине матери и плода, с оптимизацией факторов риска ещё до беременности.\src{11}

Особые вопросы --- контроль АД до и во время беременности, выбор лекарств с благоприятным профилем безопасности для беременности и отказ от терапии, способной ухудшать функцию почек, --- рассмотрены в руководстве AHA/ACC 2025 года по АД и связанных рекомендациях ACOG.\src{11,28} Если такая команда недоступна, следует рассмотреть направление к соответствующим узким специалистам.

\textbf{4. Скрининг после APO.} AHA и Американская коллегия акушеров и гинекологов (ACOG) рекомендуют частый скрининг факторов риска ССЗ и консультирование по образу жизни после родов при APO.\src{40,41} Скрининг и консультирование по ССПМ-риску рекомендуются в течение первого послеродового года.\src{12} Скрининг включает АД, массу тела/ИМТ/окружность талии, HbA1c, липиды, сывороточный креатинин и UACR.

Наблюдательные данные предполагают оценку отдельных факторов, например ИМТ и окружности талии, через 6--12 месяцев после родов;\src{42} другие требуют более частого контроля, например более ранней оценки АД для выявления и лечения гипертензии.\src{13,43--45} Без ССЗ в анамнезе следует применять рекомендации ACC/AHA по первичной профилактике,\src{46} а при установленном ССЗ --- по вторичной.\src{47,48}

Доказательств для оптимальных порогов начала лечения АД после родов мало. По гипертензии и гипертензивным расстройствам беременности следует обратиться к разделу 5.5 руководства AHA/ACC 2025 года по АД,\src{28} а по СН --- к разделу 11.3 руководства AHA/ACC 2022 года.\src{49} После медицинского допуска всем женщинам после родов рекомендуется не менее 150 минут умеренной физической активности в неделю, в соответствии с рекомендациями для небеременных взрослых.\src{15,46} Лактация связана с более благоприятным кардиометаболическим профилем: меньшими глюкозой натощак, инсулинорезистентностью, триглицеридами и АД.\src{12}

\textbf{5. Переход к длительному наблюдению.} Первые 12 недель после родов, иногда называемые четвёртым триместром, --- важнейшее окно для оптимизации ССПМ-здоровья, особенно после APO.\src{40} AHA и ACOG призывают улучшать передачу пациентки от акушерской службы в первичное звено\src{16,50} и при необходимости к кардиологу и другим специалистам, например нефрологу и эндокринологу. Это необходимо для раннего выявления и коррекции гипертензии, массы тела, гликемии, стресса и настроения и для уменьшения долгосрочного сердечно-сосудистого риска.

Небольшие локальные исследования изучали переходные послеродовые клиники после гипертензивных расстройств беременности или ГСД, интеграцию помощи матери в педиатрические посещения, навигационную поддержку, мониторинг АД и электронные напоминания для своевременного перехода в первичное звено.\src{17} Строгих крупных РКИ с оценкой сердечно-сосудистых исходов пока не было.\src{51}

Посещаемость первичного звена низкая: в одном исследовании лишь 60\% женщин после APO посетили его в течение 12 месяцев после родов.\src{52} При неблагоприятных SDOH посещаемость ещё ниже. Поэтому нужны стратегии внедрения и распространения программ переходного ведения, включая междисциплинарные модели, оценивающие и учитывающие SDOH.\src{40}

\section{Мониторинг и наблюдение после начала терапии}
\heading{Рекомендации}
Исследования, на которых основаны рекомендации, обобщены в таблице доказательств оригинала.
{\small
\begin{longtable}{P{0.07\textwidth}P{0.11\textwidth}P{0.75\textwidth}}
\toprule COR & LOE & Рекомендация\\\midrule
\endfirsthead
\toprule COR & LOE & Рекомендация (продолжение)\\\midrule
\endhead
1 & B-NR & \textbf{1.} У пациентов с ССПМ-синдромом, начавших GLP-1-терапию для снижения массы тела, следует периодически, в пределах 3--6 месяцев, повторно оценивать недостаточный ответ. При снижении массы тела менее чем на 5\% необходимо повышать дозу, переходить на другой препарат со сходными преимуществами при ССПМ-синдроме и/или направлять к специалисту по контролю массы тела для обсуждения дополнительных подходов к её оптимальному снижению.\src{1--4}\\[5pt]
1 & A & \textbf{2.} У пациентов со стадиями ССПМ 2--4 и СД2, начавших кардиопротективную сахароснижающую терапию SGLT2i или препаратом на основе GLP-1, следует оценивать гликемию по HbA1c, мониторированию глюкозы крови и/или показателям непрерывного мониторирования каждые 3--6 месяцев, а при недостижении целей --- чаще, для выбора дальнейшей тактики.\src{5--7}\\[5pt]
2a & B-R & \textbf{3.} У пациентов со стадиями ССПМ 2--4 и альбуминурией (UACR $\geq30$~мг/г), начавших нефропротективную терапию, целесообразно повторно измерить альбуминурию через 3--6 месяцев для оценки остаточного ССПМ-риска и показаний к дополнительной нефропротективной терапии, предотвращающей события ССЗ и утрату функции почек.\src{8--11}\\[5pt]
2a & B-R & \textbf{4.} У пациентов со стадиями ССПМ 2--4, начавших препарат, действующий на РААС (RASi или MRA), целесообразно повторно рассчитать рСКФ и определить калий через 2--4 недели для оценки переносимости и безопасности схемы лечения.\src{12--15}\\
\bottomrule
\end{longtable}
}

\heading{Обзор}
Эффективная лекарственная терапия ССПМ-синдрома требует устойчивой переносимости, приверженности и доступности лечения. Возможные нежелательные эффекты выявляют по анамнезу, физикальному обследованию и/или лабораторным исследованиям и уменьшают коррекцией схемы. Например, RASi и nsMRA могут вызывать гипотензию, изменения функции почек и гиперкалиемию, которые можно уменьшать изменением доз, сопутствующими поведенческими или лекарственными мерами либо отменой терапии.

Регулярная оценка нежелательных эффектов и других препятствий лечению важна для безопасного длительного применения и приверженности взаимодополняющим изменениям образа жизни и лекарствам. Регулярная оценка ответа позволяет титровать дозы, выбирать альтернативные методы, выявлять остаточный ССПМ-риск и своевременно добавлять другие вмешательства, включая комбинации. Поэтому мониторинг и безопасности, и ответа имеет решающее значение для оптимальной индивидуальной схемы. По наблюдению при лекарственном лечении гипертензии и нарушений холестерина следует обращаться к руководствам AHA/ACC по высокому АД\src{16} и дислипидемии.\src{17}

\heading{Обоснование отдельных рекомендаций}
\textbf{1. Ответ на лекарственную терапию ожирения.} Исторически эффективным ответом считали снижение массы тела не менее чем на 5\% после 3 месяцев лечения.\src{18} Более новые GLP-1-препараты обеспечивают продолжающееся снижение массы тела до 52 недель, а нередко до 68--72 недель лечения.\src{3,19} Поэтому окно оценки пользы может быть продлено с 3 до 6 месяцев.

Оптимальное снижение массы тела требует постоянного лечения и надлежащего повышения доз; прекращение приводит к возврату массы тела и ухудшению профиля ССПМ-риска.\src{2} Следует заранее объяснять фазу плато при лекарственной терапии ожирения и необходимость продолжать работу над образом жизни. Пациенты должны понимать, что GLP-1-терапия --- долгосрочная стратегия, требующая постоянной приверженности.

В исследованиях примерно 10--15\% пациентов недостаточно снижали массу тела на GLP-1-терапии. Если ожидаемого результата нет несмотря на повышение дозы, рекомендуется переход на другой GLP-1-препарат, обеспечивающий сходную защиту при других ССПМ-состояниях: показано, что такой переход улучшает снижение массы тела и гликемический контроль.\src{1,20} Если желаемый результат всё ещё не достигнут, можно рассматривать MBS. Пациентов следует консультировать о возможных нежелательных эффектах (таблица 15) и использовать меры уменьшения частых желудочно-кишечных реакций (таблица 19).

\heading{Таблица 19. Ведение желудочно-кишечных нежелательных эффектов терапии на основе GLP-1}
\heading{Изменение образа жизни}
\begin{itemize}
\item \textbf{Пищевые привычки:} чаще есть небольшими порциями, медленно; прекращать приём пищи при насыщении; избегать жирной и острой пищи; ограничивать алкоголь или отказываться от него.
\item \textbf{Гидратация:} в течение дня пить достаточное количество воды и напитков без сахара.
\item \textbf{Активность:} избегать интенсивной нагрузки после еды; предпочитать лёгкую активность, например ходьбу.
\item \textbf{Сон:} после еды выждать не менее 2--3 часов, прежде чем ложиться.
\item \textbf{Расслабление:} методы, например направленное воображение, могут помогать контролировать тошноту.
\end{itemize}
\heading{Дополнительные средства}
Рассмотреть временное применение противорвотных и/или прокинетических препаратов либо пробиотических и/или противодиарейных средств.
\heading{Повышение дозы}
Более медленная титрация и более длительное ожидание перед следующим повышением могут уменьшить желудочно-кишечные нежелательные эффекты, обычно возникающие при изменении дозы.
\begin{itemize}
\item Продлить текущий этап ещё на 2--4 недели перед переходом к следующей дозе.
\item Вернуться к предыдущей дозе на несколько дней, затем постепенно увеличивать её.
\end{itemize}
\heading{Переход на альтернативную терапию}
Рассмотреть другой препарат на основе GLP-1. Например, дулаглутид может переноситься лучше семаглутида или тирзепатида, хотя эффективность снижения массы тела может быть меньше.

\textbf{Примечание исходной таблицы 19.} Многие желудочно-кишечные эффекты со временем уменьшаются при постоянной дозе, но переносимость существенно различается. GLP-1 --- глюкагоноподобный пептид-1; GLP-1 RA --- агонист его рецептора.

\textbf{2. Контроль гликемии.} Оптимальная оценка --- HbA1c, измерение капиллярной глюкозы из пальца или непрерывное мониторирование каждые 3--6 месяцев либо чаще, например ежемесячно, при изменении лечения, недостижении целей, частой или тяжёлой гипогликемии либо гипергликемии (таблица 20).\src{5--7,21--23} Наибольшее снижение HbA1c дают схемы, включающие GLP-1-терапию, особенно семаглутид и тирзепатид, в дополнение к инсулину.\src{24,25}

При СД2 с симптомами гипергликемии либо очень высоким HbA1c или глюкозой --- HbA1c $>10\%$ или глюкоза крови $\geq300$~мг/дл --- рекомендуется направление к эндокринологу для дальнейшего ведения согласно ADA.\src{26} При СД2 с ССЗ или высоким его риском и лёгкой/умеренной остаточной гипергликемии, например HbA1c $<8{,}5$--9,0\%, к кардиопротективной терапии SGLT2i или GLP-1 можно добавить метформин.\src{27,28} При значительной остаточной гипергликемии, HbA1c $\geq8{,}5$--9,0\%, рекомендуется добавить семаглутид или тирзепатид, если пациент уже получает SGLT2i.\src{26,29,30}

При добавлении GLP-1-терапии или повышении её дозы следует пересмотреть дозу инсулина из-за риска гипогликемии. Совместное применение ингибиторов дипептидилпептидазы-4 с GLP-1-терапией не рекомендуется, поскольку дополнительной сахароснижающей пользы нет.\src{26,31}

\heading{Таблица 20. Достижение гликемических целей при СД2 и ССЗ или высоком риске ССЗ после начала кардиопротективной терапии}
{\small
\begin{longtable}{P{0.30\textwidth}P{0.64\textwidth}}
\toprule HbA1c & Тактика\\\midrule
\endfirsthead
\toprule HbA1c & Тактика (продолжение)\\\midrule
\endhead
7--8,5\%: лёгкая или умеренная гипергликемия & К SGLT2i либо GLP-1-терапии можно добавить метформин для достижения гликемических целей.\\[5pt]
$>8{,}5\%$: тяжёлая гипергликемия & Добавить семаглутид или тирзепатид, если пациент уже получает SGLT2i. При симптомах гипергликемии либо очень высоких HbA1c или глюкозе --- HbA1c $>10{,}0\%$ или глюкоза крови $\geq300$~мг/дл --- рекомендуется направление к эндокринологу для дальнейшего ведения согласно ADA.\\
\bottomrule
\end{longtable}
}
CVD --- ССЗ; GLP --- глюкагоноподобный пептид-1; HbA1c --- гликированный гемоглобин; SGLT2i --- ингибиторы натрий-глюкозного котранспортера 2 типа; T2D --- СД2.
\editorial{В пояснительном тексте оригинала граница выраженной остаточной гипергликемии записана как $\geq8{,}5$--9,0\%, а в таблице 20 --- $>8{,}5\%$. Оба варианта сохранены; перевод не устанавливает вместо них новую единую границу.}

\textbf{3. Повторная оценка альбуминурии.} Альбуминурия --- стойкий UACR $\geq30$~мг/г --- является определяющим признаком ХБП и сильным модифицируемым фактором риска ССЗ и прогрессирования ХБП.\src{32} Многие РКИ ССПМ-терапии включали альбуминурию в критерии отбора; абсолютное снижение риска сердечно-сосудистых и почечных исходов обычно было больше при более выраженной исходной альбуминурии.\src{12,15,33--38}

Поэтому альбуминурия служит показанием к препаратам для ХБП и ССПМ-синдрома: RASi, SGLT2i, GLP-1 RA и nsMRA (разделы 5.5.4, 6.2.3 и 6.3.3). Каждый из них уменьшает экскрецию альбумина с мочой.\src{8} Остаточная альбуминурия на лечении остаётся сильным фактором риска АССЗ, СН и прогрессирования ХБП.\src{9--11} Поэтому её повторная оценка позволяет определить ответ, выявить остаточный риск и своевременно добавить другие лекарства.\src{39}

Например, при СД2 остаточная альбуминурия, несмотря на RASi и SGLT2i, служит показанием к добавлению GLP-1 RA или nsMRA (раздел 5.5.4). Эффекты RASi, SGLT2i и nsMRA проявляются в течение нескольких недель,\src{40--42} поэтому повторное измерение уместно через 3--6 месяцев после начала или достижения максимальной запланированной дозы. Полный эффект GLP-1 RA\src{38} может проявиться лишь ближе к 6 месяцам, возможно, отчасти из-за времени, необходимого для титрации.

\textbf{4. Функция почек и калий.} Препараты, влияющие на РААС, включая RASi, MRA и SGLT2i, изменяют внутрипочечную гемодинамику и почечный обмен электролитов, потенциально меняя сывороточный креатинин и электролиты.\src{43} Повышение креатинина и калия на RASi и MRA возникает быстро и может оцениваться через 2--4 недели после начала.

Снижение рСКФ до 30\% после начала RASi или MRA ожидаемо и допустимо.\src{44} Если рСКФ снижается более чем на 30\%, необходимо искать другие причины: уменьшение объёма циркулирующей жидкости, НПВП, взаимодействия с лекарствами, влияющими на объём или внутрипочечную гемодинамику. Следует рассмотреть коррекцию терапии --- уменьшение дозы, отмену препарата, изменение сопутствующих лекарств либо тактики контроля объёма --- и последующий дополнительный мониторинг.

Аналогично умеренное повышение сывороточного калия до $\leq5{,}5$~мЭкв/л может быть допустимым; при $>5{,}5$~мЭкв/л могут потребоваться лекарства для уменьшения гиперкалиемии, например тиазидные или петлевые диуретики и средства, связывающие пищевой калий, изменения питания, снижение дозы или отмена препаратов, повышающих калий.

Низкая рСКФ и сопутствующие препараты, влияющие на клубочковое давление, например диуретики, SGLT2i и НПВП, являются факторами риска повышения креатинина и калия на RASi/MRA. Снижение рСКФ до 30\%, вызванное SGLT2i, не связано с неблагоприятными исходами.\src{13,14} Более выраженное снижение чаще встречается при низкой исходной рСКФ, СН, перенесённом остром повреждении почек и применении диуретиков;\src{13,14,45} у таких пациентов можно рассматривать контроль рСКФ после начала SGLT2i. SGLT2i уменьшают частоту гиперкалиемии на RASi.\src{46}

\section{Пробелы в доказательной базе и будущие направления}
Доказательная база оценки риска и клинического ведения ССПМ-синдрома быстро развивается. Однако сохраняется несколько ключевых пробелов, которые должны стать предметом будущих исследований; они подробно описаны далее и обобщены в таблице 21.

Калькулятор PREVENT существенно продвинул количественную оценку риска: он учитывает компоненты ССПМ-синдрома, добавляет исходы СН и всех ССЗ, определяет риск на горизонтах 10 и 30 лет и точнее прежних калькуляторов. Тем не менее прогнозирование требует дальнейшего уточнения. Необходимо количественно оценить вклад систематического учёта субклинических ССЗ, генетической предрасположенности и новых маркеров риска, включения многоуровневых SDOH и методов искусственного интеллекта.

Хотя 30-летний горизонт может лучше описывать риск у молодых взрослых, оптимальное использование таких оценок для выбора профилактики пока неясно. Требуют развития и стратегии объяснения риска с индивидуализированными инструментами поддержки решений. Необходимо уточнить, как включать в клинические решения усилители ССПМ-риска, связанные с большим риском прогрессирования синдрома и неблагоприятных сердечных и почечных событий. Среди них особенно важно для клиники и общественного здоровья понять влияние ключевых периодов жизни --- беременности и менопаузы --- на траектории ССПМ-здоровья.

Частая проблема --- позднее распознавание факторов и состояний, существенно повышающих риск неблагоприятных исходов. Особенно это относится к ХБП, о которой многие пациенты не знают, и MASLD, которая в настоящее время является ведущей причиной криптогенного цирроза и трансплантации печени в США. Более раннее выявление и своевременное начало доказанной терапии связаны с лучшими исходами. Ключевой клинический и популяционный вопрос: могут ли систематическая оценка ХБП одновременно по рСКФ и UACR и проверка фиброза печени по FIB-4 улучшить выявление, раннее лечение и исходы ХБП и MASLD во всей популяции?

На фоне растущей распространённости СН и ухудшения тенденций смертности желательно выявлять раннюю сердечную дисфункцию до клинической СН для целевой профилактики. В исследованиях определение сердечных биомаркеров с последующим согласованным ведением связано с меньшим риском СН. Но основные элементы междисциплинарной профилактической помощи при доклинической СН определены недостаточно. Эхокардиография уточняет абсолютный риск при повышенных биомаркерах, однако неизвестно, какие именно отклонения наиболее важны для прогноза. В целом нужны специализированные исследования профилактики СН, направленной на факторы и патологические процессы ССПМ, определяющие значительную долю риска.

Из-за полиорганной патологии пациенты часто наблюдаются у нескольких узких специалистов; поэтому стратегии уменьшения разобщённости помощи критически важны. Междисциплинарное ведение желательно, но роли и обязанности членов команды требуют уточнения. Особенно нужны доказательства влияния и оптимального использования выделенных координаторов ССПМ-помощи для обучения пациентов, помощи в прохождении маршрута лечения и междисциплинарного взаимодействия. Необходимо также определить, как работа с выявленными неблагоприятными SDOH влияет на помощь и клинические исходы. Следует уточнить эффективное внедрение междисциплинарных моделей в условиях разных ресурсов и потребностей.

Контроль массы тела --- основополагающая цель ССПМ-помощи, уменьшающая вероятность прогрессирования и способствующая регрессу синдрома. Пациентам более высокого риска рекомендуются комплексные команды по контролю массы тела для пациент-ориентированного выбора среди расширяющегося набора методов. Оптимальная организация и клинический эффект таких программ требуют дальнейшего изучения.

Ограниченное возмещение расходов на длительную поведенческую поддержку препятствует изменению образа жизни; стоимость и ценность инвестиций в эти программы --- важный вопрос организации здравоохранения. При GLP-1-терапии отмена часта и связана с возвратом массы тела и утратой нескольких благоприятных изменений профиля ССПМ-риска. Необходимо лучше определить эффективные стратегии долгосрочного поддержания сниженной массы тела на этой терапии.

СД2 и ХБП --- важнейшие факторы неблагоприятных сердечно-сосудистых и почечных событий. К счастью, растёт число методов с доказанным улучшением исходов. Однако интервенционных данных о дополнительной пользе и оптимальном сочетании этих методов пока мало. Остаются вопросы предпочтений пациентов и влияния полипрагмазии на приверженность и удовлетворённость. Учитывая обычно хроническое лечение СД2 и ХБП, экономические оценки комбинаций также помогли бы принимать решения в системе здравоохранения.

При СН сохраняются отдельные нерешённые вопросы. Польза лечения ожирения при HFmrEF/HFpEF теперь ясна: РКИ GLP-1-терапии показывают улучшение функциональных возможностей и качества жизни и снижение риска ухудшения СН. Напротив, подход при HFrEF неясен; некоторые исследования предполагают вред GLP-1-терапии. Необходимо уточнить лечение избытка массы тела при HFrEF, который часто препятствует применению методов лечения далеко зашедшей СН.

Финеренон показал пользу при диабетической болезни почек в FIDELIO и FIGARO, а также при HFmrEF/HFpEF в FINEARTS-HF, но данных при сочетании СН и ХБП мало. Дополнительные интервенционные исследования помогут определить его роль в этой частой и клинически важной подгруппе. Неясно и то, можно ли получить сходную ССПМ-пользу, наблюдаемую у нестероидных MRA, например финеренона, с помощью стероидных MRA --- эплеренона или спиронолактона.

Особенно мало интервенционных данных при далеко зашедшей ХБП и диализе. Такие пациенты преимущественно исключались из РКИ, поэтому убедительной базы для ведения этой подгруппы нет. В испытаниях нефропротективных препаратов пациентов обычно продолжали наблюдать после снижения рСКФ ниже порога начала терапии; сердечно-сосудистая и почечная польза сохранялась при приемлемой безопасности. Поэтому пришло время специализированных исследований при далеко зашедшей ХБП, чтобы уточнить лечение этой популяции высокого риска, для которой ССЗ являются ведущей причиной смерти.

Наконец, модель ССПМ-помощи означает переход от разобщённого к более интегрированному ведению нескольких взаимосвязанных состояний. Важно пересмотреть образование и подготовку кадров в соответствии с этим изменением. Необходимо определить и поддержать организационные условия: автоматизированные технологии, персонал, правила выставления счетов и возмещения расходов, телемедицину, общее физическое и виртуальное клиническое пространство, платформы для коммуникации. Это позволит внедрять оптимальную ССПМ-помощь в различных клинических и географических условиях.

\heading{Таблица 21. Будущие направления при ССПМ-синдроме: сводка исследовательских задач}
\editorial{Широкая четырёхколоночная таблица передана последовательными тематическими блоками. В каждом сохранены область, пробелы, ключевые вопросы и клиническое/организационное значение.}

\heading{1. Прогнозирование и стратификация риска}
\textbf{Пробелы в доказательствах.} PREVENT требует уточнения; использование 30-летнего риска неясно; инструменты объяснения риска ограничены.

\textbf{Ключевые вопросы.} Как биомаркеры, генетика и искусственный интеллект улучшают прогноз ССПМ-риска? Какова роль 30-летних оценок?

\textbf{Клиническое и организационное значение.} Повышение точности прогноза, поддержка индивидуализированной профилактики, обоснование стратегий вовлечения пациентов.

\heading{2. Раннее выявление ХБП и MASLD}
\textbf{Пробелы в доказательствах.} Позднее распознавание ХБП/MASLD; неясная роль систематической оценки; недостаточное применение скрининга фиброза.

\textbf{Ключевые вопросы.} Может ли систематический скрининг ХБП/MASLD улучшать выявление и исходы?

\textbf{Клиническое и организационное значение.} Более раннее лечение, уменьшение прогрессирования, обоснование правил обследования и планирования популяционного здоровья.

\heading{3. Профилактика СН}
\textbf{Пробелы в доказательствах.} Растущая распространённость СН; не определено междисциплинарное ведение; неясны прогностические эхокардиографические маркеры; мало данных о профилактической терапии.

\textbf{Ключевые вопросы.} Какие изменения при визуализации лучше всего предсказывают СН у лиц с повышенными биомаркерами? Каковы обязательные компоненты профилактической помощи?

\textbf{Клиническое и организационное значение.} Раннее вмешательство, снижение заболеваемости и смертности, определение структуры профилактической помощи.

\heading{4. Междисциплинарная и согласованная помощь}
\textbf{Пробелы в доказательствах.} Разобщённость помощи узких специалистов; неопределённые роли команды; не проверено влияние координатора; неясен эффект работы с SDOH.

\textbf{Ключевые вопросы.} Каковы исходы работы междисциплинарных ССПМ-команд? Как координаторы могут улучшать навигацию, обучение пациентов и ведение синдрома?

\textbf{Клиническое и организационное значение.} Улучшение координации, уменьшение разобщённости, поддержка возмещения расходов на командное ведение.

\heading{5. Стратегии контроля массы тела}
\textbf{Пробелы в доказательствах.} Неясна организация комплексных программ; существуют барьеры возмещения расходов; часты отмена GLP-1-терапии и возврат массы тела.

\textbf{Ключевые вопросы.} Какие комплексные программы наиболее эффективны? Что поддерживает долгосрочное сохранение снижения массы тела на GLP-1-терапии?

\textbf{Клиническое и организационное значение.} Улучшение долгосрочных исходов, обоснование возмещения расходов, обеспечение устойчивого снижения массы тела.

\heading{6. Сочетания терапии при СД2 и ХБП}
\textbf{Пробелы в доказательствах.} Мало интервенционных данных о дополнительной пользе; недостаточно изучены эффекты полипрагмазии; не хватает экономических исследований.

\textbf{Ключевые вопросы.} Каковы исходы комбинаций ССПМ-терапии? Как полипрагмазия влияет на приверженность и исходы?

\textbf{Клиническое и организационное значение.} Оптимизация лечения и приверженности, обоснование экономической эффективности и страхового покрытия.

\heading{7. ССПМ-синдром и фенотипы СН}
\textbf{Пробелы в доказательствах.} Польза лечения ожирения при HFpEF/HFmrEF установлена; при HFrEF подход неясен; сохраняется некоторая неопределённость для MRA.

\textbf{Ключевые вопросы.} Как лечить ожирение при HFrEF? Могут ли стероидные MRA воспроизвести пользу финеренона? Улучшает ли финеренон исходы ХБП при СН?

\textbf{Клиническое и организационное значение.} Безопасное лечение ожирения при HFrEF, оптимальное применение MRA, улучшение исходов сочетания СН и ХБП.

\heading{8. Далеко зашедшая ХБП и диализ}
\textbf{Пробелы в доказательствах.} Исключение из крупных РКИ; отсутствие доказательств терапии при далеко зашедшей ХБП/диализе.

\textbf{Ключевые вопросы.} Какие методы безопасны и эффективны? Сохраняется ли польза при продолжении терапии на низкой рСКФ?

\textbf{Клиническое и организационное значение.} Ответ на потребности популяции высокого риска, обоснование дизайна исследований и применения терапии на диализе.

\heading{9. Образование, подготовка кадров и инфраструктура}
\textbf{Пробелы в доказательствах.} Обучение остаётся разобщённым; недостаточны организационные механизмы поддержки --- автоматизация, телемедицина, возмещение расходов.

\textbf{Ключевые вопросы.} Как должны меняться учебные программы с учётом ССПМ-синдрома? Какая инфраструктура лучше поддерживает интегрированную помощь?

\textbf{Клиническое и организационное значение.} Подготовка кадров, поддержка инвестиций в инфраструктуру и организационных преобразований.

\heading{10. Сквозные темы}
\textbf{Пробелы в доказательствах.} Неравенство; ограниченность исследований внедрения; недостаточный акцент на значимых для пациентов исходах; недостаточно разработана глобальная применимость.

\textbf{Ключевые вопросы.} Как встроить обеспечение справедливости, науку о внедрении, пациент-ориентированные исходы и глобальное здоровье во все направления ССПМ-помощи?

\textbf{Клиническое и организационное значение.} Уменьшение различий, масштабирование вмешательств в мире, улучшение значимых для пациентов исходов, обоснование мер по обеспечению справедливости в здравоохранении.

\textbf{Сокращения исходной таблицы 21.} CKD --- ХБП; CKM --- ССПМ; GFR --- СКФ; GLP-1 --- глюкагоноподобный пептид; HF --- СН; HFmrEF --- СН с умеренно сниженной фракцией выброса; HFpEF --- СН с сохранённой фракцией выброса; HFrEF --- СН со сниженной фракцией выброса; MASLD --- стеатотическая болезнь печени, ассоциированная с метаболической дисфункцией; MRA --- антагонист минералокортикоидных рецепторов; PREVENT --- Predicting Risk of Cardiovascular Disease Events, прогнозирование риска сердечно-сосудистых событий; SDOH --- социальные детерминанты здоровья; T2D --- СД2.

\section*{Президенты и сотрудники организаций}
\addcontentsline{toc}{section}{Президенты и сотрудники организаций}
\editorial{Имена, квалификации и должности переданы по исходной публикации; этот перечень не является независимо обновлённым справочником руководства организаций на дату перевода.}

\heading{Американская коллегия кардиологов (ACC)}
\begin{itemize}
\item Roxana Mehran, MD, FACC --- президент.
\item Cathleen C. Gates --- главный исполнительный директор.
\item Richard J. Kovacs, MD, MACC --- главный медицинский советник.
\item Justine Varieur Turco --- вице-президент подразделения научных публикаций и клинических рекомендаций.
\item Mindy J. Saraco, MHA --- директор по клинической политике и рекомендациям.
\item Grace D. Ronan --- старший менеджер по производству и операционной деятельности в области публикаций по клинической политике.
\item Lauren N. Prestera, FNP-BC --- менеджер проектов по клиническому содержанию.
\item Leah Patterson --- менеджер проектов по клиническому содержанию.
\end{itemize}

\heading{Совместная структура ACC/AHA}
\begin{itemize}
\item Thomas S.D. Getchius --- национальный старший директор по клиническим рекомендациям.
\item Abdul R. Abdullah, MD --- директор по научным основам и методологии рекомендаций.
\item Shae Martinez, MLS --- консультант по справочно-библиографическому обеспечению, медицинский библиотекарь.
\item Rachel L. Schunder, MA --- советник по науке и здравоохранению в области рекомендаций.
\end{itemize}

\heading{Американская кардиологическая ассоциация (AHA)}
\begin{itemize}
\item Stacey E. Rosen, MD, FAHA --- президент.
\item Nancy Brown --- главный исполнительный директор.
\item Mariell Jessup, MD, FAHA --- главный научный и медицинский руководитель.
\item Nicole Aiello Sapio, EdD --- исполнительный вице-президент управления научной стратегии и операционной деятельности.
\item Radhika Rajgopal Singh, PhD --- старший вице-президент управления науки и медицины.
\item Prashant Nedungadi, BPharm, PhD --- национальный вице-президент по науке и медицине в области клинических рекомендаций.
\item Mu Huang, PhD, DPT --- советник по науке и медицине.
\item Heather Hartzer, MSN, APRN, AGACNP-BC --- младший советник по науке и медицине.
\item Joseph W. Loftin III --- национальный директор по научным заявлениям и рекомендациям.
\item Courtney Goodwin, MPH --- менеджер программы клинических рекомендаций.
\end{itemize}

\section*{Члены комитета рецензентов}
\addcontentsline{toc}{section}{Члены комитета рецензентов}
{\small
Nicole Bhave, MD, FAHA, FACC --- сопредседатель; Amit Khera, MD, MSC, FAHA, FACC --- сопредседатель; Kevin M. Alexander, MD, FAHA, FACC; Alison L. Bailey, MD, FAHA, FACC; Marc P. Bonaca, MD, MPH, FAHA, FACC; Tara I. Chang, MD, MS, FASN, FAHA --- представитель ASN; Cyrille K. Cornelio, PharmD, FAHA; Maria Rosa Costanzo, MD, FAHA, FACC; Colette DeJong, MD, FAHA;

Adam DeVore, MD, FAHA, FACC --- представитель JCPM; Patrick O. Gee, Sr., PhD --- представитель интересов пациентов; Beverly B. Green, MD, MPH, FAHA; Jennifer B. Green, MD, FAHA --- представитель ADA; Scott Kahan, MD, MPH, FAHA --- представитель Ассоциации по вопросам ожирения; Anuradha Lala, MD, FAHA, FACC; Melissa Magwire, RN, MSN, FAHA; Michael Miller, MD, FAHA, FACC;

Katie Moegenberg, LPN --- представитель интересов пациентов; Elif Oral, MD, FAHA; Carl E. Orringer, MD, FAHA, FACC; Kershaw V. Patel, MD, MSCS, FAHA; Sonali Patel, MD, PhD, FAHA; Jane E.B. Reusch, MD, FAHA; Philip Schauer, MD, FAHA; Daniel E. Weiner, MD, MS, FASN, FAHA.
}
\editorial{Обозначения представителей организаций и пациентов раскрыты согласно легенде сносок на странице e50 оригинала.}

\section*{Объединённый комитет ACC/AHA по клиническим практическим рекомендациям}
\addcontentsline{toc}{section}{Объединённый комитет ACC/AHA по клиническим практическим рекомендациям}
{\small
Catherine M. Otto, MD, FAHA, FACC --- председатель; Sunil V. Rao, MD, FACC, FSCAI --- избранный следующий председатель; Anastasia Armbruster, PharmD, FACC*; Vanessa Blumer, MD, FACC; Leslie L. Davis, PhD, ANP-BC, FAHA, FACC; Sharlene M. Day, MD; Dave L. Dixon, PharmD, FAHA, FACC, BCACP, CLS; Victor A. Ferrari, MD, FAHA, FACC;

Stephen E. Fremes, MD, FACC; Mario F.L. Gaudino, MD, PhD, FAHA, FACC; Hani Jneid, MD, FAHA, FACC; Heather M. Johnson, MD, MS, FAHA, FACC*; William Schuyler Jones, MD, FACC; Sadiya S. Khan, MD, MSc, FAHA, FACC; Prateeti Khazanie, MD, MPH, FAHA; Michelle M. Kittleson, MD, PhD, FAHA, FACC*;

Venugopal Menon, MD, FAHA, FACC; Debabrata Mukherjee, MD, MS, FAHA, FACC*; Daniel Munoz, MD, MPA, FACC; Kristen K. Patton, MD, FAHA, FACC; Garima Sharma, MBBS, FACC; Daichi Shimbo, MD; Joseph Woo, MD, FAHA, FACC.
}
*Бывший член рабочей группы; входил в её состав во время подготовки документа.

% --- ИСПРАВЛЕНИЕ: закомментированы строки \input, так как файлы отсутствуют ---
% \input{bibliography-ru}
% \input{appendices-ru}

% Вместо них добавлены пояснительные заголовки
\section*{Библиография}
\addcontentsline{toc}{section}{Библиография}
(Библиография не включена в данный файл перевода; см. оригинальное руководство.)

\section*{Приложения}
\addcontentsline{toc}{section}{Приложения}
(Приложения не включены в данный файл перевода; см. оригинальное руководство.)

\clearpage
\section*{Завершение перевода и статус проверки}
\addcontentsline{toc}{section}{Завершение перевода и статус проверки}
Перевод охватывает документ целиком: страницы e50--e158 журнала \textit{Circulation} (109 страниц PDF). Включены вводные материалы, все основные разделы 1--9, таблицы 1--21, текстовые адаптации рисунков 1--20, сведения об авторах, сотрудниках и комитетах, сведения о публикации, все 1054 библиографические записи с русскими переводами названий и оба приложения с раскрытием связей. В приложении 1 сохранены сведения о 36 членах авторского комитета, в приложении 2 --- о 25 рецензентах.

Адаптация рисунков и широких таблиц меняет форму представления, но не заменяет их содержательные элементы пересказом выводов. Библиографические записи сохранены по разделам оригинала; повторное цитирование одной работы в разных разделах не объединено. Внешние источники, полные тексты цитируемых публикаций, подкаст, таблицы доказательств и иные дополнительные материалы, не входящие в предоставленный PDF, в этот перевод не включены.



\end{document}
